\section{Setup}
\label{sec:1}

We take unimodal maps with critical point $0$ and exponent $\ell>1$ in the
form $g(x) = G(|x|^\ell)$, $G(0) = 1$, $G$ analytic on $D(c,r) \subset \CC$. The
Feigenbaum--Cvitanović equation \cite{feigenbaum1978quantitative, feigenbaum1979universal, coullet1978iterations} $g(x) = -\alpha \cdot g(g(x/\alpha))$, $\alpha = -1/g(1)$, in the
$t$-coordinate $t = |x|^\ell$ takes the form $F(G) = 0$, where

\begin{gather*}
F(G)(t) := G(t) + \alpha_G \cdot G(W_G(t)), \\
\alpha_G = -1/G(1), \\
W_G(t) = U_G(t)^{\ell}, \\
U_G(t) = G(s_G(t)), \\
s_G(t) = t \cdot \alpha_G^{-\ell}.
\end{gather*}

Well-definedness of the composition for $\ell \in [2, 4]$: on the disk $D(c,r)$ it is rigorously
verified that $\sup |U_G - 1| \le 0.379 < 1$ (the domain lemma for the center $P(\ell)$ --- \S\ref{sec:6.4.2}, its
transfer to the solution $\Gs$ --- \S\ref{sec:6.4.3}, certified over 46 boxes --- the acceptance
flag \texttt{nonzero\_ok}, checked rigorously by
\texttt{04-\allowbreak CODE/\allowbreak tests/\allowbreak \_\allowbreak nozeros\_\allowbreak check.py}), so that $U_G$ never vanishes and its image lies in the right half-plane, so
that $U^{\ell} = \exp(\ell \cdot \Log U)$ is single-valued.

The solution $\Gs$ of the equation $F(G) = 0$ is the ``universal scaling
function.'' For the origins of the subject and the early literature see the reprint volume \cite{cvitanovic1989universality}. What is known about it to date:

\textbf{Existence, any $\ell > 1$.} Proved twice, non-constructively:
Epstein (1989) \cite{epstein1989fixed} --- by a holomorphic method, for every
exponent $\ell > 1$; the earlier Epstein (1986) \cite{epstein1986new} proves
existence at a prescribed exponent only for $\ell \le 14$ (its main construction is
parametrized by $\lambda$, with the exponent an output), which already covers the
segment considered here --- a contemporary reading of it, Briggs (1991)
\cite{briggs1991precise}, p.~436, says exactly ``for $z < 14$ by Epstein in 1985'';
Martens (1998) \cite{martens1998periodic} --- by a real method, for an arbitrary exponent.
Neither proof gives an explicit approximation to $\Gs$, nor explicit
bounds.

\textbf{Hyperbolicity, even integer $\ell$.} Lanford (1982) \cite{lanford1982computer} --- a computer-assisted
proof at $\ell = 2$ in a fixed layer of analytic functions, later
fully closed by Eckmann--Wittwer (1987) \cite{eckmann1987complete}; the 1984 paper of the
same author \cite{lanford1984shorter} is a shorter proof of \emph{existence}, not of
hyperbolicity; conceptual proofs by Sullivan,
McMullen, and Lyubich for the quadratic case
\cite{sullivan1992bounds, mcmullen1994complex, lyubich1999feigenbaum}, extended by
Avila--Lyubich (2011) \cite{avila2011full} to an arbitrary even degree.
Computer certification at $\ell = 4$ with precision beyond 320 digits ---
Burbanks--Osbaldestin--Thurlby (2021) \cite{burbanks2021rigorous}; a separate line of work by
Breden--Gonzalez--Mireles James (2026) \cite{breden2024validated} varies the renormalization order
$m = 2\dots 10$ at degree $d=2$; separately, the same authors extend the result
to degree $d=4$, but only for $m=2,3$ (Theorem 1.6 of the same work).
Beyond period doubling of unimodal maps, the same computer-assisted style has
been applied to other renormalization fixed points --- Stirnemann (1994)
\cite{stirnemann1994siegel} to the Siegel disc renormalization, and
Arioli--Koch (2010) \cite{arioli2010critical} to the critical fixed point of
MacKay's renormalization for commuting pairs of area-preserving maps.

\textbf{Hyperbolicity, non-integer $\ell$.} The first result of this kind is older than the
interval treated here and lies below it: Collet--Eckmann--Lanford (1980)
\cite{collet1980universal} work with the family $1-(1+a)|x|^{1+\varepsilon}$ and obtain, for
every sufficiently small $\varepsilon > 0$, a fixed point of the doubling operator with a
one-dimensional unstable manifold transversal to a codimension-one stable manifold. The
exponent $1+\varepsilon$ is non-integer, but it is near $1$, and the authors state plainly
that their argument is perturbative and valid only for small $\varepsilon$; nothing in it
reaches $[2,4]$. Closer to the present setting,
Gorbovickis--Yampolsky (2018) \cite{gorbovickis2018renormalization} construct
an analytic setup for renormalization at an arbitrary exponent and
prove hyperbolicity with a one-dimensional unstable direction for
maps of bounded type at an exponent sufficiently close to an
even integer; independently, Cruz--Smania (2010) \cite{cruz2010renormalization} obtain the same for
neighborhoods of even integers. Both proofs are complex-analytic, computer-free,
and \textbf{the width of the neighborhood is not written out} in them: for a given specific
$\ell$ other than an even integer, one cannot conclude from these works whether it is
covered or not.

This is exactly what defines what has not existed until now: an explicit
statement, applicable to a specific exponent, covering an entire interval of
non-integer $\ell$ together with its interior --- and quantitative bounds instead
of pure existence.


\subsection{The t-coordinate and the invariant class}
\label{sec:1.1}

In the $t$-coordinate $t=|x|^\ell$ the operator $T$ is \textbf{fully real-analytic and
complexifiable}: $U(t) = G(t\cdot\alpha_G^{-\ell})$ lies in $[G(\alpha_G^{-\ell}), 1] \subset (0, 1]$ --- strictly
bounded away from zero. Hence $W = U^{\ell} = \exp(\ell \cdot \Log U)$ is holomorphic on a suitable
complex neighborhood avoiding the branch cut (this is an observation; a detailed
treatment --- \S\ref{sec:6.3} ``Complexification of $X$ and of the Operator $F$'' and \S\ref{sec:6.4} ``The Domain
Lemma''). Non-integrality of $\ell$ \textbf{introduces no singularities} into the domain --- the
only singularity of $|x|^\ell$ sits at $x=0$ and is absorbed by the substitution $t=|x|^\ell$.

The working space:

\[
X = \bigl\{\, h(t) = \sum h_k (t-c)^k \; : \; \norm{h}_X = \sum |h_k| \cdot r^k < \infty \,\bigr\},
\]

the disk $D(c,r)$ with center $c=3/5$ and radius $r=4/5 \supset [0,1]$.
This is a space of the same \emph{type} as the one in which Lanford (1982)
\cite{lanford1982computer} stated his historic result at $\ell=2$ --- functions of $t$,
analytic on a disk in $t$, with a weighted $\ell^{1}$ norm on the Taylor coefficients ---
but it is not his disk and does not contain it: Lanford works on $\{z : |z^{2}-1| < 5/2\}$,
that is on $|t-1| < 5/2$ in the $t$-coordinate, and $D(3/5, 4/5)$ lies strictly inside that.
His layer is also far outside the family of layers exhibited in \S\ref{sec:10.6}, so
Theorem D says nothing about it. The
affine slice $\Aff = \{G \in X : G(0) = 1\}$ (the normalization $G(0)=1$) is preserved by the
operator $T$ (from $F(G)(0) = 0$); on the tangent subspace $X_{0} = \{h \in X : h(0) = 0\}$ its
derivative acts, and $DF(\Gs)|_{X_{0}}$ is invertible (proved below: \S\ref{sec:6.6}) --- it is on
precisely this linear slice that all further linear analysis is conducted (exact
notation --- \S\ref{sec:2}).


\subsection{Polynomial vs. Irrational Case}
\label{sec:1.2}

At \textbf{odd integer $\ell=3$} the power $W = U^3$ is a polynomial, the entire composition $F(G)$
for a polynomial $G$ remains a polynomial, and the residual/derivative are computed by
direct polynomial operations (without tails). Sources --- all in
\texttt{03-\allowbreak CERTIFICATES/\allowbreak pointwise/\allowbreak }: \texttt{kant\_l3.json} (the Newton--Kantorovich gate: \texttt{Y0},
\texttt{A\_norm1}), \texttt{eigen\_\allowbreak l3\_\allowbreak fixed.json} (the simple eigenvalue \texttt{delta\_bar} and its box),
\texttt{fixL\_\allowbreak result.json} (the radii and norms \texttt{rho\_prime}, \texttt{rho\_X}, \texttt{m1}, \texttt{Gprime\_norm}),
\texttt{df\_\allowbreak package.json} (norms of the derivative: \texttt{q}, \texttt{q\_s}, \texttt{norm\_\allowbreak Gprime\_\allowbreak W}). At
\textbf{fractional $\ell=2.5$} $\Wbar = \Ubar^{5/2}$ is no longer a polynomial; the Taylor coefficients
are computed by J.~C.~P. Miller's formula for the coefficients of a power of a
power series \cite[\S4.7]{knuth1997taocp2} (derived in \S\ref{sec:5.3.2}, item 5,
from $W'U = aWU'$; this is a forward $n$-term convolution, not the backward
three-term recurrence of \cite{gautschi1967computational}); the well-definedness of the branch ---
\texttt{03-\allowbreak CERTIFICATES/\allowbreak pointwise/\allowbreak reu\_\allowbreak min\_\allowbreak l25.json} (\S\ref{sec:5.3.2}); the existence certificate at
$\ell = 5/2$ --- \texttt{kant\_\allowbreak l25.json} in the same location.


\section{Main results}
\label{sec:2}

Throughout what follows the space is fixed: analytic functions on the disk
$D(c, r)$ with $c = 3/5$, $r = 4/5$, norm $\norm{h}_X = \sum_k |h_k| r^k$; coordinate
$t = |x|^\ell$; $X_{0} = \{h \in X : h(0) = 0\}$ --- the linear subspace serving as the TANGENT
direction to the normalization $G(0) = 1$: the derivatives ($DF|_{X_{0}}$, $DT|_{X_{0}}$; the
operator $T(G) := -\alpha_G \cdot G \circ W_G$, so that $F = I - T$ and
$DT = I - D_G F$; the slice is preserved because $W_G(0) = 1$ and hence
$T(G)(0) = -\alpha_G\, G(1) = 1$, and the holomorphy of $DT$ is where this identity is
used, \S\ref{sec:6.5.1})
and the residual ($F(G) \in X_{0}$, since $F(G)(0) = 0$ identically on $\{G(0)=1\}$) do
indeed lie in $X_{0}$. The normalization itself, however, defines not a linear subspace
but an \textbf{affine slice} $\Aff := 1 + X_{0} = \{G \in X : G(0) = 1\}$ --- it is there (not in
$X_{0}$) that the solution $\Gs(\ell)$ lives. The statements of Theorems A and C below and
the pointwise theorems of \appref{sec:5} write this explicitly as $\Gs(\ell) \in \Aff$. If a shorter
notation $\Gs(\ell) \in X_{0}$ nonetheless appears somewhere in the text, this is the
shorthand adopted here for the same thing, that is, for the coordinates
$h = \Gs(\ell) - P(\ell) \in X_{0}$ relative to the center $P(\ell)$; the analytic implicit
function theorem in \S\ref{sec:6.6} is applied on the complexification of $\Aff$ (an affine shift of
the complexification of $X_{0}$), and not on the complexification of $X_{0}$ as the set of
the solutions themselves --- details and formal notation $\Aff_{\CC}$, $X_{0,\CC}$ --- \S\ref{sec:6.1}, \S\ref{sec:6.6.3}.
The projection onto the first $N_{p} = 80$ coordinates is used as a computational
device; the contribution of the tail $k > N_{p}$ is estimated explicitly and enters
the bounds rather than being discarded.

\begin{namedtheorem}[Existence and local uniqueness]
For every $\ell \in [2, 4]$ the Feigenbaum--Cvitanović equation
$F(G) = G + \alpha_G \cdot G \circ W_G = 0$, $\alpha_G = -1/G(1)$, has a solution
$\Gs(\ell) \in \Aff$ in the ball $\norm{G - P(\ell)}_X \le \rhostar_X(\ell)$ around an explicitly constructed
parametrized center $P(\ell)$, where $\rhostar_X(\ell) \le 8.49\cdot 10^{-13}$; this solution is
unique in the affine ball $P(\ell) + B_{X_{0}}(0, 1.75\cdot 10^{-7}) \subset \Aff$, that is,
on the normalization slice $\Aff$, not in the whole of $X$.
\end{namedtheorem}

Proof --- \S\ref{sec:3}: 46 boxes, the Newton--Kantorovich gate \cite{kantorovich1948newton}
$Y0 + B_A \cdot d_{\mathrm{chk}} + \Zhat \rho' \le \rho'$ with $Y0 \le 4.31\cdot 10^{-13}$, $d_{\mathrm{chk}} \le 7.98\cdot 10^{-15}$,
$\Zhat \le 0.1705$; the union of the boxes covers $[2.0000, 4.0075]$.

\begin{namedtheorem}[Hyperbolicity, uniform in $\ell$]
For every $\ell \in [2, 4]$ the derivative $DT(\Gs(\ell))|_{X_{0}}$ has \textbf{exactly
one} eigenvalue outside the unit circle; the rest of the spectrum lies
strictly inside. Consequently the unstable direction is one-dimensional, and
this eigenvalue $\delta(\ell)$ itself is real and simple.
\end{namedtheorem}

Proof --- \S\ref{sec:4}: 706 boxes, Rouché's theorem in operator form on the
unit circle, $\sup_{|z|=1}\norm{(DT_{0}-z)^{-1}}_1 \le 3.7961$,
$\norm{\Delta}_1 \le 0.280$, the worst value of the full acceptance test $0.9919 < 1$;
a chain of boxes with an anchor at $\ell = 3$; the union covers
$[2.0000, 4.0012]$. The anchor is confirmed twice by methods independent in
how they count the spectrum (\S\ref{sec:4.7} and \S\ref{sec:5.2}), but applied to a single assembly
of the matrix $DT$; the assembly itself is cross-checked by a second path
independent of the analytic derivative (\S\ref{sec:3.9}).

\begin{namedtheorem}[Analytic dependence on $\ell$]
The maps $\ell \mapsto \Gs(\ell) \in \Aff$ ($\Aff = 1 + X_{0}$ --- the affine slice of the normalization
$G(0)=1$, introduced above in this section) and $\ell \mapsto \delta(\ell) \in \RR$ are real-analytic
on $[2, 4]$: for every $\ell_0 \in [2,4]$ there exists a circular neighborhood
$U \subset \CC$ of the point $\ell_0$ and unique analytic continuations, coinciding
with $\Gs$ and $\delta$ on $U \cap [2, 4]$.
\end{namedtheorem}

Proof --- \S\ref{sec:6}, with the proof itself in
\S\S\ref{sec:6.3}--\ref{sec:6.8}: complexification of $X$, holomorphy of $F$ in $(G, \ell)$
based on the rigorous majorant $\sup |U_G - 1| \le 0.379 < 1$ over the disk,
the analytic implicit function theorem under a uniformly invertible
$D_G F$ ($Z1 \le 0.1501$), gluing of the branches through the uniqueness ball
$1.75\cdot 10^{-7}$, and the Riesz projector over the unit circle, where
$\sup \norm{(DT_{0}-z)^{-1}}_1 \le 3.7961$ for the leading matrix, and for the full operator ---
finiteness of the resolvent from the same certificates (\S\ref{sec:6.8}). The transfer of
the majorant from the center $P(\ell)$ to the solution is justified by the
explicit constant $\kappa_G \le 4.14$.

Both statements are \textbf{uniform in $\ell$}: this is not a count over a finite set of points
and not an argument ``by continuity.'' For each box an inequality is verified that is
valid simultaneously for all $\ell$ in that box, and the boxes overlap.

\textbf{What is new here.} Neither the existence of the fixed point at non-integer $\ell$,
nor the fact of hyperbolicity itself at $\ell$ close to an even integer, is new
(\S\ref{sec:1}). What is new is:

(i) \textbf{an explicit covering of the entire interval}, including its interior
($\ell = 2.5$, $\ell = 3$, and so on) --- points for which the proofs \cite{gorbovickis2018renormalization,cruz2010renormalization}
themselves do not write out the width of the neighborhood and therefore do not
allow one to conclude whether they are covered or not (see \S\ref{sec:1});

(ii) \textbf{quantitativeness}: for each $\ell$ a computable center
$P(\ell)$ and a ball containing the solution are exhibited, and for hyperbolicity ---
an explicit bound on the resolvent on the unit circle and an explicit gap in the
Rouché inequality. Prior proofs of hyperbolicity at non-integer $\ell$ give
not a single verifiable constant and do not allow one to point to a single
specific $\ell$ for which the statement is established. A caveat on external
verification: numerical values of $\alpha$ and $\delta$ as functions of a
\emph{continuous} critical-point order have been published since the early 1980s ---
Hu--Mao (1982) \cite{humao1982period} and van der Weele--Capel--Kluiving (1987)
\cite{vanderweele1987period}, for the same family $1 - \lambda|x|^{z}$ --- but the
comparison made here (\S\ref{sec:3.9}) is confined to the endpoints of the interval, where
values known to ten or more digits are available, and no comparison against those two is
attempted. At the interior point $\ell = 3$ a published pair for this family does exist, and
in the primary literature: Briggs \cite{briggs1991precise}, Table~2, prints
$\alpha(3) = -1.92769096384764084494999435296631905189265896703673$ and
$\delta(3) = 5.96796870377745104099419301997967232351260291982742$ to fifty
digits (the same values appear in his thesis, Table~2.3). The two halves of that
pair behave differently against this work, and both facts are worth stating.

The value of $\alpha(3)$ \textbf{agrees}: the scan here gives
$\alpha(3) = 1.9276909638476356$, which shares $14$ significant digits with Briggs
(relative difference $2.73\cdot 10^{-15}$). This is a genuine external comparison at
an \emph{interior} point of the interval, against a primary source computed by a
different method (collocation for the fixed point, $\alpha = 1/g(1)$), and it is the
only such comparison this work has.

The value of $\delta(3)$ \textbf{does not agree}: it differs in the second
significant figure from the Krawczyk enclosure $\deltastar(3) = 6.084689074262122$
of box width $6.226\cdot 10^{-19}$ quoted in \S\ref{sec:3.9}. Three independent
observations place the discrepancy on the side of the tabulated value rather than
ours. First, recomputing the superstable cascade of $1-\mu|x|^{d}$ from Briggs's
own definitions --- $\delta = \lim(\mu_j-\mu_{j-1})/(\mu_{j+1}-\mu_j)$ with $\mu_j$
the root of $f^{2^{j}}(0)=0$ --- in independent multiprecision arithmetic
reproduces $4.6692$ at $d=2$ and $7.2847$ at $d=4$, that is, Briggs's own numbers,
and at $d=3$ gives $6.08469$, that is, ours. Second, in Briggs's method $\alpha$ and $\delta$ are
obtained by two separate steps --- $\alpha$ read off the fixed point, $\delta$ as the
leading eigenvalue of a finite-dimensional approximation of $DT$ computed for fixed
$\alpha$ --- so agreement of $\alpha$ to $14$ digits shows that the fixed point is
the same in both works, and localises the disagreement in the spectral step.
(A third argument, from concavity of $\delta$ in $d$, suggests itself and does
\emph{not} work: across the eleven rows of that table the discrete concavity
$\delta(d) \ge (\delta(d-1)+\delta(d+1))/2$ fails not only at $d = 3$ (by
$0.009$) but also at $d = 8$ (by $0.047$) and $d = 10$ (by $0.081$), so no
conclusion about a single row can be drawn from it.)
Nothing below rests on the tabulated $\delta(3)$; it is recorded here because a
reader comparing the two numbers is owed the reason they differ. What additionally
stands at $\ell = 3$ is the cross-assembly of $DT$ by two different routes
(\S\ref{sec:3.9});

(iii) \textbf{uniformity in $\ell$}: the statement is proved not at points but
on boxes --- on each one an inequality is verified that is valid simultaneously
for all $\ell$ in that box, and the boxes overlap;

(iv) \textbf{method}: a parametrized center as a function of $\ell$ plus suprema over an
$\ell$-box, computed intervally with adaptive subdivision; transfer of the
count of eigenvalues along a chain of overlapping boxes by the operator Rouché
theorem from a single certified anchor.

The relation of this work to its predecessors is the same as that which Lanford's
computer-assisted proof has to the conceptual proofs of Sullivan--McMullen--Lyubich
that appeared later: a different type of statement about the same object, one
that is verifiable, including with numbers.

\textbf{What is not proved} is listed in \S\ref{sec:9}; briefly: outside $[2, 4]$ nothing
is asserted, and independence from the layer $(c, r)$ is proved (Theorem D,
\S\ref{sec:10}) not for an arbitrary layer, but within the family $\Ffam$, whose
exhibited members are
14 pointwise layers, all of $\Qcell$, and the connecting corridor of \S\ref{sec:10.13}, on all
of which the transfer $P \to \Gs$ of \S\ref{sec:10.7} is certified. A systematized summary table of
assumptions --- \S\ref{sec:10.11}.

\ifx\JOURNALVERSION\undefined
\textbf{A note on numbering.} A bare letter on a \emph{statement} --- Theorem
A, B, C, D --- is always one of the four main results: A, B and C are stated
just above, and D in \S\ref{sec:10.5}. No other
statement in this article carries a bare letter. The single appendix is
Appendix~A, and its letter is unrelated to Theorem~A. Everything else numbered here carries a digit and
belongs to the section in which it stands: Remark 3.1 in \S\ref{sec:3},
Theorem 4.1 in \S\ref{sec:4}, and so on. Two conventions cut across this. The
four lemmas that carry the fixed point from one layer of the family to the next
are named Lemma 1 to Lemma 4 throughout: they stand outside every counter, and
a bare \emph{Lemma} with a digit always means one of those four. And a
subsection number may coincide with the number of a statement --- \S5.2 is a
subsection, Remark 5.2 is a statement standing elsewhere --- so the section
sign is what distinguishes them.
\fi



\section{Parametric existence certificate: [2,4] closed by forty-six boxes}
\label{sec:3}

The earlier approach --- the \emph{march}: a chain of micro-intervals, on each of which
the center of the ball is static and $\ell$ is interval-valued (comparison with it --- \S\ref{sec:3.8}).
The march works, but pays for every micro-step: in 8 hours on eight machines ---
that is, 64 core-hours --- it covered 15.15\% of the interval, from which the total
cost extrapolates linearly to \textbf{$\approx 422$ core-hours} (\texttt{64 / 0.1515}; the fraction is
computed from \texttt{05-\allowbreak DATA/\allowbreak coverage.json::\allowbreak covered} --- the total length of the march's
coverage over an interval of length 2). Described here is a method that closes the
same interval in \textbf{42.2 core-minutes} on a single machine and --- more importantly
--- that yields a statement uniform in $\ell$ rather than a chain of pointwise ones.
An earlier version of this section stated ``$\sim$1200 core-hours'': the figure was not derived
from the very foundation given and overstated the cost of the earlier approach by
a factor of 2.8.


\subsection{The idea: parametrize the center rather than move it}
\label{sec:3.1}

In the march, each subinterval needs its own center, and the whole difficulty is in
dragging it along: polishing, drift, re-seeding. The parametric certificate
removes this dynamics: the center is taken directly as a \textbf{function of $\ell$},

\[
P(\ell) = \sum_j c_j \cdot T_j(s), \qquad s = 2(\ell - \ell_0)/w, \qquad \ell \in [\ell_0 - w/2, \ell_0 + w/2],
\]

a Chebyshev interpolant over $m+1$ nodes inside the box, whose coefficients are
obtained from high-precision solutions at the nodes. The approximate inverse
operator $\Ahat(\ell)$ is parametrized analogously. The Newton--Kantorovich gate
is verified \textbf{for the whole box at once}:

\[
Y0 + B_A \cdot d_{\mathrm{chk}} + \Zhat \cdot \rho' \le \rho', \qquad \Zhat = Z1 + 1.75^2 \cdot B_A \cdot L \cdot \rho' + Z_{\mathrm{conv}},
\]

where $Y0 \ge \sup_\ell \norm{\Ahat(\ell) \widetilde{F}(P(\ell),\ell)}$, $Z1 \ge \sup_\ell \norm{I - \Ahat(\ell) M'(\ell)}_1$, the suprema
being over the whole box, not over the nodes, and $Z_{\mathrm{conv}} = B_A \cdot L \cdot \norm{\Delta P_{\mathrm{conv}}}$
is the correction for the center conversion (Taylor center versus Chebyshev center;
for $Y0$ the same discrepancy is accounted for by the term $B_A \cdot d_{\mathrm{chk}}$).
Numerically $Z_{\mathrm{conv}}$ is negligible (on box \texttt{g01}, $Z_{\mathrm{conv}} \approx 1.1 \cdot 10^{-12}$ in the deposited run against $\Zhat \approx 8.4 \cdot 10^{-2}$; a rerun of that box gives $5.95\cdot 10^{-13}$, since $Z_{\mathrm{conv}}$ inherits the float-seed non-determinism of $d_{\mathrm{chk}}$), but it enters the gate explicitly:
\texttt{rig2fix3.py:764-765}. It is not deposited as a field of any certificate, and the lane logs of the
deposited run do not contain it either: the run driver
(\texttt{runplan6.sh:24}) filters stdout through a \texttt{grep} whose
pattern does not match the line that carries it. To see the number one has
to rerun \texttt{rig2fix3.py}, which prints it at line 774 as
\texttt{(Zconv=...)}.

Here and everywhere below $\widetilde{F}$ is the
same residual $F$, written in the coordinates of the slice $\Xz$ (the coefficients
$k \ge 1$ of the b-basis); the zeroth coordinate is dropped not for convenience but
because it is identically zero: $F(G)(0) \equiv 0$ on $\{G(0) = 1\}$ (\S\ref{sec:2}). Hence the
zeros of $\widetilde{F}$ and of $F$ on the slice are one and the same set, and the existence
certificate, formulated for $\widetilde{F}$, gives the zero of $F$ that Theorem A speaks of.
The factor $1.75$ is not a fudge: it is $\max_{k \ge 1} \norm{b_k}_X = \max_k (1 + (c/r)^k) = 1 + 3/4$, the norm of the basis vector of the slice $b_k = e_k - (-c/r)^k e_0$,
that is, the conversion of the radius from b-basis coordinates into $X$. The
basis here is \textbf{normalized}: $e_k(t) = ((t-c)/r)^k$, so that $\norm{e_k}_X = 1$ for
all $k$, while $e_k(0) = (-c/r)^k$ --- hence the subtracted term. (With the
unnormalized $e_k = (t-c)^k$, where $\norm{e_k}_X = r^k$, the coefficient of $e_0$
would be $(-c)^k$, not $(-c/r)^k$; mixing the two normalizations in one triangle
inequality is not allowed.)

In the gate the factor is taken \textbf{squared}, and this is not an extra margin but
the \textbf{load-bearing part of the derivation} --- it is precisely what buys the ball
$1.75\rho'$ in the X-norm, claimed at the end of the paragraph. The check ``with a
single power, $Z1 + 1.75 \cdot B_A \cdot L \cdot \rho' \le 0.162$, all 46 boxes pass'' is true as a
fact, but proves something else. That bound is not deposited as a field, and the
one-off run that produced it is not kept --- but it does reassemble from the
package in a dozen lines: all four factors are fields of the certificates
(\texttt{Z1}, \texttt{B\_A}, \texttt{best.L}, \texttt{best.rho\_prime}), and their
maximum over the 46 boxes is $0.161733946474$ on \texttt{g43} ($0.161734679934$ if
every field is read as ``printed value plus one unit of its last digit''), i.e.
$\le 0.162$ under either reading.

The square is accounted for as in the paragraph ``Where exactly $L$ is certified''
below and in \S\ref{sec:6.6.2}: one factor of $1.75$ converts the radius from the
b-norm to the X-norm, the second converts the input vector in
$\norm{T}_{1\to1}\le1.75\,\norm{T}_{X\to X}$. On a b-ball of radius $s$ the Lipschitz
term is therefore $1.75\cdot B_A\cdot L\cdot(1.75 s) = 1.75^{2}B_A L s$, so with
$s=\rho'$ it is exactly the printed $\Zhat$; a b-ball of radius $1.75\rho'$ would
require $1.75^{3}$. \emph{Uniqueness on the X-ball does not go through the larger
b-ball at all, and needs no third factor.} Let $u,v$ be zeros of
$\widetilde{F}$ in $B_X(P,1.75\rho')\cap X_{0,s}$. The X-ball is convex, so the whole
segment $w(\theta)=v+\theta(u-v)$, $\theta\in[0,1]$, lies in it and
$\sup_\theta\norm{w(\theta)-P}_X\le1.75\rho'$; the domain on which $L$ is certified
is the X-ball of radius $3.0625\rho'+dP$ around $P_s$ (see below), which contains
this segment with room to spare. The mean value inequality along the segment gives
\[
  \norm{u-v}_1 \;\le\; \sup_\theta\norm{I-\widehat{A}M'(w(\theta))}_{1\to1}\cdot\norm{u-v}_1
  \;\le\; \bigl(Z1 + 1.75\cdot B_A\cdot L\cdot 1.75\rho'\bigr)\norm{u-v}_1
  \;=\; \Zhat\,\norm{u-v}_1 ,
\]
and $\Zhat<1$ forces $u=v$. So the printed $\Zhat$ --- with the square and nothing
more --- already gives uniqueness on the whole of $B_X(P,1.75\rho')$, which is the
ball the statement of Theorem~A claims.

(An earlier version of this paragraph derived the same conclusion by inflating the
gate from $\rho'$ to the b-ball of radius $1.75\rho'$ and called the square ``a double
conversion of the radius''. That accounting is inconsistent with the one used
everywhere else --- and with the code --- where the two factors are the radius and
the input vector; by that accounting the larger b-ball needs $1.75^{3}$. Nothing
numerical depended on it: the cube also passes, $\max_{46}(Z1+1.75^{3}B_A L\rho')
= 0.1858$ on \texttt{g43}, but the segment argument above is the correct route and
uses exactly the certified $\Zhat$.)

\textbf{Where exactly $L$ is certified.} Here it is important not to confuse two
centers. $L$ is computed around the center of the SUB-BOX $P_s$ (the function
\texttt{\_compute\_L}, which certifies the Lipschitz chain on the X-ball of radius $1.75x$
from the passed argument $x$), and it is certified as a constant of the map
$X \to X$: both norms are $\norm{\cdot}_X$. That is why the term carries
two factors of $1.75$ and not three --- one converts the radius, one the input
vector, and the output costs nothing, since $\norm{\cdot}_1 \le \norm{\cdot}_X$
on the slice.
The chain is two-point, not radial: every $P$-dependent factor enters it
inflated over the whole ball --- each majorant $\bar A$ appears in the products
as $\bar A + \Delta A$, with $\Delta A$ its own variation over that ball --- and
each $\Delta$ is monotone in the increment and vanishes with it. Since
$G\mapsto DT(G)$ is analytic and the ball is convex, the mean value inequality
along the segment joining $G_1$ to $G_2$ turns the chain into the two-point
bound stated as premise a4 of Table~\ref{tab:certified-premises}. So $D_{\mathrm{total}}/\rho_X$ measures the variation per unit
increment over the ball, not the displacement from $P_s$ alone --- which is
exactly what the next sentence needs, since the uniqueness ball is centered at
the
MOVING center $P(\ell)$. Within the sub-box $\norm{P(\ell) - P_s}_X \le dP$, so around $P_s$
an X-radius is required of

\[
3.0625 \cdot \rho' + dP \; = \; 1.75^2 \rho' \;\; \text{(the b-ball of radius } 1.75\rho'\text{, which contains it)} \;+\; dP \;\; \text{(center travel)},
\]

and the argument of \texttt{\_compute\_L} must be $(3.0625\rho' + dP)/1.75 = 1.75\rho' + dP/1.75$. This is exactly how \texttt{\_compute\_L} is called in \texttt{rig2fix3.py} (marked
\texttt{L-DOMAIN}); if the sub-box's package fails its own internal domain checks, the
box is rejected, not accepted on faith.

\begin{remark}[Fixed]
Previously the justification here relied on a gate:
$dP \ge 1.3125\rho' = (1.75^2-1.75)\rho'$, with $L$ computed at the argument $\rho' + dP/1.75$, that is, at the X-radius $1.75\rho' + dP$. This is $1.3125\rho'$ less
than the required $3.0625\rho' + dP$, and the gap does not depend on $dP$: the
gate compared the certified radius against $3.0625\rho'$ as if the uniqueness
ball were centered at $P_s$ itself, and did not add the center's travel to
the needed value. The justification did not hold. The argument has been
fixed, all 46 boxes recomputed; since $dP$ is in fact three to four orders of
magnitude larger than $\rho'$, the radius grew by a fraction of a percent.
Across all 46 boxes: the worst $L$ --- from $4.2512 \cdot 10^4$ to $4.2518 \cdot 10^4$
(+0.014\%), the worst $\Zhat$ --- from $0.170473$ to $0.1704744$, the worst $\rhostar_X$ ---
from $8.48288 \cdot 10^{-13}$ to $8.482891 \cdot 10^{-13}$, the worst \texttt{extra} --- from
$2.851703 \cdot 10^{-6}$ to $2.852127 \cdot 10^{-6}$, the worst full tube test --- from
$0.991843707425$ to $0.991843707457$. The gate is passed on all 46 existence
boxes and all 706 tube boxes; not a single published bound ceased to be valid
as a result.
\textbf{The ``from'' halves above are the values of the superseded run: they
are not in the package, not in \texttt{05-\allowbreak DATA-\allowbreak INDEX.md},
and not reproducible from it.} They are printed only so that the size of the
change can be seen; the ``to'' halves are the certified values, and no step of
any proof rests on the comparison.
\end{remark}

\begin{remark}[The same defect one section away]
A reference to this same removed gate remained in \S\ref{sec:6.7}
(step 1) --- there it was presented as an active check ``per box.'' It was
replaced with the active justification. The local fix left a dangling
reference in another section --- a class of defect, not just one instance.
Since that revision it is found automatically: \texttt{04-\allowbreak CODE/\allowbreak tests/\allowbreak \_\allowbreak refs\_\allowbreak check.py}
checks every field, flag, function, and file name mentioned in the package
documentation against what actually exists in the code and in the json.
\end{remark}

\begin{remark}[The same class in a third guise]
Not ``a reference to
something that does not exist,'' but ``a reference to something that exists,
but holds something else.'' The watchdog just described checked for the existence
of the field and was satisfied with that, so past it slipped: line 24 of
Appendix~\ref{sec:A} ($\beta_{81}$ referenced to \texttt{kant\_\allowbreak l3.json::\allowbreak beta\_\allowbreak Np1}, where a different
quantity resides --- a tail estimate, now a separate line 24a); the K-mode
(\S\ref{sec:5.1.2} --- the numerator held the norm $\norm{\Gbar'}_X$ instead of the value
$\Gbar'(1)$); the run cost given as three inconsistent numbers, none of which
matched the sum of \texttt{time\_s} over the certificates. The lesson:
what must be cross-checked is not the existence of a reference but its
value, and not against a list of rules, but exhaustively. This is what
\texttt{04-\allowbreak CODE/\allowbreak tests/\allowbreak \_\allowbreak appendix\_\allowbreak check.py} does: it goes through every line of
that appendix as it stands in the package's own copy of the article
(\texttt{01-\allowbreak ARTICLE.md}; see Reproducibility), checks each number
there against the artifact, and \textbf{fails
on the line it could not verify} unless that line has been entered into the
exceptions list with a reason. The silent pass-through that the defect lived
on has been closed as a mechanism.
\end{remark}

Conclusion: \textbf{for every $\ell$ in the box, $\widetilde{F}(\cdot,\ell)$ has a zero in the ball
$B_X(P(\ell), 1.75\rho')$, and it is unique there.}

Beyond speed, this gives something the march did not give in principle: an
explicit analytic dependence of the center on $\ell$, ready for differentiation ---
that is, the object needed for the next step (the spectrum along the branch,
\S\ref{sec:7}).


\subsection{How the suprema over a box are actually computed}
\label{sec:3.2}

Precision matters here, because this is exactly the point at which two errors
were made and fixed (\S\ref{sec:3.3}).

There are two ways to obtain the supremum of a quantity over an $\ell$-box.

\textbf{The $\sigma$-path.} Everything is expanded in a series in $\sigma = \ell - \ell_0$, the series is
truncated at degree D, and the discarded mass is bounded by a quantity $\tau$. The
supremum is then majorized by the sum $\sum_i |\cdot| \cdot \sigma_{\max}^i + \tau$. The path is tight,
but is correct only to the extent that the bookkeeping of $\tau$ is correct.

\textbf{The interval path.} The quantity is computed once, at an \textbf{interval} $\ell$ and an
\textbf{interval} center: $P$ is substituted as an interval of coefficients, $\ell$ as
the interval $[\ell_0-w/2, \ell_0+w/2]$, and arb arithmetic itself returns an inclusion
valid for all points of the box \cite{moore1966interval, tucker2011validated, rump2010verification}. No $\tau$, no truncation --- but the interval widths
inflate, and on a wide box the bound explodes: composition of power series
produces a Miller-recurrence wrapping $e^{O(w \cdot \NW)}$ ($\NW$ being the truncation order
of Miller's series, see \S\ref{sec:3.6}).

\textbf{In the final certificate the tail quantities are computed by the interval path
only} (the reason is in \S\ref{sec:3.3}), and the heads $Y0$ and $Z1$ by the $\sigma$-path only
(the reason is in \S\ref{sec:3.4}). The explosion is cured not by choice of path but by
\textbf{subdivision}: the supremum over the box equals the maximum of the suprema over
its pieces, so the box is cut into K parts, each with a narrow center and a small
wrapping, and the maximum is taken. A measurement on the heaviest box computed
for this check ($\ell_0=3.60$, $w=0.05$), for the quantity Z1c (a component of the
decomposition $Z1=Z1a+Z1b+Z1c$, formally introduced in \S\ref{sec:4.6}). It is a
measurement, not a certified quantity, and it is reproduced by
\texttt{04-\allowbreak CODE/\allowbreak tests/\allowbreak \_\allowbreak z1bc\_\allowbreak sub.py}:

\begin{center}
\begin{tabular}{|c|c|c|c|c|}
\hline
K & 1 & 2 & 4 & 8 \\
\hline
Z1c & 1.199 (gate fails) & $7.5 \cdot 10^{-5}$ & $7.8 \cdot 10^{-7}$ & $2.5 \cdot 10^{-7}$ \\
\hline
\end{tabular}
\end{center}

K is chosen adaptively, by doubling, until the quantity settles; on $[2,4]$ $K \le 4$
sufficed. The statement itself does not change --- only the way of obtaining the
same upper bound changes.


\subsection{Two errors, and why the $\sigma$-path was removed from the tails}
\label{sec:3.3}

The method came from outside, and a ``numerical'' cross-check would not have
proved anything here: agreement to ten digits only says that the two
implementations are similar. Therefore tools capable of returning a negative
answer were applied.
\textbf{One-off mutation battery, 4/4.} Deliberate defects were introduced into the
code, and it was checked that the gate catches them: shifting the center by
$2 \cdot 10^{-6}$ in coefficient 5 --- FAIL; an operator built at $\ell_0+10^{-3}$ but checked at
$\ell_0$ --- FAIL; $\Ahat$ replaced by $2 \cdot I$ --- FAIL; the baseline --- PASS. A certificate
that cannot fail certifies nothing. (This is a one-off experiment outside the
package --- there is no script in \texttt{04-CODE} that reproduces exactly these four
mutations; this result cannot be reproduced from the package as it stands. The
outcome \textbf{4/4} is therefore the author's statement about that one run, not
a certificate; it is the same status as the five-run battery of
\S\ref{sec:4.10}, and no step of any proof rests on it.)

\textbf{Exact recomputation of float tolerances in arb.} It was this step --- not
code reading --- that found the first error.

\begin{remark}[Error 1 (Z1 was underestimated by 4.69 \%)]
The majorant of the quantity
$\sup_{|\sigma| \le \sigma_{\max}} \norm{I - \Ahat(\sigma) M'(\sigma)}_1$ was taken as $\max_i (\max_k \mathrm{col}_i[k] \cdot \sigma_{\max}^i)$,
whereas the correct estimate is $\max_k \sum_i \mathrm{col}_i[k] \cdot \sigma_{\max}^i$: the column norm
of a sum of slices does not exceed the \textbf{sum} of the column norms, not the
maximum over them. At $\ell_0=3$, $w=0.2$ the exact recomputation gives
$0.41641364780096$ versus $0.39777452179569$ in the code.
\textbf{The second number, and the $4.69\,\%$ of the heading computed against
it, are outputs of the superseded majorant: they are not in the package and
are not reproducible from it.} They are printed to record the size of the
error, and no step of any proof rests on them.
\end{remark}

Independent verification confirmed the fix and produced a counterexample to
the original rule: for $Z_0 = Z_1 = [[1]]$, $\sigma_{\max} = 0.5$ the true supremum
equals 1.5, while the \texttt{max\_i} rule gives 1.0 --- an underestimate by a third.

\textbf{A second fix was rejected by the same check, and that turned out to matter
more.} The tail $\sum_{k>N_p} |F_k| r^k$ is estimated by both routes, and I took
the \texttt{min} of them, reasoning ``the minimum of two valid upper bounds is a
valid upper bound.'' The premise did not hold.

\begin{remark}[Error 2 (the $\sigma$-path is not a proven bound)]
Its scalar constants
\texttt{c1\_sup}, \texttt{alpha\_sup} are computed as the sum of the moduli of the
coefficients of the \textbf{truncated} $\sigma$-series, i.e.\ they majorize a
degree-D polynomial, not the function itself; the remainder beyond degree D
is not covered by anything, because $\tau$ from \texttt{sig\_lib} majorizes a different
truncation and misses by 13--18 orders of magnitude. The underestimate is
0.2--5.6 \% across boxes. Both figures are measurements of the withdrawn
$\sigma$-path, which is not in the package. And \texttt{min} chose exactly this path precisely on the
boxes where the interval path blew up.
\end{remark}

In the process, a defect was exposed that no one had seen: on an exception in the
interval path, the code set the tail to zero, i.e.\ it \textbf{silently cheapened
the certificate} instead of failing the box.

An attempt was made to fix the $\sigma$-path locally, and it \textbf{failed on the
merits}. The review went through all 13 sites of the form ``sum of moduli of
coefficients $\times \sigma_{\max}^i$'': six were already rigorous (where the $\sigma$-series is
finite because the center is a polynomial), three were made rigorous by
direct interval computation over a subdivided $\sigma$, and the remaining four are
three categories over eight objects (\texttt{poly\_norm[k]}, \texttt{b0star}, and six
bivariate norms), which remained non-rigorous because they rely on the same
$\tau$-accounting. All of them enter Z1b/Z1c (the formal definition of both
quantities is in \S\ref{sec:4.6}, formula $Z1b := B_A \cdot \hat{c}$, $Z1c := (1+B_A) \cdot \hat{b}$).
Conclusion: the $\sigma$-path is rigorous for the Y0 tail, but \textbf{not for Z1}, and a
local fix does not cure this.

The $\sigma$-path is removed from the Y0 tail not as a result of the analysis
above (there it is recognized as rigorous for Y0), but for an earlier
reason --- the withdrawal of the unproven \texttt{min(interval, sigma)} after Error 2
(FIX-3, \texttt{04-\allowbreak CODE/\allowbreak rig2fix3.py:596-611}): its constants \texttt{c1\_sup}/\texttt{alpha\_sup}
are underestimated by 0.2--5.6\% (again a measurement of the withdrawn path,
not a package quantity). From Z1b/Z1c it is removed for the reason
discussed above. In both cases the $\sigma$-path was the sole support of the tail
quantities, and wherever it is removed it is replaced by the interval
\texttt{\_\allowbreak z1bc\_\allowbreak subdivided} with subdivision. The cost --- the boxes became narrower.

The $\sigma$-model \textbf{cannot} be dropped entirely, and this is a matter of
principle, not of diligence: the heads of Y0 and Z1 cannot be computed
without it (\S\ref{sec:3.4}). So the question of the rigor of $\tau$ is not removed by
switching paths --- it has to be settled on the merits, which is done in \S\ref{sec:3.5}.

Both \textbf{main} errors were underestimates in sign (not overestimates: Z1 and
the $\sigma$-path gave values smaller than the true ones --- see the descriptions
above), i.e.\ anti-conservative; nevertheless no previously issued
certificate is refuted by them --- not because of conservativeness, but
because all 46 boxes were recomputed with the fixed engine (\S\ref{sec:3.6}). The side
defect, however --- zeroing the tail in the \texttt{except} branch --- was \textbf{not}
conservative: it made the estimate more optimistic and could give a false
PASS on a box where the interval path had failed. It did not in fact give
any false PASS (verified against the \texttt{tail\_\allowbreak interval} field across all boxes
of both runs: a zero tail never occurs --- so the branch never fired), but it
was removed precisely as anti-conservative: now a failure of the interval
path fails the box (\texttt{rig2fix3.py}, tag FIX-3). But both had passed through a
prior code review with the verdict ``the method is correct.'' Hence a
methodological conclusion worth stating separately: \textbf{reading code does not
catch majorant errors.} Only exact recomputation of the same quantity by an
independent route catches them --- together with the explicit question, ``does
this formula bound the function, or its model?''


\subsection{Why the $\sigma$-model cannot be replaced by interval computation}
\label{sec:3.4}

The temptation is obvious: if the interval path is rigorous by construction,
why not use it for the heads $Y0$, $Z1$ as well? Checked by direct
measurement and \textbf{refuted} (\texttt{04-\allowbreak CODE/\allowbreak tests/\allowbreak pure\_\allowbreak interval.py}, box \texttt{g01},
$\ell_0=2.07$, $w=0.05$; in this variant $A$ is taken as a single matrix --- the value
of $\Ahat$ at the box center):

\begin{center}
\begin{tabular}{|c|c|c|c|}
\hline
K pieces & piece width & head $Y0$, interval & $Z1a$, interval \\
\hline
1 & 0.05 & $2.584\ldots \cdot 10^{-1}$ & 5.252 \\
\hline
2 & 0.025 & $1.224\ldots \cdot 10^{-1}$ & 2.436 \\
\hline
4 & 0.0125 & $5.966 \cdot 10^{-2}$ & 1.192 \\
\hline
\end{tabular}
\end{center}

Both quantities fall \textbf{exactly linearly} with piece width. This is not a
true variation but an interval-arithmetic dependency: the remainder itself
$\widetilde{F}(P(\ell),\ell)$ is of order $10^{-13}$ for any $\ell$ in the box (independently visible
both from the 28-point cross-check against the march engine and from the
self-check ``$\sigma$-model versus engine'' built into an early version of $d_{\mathrm{chk}}$,
whose discrepancy came out at $1.3 \cdot 10^{-15}$ --- a float diagnostic of
that early version, not in the package and not reproducible from it; the
mechanism's current description is in \S\ref{sec:3.7}).
Interval arithmetic with an interval center loses exactly the cancellation
for the sake of which the center was constructed in the first place: to get
$Y0 \approx 10^{-7}$, the piece would have to be of width around $10^{-8}$, i.e.\ about $10^{6}$
pieces per box.

Conclusion: here the $\sigma$-model is not an optimization but the only known way
to preserve the cancellation. Consequently its $\tau$ must be rigorous.


\subsection{$\tau$: defect and fix}
\label{sec:3.5}

All four formulas for $\tau$ in \texttt{sig\_lib.py} \textbf{were not upper bounds}. The
error is the same in all four: only terms $h^n$ with $n > D$ were counted
toward the tail, whereas $h^n$ with $n \le D$ also contributes to $\sigma$-powers above
$D$ ($h$ is not a monomial). Plus a fifth site: when multiplying by a scalar
$\sigma$-series, the intrinsic $\tau$ of the multiplied object was lost.

A counterexample settling the question: $h(\sigma) = 0.1\sigma + 0.1\sigma^2$, $D = 1$,
$\mathit{smax} = 1$. Then $[1/(1+h)]_2 = h_1^2 - h_2 = -0.09$, i.e.\ the true tail is
$\ge 0.09$, while the old formula gives $h_{\mathrm{abs}}^{D+1}/(1-h_{\mathrm{abs}}) = 0.05$.

The falsifying test (\texttt{\_\allowbreak tau\_\allowbreak test.py}) compares $\tau$ against the true tail,
computed by expansion up to degree 80, on 15 pairs (function $\times$ case): \textbf{the
old version failed 15 of 15}, underestimating by up to $2.8 \cdot 10^5 \times$; the fixed
version passes all 15 and in four cases matches the true tail exactly.

The fix is the same for all four sites: instead of a Lagrange or geometric
remainder, one takes the \textbf{exact tail of the majorant series} --- the
closed-form value minus the computed head:

\begin{center}
{\small\setlength{\tabcolsep}{4pt}
\begin{tabular}{|l|l|>{\raggedright\arraybackslash}p{0.291\textwidth}|}
\hline
function & majorant & tail \\
\hline
$1/g$, $g = g_0(1+h)$ & $1/(1-H)$, $H = \sum |h_k| \sigma^k$ & $|1/g_0| \cdot (1/(1-H(\sigma_{\max})) - \sum_{i \le D} u_i \sigma_{\max}^i)$ \\
\hline
$\Log g$ & $-\log(1-H)$ & $-\log(1-H(\sigma_{\max})) - \sum_{i \le D} v_i \sigma_{\max}^i$ \\
\hline
$\exp M$ & $e^{M_0} e^R$, $R = \sum_{i \ge 1} |M_i| \sigma^i$ & $|e^{M_0}| \cdot (e^{R(\sigma_{\max})} - \sum_{i \le D} e_i \sigma_{\max}^i)$ \\
\hline
$X = X_0 e^{M-M_0}$ (bivariate) & $\norm{X_0} e^R$ & same, via its own recursion + carry-over of $\tau_M$ \\
\hline
\end{tabular}}
\end{center}

Legitimacy of the technique: $X$ is a weighted $\ell^1$ with weights $r^k$, a
Banach algebra under convolution, and truncating the t-series only decreases
the norm; hence any composition of series is majorized by a scalar
recursion over moduli.

The matter did not end with a single library. Later \textbf{four more sites of
the same kind} were found --- where \texttt{op\_sig.py} multiplies scalar $\sigma$-series by
hand, bypassing \texttt{Sig.mul}:

\begin{center}
{\small\setlength{\tabcolsep}{4pt}
\begin{tabular}{|>{\raggedright\arraybackslash}p{0.34\textwidth}|>{\raggedright\arraybackslash}p{0.533\textwidth}|}
\hline
site & what was being lost \\
\hline
\texttt{\_smul} in \texttt{c2}, \texttt{Da}, \texttt{lc2}, \texttt{coef0}, \texttt{coef} & head\,$\times$\,head terms of degrees $D+1\dots 2D$ \\
\hline
\texttt{tau\_c1} after \texttt{exp(h)} & the tail of the argument $h$ itself (\texttt{tau\_h} was computed and not used) \\
\hline
$\tau_L$ (the article's label, not a name in the code) & the $\sigma$-tail of $\Log(U)$ itself: the logarithm of a polynomial is not a polynomial \\
\hline
\texttt{tau\_M}, and $\tau_{Mlm}$ (again a label, not a name) & the term $L_D \cdot \sigma^{D+1}$, cut off by multiplication by $(\ell_0+\sigma)$ \\
\hline
\end{tabular}}
\end{center}

All four are closed by the same technique (either the exact mass of what
was discarded, or the tail of the majorant series). The overall effect on
the numbers (box \texttt{g01}): $\tau_F$ grew from $6.39 \cdot 10^{-16}$ to $4.62 \cdot 10^{-14}$, $Y0$ --- from
$3.3238 \cdot 10^{-13}$ to $4.6343 \cdot 10^{-13}$ (+39 \%; these are the values of
that one run, and $Y0$ varies by tens of percent between runs --- the final
figures over all 46 boxes are in Table~\ref{tab:3.6-ranges}), $Z1$ did not change in seven
significant digits. The gate requires $Y0 \le (1-\Zhat)\rho'$, which over the 46 boxes runs from $8.2953\cdot 10^{-8}$ (\texttt{g43}) to $9.1555\cdot 10^{-8}$ (\texttt{g00}) and equals $9.1540\cdot 10^{-8}$ on \texttt{g01} --- so the honest round figure for the threshold is $8.29\cdot 10^{-8}$ (rounded DOWN: the minimum itself is $8.295256\cdot 10^{-8}$), not $9.2\cdot 10^{-8}$, and it still leaves a margin
of five orders of magnitude. \textbf{The defect was real, but not fatal}; all 46
boxes were recomputed with the fixed engine (\S\ref{sec:3.6}).

A lesson worth stating separately: the error lives not in a function but in
a technique --- ``truncated the series and forgot about what was truncated.''
It has to be sought at every site where something is truncated, not in a
single file. Nine such sites were found at that time.

A tenth was found later, by an outside review of the finished package, and it
is recorded here because of what it says about the search rather than about
the number. Where the engine passes to the logarithm --- $\alpha=-1/G_{1}$,
$\Log\alpha$ --- it called the majorant of $\Log$ on the \emph{truncated}
series for $1/G_{1}$, whose own tail then entered nothing downstream: the
routine bounds the tail of $\Log$ of the finite polynomial it is handed, and
the tail of its argument is a separate term. This is precisely the pattern
named one row above for $\exp$ (``the tail of the argument $h$ itself''),
missed at the neighbouring site. The falsifying test did not catch it because
it checked $1/g$ and $\Log g$ separately, each on a finite polynomial, and
never their composition; a composition case has since been added, and on it
the naive route fails to be an upper bound in three of the four test cases,
understating the true tail by up to a factor of $21$ (exactly $20.4986$ on the adversarial case $h = 0.1\sigma + 0.1\sigma^{2}$ of that same test).

The fix is not a Lipschitz term but an identity: $\Log(-1/G_{1})=-\Log(-G_{1})$,
and $-G_{1}$ is an exact finite polynomial with no input tail by construction.
The head is unchanged --- the coefficients of $\Log\alpha$ up to degree $D$
depend only on the coefficients of $1/G_{1}$ up to degree $D$, and those are
exact --- so no printed number moves. The $\tau$ the identity yields is
\emph{smaller} than the one on which the deposited certificates were built, by
a factor of at least $1.5$ on every one of the 46 boxes (worst case \texttt{gtop}, where the factor is $1.53$; the median over the 46 is $6.7436$, and the factor exceeds $3.5$ on 43 of them); this was checked box by box, rebuilding each from its seed, and
it is what leaves the deposited certificates standing without recomputing them.
The check is documented and reproducible; it is not one of the 26 steps of the
acceptance driver only because it takes about three quarters of an hour.



\subsection{Covering [2,4]: 46 boxes}\label{sec:3.6}

\begin{table}[h]
\centering
\caption{Grid of the covering boxes.}\label{tab:3.6-grid}
{\footnotesize\setlength{\tabcolsep}{3pt}
\begin{tabular}{llllll}
\toprule
boxes & $\ell_{0}$ & $w$ & $m$ & $D$ & $\rho'$ \\
\midrule
\texttt{g00} \ldots \texttt{g43} & 2.025 \ldots 3.960, step 0.045 & 0.05 & 6 & 8 & $10^{-7}$ \\
\texttt{ganchor} & 3.000 (cross-check with the pointwise certificate at $\ell=3$) & 0.05 & 6 & 8 & $10^{-7}$ \\
\texttt{gtop} & 3.9925 & 0.03 & 6 & 8 & $10^{-7}$ \\
\bottomrule
\end{tabular}}
\end{table}

Ranges over all 46 boxes (cost: 42.2 core-minutes):

\begin{table}[h]
\centering
\caption{Ranges of quantities over the boxes.}\label{tab:3.6-ranges}
{\small\setlength{\tabcolsep}{4pt}
\begin{tabular}{>{\raggedright\arraybackslash}p{0.608\textwidth}ll}
\toprule
quantity & from & to \\
\midrule
\texttt{Y0} & $6.24\cdot 10^{-14}$ & $4.31\cdot 10^{-13}$ \\
\texttt{Z1} & 0.0828 & 0.1501 \\
$\widehat{Z}$ & 0.0844 & 0.1705 \\
\texttt{d\_chk} (derived bound on the discrepancy between the $\sigma$-model and the engine; derivation --- \texttt{rig2fix3.py}, tag DCHK-CONV, discussion --- \S\ref{sec:3.7}) & $1.07\cdot 10^{-16}$ & $7.98\cdot 10^{-15}$ \\
\texttt{Y0} tail (interval path, with subdivision) & $2.9534\cdot 10^{-33}$ & $1.08\cdot 10^{-13}$ \\
\texttt{Z1b} (honest column tail; interval recomputation --- \S\ref{sec:3.3}) & $8.85\cdot 10^{-13}$ & $1.32\cdot 10^{-3}$ \\
\texttt{Z1c} & $5.47\cdot 10^{-7}$ & $8.96\cdot 10^{-4}$ \\
$\tau_{F}$ (after the fixes, \S\ref{sec:3.5}) & $3.98\cdot 10^{-19}$ & $6.93\cdot 10^{-14}$ \\
$L$ & $1.85\cdot 10^{3}$ & $4.26\cdot 10^{4}$ \\
$\rho^{\star}_{b} = (Y0 + B_{A}\cdot d_{\mathrm{chk}})/(1-\widehat{Z})$ --- radius in the $b$-basis & $8.1\cdot 10^{-14}$ & $4.85\cdot 10^{-13}$ \\
$\rho^{\star}_{X} = 1.75\cdot \rho^{\star}_{b}$ --- where $\Gs(\ell)$ actually lies in the norm $X$ & $1.41\cdot 10^{-13}$ & $8.49\cdot 10^{-13}$ \\
guard criterion $B_{A}\cdot d_{\mathrm{chk}} / (Y0/2)$ & 0.008 & 0.352 \\
\texttt{U\_min} (domain lemma; on \texttt{g37}--\texttt{g43} anti-conservative by $\le 4.45\cdot10^{-4}$, \S\ref{sec:6.2.a}) & 0.626 & 0.663 \\
$\min \operatorname{Re} U$ on the circle of radius 1.2 around $c$ (non-vanishing of $U$, \S\ref{sec:10.2}/\S\ref{sec:10.7}) & 0.526 & 0.572 \\
geometry bases $q$ (defined in \S\ref{sec:6.5}) & 0.736 & 0.9911 \\
geometry bases $q_s$ (defined in \S\ref{sec:10.2}) & 0.791 & 0.8112 \\
\bottomrule
\end{tabular}}
\end{table}
\FloatBarrier

$L$ and $\widehat{Z}$ in this table are the \textbf{upper ends} of the intervals, not their midpoints. The distinction matters: $L = D_{\mathrm{total}}/\rho_{X}$ is assembled from interval quantities and remains an arb-ball whose upper end exceeds the midpoint by up to $79\,\%$ (worst box \texttt{gtop}: the ratio of the upper end to the midpoint is $1.7824$), and it is its worst end that goes into the gate. While the midpoint was being written into the certificate, the table underestimated $L$ by nearly a factor of two, and with it $\widehat{Z}$ and the displacement term $\mathrm{extra}$ in \S\ref{sec:4.5}, which is computed via $L$. The gate itself was checked against the worst end the whole time and therefore remained rigorous --- but the published number was not the one used in the proof. The fields \texttt{gate\_num} and \texttt{gate\_den} in the json now store the numerator and denominator of the gate separately, so that the identity $\rho^{\star}_{b}\cdot(1-\widehat{Z}) = Y0 + B_{A}\cdot d_{\mathrm{chk}}$ can be checked externally, exactly (\texttt{tests/\allowbreak \_\allowbreak numbers\_\allowbreak check.py}).

One row of the table deserves a word here rather than later. The discrepancy
$d_{\mathrm{chk}}$ between the $\sigma$-model and the engine does not fall below
a floor of order $10^{-16}$, independent of $D$ and of the box width; the guard
criterion $B_{A}\cdot d_{\mathrm{chk}} \le \tfrac{1}{2}Y0$ is what keeps that
floor out of the gate --- it is filtered, not eliminated. Its source, and how
far it can be pushed down, are in \S\ref{sec:3.10}.

Like \texttt{Y0} (\texttt{02-\allowbreak REPRODUCE.md}, ``Note on reproducibility''), $\tau_{F}$, $\rho^{\star}_{b}$, and the $\mathrm{extra}$ (\S\ref{sec:4.8}) that inherits from it are built through the non-deterministic float seed \texttt{solve\_\allowbreak fixed\_\allowbreak point} --- on an independent regeneration (a different flint/gmp/mpfr/LAPACK build, whose output is not part of the deposited package), for the degenerate box \texttt{gtop} values were observed up to $\tau_{F} = 5.60\cdot 10^{-19}$, $\rho^{\star}_{b} = 3.53\cdot 10^{-13}$, $\mathrm{extra} = 2.96\cdot 10^{-6}$ (these three are observations of a non-deterministic float seed on that other build: they are not interval certificates, they are not persisted in the package, and no step of the proof rests on them) --- 4--5 orders of magnitude below what actually bounds the conclusion (gate $Y0\le 8.29\cdot 10^{-8}$, the worst of the 46 boxes, \texttt{g43}; the range is $8.2953\cdot 10^{-8}\ldots9.1555\cdot 10^{-8}$; inequality \S\ref{sec:4.8}, where recomputing with a larger $\mathrm{extra}$ changes the margin in the 6th digit). What must be compared is whether the gate passes, not bitwise agreement of these quantities.

This is exactly the same class of error as the others found (\S\ref{sec:3.3}, \S\ref{sec:4.10}): not an overestimate but an \textbf{underestimate}, which gates of the form ``$< 1$'' cannot by construction see.

\textbf{Fractional power inside the chain for $L$.} The chain passes through the quantity $\norm{U_{G}^{\ell-1} - U_{P}^{\ell-1}}_{X}$ and through $\norm{W_{G} - W_{P}}_{X} = \norm{U_{G}^{\ell} - U_{P}^{\ell}}_{X}$, and it is easy to write an incorrect inequality here. For an \textbf{integer} power everything is simple: $U^{n} - V^{n} = (U-V)\cdot \sum_{k} U^{k} V^{n-1-k}$, whence by submultiplicativity $\norm{U^{n} - V^{n}} \le n\cdot \max(\norm{U},\norm{V})^{n-1}\cdot \norm{U-V}$. For a \textbf{non-integer} power this inequality is false, and it cannot be obtained by a ``mean value theorem'': $X$ is a Banach algebra, but $\norm{U^{a}} \le \norm{U}^{a}$ does not hold for non-integer $a$. The generalized binomial series has coefficients that alternate in sign, and after taking moduli the sum can exceed the power of the norm: for $U = 1 + z/2$ (in the normalized variable $z = (t-c)/r$) and $a = 3/2$
\[ \norm{U^{a}}_{X} = 1.85355\ldots \; > \; \norm{U}_{X}^{a} = 1.83712\ldots \]

The correct route is the mean value theorem in \textbf{integral} form, which is valid in a Banach algebra:
\[ U_{G}^{a} - U_{P}^{a} = a \int_{0}^{1} (U_{P} + \theta H)^{a-1} H \,d\theta, \qquad H = U_{G} - U_{P}, \]
it is applicable because the spectrum of every $U_{P} + \theta H$ lies in the disk $|w - 1| \le \kappa_{U} < 1$, i.e. in the right half-plane, where $z^{a}$ is single-valued and holomorphic ($\kappa_{U} := \norm{U_{P} - 1}_{X} + \Delta U$; over all 46 boxes $\kappa_{U} \le 0.39$ --- the same quantity is certified more sharply in \S\ref{sec:6.5} item 5 as $\kappa_{U} \le 0.363184$; here a cruder but sufficient bound is used).
Hence
\begin{align*}
\norm{U_{G}^{a} - U_{P}^{a}}_{X} &\le a\cdot \bigl( \norm{U_{P}^{a-1}}_{X} + S(a-1, \kappa_{U})\cdot \Delta U \bigr)\cdot \Delta U, \\
S(b, \kappa) &:= \sum_{n\ge 1} n\cdot |C(b,n)|\cdot \kappa^{n-1},
\end{align*}
where $\norm{U_{P}^{a-1}}_{X}$ is the genuine norm of the series (the engine already builds it via Miller), and $S$ is the Lipschitz constant of the map $v \mapsto (1+v)^{b}$ on the ball $\norm{v} \le \kappa$, computed as the sum of a series of the \textbf{moduli} of the binomial coefficients with an explicit tail: for $n > b$ it holds that $|C(b,n+1)|/|C(b,n)| = (n-b)/(n+1) < 1$, hence $\sum_{n\ge N} n|C(b,n)|\kappa^{n-1} \le |C(b,N)|\cdot (N\kappa^{N-1} - (N-1)\kappa^{N})/(1-\kappa)^{2}$. For \textbf{integer} $a$ the two formulas coincide identically --- $S(m,\kappa) = m(1+\kappa)^{m-1}$ --- so nothing changes for $\ell = 2, 3, 4$, while for non-integer $\ell$ the chain becomes derived rather than postulated. Numerically $L$ grows under this by no more than 14\,\% (worst case --- the upper end of the interval), $\widehat{Z}$ by 4.2\,\%, $\rho^{\star}$ by 0.7\,\%; this does not move the gate.
\textbf{These three percentages are float64 measurements against the values
before the fix; those earlier values are not in the package and the
percentages are not reproducible from it.} They say how much room the change
took, and no step of any proof rests on them.

Separately: the norms $\norm{W_{P} - c}_{X}$ and $\norm{U_{P}^{\ell-1}}_{X}$ are taken from Miller series truncated at $\NW = 192$, so Cauchy tails over the circle $|t - c| = 1.2$ are added to them (quantities of order $10^{-33}$, but without them these would not be upper bounds). The same truncated $W_{P}$ enters two more norms of the same function, $\norm{\Gbar\circ W_{P}}_{X}$ and $\norm{\Gbar'\circ W_{P}}_{X}$, where the tail is NOT added (\texttt{tube\_\allowbreak lib.py}, tag \texttt{TRUNC-\allowbreak TAIL}); the omission is at most
$\mathrm{Lip}(\Gbar;\nu_W)\cdot\mathtt{tail\_hW}\le 1.95\times 10^{-33}$ and
$\mathrm{Lip}(\Gbar';\nu_W)\cdot\mathtt{tail\_hW}\le 2.03\times 10^{-33}$
--- both maxima over all $46$ boxes, attained on \texttt{g43} ($\ell_0=3.96$);
at the lower end of the range, on \texttt{g01}, they are three and ten times
smaller. That is at least twenty-two orders below the radius of the interval
already carried by those two quantities, so the deposited upper ends remain
upper bounds --- but by the accumulated interval, not by construction.

The \texttt{Y0} tail also had to be computed with subdivision: with a single piece covering the whole box, the Miller wrapper inflated it at the upper end of the interval to $2.8\cdot 10^{-8}$ against a true value of $\sim 10^{-13}$
(a float64 diagnostic of the pre-subdivision run, not a certified quantity and
not reproducible from the package). This did not fail the gate (a margin of five orders of magnitude), but it inflated $\rho^{\star}$, and through it the displacement term of the fixed point in the hyperbolicity tube (\S\ref{sec:4.5}). After subdivision $\rho^{\star}_{b}$ does not exceed $4.85\cdot 10^{-13}$ over the whole range.

All 46 boxes pass the gate at $\rho' = 10^{-7}$; not a single FAIL, not a single exception. The cost of the final run is 42.2 core-minutes.

What bounds this estimate from above is not the gate but the geometry base $q$: at $q \ge 1$ both the $\widehat{\beta}$-majorant (the certified quantity $\widehat{\beta}_{N_{p}+1}$, a tail estimate in the Kantorovich chain --- Appendix~\ref{sec:A}) and the Cauchy tails stop converging. It comes closest to unity on box \texttt{g36} ($\ell \in [3.62, 3.67]$), where $q = 0.9911$; it reaches the threshold $q = 1$ nowhere. (This holds for the specific implementation of the composition-tail estimate via $q$; \S\ref{sec:10.4} gives a weaker sufficient condition, $\texttt{margin\_c} > 0$, for a related estimate in Lemma 2, where $q \ge 1$ no longer poses a problem --- see also \S\ref{sec:6.5} item 6.) The quantity is not monotone in $\ell$ and does not tend to unity as $\ell \to 4$: on \texttt{g43} ($\ell_{0} = 3.96$), $q = 0.9552$; on \texttt{gtop}, $0.9858$. The dips in this sequence are explained not by the structure of the problem but by subdivision: where $q$ approaches unity, $K = 2$ is taken, and the narrowing of the $\ell$-intervals drives the quantity down. So continuation beyond $\ell = 4$ runs into the choice of layer, not into the volume of computation.

\textbf{The acceptance run --- eight conditions at once, not just the gate:}

\begin{table}[h]
\centering
\caption{Acceptance-run checks.}\label{tab:3.6-acceptance}
\begin{tabular}{p{0.55\textwidth}p{0.35\textwidth}}
\toprule
check & result \\
\midrule
gate $Y0 + B_{A}\cdot d_{\mathrm{chk}} + \widehat{Z}\rho' \le \rho'$ & \checkmark\ 46/46 \\
guard criterion $B_{A}\cdot d_{\mathrm{chk}} \le \frac{1}{2}\cdot Y0$ (now also \textbf{inside} the gate) & \checkmark\ worst share $B_{A}\cdot d_{\mathrm{chk}}/Y0$ = \textbf{0.176} against threshold \texttt{GUARD\_FRAC} = 0.5, i.e. $B_{A}\cdot d_{\mathrm{chk}} \le 0.5\cdot Y0$ --- the same condition as the guard criterion on the left (box \texttt{g31}) \\
domain lemma on every box (interval center, interval $\ell$, subdivision) & \checkmark\ 46/46, $U_{\min} \ge 0.626$ (on \texttt{g37}--\texttt{g43} the printed field is anti-conservative by $\le 4.45\cdot10^{-4}$, \S\ref{sec:6.2.a}), both margins $> 0$ \\
non-vanishing of $U$ in the disk: $\min \operatorname{Re} U > 0$ on both circles & \checkmark\ 46/46, $\min \operatorname{Re} U \ge 0.526$ \\
geometry bases $q, q_{s} < 1$ on every box & \checkmark\ 46/46, $q \le 0.9911$, $q_{s} \le 0.8112$ \\
seams: the ZERO of every box lies in the uniqueness ball of its neighbor --- $\mathrm{dist}(P_{i}, P_{j})(\ell^{*}) + 1.75\cdot \rho^{\star}_{i} < 1.75\cdot \rho'_{j}$ in both directions, at a rational point $\ell^{*}$ of overlap;
mere intersection of the balls is NOT sufficient for this & \checkmark\ 45/45, minimum margin $\ge 1.749\cdot 10^{-7}$, that is, more than 0.999$\cdot$(uniqueness radius): it is $1.75\cdot 10^{-7}$ less $\mathrm{dist} + 1.75\cdot \rho^{\star} \le 1.07\cdot 10^{-12}$ (worst seam \texttt{g41}/\texttt{g42}; the distance alone stays below $2.14\cdot 10^{-13}$), truncated downward \\
gaps & \checkmark\ none \\
\textbf{union} & \textbf{[2.0000, 4.0075] $\supseteq$ [2,4]} (the left end is tight: the certified $\ell$-interval of \texttt{g00} begins at $1.9999999999999999086$, i.e.\ $9.14\cdot 10^{-17}$ below $2$, whereas the printed \texttt{ell\_lo} is exactly $2.0$; \texttt{glue.py} now compares the certified end, not the printed one. The right end has slack $0.0075$) \\
\bottomrule
\end{tabular}
\end{table}
\FloatBarrier

A seam is not cosmetic, but it is also not what it may seem. Each certificate asserts existence and \textbf{uniqueness of its own} zero in \textbf{its own} ball, and two boxes with different zeros in non-intersecting balls do not contradict each other: both certificates remain true. The seam answers a different question --- whether all 46 boxes certify ONE AND THE SAME zero, i.e. a single branch $\Gs(\ell)$ over the whole \texttt{[2,4]}. Its failure is not ``proof of an error'' but a failure of the covering to glue into one branch. Independently of this argument, the seams are also invisible in the scan: the exploratory scan crosses them dozens of times and no jump in $\delta$ appears at any of them (\S\ref{sec:7.2}).

The check is done at a single rational point $\ell^{*}$ of overlap, and that suffices, but the argument is worth writing out rather than left implicit. Let $G_{i}(\ell)$ and $G_{j}(\ell)$ be the zeros selected by the certificates of neighboring boxes on their common (connected) interval $J$. Both functions are continuous in $\ell$: the center $P(\ell)$ is polynomial, and the zero is obtained by a uniformly contracting Newton operator whose parameters are continuous in $\ell$. The set $E = \{\ell \in J : G_{i}(\ell) = G_{j}(\ell)\}$ is closed in $J$ by continuity. It is also open: if $G_{i}(\ell_{0}) = G_{j}(\ell_{0})$, then by continuity, for $\ell$ close by, the point $G_{i}(\ell)$ remains in the uniqueness ball of the $j$-th certificate, and there is exactly one zero there --- hence $G_{i}(\ell) = G_{j}(\ell)$. The check at $\ell^{*}$ gives $E \ne \varnothing$; connectedness of $J$ gives $E = J$. This is exactly why one checks the nesting of the balls with a margin, rather than mere intersection: nesting is precisely the condition from which openness follows immediately.

\textbf{Cost: 42.2 core-minutes} (46 boxes at 45.7--76.2 s each) versus $\approx 422$ core-hours for the march (64 core-hours for 15.15\,\% of the interval) --- roughly \textbf{$600\times$}, for a stronger statement. The cost is not a bare declaration: it is the sum of the \texttt{time\_s} fields across the 46 certificates \texttt{03-\allowbreak CERTIFICATES/\allowbreak existence/\allowbreak box\_\allowbreak g*.json}, each the full wall-clock time to certify its box (\texttt{04-\allowbreak CODE/\allowbreak rig2fix3.py}, lines $291$, $794$ and $864$),
and the same number is checked against that sum by \texttt{04-\allowbreak CODE/\allowbreak tests/\allowbreak \_\allowbreak appendix\_\allowbreak check.py} in the package's copy of the article.


\subsection{Margin on $\tau$}\label{sec:3.7}

The dependence on $\tau$ remains in exactly one place: \texttt{Y0\_head} and \texttt{Z1\_head} are built from the $\sigma$-model of the operator, whose discarded mass is described by \texttt{tau\_F} and \texttt{taus\_cols} (now rigorous, \S\ref{sec:3.5}). It is useful to know not only that the estimate is valid, but also how load-bearing it is: \textbf{by what factor $\tau$ would have to grow for the gate to stop passing}. The measurement is on the current, already-fixed values of $\tau$ (\texttt{tests/\allowbreak \_\allowbreak tau\_\allowbreak share.py}); in parentheses, for comparison, is what the pre-fix accounting gave --- those parenthesised values are outputs of the superseded accounting, are not in the package and are not reproducible from it:

\begin{table}[h]
\centering
\caption{Margin on $\tau$ for three representative boxes.}\label{tab:3.7-tau}
{\small\setlength{\tabcolsep}{4pt}
\begin{tabular}{lll>{\raggedright\arraybackslash}p{0.609\textwidth}}
\toprule
box & $\ell_{0}$ & $Z1_{\tau} / Z1$ & gate breaks at growth of $\tau$ by (margin divided by $B_{A}\cdot \tau_{F}$ --- the factor $B_{A}$ from premise a3 of Table~\ref{tab:certified-premises} is mandatory) \\
\midrule
\texttt{g00} & 2.025 & $2.8\cdot 10^{-11}$ & $4.6\cdot 10^{5}\times$ (was $3.4\cdot 10^{7}$) \\
\texttt{ganchor} & 3.000 & $9.6\cdot 10^{-14}$ & $1.3\cdot 10^{8}\times$ (was $1.4\cdot 10^{10}$) \\
\texttt{g43} & 3.960 & $3.7\cdot 10^{-15}$ & $3.5\cdot 10^{9}\times$ (was $9.8\cdot 10^{10}$) \\
\bottomrule
\end{tabular}}
\end{table}

That is, the margin amounts to \textbf{five to nine orders of magnitude} (exact measured values --- \texttt{03-\allowbreak CERTIFICATES/\allowbreak pointwise/\allowbreak tau\_\allowbreak share.json}; the shares $Z1_\tau / Z1$ are rounded UP, the margin factors are rounded DOWN). This is a single measurement and depends on the seed: $Y0$ varies by tens of percent between runs (\texttt{02-\allowbreak REPRODUCE.md}), and with it the factors also drift, so the table gives directional rather than exact numbers. Two separate causes make the drift larger than run-to-run noise, and both push the same way. The seed accounts for the smaller part. The larger part is the correction to the tail of $\Log\alpha$ described in \S\ref{sec:3.5}: it makes $\tau$ smaller, so re-running \texttt{tests/\allowbreak \_\allowbreak tau\_\allowbreak share.py} against the current engine returns shares smaller by a factor of $1.7$ to $4.2$ and margin factors larger by a factor of $2.0$ to $8.1$ than the deposited \texttt{tau\_\allowbreak share.json}, which was produced before that correction and describes the deposited certificates. The recipe in \texttt{05-\allowbreak DATA-\allowbreak INDEX.md} therefore does not reproduce that file byte for byte, and is not meant to. The numbers printed above are conservative against both: every published share is an upper bound for the deposited and for the current value alike, and every published factor is a lower bound for both, so nothing in the paragraph below depends on which of the two a reader recomputes. The fix to $\tau$ cost more than it looked at a glance: at the left end the margin dropped by nearly two orders of magnitude (more than $70\times$), not ``by half an order of magnitude'' as the draft had said. Even so, for the gate to fail, the entire discarded mass of the $\sigma$-model would have to be underestimated by hundreds of thousands of times. Historical note: in the EARLY version, $d_{\mathrm{chk}}$ was a built-in self-check with an inconsistent power of the weight factor, and therefore overestimated the discrepancy between the model and the engine by $5.6\times$ (the measurement of the code of that time: $1.62\cdot 10^{-15}$ versus $2.89\cdot 10^{-16}$ at $\ell_{0}=2.025$ --- outputs of the superseded implementation, not persisted in the package and not reproducible from it, quoted here only to record the size of the old bias; the ratio itself is still visible in the \texttt{\_\allowbreak dchk\_\allowbreak check.py} diagnostic) --- the error was conservative, and the gate did not suffer. Since 18.08, $d_{\mathrm{chk}}$ in the gate is DERIVED differently (tag DCHK-CONV, \S\ref{sec:3.6}) and does not inherit this defect; the current value at $\ell_{0}=2.025$ is $8.32\cdot 10^{-16}$ (\texttt{box\_\allowbreak g00.json::\allowbreak dchk}), not $1.62\cdot 10^{-15}$.


\subsection{Independent cross-check against the march at scale}\label{sec:3.8}

By the time it was stopped, the march had accumulated \textbf{81,479 certified records} (all with $\texttt{gate\_pass}=\mathrm{True}$, 141 MB, 59 files, none unreadable), covering a total of 0.303 in $\ell$ across 55 disconnected components within $[2.0700,\, 3.5808]$ (the union reaches $3.580716250000099$; both ends are rounded OUTWARD, as \appref{sec:A} requires). As a proof of existence this is no longer needed; as an \textbf{independent oracle} --- it is exactly what is required.

The key test is set up so that, at the decisive step, it does not depend on a single line of the parametric code. For 28 values of $\ell$ --- one from each of 28 components of the march's covering, chosen uniformly by component number, from the first to the last --- the parametric center $P(\ell)$ is taken, and its residual is computed by the \textbf{march engine} --- \texttt{compute\_\allowbreak F\_\allowbreak poly} in arb at 256 bits --- in the related norm $\norm{\widetilde{F}}_{1} = \sum_{k\ge 1}|F_{k}|\cdot r^{k}$ (the head of this sum enters the march's own $Y0$ via the $\widetilde{A}$-transform and explicit tails, not literally in this form), namely $Y0$. Alongside it, the same quantity is computed for a deliberately shifted vector.

\begin{table}[h]
\centering
\caption{Results of the march cross-check.}\label{tab:3.8-march}
{\small\setlength{\tabcolsep}{4pt}
\begin{tabular}{l>{\raggedright\arraybackslash}p{0.387\textwidth}}
\toprule
quantity & result over 28 points \\
\midrule
$\norm{\widetilde{F}(P(\ell),\ell)}_{1}$ & $1.81\cdot 10^{-15} \ldots 3.48\cdot 10^{-14}$ \\
control (shift of $10^{-8}$ in coeff.\ 5) & \textbf{$2.6\cdot 10^{6} \ldots 6.1\cdot 10^{7}$ times} larger \\
residual / the march's own $Y0$ & $\le 1.57\cdot 10^{-4}$ (\texttt{05-\allowbreak DATA/\allowbreak COVERAGE.md}: $1.564\cdot 10^{-4}$) \\
partial center distance (20 leading coeffs.) & $6.4\cdot 10^{-9} \ldots 5.1\cdot 10^{-7}$ at a tolerance of $1.75\cdot \rho'_{\mathrm{march}} \approx 5\cdot 10^{-5}$ \\
\bottomrule
\end{tabular}}
\end{table}

Here is how to read this. The parametric center is a genuine zero of the very operator that certifies the march, and it is four orders of magnitude more accurate than the march's own polished seed. The sensitivity of the test is measured: a shift 3,000 times smaller than the radius of the march's ball raises the residual by millions to tens of millions of times ($2.6\cdot 10^{6}\ldots 6.1\cdot 10^{7}$, the range of Table~\ref{tab:3.8-march}), so a coincidence is ruled out. Finally, the partial center distance (over the 20 leading coefficients --- a lower bound on the full $X$-distance; the tail of coefficients 21+ is not separately estimated) fits within the tolerance of the march's uniqueness ball ($6.4\cdot 10^{-9} \ldots 5.1\cdot 10^{-7}$ at a tolerance of $1.75\cdot \rho'_{\mathrm{march}} \approx 5\cdot 10^{-5}$, where $\rho'_{\mathrm{march}}$ is the radius of the MARCH's uniqueness ball, not the parametric $\rho' = 10^{-7}$ of Theorem A; Table~\ref{tab:3.8-march}) --- strong, though not rigorous (owing to the unestimated tail), evidence that the march asserts uniqueness within this same ball, i.e. the two methods probably refer to \textbf{one and the same zero}. The reverse inclusion does not hold and is not needed --- the uniqueness ball of the parametric certificate ($1.75\cdot 10^{-7}$, Theorem A) is narrower than the largest partial center distance.


\subsection{External cross-check: the classical Feigenbaum constants}\label{sec:3.9}

All checks in the preceding subsections are internal. They share a common point of failure: the existence certificates, the tube of \S\ref{sec:4}, and both pointwise anchors of \appref{sec:5} all rely on one and the same derivation of the derivative $DF$. A systematic error in that derivation would be inherited by every check at once --- the certificates would faithfully confirm the wrong quantity, and no internal inequality would notice. Reading the code does not catch this either: the reviewer reads the same derivation.

Against this there is one cheap and convincing remedy --- compare our numbers with numbers computed by \textbf{other people using other methods}. Both ends of the interval qualify for this, and both lie inside the certificate covering: $\ell = 2$ is the classical quadratic case, $\ell = 4$ is the degree-four critical point.

Our values are obtained by a path that depends neither on the scan of \S\ref{sec:7} nor on the pointwise anchors of \appref{sec:5}: $\alpha = -1/G(1)$ --- directly from the coefficients of the parametric center (they are stored in the monomial basis, so $G(1)$ is their sum); $\delta$ --- by the power method in arb on the matrix $\DT$ ($|\lambda_{2}|/\delta \approx 0.03$, convergence within three to four dozen iterations). The comparison is a cross-check, not a premise: no statement of Theorems A--D uses these digits, and the check itself is in the acceptance run (\texttt{tests/\allowbreak \_\allowbreak classic\_\allowbreak constants.py}).

\begin{table}[h]
\centering
\caption{Comparison of $\delta$ and $\alpha$ with the published classical values.
The published values at $\ell = 4$ are those of Table~II of BOT ($325$ digits of $\delta$, $331$ of $\alpha$), quoted here to twenty; the count is the
number of digits that agree once our value is rounded to the precision at which
the published one is printed.}\label{tab:3.9-classic}
\begin{tabular}{lllll}
\toprule
$\ell$ & & our value & published & matching digits \\
\midrule
2 & $\delta$ & 4.669201609102991531 & 4.669201609102990672 & \textbf{15} \\
2 & $\alpha$ & 2.502907875095880978 & 2.502907875095892822 & \textbf{14} \\
4 & $\delta$ & 7.284686217073389754 & 7.28468621707334336431 & \textbf{14} \\
4 & $\alpha$ & 1.690302971405251591 & 1.69030297140524485334 & \textbf{14} \\
\bottomrule
\end{tabular}
\end{table}

Sources of the external numbers: for $\ell = 2$ --- OEIS A006890 and A006891 \cite{oeisA006890, oeisA006891} (hundreds of digits, computed independently by Briggs (1991) \cite{briggs1991precise} and Broadhurst (1999) \cite{broadhurst1999feigenbaum}); for $\ell = 4$ --- Burbanks--Osbaldestin--Thurlby (2021) \cite{burbanks2021rigorous}, below BOT: Table~II of that paper (p.~8) prints $331$ proven digits of $\alpha$ and $325$ of $\delta$, and the ``published'' column quotes them --- the eleven-digit figure that appeared here in earlier versions was the crude intermediate bound of their \S IV\,E, not their result; Briggs independently computed $\ell = 4$ as well --- Table 2.3 of his dissertation \cite{briggs1997thesis} contains $\alpha(d)$, $\delta(d)$ for $d = 2\ldots 12$ and agrees with BOT on all shared digits.

\textbf{Normalization of $\alpha$.} The comparison for $\alpha$ is a comparison of ABSOLUTE VALUES, and this must be stated plainly. The external sources define $\alpha := 1/G(1)$ without the minus sign, and for them $\alpha$ is negative: Briggs has $\alpha(2) = -2.5029078750958928\ldots$, $\alpha(4) = -1.6903029714052448\ldots$ (Table 2.3); Burbanks--Osbaldestin--Thurlby have $\alpha = 1/a = -1.6903029714\ldots$ with $a = \Gs(1)$ (caption to Fig. 1). Our $\alpha_{G} = -1/G(1)$ (\S\ref{sec:1}) gives the same number with a plus sign. The matching digits in Table~\ref{tab:3.9-classic} are a match of $|\alpha|$, not a literal match of the definitions.

The match cannot go deeper than this: the center $P(\ell)$ is at most $\rho^{\star}_{X} \approx 10^{-12}$ away from the solution $\Gs(\ell)$, and it is this distance, not the precision of the arithmetic, that sets the limit. The observed discrepancy of $8.6\cdot 10^{-16}$ for $\delta(2)$ and $1.2\cdot 10^{-14}$ for $\alpha(2)$ lies within this limit.

Check: \texttt{04-\allowbreak CODE/\allowbreak tests/\allowbreak \_\allowbreak classic\_\allowbreak constants.py}. The test fails if the match is worse than fifteen digits for $\delta(2)$ and fourteen for $\alpha(2)$, $\delta(4)$ and $\alpha(4)$ --- in every case the full published reference is used, not a truncation of it.

\textbf{What this closes and what it does not.} It closes the common point of failure \textbf{at the ends of the interval}: an error in the derivation of $DF$ gross enough to move the leading eigenvalue would show up at the fifteenth digit of $\delta(2)$. But only at the ends, along the leading mode, and no more than that. The comparison sees two numbers, and only two: $\delta$, the leading eigenvalue of $\DT$, and $\alpha$, read off the centre $P(\ell)$. An error confined to the rest of the spectrum is invisible to it --- and such an error is by itself enough to change the count of eigenvalues outside the unit disc, on which Theorem B stands;
what compares the assembly of $\DT$ entry by entry is the cross-check described next --- inside the acceptance run at $\ell = 3$ alone, and by hand at two further points. An error with a parametric dependence that vanishes at $\ell = 2$ and $\ell = 4$ --- say, with a factor $(\ell-2)(\ell-4)$ --- is not caught by this cross-check at all. Such an error looks artificial, but this check gives no grounds to claim that the common point of failure is closed over the whole \texttt{[2,4]}.
It closes nothing about the \textbf{interior} of the interval either. Numerical values of
$\alpha$ and $\delta$ for a continuous critical-point order, and hence at non-integer
$\ell$, have been published for this family since the early 1980s
\cite{humao1982period, vanderweele1987period}; no comparison against them is made here.
One published interior value is used: $\alpha(3)$ of Briggs \cite{briggs1991precise},
Table~2, printed to fifty digits. The float64 scan here agrees with it to $14$
significant digits (\S\ref{sec:2}), and that number is the resolution of the
scan rather than the strength of the comparison: against
$\alpha_{\Gbar} = -1/\Gbar(1)$ recomputed from the deposited seed the agreement is
to $50$ digits, and the pointwise certificate of \S\ref{sec:5.2} pins the
certified value, $|\alpha(\Gs(3)) - \alpha_{\Gbar}| \le \rho_{X}/\Gbar(1)^{2}
\le 6.51\cdot 10^{-20}$ with $\rho_{X} = 1.75\cdot 10^{-20}$
(\texttt{kant\_\allowbreak l3.json::\allowbreak rho\_\allowbreak X}), so the
published value lies inside a certified interval whose two ends share their first $18$ significant digits (the width sits in the 20th place, but the 19th digit of the center is a nine, so adding the half-width carries). It is a comparison of $\alpha$, not of the spectral quantity
$\delta$ this work certifies, and the $\delta(3)$ of that same table disagrees ---
see \S\ref{sec:2} for why the disagreement is placed on the side of the table. So the
interior is not wholly unverified externally, but what is verified there is the fixed
point, not the eigenvalue count that is the novelty of the work (\S\ref{sec:2}).

\textbf{An interior point: cross-assembly of $\DT$ by finite differences.} The common point of failure --- the derivation of $DF$ --- is now closed at the interior point $\ell = 3$ as well, and without appeal to external constants. The acceptance-run step \texttt{04-\allowbreak CODE/\allowbreak tests/\allowbreak \_\allowbreak dt\_\allowbreak fd\_\allowbreak cross.py} assembles the same matrix $\DT$ by a second route: the columns are obtained by a central finite difference of the operator's VALUE (\texttt{tube\_\allowbreak lib.compute\_\allowbreak F\_\allowbreak poly}), without a single call to the hand-written analytic derivative \texttt{DF\_at\_ell} (\texttt{tube\_\allowbreak lib.py:489-506}) or to \texttt{M\_full}. The two assemblies are compared element by element and by spectrum: relative discrepancy $1.18\cdot 10^{-12}$, the leading eigenvalue agrees to twelve digits ($\delta = 6.084689074262$), and the ``exactly one outside'' count is the same in both. The strength of the check is measured, not merely claimed: artificially corrupting \texttt{DF\_at\_ell} by a factor $1 + 10^{-4}$ turns the step red (discrepancy $1.34\cdot 10^{-4}$ against a threshold of $10^{-6}$). This closes the class of ``error in the analytic derivative'' at the point $\ell = 3$ --- including an error with the factor $(\ell-2)(\ell-4)$, which is invisible to the cross-check at the ends.

The check was carried out at three points, including a non-integer one, where no external values for comparison exist at all. Table~\ref{tab:3.9-dtfd} prints two diagnostics per point and one fact. The diagnostics --- the discrepancy of the two assemblies and the two values of $\delta$ --- are not certificates, and no statement of this article rests on them. The last column is different: the counts in it are computed in float, but what they agree on is exactly what Theorem B certifies for the whole of $[2,4]$, so that column carries no new claim, only a second route to an established one. The columns of both assemblies are built in \texttt{arb} at 256 bits, but the finite-difference truncation is not enclosed rigorously; the discrepancy is a comparison against a fixed threshold of $10^{-6}$; and the two $\delta$'s of each row are eigenvalues returned by \texttt{numpy.\allowbreak linalg.\allowbreak eigvals} on the float64 midpoints, so their digits carry no interval guarantee. Two of the three points do have a certified $\delta$, established elsewhere and by a different route: at $\ell = 3$ the Krawczyk enclosure $\deltastar = 6.084689074262122$ of box width $6.226\cdot 10^{-19}$, and at $\ell = 5/2$ the separate pointwise certificate $\deltastar = 5.412736346083841$. Neither is what this table prints. Its digits are the two float64 spectra agreeing with each other, nothing more --- and at $\ell = 3.9$ no enclosure of $\delta$ is claimed anywhere in this article.

\begin{table}[h]
\centering
\caption{Cross-assembly of $\DT$ at three interior points.}\label{tab:3.9-dtfd}
{\footnotesize\setlength{\tabcolsep}{4pt}
\begin{tabular}{lllll}
\toprule
$\ell_{0}$ & branch & discrepancy of assemblies & $\delta$ (both routes) & ``outside'' \\
\midrule
2.5 & \texttt{branch\_\allowbreak g11.npz} & $7.36\cdot 10^{-14}$ & 5.412736346084 & 1 = 1 \\
3.0 & \texttt{branch\_\allowbreak ganchor.npz} & $1.18\cdot 10^{-12}$ & 6.084689074262 & 1 = 1 \\
3.9 & \texttt{branch\_\allowbreak g42.npz} & $1.12\cdot 10^{-12}$ & 7.171724869293 & 1 = 1 \\
\bottomrule
\end{tabular}}
\end{table}
\FloatBarrier

A caveat remains: these are three points, not an interval --- the second assembly was not run over the whole \texttt{[2,4]} (about 100 s per point). And the second route uses the same code for the operator's value, \texttt{compute\_\allowbreak F\_\allowbreak poly}: it is independent of the derivation of the DERIVATIVE, but not of the engine as a whole. The acceptance run includes the point $\ell = 3$; the others are reproduced with a command-line argument (\texttt{python tests/\allowbreak \_\allowbreak dt\_\allowbreak fd\_\allowbreak cross.py 2.5}).


\subsection{What the existence certificate does not prove}\label{sec:3.10}

\begin{itemize}
\item \textbf{Hyperbolicity on [2,4] is not covered by this section} --- only existence is covered here. The spectral part is discussed separately.
\item \textbf{A $\sigma$-independent floor of $\sim 7\cdot 10^{-16}$} remains in the model self-check; it does not decrease either with growing $D$ or with a narrower box. Its source has been identified --- the Chebyshev $\to$ Taylor conversion in float64 --- and removing that conversion lowers the floor by a factor of 36--71, but not to zero: a residue of $\sim 9\cdot 10^{-18}$ (the top of that range applied to the floor: $7\cdot 10^{-16}/71 \approx 9.9\cdot 10^{-18}$) is produced by arb rounding within the $\sigma$-operator assembly itself. \textbf{Both figures here --- the factor $36$--$71$ and the residue --- are float64 measurements, not certified quantities, and no step of any proof rests on them.} (A one-off experiment outside the package: \texttt{04-CODE} contains no assembly that replaces the float64 conversion with arb entirely --- \texttt{\_\allowbreak cheb\_\allowbreak to\_\allowbreak taylor\_\allowbreak arb} in \texttt{rig2fix3.py:367-392} is used only to estimate the conversion error; it is not reproducible from the package as shipped). In the current certificate the floor is bypassed by the guard criterion (\S\ref{sec:3.6}), not eliminated.
\item \textbf{The upper end of the interval is more demanding than the lower one.} The interval \texttt{Z1c} grows with $\ell$ over most of the segment, though not from the very start: on \texttt{g00}--\texttt{g11} ($\ell\le 2.52$) it \emph{falls}, by a factor of $1.4832$, and its minimum over the 46 boxes sits on \texttt{g29}, not at the left end. Over the 46 boxes deployed the largest GROWTH from one neighbour to the next (step 0.045) is $2.9425$ (\texttt{g33}$\to$\texttt{g34}); in the opposite direction the drop reaches $1059.3$ (\texttt{g28}$\to$\texttt{g29}), three orders of magnitude --- so ``ratio of neighbours has to name its direction. From the minimum to the maximum across all boxes \texttt{Z1c} grows by more than three orders of magnitude, $1635.6\times$ (field \texttt{Z1c} in \texttt{box\_g*.json}); the maximum is on \texttt{g43}, while the topmost box \texttt{gtop} --- narrow by construction --- carries a \texttt{Z1c} smaller by a factor of $5.84$. Hence both the adaptive $K$ and the narrower box at the very top.
\end{itemize}

*Code, certificates, and logs for this section --- \texttt{04-CODE/}, \texttt{03-\allowbreak CERTIFICATES/\allowbreak existence/\allowbreak }, index in \texttt{05-\allowbreak DATA-\allowbreak INDEX.md}. Certificates: \texttt{box\_g*.json}, \texttt{glue.json} (in \texttt{existence/}), \texttt{branch\_g*.npz} (in \texttt{branches/}). Diagnostics for this section (they print numbers, are not included in the acceptance run \texttt{verify\_\allowbreak pkg.sh} --- the acceptance-run script itself is not part of the package, it is supplied to the referee separately together with instructions --- and are NOT a verdict): \texttt{04-\allowbreak CODE/\allowbreak tests/\allowbreak \_\allowbreak z1bc\_\allowbreak sub.py}, \texttt{\_\allowbreak dchk\_\allowbreak check.py}, \texttt{\_\allowbreak tau\_\allowbreak share.py}. The section's gates, which give a verdict and a return code, are included in the acceptance run separately (\texttt{\_\allowbreak tau\_\allowbreak test.py}, \texttt{\_\allowbreak dchk\_\allowbreak bound\_\allowbreak test.py}, \texttt{\_\allowbreak arc\_\allowbreak cover\_\allowbreak test.py}, \texttt{\_\allowbreak nozeros\_\allowbreak check.py}, \texttt{glue.py}). Data: \texttt{05-\allowbreak DATA/\allowbreak COVERAGE.md}, \texttt{05-\allowbreak DATA/\allowbreak xcheck\_\allowbreak scale.py} (the script itself; its input --- the march bank --- is not part of the package, so the 28-point cross-check table of \S\ref{sec:3.8} is not reproducible from the package, though its four rows are written out in \texttt{05-\allowbreak DATA/\allowbreak COVERAGE.md} \S5. The coverage statistics of the same section do reproduce, and not only by reading: \texttt{05-\allowbreak DATA/\allowbreak coverage.json} gives $n_{\mathrm{total}} = n_{\mathrm{pass}} = 81\,479$, $n_{\mathrm{fail}} = 0$, $55$ connected components inside $[2.07, 3.580716250000099]$, and the sum of their lengths reproduces the stored \texttt{covered} $= 0.30302591324252637$ to the printed precision; the same union is the union of the eight machine lists, whose counts add up to the same $81\,479$. See \texttt{05-\allowbreak DATA-\allowbreak INDEX.md}, which states this split more precisely than the sentence above used to.)*



\section{Hyperbolicity on the whole interval: the Rouché tube}\label{sec:4}


\subsection{What is claimed}\label{sec:4.1}

\begin{theorem}[Theorem B, restated]
For every $\ell \in [2, 4]$, the operator
$DT(\Gs(\ell))|_{X_{0}}$ on the working layer ($c = 3/5$, $r = 4/5$) has
\textbf{exactly one} eigenvalue outside the unit disk; the rest of the spectrum
lies strictly inside. In particular, the unstable direction is
one-dimensional, the eigenvalue $\delta(\ell)$ is real and simple, and the spectrum
never touches the unit circle for any $\ell$ in the interval.

The count is carried out first for the projection onto the first
$N_p = 80$ coordinates of the b-basis, and is then transferred to the full
operator (\S\ref{sec:4.6}): the tail enters via separate explicit bounds, rather
than being discarded.
\end{theorem}

This claim is uniform in $\ell$: it is not derived from counting at individual
points, and does not rely on continuity of the spectrum as an argument.
Below is how it is obtained.


\subsection{Why not ``deflation + power of the norm''}\label{sec:4.2}

At $\ell = 3$ and $\ell = 5/2$, hyperbolicity is proved by deflation (\S\ref{sec:5.2} for $\ell=3$,
\S\ref{sec:5.3.4} for $\ell=5/2$): the unstable direction is split off from the operator,
and for the remainder $Q^{\star}$ one checks $\norm{(Q^{\star})^{n}}_{1} < 1$. Carrying this device
over to a box did not succeed, and this is measured, not assumed --- the
figures in the table are float64 diagnostics of the failed routes, not
interval certificates, and nothing in the article rests on them:

\begin{table}[ht]
\centering
\caption{Routes attempted for the uniform eigenvalue count over a parameter box.}
\label{tab:routes}
\begin{tabular}{@{}p{0.32\textwidth}p{0.60\textwidth}@{}}
\toprule
route & what happened \\
\midrule
deflation + $\norm{Q^{n}}$ on the interval matrix &
the interval $Q^{16}$ blows up to $2.5\cdot 10^{20}$ even at box width
$10^{-5}$: the sixteen multiplications themselves are to blame, not the
width \\
\addlinespace
Bauer--Fike ($\norm{\Delta}$ against the spectral gap) &
requires normality; $\kappa(V) = 4.12\cdot 10^{9}$, and the half-width
drops to $4\cdot 10^{-12}$ \\
\addlinespace
quadratic form, inertia of $A^{\mathsf{T}}HA - H$ &
\textbf{the premise is false}: for arbitrary $H \succ 0$, the inertia of
$A^{\mathsf{T}}HA - H$ does not count the spectrum outside the disk
(counterexample: $A = [[\tfrac{1}{2},10],[0,\tfrac{1}{2}]]$,
$H = I$ --- spectrum inside, inertia $(1,1,0)$). The correct statement is
the Stein--Taussky--Wimmer inertia theorem (reference \cite{stein1952theorems, taussky1964matrices, wimmer1973ostrowski, horn1991topics}): when
$A^{\mathsf{T}}HA - H \prec 0$, the spectrum is counted by the inertia of
$\mathbf{H}$ \textbf{itself} \\
\addlinespace
\textbf{resolvent + argument principle (Rouché)} &
works; below \\
\bottomrule
\end{tabular}
\end{table}

The reason the first three fail is the same in each case: they estimate a
\textbf{magnitude}, which is sensitive to non-normality ($\kappa(V) \approx 4\cdot 10^{9}$), whereas
the count of eigenvalues outside the contour is a \textbf{discrete} quantity and
does not depend on the conditioning of the basis at all.


\subsection{Scheme: Rouché in operator form}\label{sec:4.3}

Let $\Gamma$ be the unit circle, $DT_{0} = DT(P(\ell_{0}), \ell_{0})$. If $\Gamma$ contains no
spectrum of $DT_{0}$ and
\[
\sup_{z \in \Gamma} \norm{(DT_{0} - z)^{-1}}_{1} \cdot \norm{\Delta}_{1} < 1,
\]
then $DT_{0} + \Delta$ has the same number of eigenvalues \emph{outside} $\Gamma$ as
$DT_{0}$ (with multiplicities). The count is taken outside, not inside, and this is
not a matter of taste: the operator compared against annihilates the whole tail, so
inside $\Gamma$ the total multiplicity is infinite and the hypotheses of the argument
principle fail there outright. What is preserved is the rank of the \emph{outer}
Riesz projector, finite by construction; the operator form of the argument principle
(Gohberg--Sigal) is applied on an annulus whose inner boundary is $\Gamma$ and whose
outer circle is taken large enough. The full account --- the choice of domain and the
verification of the outer boundary --- is \S\ref{sec:4.6}. Taking $\Gamma =$ the unit circle therefore gives ``exactly one outside''
for the whole box at once, and moreover \textbf{the simplicity of $\boldsymbol{\delta}$ comes for
free}: if $\delta$ were to split, or a second eigenvalue were to approach it, the
outside count would become greater than one.

Denote $R = \sup_{z\in\Gamma} \norm{(DT_{0} - z)^{-1}}_{1}$. The entire substance lies in
obtaining $R$ and $\norm{\Delta}$ as \textbf{enclosures}, not as approximations.

One point about the norm should be made right away: $R$ and $\norm{\Delta}$ are
taken in \textbf{one and the same} column 1-norm of the coordinates of the
b-basis $b_k = e_k - (-c/r)^k e_{0}$, $k \ge 1$. These coordinates are honest
coordinates on $X_{0}$: the functional $h \mapsto h(0)$, in the scaled coefficients,
is $\sum \hat{h}_k (-c/r)^k$, and it is precisely $b_k$ that annihilates it. The
norm $\sum_{k\ge 1}\abs{\hat{h}_k}$ is equivalent to $\norm{\cdot}_X$ on $X_{0}$, and the count of
eigenvalues outside the contour does not depend at all on the choice of an
equivalent norm --- so no equivalence constant enters the estimates; all
that matters is that $R$ and $\norm{\Delta}$ are measured in the same norm.


\subsection{Rigorous resolvent}\label{sec:4.4}

At each of the $N$ nodes $z_k = e^{2\pi i k/N}$ one takes an APPROXIMATE
inverse $A_k$ (an ordinary float inverse of the matrix $DT_{0} - z_k$). It is
\textbf{not part of the proof} --- it is merely a preconditioner. Next, in complex
interval arithmetic (\texttt{acb}), one computes $E_k = I - A_k(DT_{0} - z_k)$, and
if $\norm{E_k}_{1} < 1$, then
\[
\norm{(DT_{0} - z_k)^{-1}}_{1} \le \norm{A_k}_{1} / (1 - \norm{E_k}_{1}) \qquad \text{(Neumann series)}
\]
--- and this is an enclosure. Between nodes, the resolvent identity is used:
\[
\norm{(DT_{0} - z)^{-1}} \le R_k / (1 - \abs{z - z_k}\cdot R_k), \qquad \abs{z - z_k} \le \pi/N,
\]
where the chord is majorized by the arc. At $N = 180\dots 240$ the arc markup is
5.2--7.1\,\% --- this is against the global bound $R \le 3.7961$ (premise a7 of
Table~\ref{tab:certified-premises}); for a
typical node with $R_k \approx 3.2$ it is smaller, 4.3--6.0\,\%. Both brackets are
rounded OUTWARD (the lower end down, the upper end up), as are all other
numbers in the package. The actual count uses the per-node $R_k$, not this
illustration.

A separate pitfall, which cost a false PASS on the first run: comparing
arb intervals. The operator \texttt{>} is true only under \textbf{guaranteed}
superiority, so accumulating a maximum via \texttt{if col > best} silently leaves
zero in place if the interval contains zero --- and $\norm{\Delta}$ comes out as zero.
The maximum is taken over the upper endpoints of the balls, which stay in
\texttt{arb} to the end; the cast to \texttt{double} happens once, at the
output, and rounds upward (\texttt{up\_\allowbreak float} in
\texttt{hyp\_\allowbreak rouche.py}). Taking it through
\texttt{float(x.upper())} would not do: \texttt{float} rounds to nearest, not
upward, and about half of the column quantities on the anchor box then stay
below their true upper bound.


\subsection{Perturbation over the box}\label{sec:4.5}

$\Delta$ must cover everything by which $DT$ at a point $\ell$ of the box differs
from $DT_{0}$. There are two terms, and the second is easy to lose:

\textbf{(1) Motion of the parameter and the center.} $DT$ is computed with an
INTERVAL $\ell$ and an INTERVAL center $P(\ell)$ simultaneously; against the
Miller wrapping, the box is cut into $K$ parts, and the maximum is taken
over the parts.

\textbf{(2) Displacement of the true fixed point.} The tube computes $DT$ at the
parameterized center $P(\ell)$, while the theorem is about $\Gs(\ell)$. Between
them $\norm{\Gs - P}_X \le \rhostar_X = 1.75\cdot\rhostar_b$ --- this is the radius of the EXISTENCE
BALL (premise a1 of Table~\ref{tab:certified-premises}), which is smaller by ${\sim}5.07$ orders of magnitude (maximum over the
46 boxes: $8.49\cdot 10^{-13}$ against $\rho' = 10^{-7}$) than the radius of the
Kantorovich gate $\rho' \sim 10^{-7}$ (that radius is the one INSIDE which
uniqueness and Newton convergence are guaranteed, not the distance to the
solution found). This term was absent in the first version; it is
recovered from quantities already certified in the existence certificate,
with no new assumptions:
\begin{gather*}
\norm{\Ahat\cdot(M'(u) - M'(P))}_{1} \le 1.75^{2}\cdot B_A\cdot L\cdot\rhostar_b \qquad \text{(second term of $\Zhat$)} \\
\Ahat M' = I - Z, \quad \norm{Z}_{1} = Z1 < 1 \quad\then\quad \Ahat^{-1} = M'(I - Z)^{-1} \\
\then\quad \norm{M'(\Gs) - M'(P)}_{1} \le \bigl( \norm{M'}_{1} / (1 - Z1) \bigr) \cdot 1.75^{2}\cdot B_A\cdot L\cdot\rhostar_b\,.
\end{gather*}

(The factor $1.75^{2}\cdot B_A\cdot L\cdot\rhostar_b$ here is not a fresh derivation from the
line above (that line gives only a single factor of $1.75$), but a REUSE
of what was already computed in \S\ref{sec:3.1} as the ``second term of $\Zhat$'' for this
same box, where $1.75^{2}$ is justified separately --- one factor converts the
radius, one the input vector (\S\ref{sec:3.1}, \S\ref{sec:6.6.2}). Numerically it does not affect the outcome: \texttt{extra}
amounts to $1.4346\cdot 10^{-8}\dots 2.86\cdot 10^{-6}$ (\S\ref{sec:4.8}) against a dominant
$\norm{\Delta}\approx 0.13\dots 0.28$.)

Here $\norm{M'}_{1}$ is the rigorous supremum over the box (interval $\ell$, interval
center), and $Z1, B_A, L, \rhostar_b$ come from the json of the existence box
that covers this piece of the tube ($\rhostar_b$ is the field \texttt{best.rho\_\allowbreak star\_\allowbreak b},
the same existence radius as in premise a1 of Table~\ref{tab:certified-premises};
NOT the gate radius $\rho'$ --- the symbol
$\rho'$ in the draft was left there by mistake and substituted a completely
different quantity in the formula, although the code (\texttt{tube\_\allowbreak offset.py})
always computed via \texttt{rho\_star\_b}, and the final number \texttt{extra} did not
suffer from this).

About $\norm{M'}_{1}$ one must be precise, because it is easy to substitute the
wrong quantity here. The inequality derived above contains the norm of the
\textbf{full} operator on $X_{0}$, while the engine naturally computes the norm of
its $80\times 80$ head projection. The norm of the projection does \textbf{not}
majorize the full norm, and the gate $< 1$ would not notice such a
substitution. Therefore the full column norm is used, assembled from the
same blocks as in \S\ref{sec:4.6}:
\begin{gather*}
k \le N_p:\quad \sum_j \abs{M'_{jk}} = \text{(head mass)} + \text{(tail mass)} \le \norm{M'_{\mathrm{head}}}_{1} + \norm{C}_{1} \\
k > N_p:\quad \sum_j \abs{M'_{jk}} \le 1 + \norm{[B;\,D]}_{1} \\
\then\quad \norm{M'}_{1} \le \max\bigl( \norm{M'_{\mathrm{head}}}_{1} + \norm{C}_{1}\,,\; 1 + \norm{[B;\,D]}_{1} \bigr)\,.
\end{gather*}

Here $C$ and $[B;\,D]$ are blocks of $\mathbf{DT}$, not of $M'$ (definitions ---
\S\ref{sec:4.6}). Off the diagonal, $M' = I - DT$ differs from $DT$ only in sign, so
for block $C$ (tail rows of head columns) the masses coincide. But on the
tail columns the diagonal also comes into play: $D_{M'} = I - D_{DT}$, and
the unity from $I$ gives the term $1 +$ in the second line. The earlier
draft simply wrote $\norm{[B;\,D]}_{1}$ there, that is, it substituted the mass of
the $DT$ block into the estimate of the $M'$ norm. This does not affect
the outcome --- the maximum is taken on the first line regardless
($\norm{M'_{\mathrm{head}}} \ge 6.34$ against $1 + 1.95\cdot 10^{-4}$, a margin of more than a
factor of six), and the code (\texttt{tube\_\allowbreak offset.py}) takes the same first
branch on all 46 existence boxes, from which the 706 tube boxes inherit their
value of \texttt{extra} --- but without the unity, the derivation is not
formally closed.

The maximum of the sum of maxima majorizes the maximum of the sums, so the
estimate is valid and slightly conservative. Numerically the correction is
small --- $\norm{C} \le 4.07\cdot 10^{-4}$ against $\norm{M'_{\mathrm{head}}} \in [6.34, 18.66]$, i.e.\ \texttt{extra}
grows by no more than a factor of $1.0000318$ --- but without it the chain
is not formally closed. The factor is a maximum over boxes of a ratio taken
\emph{within} each box, and it is attained on \texttt{g28}, where
$\norm{C} = 4.0682\cdot 10^{-4}$ and $\norm{M'_{\mathrm{head}}} = 12.80$; the two
extremes of the printed range fall on different boxes, so pairing the largest
$\norm{C}$ with the smallest $\norm{M'_{\mathrm{head}}}$ would give the weaker and
unattained $1.0000642$.

The strengthened inequality that is checked is
\[
R \cdot \bigl( \norm{\Delta} + \mathrm{extra} \bigr) < 1\,.
\]



\subsection{From the projection to the operator}
\label{sec:4.6}

Everything above is computed for the matrix $A = \DT|_{\mathrm{head}}$ of
size $N_{p} \times N_{p}$. By itself, this is a statement about the
projection, not about the operator. The transfer is made using the same
Rouch\'e device, and the existence certificate already supplies all the
numbers needed.

In the b-basis it is $\mathbf{DT}$ itself that is split into blocks --- not
$M' = I - \DT$:

\begin{align*}
\DT &= \begin{bmatrix} A & B \\ C & D \end{bmatrix}
      \qquad
      \begin{aligned}
         &\text{head: } k, j \le N_{p} , \\
         &\text{tail: } k, j > N_{p} .
      \end{aligned}
\end{align*}

It is specifically $\DT$, and this is not nitpicking over notation but a
difference in order of magnitude. The tail column of $\DT$ is small:
$\norm{[B; D]}_{1} \le 1.95\cdot 10^{-4}$ (certified below). The tail
column of $M' = I - \DT$, however, also contains the diagonal of the
identity operator, i.e.\ it equals $1 \pm 1.95\cdot 10^{-4}$. Hence the
comparison with $A(\ell_{0}) \oplus 0$ below is valid for $\DT$ and would
be FALSE for $M'$: for $M'$ the tail perturbation would be approximately
unity, and the product with $R_{A} \approx 3$ would not pass the
gate $< 1$. Later in the text (\S\ref{sec:5.1.2}, where the matrix $M' = DF|_{X_{0}}$
of the existence certificate is introduced), the same letter denotes
$DF = I - \DT$ --- a different object; here and in \S\ref{sec:4.5} the blocks
$A$, $B$, $C$, $D$ are always blocks of $\DT$.

The engine computes rigorous majorants of the column masses of these
blocks, and in the certificate they are stored not by themselves but
multiplied by the preconditioner constants --- and it is in this form that
they enter as terms in $Z1 = \norm{I - \widetilde{A} M'}_{1}$:

\begin{align*}
\hat{c} &:= \max_{k \le N_{p}} \Bigl[ \sum_{j > N_{p}} |DT_{jk}|\, r^{j}
      + (c/r)^{k}\cdot \sum_{j > N_{p}} |DT_{j0}|\, r^{j} \Bigr] \\
   &\qquad\qquad \text{(mass of block } C \text{, b-basis correction --- see (iii))} , \notag \\
\hat{b} &:= \hat{\beta}_{N_{p}+1} + (c/r)^{N_{p}+1}\cdot \beta_{0}^{\star}
      \qquad \text{(mass of block } [B; D] \text{)} , \\
Z1b &:= B_{A} \cdot \hat{c} ,\qquad Z1c := (1 + B_{A}) \cdot \hat{b} .
\end{align*}

(\code{04-CODE/tube\_lib.py:880} and \code{04-CODE/tube\_lib.py:904} --- this
is the route whose values go into the certificate; the diagnostic
$\sigma$-route, \code{04-CODE/rig2fix3.py:174} and
\code{04-CODE/rig2fix3.py:205}, has the same form.) Hence

\begin{align*}
\norm{C}_{1} &\le \mathrm{Z1b} / B_{A} , \\
\norm{[B; D]}_{1} &\le \mathrm{Z1c} / (1 + B_{A}) .
\end{align*}

Division by itself proves nothing --- from
$\norm{\widetilde{A}X}_{1} \le \eta$ and $\norm{\widetilde{A}}_{1} \le B_{A}$
one CANNOT derive $\norm{X}_{1} \le \eta/B_{A}$ (with
$\widetilde{A} = \operatorname{diag}(100, 0.01)$ and
$X = (0, 1)^{\mathsf{T}}$, the left side equals $0.01$, $B_{A} = 100$, and
$\norm{X}_{1} = 1$, i.e.\ the ``estimate'' would give the false
$1 \le 10^{-4}$). The validity rests on three links, and it is worth
spelling them out, because each of them was once tacitly assumed.

\textbf{(i) The multipliers are built into the construction itself}, not
appended afterward: this is visible in the two lines of code above.

\textbf{(ii) Where $1 + B_{A}$, not $B_{A}$, comes from.} The preconditioner
is defined as $\widetilde{A}_{\mathrm{full}} :=
\widetilde{A}_{\mathrm{head}} \oplus I_{\mathrm{tail}}$
($\widetilde{A}_{\mathrm{head}}$ is an $N_{p} \times N_{p}$ matrix; outside
the head, the continuation is the identity). For a tail basis column this
gives

\begin{align*}
\norm{\widetilde{A}_{\mathrm{full}}\cdot \DT \cdot b_{k}}_{1}
   &\le \norm{\widetilde{A}_{\mathrm{head}}\cdot B\, b_{k}}_{1} + \norm{D\, b_{k}}_{1} \\
   &\le B_{A}\cdot \norm{B\, b_{k}}_{1} + \norm{D\, b_{k}}_{1}
      \le (1 + B_{A})\cdot \hat{b} ,
\end{align*}

that is, the head output passes through $\widetilde{A}_{\mathrm{head}}$,
while the tail output passes through untouched; the sum gives
$1 + B_{A}$. (This is conservative: for a majorant on the COMBINED norm
$[B; D]$, $\max(B_{A}, 1)$ would suffice. That is why it is written
$\le \mathrm{Z1c}$, and not ``equals''.)

\textbf{(iii) Why $\hat{b}$ majorizes precisely $\norm{[B; D]}_{1}$.} The
tail column of the b-basis is $b_{k} = e_{k} - (-c/r)^{k} e_{0}$, so
$\DT b_{k} = \DT e_{k} - (-c/r)^{k}\cdot \DT e_{0}$, and with
$\beta_{k} \ge \norm{\DT e_{k}}_{1}$,
$\beta_{0}^{\star} \ge \norm{\DT e_{0}}_{1}$, and monotonicity
$\beta_{k} \le \hat{\beta}_{N_{p}+1}$ for $k > N_{p}$, we get
$\norm{\DT b_{k}}_{1} \le \hat{\beta}_{N_{p}+1} + (c/r)^{N_{p}+1}\cdot
\beta_{0}^{\star} = \hat{b}$.
Monotonicity of $\beta_{k}$ for $k > N_{p}$ is not a hypothesis but a
direct consequence of its CLOSED form (not an empirical fit):
\code{04-CODE/beta\_bound\_lib.py} gives a rigorous arb-estimate of
$\beta_{k}$, decomposing into three terms, each a geometric progression in
$k$ with base $q$, $q_s$, or $\frac{1}{2}$
($\delta1_{k} = ((1-c)/r)^{k} = 0.5^{k}$); all three bases are STRICTLY
$< 1$, which is checked separately (\code{bases\_check}, the same
\code{RuntimeError} gate as the other gates \textbf{[C]} --- not a silent
assumption). Each term strictly decreases in $k$, hence so does the whole
sum, and $\sup_{k > N_{p}}\beta_{k} = \beta_{N_{p}+1}$ holds literally,
not heuristically.

This also shows why the identity term from $M'$ does not enter here:
$\widetilde{A}_{\mathrm{full}}$ is the identity outside the head, so
$I - \widetilde{A}_{\mathrm{full}}\cdot M' = I -
\widetilde{A}_{\mathrm{full}}\cdot (I - \DT) = (I -
\widetilde{A}_{\mathrm{full}}) + \widetilde{A}_{\mathrm{full}}\cdot \DT$,
and on the tail columns the first term vanishes --- leaving exactly
$\widetilde{A}_{\mathrm{full}}\cdot \DT$. For block $C$ the question does
not arise at all: it is off-diagonal, and off the diagonal $M'$ and $\DT$
differ only in sign. The acceptance step
\code{04-CODE/tests/\_numbers\_check.py} guards both the structure of the
definitions in the code and the extraction on every box, and the fact that
the maximum of the full norm is taken on the head branch.

Let us compare the full operator with the block-diagonal
$S = A(\ell_{0}) \oplus 0$, that is, with the head block at the center of
the box and a ZERO tail block:

\begin{itemize}
\item the \textbf{resolvent of $S$} on $|z| = 1$ equals the maximum over
  blocks, $\max(R_{A}, 1) = R_{A}$, since $R_{A} \approx 3 > 1$;
\item \textbf{the count of eigenvalues outside for $S$ is the same as for
  $A(\ell_{0})$}: the zero block contributes only the eigenvalue $0$.
  The same holds for $A(\ell_{0}) \oplus D$: the spectral radius of the
  tail block does not exceed its norm (a notational remark: $\norm{D}$ is
  not computed separately anywhere in the package --- wherever it enters
  an estimate, the computed $\mathrm{BD\_norm} = \norm{[B;D]}$ is
  substituted, a larger quantity, so the substitution is conservative:
  \code{tube\_offset.py:154-155}, \code{tube\_accept.py:130-136}), and that
  norm is four orders of magnitude below one
  ($\norm{D} \le \norm{[B;D]} \le 1.95\cdot 10^{-4}$), so the choice between
  $\oplus 0$ and $\oplus D$ does not affect the conclusion;
\item the \textbf{perturbation} $\DT - S$ in the column norm (specifically
  $\DT$, not $M' = I - \DT$: for $M'$ the tail column carries an extra
  identity term, and the norm $\norm{[B;D]} + \mathrm{extra}$ would not
  majorize it): the head columns are $[A' - A(\ell_{0},P); C']$, with
  norm $\le \norm{\Delta} + \mathrm{extra} + \norm{C}$; the tail columns
  are $[B'; D']$, with norm $\le \norm{[B;D]} + \mathrm{extra}$. These
  are exactly the two quantities that \code{tube\_accept.py} computes.
\end{itemize}

\paragraph{Why $\oplus 0$, and not $\oplus D$.} Against
$A(\ell_{0}) \oplus D$ the tail perturbation would be
$[B'; D' - D]$ with norm $\le \norm{B'} + \norm{D' - D}$, and the quantity
$\norm{[B;D]} + \mathrm{extra}$ (this is the norm of the WHOLE tail column
of $\DT$, together with $D'$) does not majorize it. The certified numbers
support the comparison specifically with $A(\ell_{0}) \oplus 0$; this does
not affect the conclusion, by the second point above. The stated
formulation diverged from what the code actually does, although the code
itself is conservative.

The resulting inequality, checked on every box:

\begin{multline*}
\max(R_{A},\, 1/(1-\norm{[B;D]}-\mathrm{extra}))\cdot{}\\
\max(\norm{\Delta} + \mathrm{extra} + 2\norm{C},\; \norm{[B;D]} + \mathrm{extra}) < 1 .
\end{multline*}

The first factor here is guaranteed to be no smaller than the required
$R_{A}$: $1/(1-\norm{[B;D]}-\mathrm{extra})$ is slightly above one, while
$R_{A} \approx 3$, so the $\max$ gives $R_{A}$ in any case. (The
denominator carries the norm of the WHOLE tail column $\norm{[B;D]}$,
not a single block $\norm{D}$ --- this is exactly what
\code{tube\_accept.py} takes; see the ``Why $\oplus 0$, and not
$\oplus D$'' remark above.) It is left in this form as a margin, not
because the comparison is against $\oplus D$.

Numbers: $\norm{C} \le 4.07\cdot 10^{-4}$,
$\norm{[B;D]} \le 1.95\cdot 10^{-4}$,
$\mathrm{extra} \le 2.86\cdot 10^{-6}$ --- all three much smaller than
$\norm{\Delta} \approx 0.13\ldots 0.28$, so the transition from the
projection to the operator costs about a quarter of a percent of margin:
the worst full test $0.9896$ ($R\norm{\Delta}$, still without the tail
coupling) $\to$ $\mathbf{0.9919}$ (with $2\norm{C}$). The check passes on
all $706$ boxes (\S\ref{sec:4.8}).

The factor $2$ on $\norm{C}$ in the head columns is not derived --- it is
a deliberate margin: $2\norm{C}$ instead of $\norm{C}$, kept so the
certificate need not be recomputed when $\norm{C}$ is refined; it costs
no more than
$R_{A}\cdot \norm{C} \lesssim 3.8\cdot 4.07\cdot 10^{-4} \approx
1.5\cdot 10^{-3}$ of margin, that is, it does not affect the verdict.

Separately, for the TAIL block $[B;D]$ the B-FIX fix applies
(\code{tube\_\allowbreak accept.py}): the honest bound
$\norm{[B;D]} + \mathrm{extra}$ majorizes the old formula
$2\norm{[B;D]}$ whenever $\mathrm{extra} \le \norm{[B;D]}$ --- verified on
all $46$ boxes of the certificate: no violations, the worst ratio
$\mathrm{extra}/\norm{[B;D]} = 0.888$ (box \code{g30}).

$\norm{C}$ deserves a separate remark: the tail of the operator column is
the most fragile spot in the whole bookkeeping, and an error was found
here twice.

\emph{First time.} The tail beyond degree $N_{p}$ was estimated by the
quantity $4\cdot \max(\text{tail of } \overline{W}, \text{tail of }
\overline{U}^{\ell-1})\cdot \norm{\mathrm{column}}$, that is, by the
\textbf{Cauchy tail of the series $\overline{W}$ beyond degree $192$},
whereas what is needed is the tail of the composition with $\overline{W}$
beyond degree $80$. These are different quantities: the second exceeds the
first by 13--24 orders of magnitude ($5.7\cdot 10^{-32}$ versus
$1.95\cdot 10^{-8}$ at $\ell = 3.9925$; these two are the output of a one-off
diagnostic recomputation, not certified quantities --- they establish the gap
in order of magnitude and enter no bound). Found by an independent
recomputation of the columns up to t-degree $300$.

\emph{Second time.} In the corrected scheme the column was computed as a
direct sum up to $\NMAX = 320$, and the remainder beyond
$\NMAX$ was bounded only for one of the five terms ---
$\alpha_{G}\cdot G'_{c}(\overline{W}-c)\cdot DW_{v}$ --- with the
justification ``the remaining terms have degree no higher than $80$.'' The
justification is false: $\overline{W}-c$ is a series, so both
$v_{c}(\overline{W}-c) = (\overline{W}-c)^{k}/r^{k}$ and
$G_{c}(\overline{W}-c)$ have nonzero coefficients of arbitrarily high degree
--- already for a polynomial $\overline{W}$ this is $k\cdot \NW = 15360$.
Their tails did not enter the estimate at all. On top of that,
$\overline{W}$ itself is taken truncated at $\NW = 192$, and the cost of
this truncation was discarded along with the old term.

A naive fix is not possible:
$\sup_{|u|\le R}|v_{c}(\overline{W})|$ grows like
$((M_{W}+c)/r)^{k}$ and at $\ell \approx 4$, $k = 80$ reaches $10^{231}$
--- against the Cauchy factor $(r/R)^{321} \approx 10^{-57}$ this would
blow up the gate. So the tail is estimated not through a $\sup$ on the
circle of radius $R$, but in a weighted $\ell^{1}$-norm --- in the same
Banach algebra in which everything else works (\S\ref{sec:6.5}, item 6):

\begin{align*}
A_{f}(\lambda) &:= \sum_{j} |f_{j}|(\lambda r)^{j} , \qquad \lambda \ge 1 , \\
A_{fg}(\lambda) &\le A_{f}(\lambda)\cdot A_{g}(\lambda) , \\
\sum_{j > N} |f_{j}| r^{j} &\le \lambda^{-(N+1)}\cdot A_{f}(\lambda) .
\end{align*}

The first is the submultiplicativity of the convolution norm, the second
is the identity $|f_{j}|r^{j} = |f_{j}|(\lambda r)^{j}\cdot \lambda^{-j}$.
From this the tail of a product is bounded by the product of the
$A$-norms of the factors, and the coefficient decay of $\overline{W}$ does
the work where a $\sup$ on a larger circle is useless. The tails of
the series $\overline{W}$, $\overline{U}^{\ell-1}$ themselves beyond
$\NW$ enter their $A$-norms via Cauchy ($|w_{j}| \le M_{W}/R^{j}$),
which is exactly what requires $\lambda < R/r = 1.5$. The column estimate
is the sum of three derived parts:

\begin{align*}
\mathrm{tail\_col}(k) &= \sum_{j=N_{p}+1}^{\NMAX} |res_{j}|\, r^{j}
      \qquad && \text{direct sum} \\
   &\quad + E_{\mathrm{trunc}}(k)
      \qquad && \text{cost of truncation at } \NW \\
   &\quad + \min_{\lambda} \lambda^{-(\NMAX+1)}\cdot A_{\mathrm{res}}(\lambda,k)
      \qquad && \text{tail beyond } \NMAX
\end{align*}

(The technique is standard: this is the weighted $\ell^{1}$ Wiener algebra,
where the weight $\lambda^{j}$ simultaneously provides submultiplicativity
and a geometric tail cutoff.)
\code{E\_trunc} is derived from the same algebra by the inequality
$\norm{a^{k} - b^{k}} \le k\cdot \max(\norm{a},\norm{b})^{k-1}\norm{a-b}$,
applied at the four places where $\overline{W}$ and
$\overline{U}^{\ell-1}$ enter \code{res} \textbf{as factors}. The exponent
$k$ here is ALWAYS the integer monomial index ($k = 0..N_{p}$, the same
$k$ that numbers the matrix column, or the index of the Taylor coefficient
of \code{G\_c}/\code{G'\_c}) --- NOT the fractional exponent $\ell-1$ of
$\overline{U}$ itself; \S\ref{sec:3.6} declares this very inequality false
specifically for a non-integer POWER of $a$, $b$ (there $a = \overline{U}$,
a fractional exponent), not for an integer power of the scalar norm
$\norm{\cdot}^{k}$, which is what is taken here
(\code{04-CODE/tube\_lib.py}, the line with \code{nu\_W ** (k - 1)},
$\mathrm{nu}_{W} := \norm{\overline{W}-c}$ --- a norm, a number, $k$ an
integer). The truncation of $\overline{U}^{\ell-1}$ itself
(\code{tail\_hUlm1}) is estimated by a SEPARATE technique, valid for
non-integer powers --- a geometric Cauchy tail --- without this inequality
at all.

Numerically: in the column that determines $Z1b$ (the maximum over $k$),
the two new terms amount to a fraction of about $2\cdot 10^{-9}$ --- nine
orders of magnitude below the direct sum. In columns with small $k$, on
boxes at small $\ell$, the direct sum is practically zero, and there the
new terms make up almost the whole quantity, but that quantity itself is of
order $10^{-32}$ and does not enter the maximum. The outcome is checked not
by estimation but by fact: after the fix, $\mathrm{Z1b}$ and $\mathrm{Z1c}$
shifted by no more than $10^{-6}$ relative in all $46$ certificates ---
exactly a decimal round-up. What changed was not the number but its
justification.

On the factor $2$ in front of $\norm{C}$. The derived bound on the head
column of the full perturbation is $(\norm{\Delta} + \mathrm{extra}) +
\norm{C}$ --- the 1-norm of the block column
$[A(\ell,\Gs) - A(\ell_{0},P) ; C(\ell,\Gs)]$ does not exceed the sum of
the block norms. The gate uses
$\norm{\Delta} + \mathrm{extra} + 2\norm{C}$: a deliberate margin on top
of what was derived, left in place so the number in the certificate would
not need recomputing when $\norm{C}$ is refined. It majorizes the honest
estimate, since $2\norm{C} \ge \norm{C}$.

Compactness of the operator is \textbf{not needed} here --- the block
structure suffices. For the block-diagonal $A \oplus 0$, the spectrum is
the union of the spectra, and outside the unit circle it consists only of
the points of $\spec(A)$ (the zero block contributes only the eigenvalue
$0$), that is, of finitely many points; the corresponding Riesz projector
has finite rank, equal to 1. A perturbation for which
$\sup_{\Gamma}\norm{(A \oplus 0 - z)^{-1}}\cdot \norm{\Delta} < 1$
preserves the rank of this projector (the Gohberg--Sigal theorem \cite{gohberg1971operator}, the
operator form of the argument principle). \emph{The domain to which that theorem
is applied is the annulus} $D = \{1 < \lvert z\rvert < R\}$ with
$R > \norm{A \oplus 0} + \norm{\Delta}$, not the disk and not the exterior of
$\Gamma$. This matters, and in two directions. On the disk the hypotheses fail
outright: $A \oplus 0$ annihilates the whole tail $k > N_{p}$, so
$(A \oplus 0) - z$ has an infinite-dimensional kernel at $z = 0$, is not Fredholm
there, and the total multiplicity inside would be infinite. The exterior of
$\Gamma$ is not bounded, so it is not admissible either. On the annulus, by
contrast, $(A \oplus 0) - z$ is invertible except at the finitely many
eigenvalues of $A$ outside the unit circle, each of finite multiplicity, so the
family is finitely meromorphic and Fredholm on $\overline{D}$; its boundary is
$\Gamma \cup \Gamma_{R}$, and on the outer circle the Rouch\'e condition holds
automatically because $\norm{((A \oplus 0) - z)^{-1}} \le 1/(R - \norm{A \oplus 0})$,
which is below $1/\norm{\Delta}$ for $R$ large enough. Only the condition on
$\Gamma$ therefore needs certifying, and that is what is certified. (An
alternative route avoids the outer circle altogether: the Riesz projector of
the outer part of $A \oplus 0 + \tau\Delta$ is norm-continuous in
$\tau \in [0,1]$ under the same condition on $\Gamma$, and the rank of a
continuous family of projectors is locally constant; this needs the contour
$\Gamma$ only.) Hence the full operator
$\DT|_{X_{0}}$, too,
has exactly one eigenvalue outside the unit circle, counted with
multiplicity. (The comparison is made against $\oplus 0$, not $\oplus D$,
for the reason given in the ``Why $\oplus 0$, and not $\oplus D$'' remark
above; for the spectral structure the
choice is immaterial, since $\norm{D} \le 1.95\cdot 10^{-4}$, but the
certified numbers are specifically for $\oplus 0$, and only that appears
in the rest of the text.)



\subsection{Chain and anchor}\label{sec:4.7}

Rouché gives \textbf{constancy of the count over the whole box}, equal to the
count at its center. It therefore suffices that neighboring boxes
\textbf{overlap}: at the common point $\ell^{\star}$ the count equals both the count at
the center of the left box and the count at the center of the right box,
hence they equal each other. Centers are spaced at a NOMINAL step of $0.8\cdot w$, that is, with a nominal overlap of $0.2\cdot w$ --- the neighbor's center then lies outside the
previous box, and that is fine: the requirement ``the center lies inside the
neighbor'' would be superfluous. The nominal figure is not the guarantee, though: the actual grid is
locally denser than nominal, and what is certified is the minimum over all
705 seams --- an overlap of $3.3\,\%$ of a width (\S\ref{sec:4.8}).

The anchor of the chain is $\ell = 3$. Here there are \textbf{two certificates,
independent in their method of counting the spectrum}, and this matters:
the Rouché chain propagates the number, but the number itself has to come
from somewhere.

\emph{The first} --- the pointwise certificate of \S\ref{sec:5.2}: spectral deflation, the
eigenpair verified by Krawczyk \cite{krawczyk1969newton}, $\rho(Q^{\star}) \le 0.9257 < 1$, the inflation $\varepsilon_{DT}$
covers a ball around $\Gs$.

\emph{The second, obtained by this work, is the \textbf{inertia count of the spectrum
relative to the unit circle} (a discrete analogue of Lyapunov's theorem;
Horn--Johnson \cite[Ch.~2]{horn1991topics}, in the line of
Stein, Taussky and Wimmer \cite{stein1952theorems, taussky1964matrices, wimmer1973ostrowski}).}
Two of those three are cited for the line, not for the formulation, and it is
worth saying which: Stein (1952) proves the \emph{definite} case (his Theorem~1:
$B^{n}\to 0$ iff some positive definite $H$ has $H - B^{*}HB \succ 0$) and does not
state an inertia count at all; Wimmer (1973) is the Ostrowski--Schneider theorem,
that is, inertia relative to the \emph{imaginary axis} --- the continuous variant.
The discrete statement used below is the one in Horn--Johnson. Let $A$ be the matrix $\DT|_{X_{0}}$
at $\ell = 3$, $H = H^{\mathsf{T}}$ (definiteness is not assumed). If $A^{\mathsf{T}}HA - H \prec 0$, then
$H$ is nonsingular and its inertia counts the spectrum: $n_{+}(H)$ is the
number of eigenvalues inside the unit disk, $n_{-}(H)$ --- outside. $H$ is
obtained in float by solving the Stein equation $H - A^{\mathsf{T}}HA = I$ and enters
the proof only as a set of numbers; the proof consists of two interval
checks:

\begin{gather*}
-S = H - A^{\mathsf{T}}HA \succ 0 :  \text{all 80 pivots of } LDL^{\mathsf{T}} \ge 0.9999,\ \text{radii } \le 4.76\cdot 10^{-72}, \\
\operatorname{In}(H) = (79, 1, 0) :  \text{one pivot } -0.4085,\ \text{the rest } \approx 1,\ \text{radii } \le 4.76\cdot 10^{-72}.
\end{gather*}

The condition ``no spectrum on the circle itself'' gives finiteness of the
resolvent, already computed rigorously: $\sup_{|z|=1}\|(A - z)^{-1}\|_{1} \le 3.279$
(\texttt{03-\allowbreak CERTIFICATES/\allowbreak anchor-\allowbreak inertia/\allowbreak s2\_\allowbreak resolvent.py};
the gate \texttt{R\_A} of \texttt{04-\allowbreak CODE/\allowbreak tests/\allowbreak \_\allowbreak numbers\_\allowbreak check.py}
holds it against the recorded \texttt{R\_\allowbreak rig}). All four figures of
this anchor --- the two pivot bounds, their common radius and this resolvent ---
are recorded in
\texttt{03-\allowbreak CERTIFICATES/\allowbreak anchor-\allowbreak inertia/\allowbreak REPORT-\allowbreak ANCHOR.md};
the result file of the interval run itself,
\texttt{\_\allowbreak s3s4\_\allowbreak result.json}, is written at run time and
deliberately not shipped (\texttt{00-\allowbreak README.md}).

Thus the count ``exactly one outside'' is \textbf{proved within this work},
without reference to the earlier certificate; the pointwise certificate of
\S\ref{sec:5.2} remains a second confirmation, independent in its counting method. A
caveat on independence: the methods themselves (deflation with Krawczyk
versus Stein--Taussky--Wimmer inertia) do not rely on each other in any way,
but both are applied to the same matrix $\DT|_{X_{0}}$, assembled by the same
code (\texttt{hyp\_\allowbreak rouche.dt\_\allowbreak from\_\allowbreak pkg}). The common point of failure is not the
spectrum count but the assembly of $\DT$; this is checked separately and
externally, by comparing $\delta(\ell)$ and $\alpha(\ell)$ against the
published Feigenbaum constants (\texttt{tests/\allowbreak \_\allowbreak classic\_\allowbreak constants.py}, \S\ref{sec:3.9}). The
two anchors cannot be called fully independent chains, and this is
accounted for below.
The assembly of $\DT$ itself is no longer unverified, however: the
acceptance step \texttt{tests/\allowbreak \_\allowbreak dt\_\allowbreak fd\_\allowbreak cross.py} builds it a second way --- by finite
differences of the operator's value, without the analytic derivative --- and
the two assemblies agree to a relative accuracy of $1.18\cdot 10^{-12}$ (a
float64/\texttt{arb} diagnostic against a fixed threshold, not an interval
certificate: the finite-difference truncation is not enclosed rigorously,
\S\ref{sec:3.9}). The
common point of failure that remains is the code for the operator's value
itself (\texttt{compute\_\allowbreak F\_\allowbreak poly}). From the anchor, the number propagates along the
chain in both directions.

A remark on the formulation. Initially an incorrect premise was written
down --- that the spectrum is counted by the inertia of $S = A^{\mathsf{T}}HA - H$
itself, for any $H \succ 0$. It was \textbf{refuted by a counterexample}
($A = [[\tfrac{1}{2}, 10], [0, \tfrac{1}{2}]]$, $H = I$: the spectrum lies entirely inside the
disk, yet the inertia of $S$ equals $(1,1,0)$), and no count was ever
performed with it. The correct theorem concerns the inertia of $H$ under a
definite $S$; everything here is built on that. Exactly the same mistake
had earlier gone unnoticed along a different route (\S\ref{sec:4.2}).

It is worth spelling out how the anchor connects to the box statement,
because the objects involved are different. Box $k$ proves the equality

\[
N(\DT(\Gs(\ell))) = N(A(P_{k}(\ell_{0}^{k})) \oplus 0) \qquad \text{for all } \ell \text{ in the box},
\]

where the right-hand side is the block-diagonal operator built from the
center of the box ($N$ is the number of eigenvalues outside the unit
circle, with multiplicities). The value of the right-hand side is itself
unknown --- it is precisely the constant that the chain propagates, and it
needs to be determined only once.

It is determined \textbf{directly from the right-hand side}, by the inertia
count above: the matrix $A$ there is built from the same interpolant as
the whole tube --- the branch of the box covering $\ell = 3$, the interpolant
at that point, then $\DT$. Two existence boxes cover $\ell = 3$; the loader
takes the one whose box holds the point with the larger margin
(\texttt{hyp\_\allowbreak rouche.load\_\allowbreak box}), which for both the
anchor and the tube box \texttt{t0215} is
\texttt{branch\_\allowbreak ganchor.npz}, so the two are the same interpolant
and not merely similar ones. The rigorous, interval part of the count
(LDL$^T$ in \texttt{arb}, the signs of all 80 pivots, published radii) is done by
\texttt{03-\allowbreak CERTIFICATES/\allowbreak anchor-\allowbreak inertia/\allowbreak s3\_\allowbreak prepare\_\allowbreak AB.py}; alongside it,
\texttt{03-\allowbreak CERTIFICATES/\allowbreak anchor-\allowbreak inertia/\allowbreak \_\allowbreak anchor\_\allowbreak check.py} is an INDEPENDENT float
cross-check of the same numbers (its docstring states as much: it does not
replace the interval part, but catches a gross error in the setup).
The inertia gives $\operatorname{In}(H) = (79, 1, 0)$, that is, the right-hand side equals
$1$; hence the left-hand side also equals $1$ on the entire anchor box
\texttt{t0215}, and from there, at every intersection of neighboring boxes, the
left-hand sides agree, and through them so do the right-hand sides. In this
way the number ``1'' reaches both ends of the interval.

This is worth spelling out, because via the LEFT-hand side the chain would
not close so simply. The pointwise certificate of \S\ref{sec:5.2} is built around the
numerical center $\Gbar$, while $\Gs(3)$ on the left is given by the uniqueness
of Theorem~A around $P(3)$; to cite \S\ref{sec:5.2} as a link, one needs
$\|\Gbar - P(3)\|_{X} < 1.75\cdot 10^{-7}$. That estimate \emph{is} in the package:
the certifier computes it and deposits it as
\texttt{box\_\allowbreak ganchor.json::\allowbreak containment}
($d_{X} = 2.7934\cdot 10^{-16}$ against $\rho_{X} = 1.75\cdot 10^{-7}$,
\texttt{rig2fix3.py:797--810}); no step of the acceptance used to read that
field, and \texttt{tests/\allowbreak \_\allowbreak numbers\_\allowbreak check.py}
(block K2) now does. Recomputed from the full-precision
seed of \texttt{hiprec\_seeds.json} the same distance is
$\|\Gbar - P(3)\|_{X} \le 4.02\cdot 10^{-16}$ --- eight orders below the
uniqueness radius; the deposited field measures it to the seed rounded to
double, and the polishing step performs zero iterations at $\ell = 3$, so the
two differ only by that rounding. Closing the chain via the right-hand side
does not depend on this question at all, and that is still how it is closed
here; but \S\ref{sec:5.2} may now be cited as a link and not only as a second
confirmation, independent in its counting method.



\subsection{Numbers and acceptance}\label{sec:4.8}

706 boxes, width from 0.0012 (upper end of the interval) to 0.007 (lower
end), $K = 4$ subdivisions in $\ell$, 180--240 nodes on the circle; total 8.20
core-hours, 41.8 s per box. (The original plan contained 696 boxes; one was
added at the junction of widths where two boxes touched exactly at a point,
and nine were replaced by eighteen narrower ones --- see \S\ref{sec:4.9}, defects 5 and
6.)

\begin{table}[h]
\centering
\caption{Range of tube quantities.}\label{tab:tube-quantities}
\begin{tabular}{lll}
\toprule
quantity & from & to \\
\midrule
$R$ (rigorous resolvent) & 2.9884 & 3.7961 \\
$\|\Delta\|$ (rigorous, per box) & $1.312\cdot 10^{-1}$ & $2.798\cdot 10^{-1}$ \\
$R\cdot \|\Delta\|$ & 0.4269 & 0.9896 \\
\texttt{extra} (offset $P \to \Gs$) & $1.4346\cdot 10^{-8}$ & $2.86\cdot 10^{-6}$ \\
$\|C\|$ (head $\to$ tail, honest) & $3.08\cdot 10^{-13}$ & $4.07\cdot 10^{-4}$ \\
$\|[B;D]\|$ (tail columns) & $1.288\cdot 10^{-7}$ & $1.95\cdot 10^{-4}$ \\
\bottomrule
\end{tabular}
\end{table}

Acceptance --- six conditions, not just the gate:

\begin{table}[h]
\centering
\caption{Acceptance checks.}\label{tab:acceptance}
{\small\setlength{\tabcolsep}{4pt}
\begin{tabular}{>{\raggedright\arraybackslash}p{0.395\textwidth}>{\raggedright\arraybackslash}p{0.529\textwidth}}
\toprule
check & result \\
\midrule
Rouché gate on every box & 706 of 706 \\
full test $R_{\mathrm{tot}}\cdot \max(\|\Delta\| + \mathrm{extra} + 2\|C\|, \|[B;D]\| + \mathrm{extra}) < 1$ ($R_{\mathrm{tot}} := \max(R, 1/(1-\|[B;D]\|-\mathrm{extra}))$, defined in \texttt{04-\allowbreak CODE/\allowbreak tube\_\allowbreak accept.py}) & 706 of 706, worst \textbf{0.9919} (box \texttt{t0464}, $\ell_{0} = 3.6976$; the full set of inputs is \texttt{tube/\allowbreak box\_\allowbreak t0464.json} and \texttt{tube/\allowbreak offsets.json::\allowbreak g37}, the number reproduces by running \texttt{tube\_\allowbreak accept.py}, which prints the exact $0.991843707457$ together with its short form rounded UP --- $0.9919$, the same direction as here, as befits a majorant of a quantity that must stay below $1$; the tool marks the direction in the line it prints, and the check of that direction is part of the acceptance) \\
every tube box is contained in its existence box & 706 of 706 (\texttt{t0695} lies wholly inside \texttt{gtop} $= [3.9775, 4.0075]$ with slack $6.3\cdot 10^{-3}$; \texttt{t0000} is nested after intersecting with $[2,4]$, where its slack is $+9.14\cdot 10^{-17}$; before that intersection it overhangs the left edge by $-6.87\cdot 10^{-17}$. Both numbers are computed from the certified ends $\ell_0 \mp w/2$ in exact rational arithmetic, not from their float printout --- the printout of the left end of \texttt{g00} is exactly $2.0$, which made the comparison establish nothing) \\
chain: neighboring boxes overlap & no gaps, minimum overlap 3.3\% of width \\
anchor $\ell=3$ covered by a box with a known count & yes (\texttt{t0215}) \\
coverage & $[2.0000, 4.0012] \supseteq [2,4]$ \\
\bottomrule
\end{tabular}}
\end{table}
\FloatBarrier

The coverage of the two certificates does not coincide: existence is
closed on $[2.0000, 4.0075]$, hyperbolicity on $[2.0000, 4.0012]$. The joint
statement of both theorems is formulated on the intersection, that is, on
$[2, 4.0012] \supset [2, 4]$; for simplicity of the wording, $[2, 4]$ is written
throughout the text. The point $\ell = 4$ itself is interior to the preceding box, which overlaps
\texttt{t0695} (minimum overlap 3.3\,\% of a width,
Table~\ref{tab:acceptance}) and for which
\texttt{extra} is checked, so no part of the statement about $[2,4]$ rests on
the degenerate remainder. The right endpoint 4.0012 of the covering rests on \texttt{t0695}, whose intersection with $[2,4]$ is a point; \texttt{extra} is justified on it all the same, because $[4.0000, 4.0012] \subset \texttt{gtop} = [3.9775, 4.0075]$, and the containment check covers it (see Table~\ref{tab:acceptance} and a7 in Table~\ref{tab:certified-premises}).

\textbf{The anchor box.} \texttt{t0215}: $\ell_{0} = 3.00000$, $w = 0.004$, $R = 3.3266$,
$\|\Delta\| = 1.9525\cdot 10^{-1}$, $R\cdot \|\Delta\| = 0.6496$; the count at its center in float
gives $\delta = 6.084689074$ and $\lambda_{2} = 0.237146112$ --- the same digits that are
certified by an independent method: $\delta$ in \S\ref{sec:5.2}, $\lambda_{2} \in
[0.23714611165001134 \pm 1.6\cdot 10^{-18}]$ in \S\ref{sec:5.3.5} (deflation + mean-value-form
Krawczyk). It is through this box that the statement ``exactly one outside''
enters the chain.

For reference (the count at the centers in float, \textbf{not part of the
proof}): $\delta(\ell)$ ranges over $[4.6747, 7.2854]$, $\lambda_{2}(\ell)$ over $[0.1599,
0.2919]$, and at each of the 706 centers there is exactly one value outside
the unit circle --- the same as given by the rigorous count.

\textbf{Simplicity and reality of $\delta(\ell)$ come for free.} The matrix is real: a
non-real eigenvalue would come together with its conjugate and give a
count $\ge 2$ outside, and a repeated one likewise $\ge 2$ (the count is taken
with multiplicities). Since there is exactly one outside, it is real and
simple.



\subsection{Six defects found while assembling this certificate}\label{sec:4.9}

All six are the work's own, and all six were found not by reading the
code:

\begin{table}[h]
\centering
\caption{Six defects found while assembling this certificate.}\label{tab:defects}
\begin{tabular}{p{0.04\textwidth}p{0.55\textwidth}p{0.35\textwidth}}
\toprule
\# & defect & caught by \\
\midrule
1 & maximum over arb-intervals via \texttt{>}: an interval containing zero is ``not greater'' than zero, so the maximum silently stays zero $\Rightarrow \|\Delta\| = 0$ and a false PASS & a control run gave exactly zero --- a value that is physically impossible \\
2 & $\Delta$ was built at a fixed center: the motion of the fixed point along the branch was not covered by anything & comparing the formula against what Rouché's theorem actually requires \\
3 & the resolvent was computed with \texttt{numpy.linalg.solve} in \texttt{float64} & the question ``is this an enclosure or an approximation?'' \\
4 & the chain criterion ``the neighbor's center lies inside the previous box'' was superfluous and was not satisfied & the acceptance run showed 393 ``gaps'' at step $0.8\cdot w$; analysis showed that overlap suffices \\
5 & at the junction of different widths, two boxes touched at exactly one point (\texttt{hi = lo} in float) & acceptance \\
6 & nine tube boxes extended beyond their existence box (by up to $9\cdot 10^{-4}$), while the term \texttt{extra} is proved only inside it & a separate containment check (\texttt{\_\allowbreak containment.py}) \\
\bottomrule
\end{tabular}
\end{table}
\FloatBarrier

Defect 6 deserves an explanation: existence boxes of width 0.05 overlap by
only 0.005, and a tube box of width 0.007 falling into the overlap zone did
not fit entirely into \textbf{either} of them. On the protruding piece, the
quantity \texttt{extra} (and with it the passage to the true fixed point) was
unsupported by anything. The fix: nine such boxes were replaced by
eighteen --- each entirely inside its own existence box, overlapping with
one another; the choice of branch is now made by containment, not by the
center (\texttt{hyp\_\allowbreak rouche.load\_\allowbreak box(span=\ldots)}). A containment check was added to
the acceptance run.

Defect 1 is the most telling: it produced not a failure but a \textbf{success}.
A certificate whose checked quantity vanishes passes any gate; the only
way to catch it is by comparison against the expected order of magnitude.



\subsection{What the gate does not catch}\label{sec:4.10}

It is useful to know not only what the check catches, but also what it
lets through. Five deliberate defects were introduced into the code to see
which ones defeat the gate (box g01):

\begin{table}[h]
\centering
\caption{Gate behaviour under five deliberate defects.}\label{tab:gate-defects}
\begin{tabular}{p{0.4\textwidth}p{0.1\textwidth}p{0.4\textwidth}}
\toprule
defect & gate & why \\
\midrule
center shifted by $10^{-6}$ only in the model (not consistently) & \textbf{FAIL} & $d_{\mathrm{chk}}$ (the mechanism at work --- \S\ref{sec:3.7}) catches it: $4.1\cdot 10^{-6}$ against the threshold $10^{-6}$ \\
$\widetilde{A} := 2I$ & \textbf{FAIL} & $Z_1 = 11.4$, $\widehat{Z} \ge 1$ for any $\rho'$ \\
center shifted consistently (both in the model and in the engine) & PASS & the certificate remains valid, but the radius degrades: the gate passes at $\rho' = 10^{-6}$ instead of $10^{-7}$ \\
$L$ underestimated threefold & PASS & the gate is monotone in $L$, so lowering it can only make the gate easier; the Lipschitz term $1.75^{2}B_{A}L\rho'$ contributes $1.98\%$ of $\widehat{Z}$ on \texttt{g01} and between $1.92\%$ (\texttt{g00}) and $34.3\%$ (\texttt{gtop}) over the 46 boxes, median $3.36\%$ (computed with the gate's own $L$; the earlier figure $\sim 1.6\%$ was computed with $L_{\mathrm{mid}}$) \\
\texttt{extra}, $\|C\|$, $\|[B;D]\|$ zeroed out & PASS & the full test is monotone in them \\
\bottomrule
\end{tabular}
\end{table}
\FloatBarrier

These five runs were carried out as a one-off, ad hoc experiment outside
the package: there is no script in \texttt{04-CODE} that reproduces exactly these
five mutations (it is not \texttt{layer\_\allowbreak mutation1.py}/\texttt{layer\_\allowbreak mutation2.py} --- those
concern a different certificate, \S\ref{sec:10} --- nor \texttt{\_\allowbreak dchk\_\allowbreak bound\_\allowbreak test.py} --- that one
concerns the derivation of the bound $d_{\mathrm{chk}}$ itself). The results in the
table should be taken as the author's statement about a one-off run, not
as a certificate reproducible from the package.

The pattern is simple and structural: \textbf{the gate catches what makes the
bound larger, and does not catch what makes it smaller.} Any certificate
of the form ``chain of inequalities $< 1$'' is built this way. Two consequences
follow from this, worth keeping in mind when reading work of this kind:

\begin{enumerate}
\item[1.] The rigor of each quantity ($L$, \texttt{extra}, $\tau$, the center) rests not on
  the gate but on the fact that the quantity was computed by an
  independent rigorous path. This is exactly why so much attention is
  paid in this work to the question ``does this formula bound the
  function, or its model?'' --- the gate does not answer that. What the
  numbers gate does establish is that the shipped artefact is the output
  of that path and not a desynchronised copy of it: the nine fields of
  \texttt{offsets.json} that are copies of \texttt{box\_g*.json} are compared with
  it for exact equality, and \texttt{extra}, $\|C\|$, $\|[B;D]\|$ and
  $\|M'\|_{\mathrm{full}}$ are re-derived from the certificate's own fields and
  required to be no larger than what is shipped. The expensive head norm
  $\|M'\|$ is not re-derived and is taken as shipped.
\item[2.] The only built-in cross-check --- $d_{\mathrm{chk}}$ (the mechanism at work is a
  derived majorant of the conversion error, \S\ref{sec:3.6}/\S\ref{sec:3.7}, not a direct
  comparison) --- catches only a \textbf{mismatch} between the model and the
  engine. It does not see an error that is consistent between the two
  (e.g., a wrong center in both); that is caught by the degradation of
  $\rho'$, which is recorded in the certificate of every box and visible in
  the table.
\end{enumerate}



\subsection{What the tube certificate does not prove}\label{sec:4.11}
\begin{itemize}
\item \textbf{The working layer is fixed} ($c = 3/5$, $r = 4/5$): the statement concerns the
  operator on that space of analytic functions, not on an arbitrary class. The
  result does not depend on the choice of $(c, r)$ within the presented family
  $\Ffam$, on all of which the transfer $P \to \Gs$ of \S\ref{sec:10.7} is rigorously
  certified (Theorem~D, \S\ref{sec:10}; how far $\Ffamz$ itself extends --- \S\ref{sec:10.9}); an
  arbitrary class of spaces remains uncovered. The projection $N_{p} = 80$ is
  \textbf{not} a limitation of the formulation: the tail is accounted for explicitly (\S\ref{sec:4.6}).
\item Smoothness of $\delta(\ell)$ in $\ell$ is \textbf{not} proved at this stage --- only that for each
  $\ell$ there is exactly one value outside the unit disk. (Established by
  Theorem~C, stated in \S\ref{sec:6}: real analyticity of $\delta(\ell)$ on $[2,4]$ --- rigorously stronger than
  smoothness; see also \S\ref{sec:6.10} item 1.)
\item $\lambda_{2}(\ell)$ is not certified: the numbers given are a count in float.
\item Nothing is asserted outside $[2,4]$.
\item The words ``on the whole interval'' in the heading are an abbreviation. The
  union of the boxes covers $[2.0000, 4.0012]$, i.e. $[2,4]$ entirely and with
  margin beyond the right endpoint (\S\ref{sec:4.8}); inside $[2,4]$ there are no gaps.
  About the margin itself, $(4, 4.0012]$, nothing is asserted separately: it is a
  by-product of the box grid, not a result, and its right endpoint is not
  distinguished in any way.
\end{itemize}

\emph{Code, certificates and protocols of this section --- \texttt{04-CODE/}, \texttt{03-\allowbreak CERTIFICATES/\allowbreak tube/\allowbreak }, index in \texttt{05-\allowbreak DATA-\allowbreak INDEX.md}. Certificates: \texttt{box\_t*.json}, \texttt{offsets.json}. Acceptance log: \texttt{\_accept.out}.}



\section{Pointwise anchor certificates}
\label{sec:5}

The chain of boxes in \S\ref{sec:4} transfers the number of eigenvalues outside the unit
disk, but the number itself must be proved somewhere. Below are the pointwise
certificates at $\ell = 3$ and at $\ell = 5/2$ (the enclosure differs: the anchor of the
CHAIN is only $\ell = 3$,
see \S\ref{sec:4.7}; the certificate at
$\ell = 5/2$ is not part of the count transfer along the chain, it is a separate,
independent pointwise verification). The pointwise certificate at $\ell = 3$ (below)
is duplicated by an independent proof in \S\ref{sec:4.7} (inertia count, reference \cite{stein1952theorems, taussky1964matrices, wimmer1973ostrowski, horn1991topics}),
so that the chain is closed within the work; the results given here remain a
second, independent confirmation. The word ``independent'' here means
independence in the METHOD OF SPECTRAL COUNTING (Krawczyk deflation vs.
Stein--Taussky--Wimmer inertia), not along the whole chain: both counts are
applied to the same matrix $DT|_{X_{0}}$. That assembly itself is verified
separately and by a second route --- by finite differences of the operator
value, without an analytic derivative (\S\ref{sec:3.9}, acceptance step
\texttt{tests/\allowbreak \_\allowbreak dt\_\allowbreak fd\_\allowbreak cross.py}); the code of the operator value itself is shared
(detailed caveat --- \S\ref{sec:4.7}).


\subsection{Pointwise existence certificate at $\ell = 3$}
\label{sec:5.1}

\subsubsection{Statement}
\label{sec:5.1.1}

\begin{theorem}[Hyp-5-W5-L3]
At $\ell=3$, $c=3/5$, $r=4/5$, $N_p=80$, and $\Gbar$ --- an analytic
function on $D(c,r)$ (degree 64 in $u=(t-c)/r$; the coefficients of
\texttt{04-\allowbreak CODE/\allowbreak hiprec\_\allowbreak seeds.json} are given in $t$ from center 0 --- the shift
$t=c+r\cdot u$ is performed by the verifier code; that file holds four points,
and this is its entry at $\ell=3$),
with $\hat{\rho}' = 1\cdot 10^{-20}$ there exists a \textbf{unique} function
$\Gs \in \Aff$ ($\Gs(0)=1$) in the ball
$\norm{G-\Gbar}_X \le 1.75\cdot \hat{\rho}' = 1.75\cdot 10^{-20}$, analytic on
$D(c,r) \supset [0,1]$, satisfying $F(\Gs) = 0$, i.e.
$\Gs + \alpha_{\Gs}\cdot \Gs \circ (U^{\star})^{3} = 0$, where
$U^{\star} = \Gs \circ s_{\Gs}$ and $s_{\Gs}(t) = t\cdot \alpha_{\Gs}^{-3}$ (notation \S\ref{sec:1}),
$\alpha_G = -1/G(1)$, with constants
(prec=512):
\end{theorem}

\begin{table}[htbp]
\centering
\small
\caption{Rigorous constants of the existence certificate at $\ell=3$.}
\label{tab:5.1.1-constants}
\begin{tabular}{p{0.23\textwidth} p{0.47\textwidth} p{0.24\textwidth}}
\toprule
Constant & Value & Source file \\
\midrule
$Y_0 = \norm{\Ahat\cdot \bar{F}}_1$ (head) + tails &
1.138567618961918\ldots$\times 10^{-42}$ $\pm$ 3.18e-150 &
\texttt{03-\allowbreak CERTIFICATES/\allowbreak pointwise/\allowbreak kant\_\allowbreak l3.json\allowbreak ::\allowbreak Y0} \\
$Z_1 = Z_{1a}+Z_{1b}+Z_{1c}$ &
4.4786e-15 + 2.3119e-11 + 5.1978e-8 = \textbf{5.2002e-8} (rigorous formulaic
$\hat{\beta}$, rounded outward; exactly 5.20010896072677\ldots$\times 10^{-8}$) &
\texttt{03-\allowbreak CERTIFICATES/\allowbreak pointwise/\allowbreak Z1c\_\allowbreak l3\_\allowbreak strict.json\allowbreak ::\allowbreak Z1\_\allowbreak strict} \\
$\Zhat = Z_1 + 1.75^{2}\cdot \norm{\Ahat}_1\cdot L\cdot \hat{\rho}'$ &
5.2002e-8 $\ll$ 1 (rounded outward; second term $\approx$ 2$\cdot 10^{-16}$) &
\texttt{Z1c\_\allowbreak l3\_\allowbreak strict.json} + \texttt{fixL\_\allowbreak result.json\allowbreak ::\allowbreak Z\_\allowbreak hat\_\allowbreak new\_\allowbreak strict} \\
$L$ (Lipschitz of $DT$) & $\mathbf{2.0648726\times 10^3}$ &
\texttt{03-\allowbreak CERTIFICATES/\allowbreak pointwise/\allowbreak fixL\_\allowbreak result.json\allowbreak ::\allowbreak L\_\allowbreak strict} \\
$\rho^{\star} = Y_0 / (1-\Zhat)$ &
1.1385678 $\times 10^{-42}$ (the field \texttt{rho\_\allowbreak star\_\allowbreak new} is computed from
the historical $\Zhat = 1.4742\cdot 10^{-7}$, not from the rigorous $5.2002\cdot 10^{-8}$ (rounded
outward) given above in this table --- the direction is conservative: with the
rigorous $\Zhat$ it would be slightly smaller, $1.1385677\times 10^{-42}$; no
inequality is affected, the uniqueness ball is 22 orders of magnitude larger) &
\texttt{03-\allowbreak CERTIFICATES/\allowbreak pointwise/\allowbreak fixL\_\allowbreak result.json\allowbreak ::\allowbreak rho\_\allowbreak star\_\allowbreak new} \\
uniqueness ball &
$\norm{G-\Gbar}_X \le 1.75\times 10^{-20}$ &
\texttt{03-\allowbreak CERTIFICATES/\allowbreak pointwise/\allowbreak kant\_\allowbreak l3.json\allowbreak ::\allowbreak rho\_\allowbreak X} \\
\bottomrule
\end{tabular}
\end{table}

\subsubsection{Proof outline}
\label{sec:5.1.2}

\begin{enumerate}
\item \textbf{$X_{0}$-slice.} $F(G)(0) = 0$ is identically true on $\{G(0)=1\}$
(hence $T$ maps the slice $\{G(0)=1\}$ into itself: $T(G)(0)=1$), so
$F(G) \in X_{0}$; $DF|_{X_{0}}$ is invertible (see
\texttt{03-\allowbreak CERTIFICATES/\allowbreak pointwise/\allowbreak kant\_\allowbreak l3.json\allowbreak ::\allowbreak gate2\_\allowbreak pass}).
\item \textbf{$\Ahat$ --- the inverse matrix} in the $b$-basis
$\{b_k = e_k - (-c/r)^k\cdot e_0\}$, $k\ge 1$; the matrix $M' = DF|_{X_{0}}$
($N_p\times N_p=80\times 80$), $\norm{\Ahat}_1 = 3.140$
(\texttt{03-\allowbreak CERTIFICATES/\allowbreak pointwise/\allowbreak kant\_\allowbreak l3.json\allowbreak ::\allowbreak A\_\allowbreak norm1}).
\item \textbf{$Y_0 = \Ahat\cdot \bar{F}$.} Rigorous ball at arb prec=512:
$Y_0 = 1.139\times 10^{-42}$.
\item \textbf{$Z_1 = Z_{1a}+Z_{1b}+Z_{1c}$} from 65-coefficient polynomials:
4.4786e-15, 2.3119e-11, 5.1978e-8, \textbf{$Z_1 =$ 5.2002e-8} $\ll 1$ (rounded
outward; gate 3 PASS). The value is the rigorous formulaic $\hat{\beta}$; with
the previous float fit $\beta_k \approx C\cdot \kappa^k$ the same $Z_1$
equalled 1.4742e-7 --- the estimate was correct, but 2.8 times looser.
\item \textbf{L --- Lipschitz of DT.} In
\texttt{03-\allowbreak CERTIFICATES/\allowbreak pointwise/\allowbreak fixL\_\allowbreak result.json} a step4 bug was found and
\textbf{fixed} ($\norm{\bar{U}}$ was taken as $\sqrt{\norm{\bar{U}^{2}}_X}$ ---
that is a lower bound; the correct value is $\sum |\mathrm{coef}|\cdot r^k$). L is
derived for the first time from an honest chain: $D_{\mathrm{total}} = T_2 +
T_3 + \Delta T_4 + \Delta T_5 = 3.6136\times 10^{-17}$,
$L = D_{\mathrm{total}}/\hat{\rho}_X = \textbf{2.065}\times \textbf{10}^{\textbf{3}}$.
(Equality holds exactly on non-rounded interval quantities. The sum of the
PRINTED terms, each of which is rounded outward, yields $\approx 3.629\times
10^{-17}$ --- 0.43\,\% larger, i.e. on the conservative side: the $L$ assembled
this way is only rigorously tighter.) Closed-form telescoping gives the values
of $\Delta T_4$, $\Delta T_5$ (\texttt{fixL\_\allowbreak result.json\allowbreak ::\allowbreak Delta\_\allowbreak T4},
\texttt{fixL\_\allowbreak result.json\allowbreak ::\allowbreak Delta\_\allowbreak T5}); there is no independent cross-check of
these quantities in the package.
\item \textbf{Banach theorem.} With the derived L we verify rigorously:
$\Zhat := Z_1 + 1.75^{2}\cdot \norm{\Ahat}_1\cdot L\cdot \hat{\rho}' =$ 5.2002e-8
(rounded outward) $< 1$ and $Y_0 + \Zhat\cdot \hat{\rho}' \le \hat{\rho}'$; the
operator $\tilde{T}(h) = h - \Ahat\cdot \tilde{F}(\Gbar+h)$ is a contraction of
the closed ball of radius $\hat{\rho}'$, its fixed point $h^{\star}$ is unique,
$\Ahat$ is invertible (from $\norm{I-\Ahat M'}<1$ on the block and $\Ahat=I$ on
the tail) $\Rightarrow$ $F(\Gs)=0$, $\Gs = \Gbar + h^{\star}$ is unique in the
ball; $\norm{\Gs-\Gbar}_X \le 1.75\cdot \rho^{\star} \le 1.993\times 10^{-42}$
(exactly: 1.9924936269\ldots$\times 10^{-42}$).
\end{enumerate}

\textbf{K-mode.} Exactly one \textbf{number} is certified here:
$\lambda_K := 3\cdot \Gbar'(1)/\Gbar(1) = 3\cdot 1.238664150\ldots
/0.518755349\ldots = 7.1632851\ldots$
(both values are negative, the ratio is positive; the balls lie in
\texttt{03-\allowbreak CERTIFICATES/\allowbreak pointwise/\allowbreak kmode\_\allowbreak l3.json}, lines 59--61 of Appendix~\ref{sec:A}).
The identification of $K = \Gbar - 3t\Gbar'$ with the RIGHT eigenvector of
$DT(\Gbar)$, corresponding to $\lambda_K$, is \textbf{not} certified and is not
part of the proof of the certificates of \S\ref{sec:5.1}, \S\ref{sec:5.2} and \S\ref{sec:5.3.3}. The residuals $\norm{DT(K) -
\lambda_K\cdot K}$ do not exist in the package --- and cannot, in the available
basis: $K(0) = \Gbar(0) = 1 \ne 0$, i.e. $K \notin X_{0}$, whereas the matrix
$DT$ is built precisely on the slice $X_{0}$; \texttt{04-\allowbreak CODE/\allowbreak tests/\allowbreak \_\allowbreak kmode.py}
measures both.

But there is no need to verify the right eigenvector: that $\lambda_K$ is an
eigenvalue of $DT$ on $X$ follows in two lines via the LEFT eigenvector, and
the derivation does not rely on any artefact. Let $\varphi(h) := h(0)$ --- a
bounded functional on $X$. For $G$ with $G(0) = 1$ we have $U_G(0) = G(0)$,
$W_G(0) = G(0)^\ell$, hence $\varphi(T(G)) = -\alpha_G\cdot G(G(0)^\ell)$;
differentiating in the direction $h$ and substituting
$d\alpha_G(h) = h(1)/G(1)^{2}$,

\begin{align*}
\varphi(DT(G)h) &= -\bigl(h(1)/G(1)^{2}\bigr)\cdot G(1) - \alpha_G\cdot
\bigl( h(1) + \ell\cdot G'(1)\cdot h(0) \bigr) \\
&= \bigl(\ell\cdot G'(1)/G(1)\bigr)\cdot h(0) = \lambda_K\cdot \varphi(h) ,
\end{align*}

where the first two terms cancel exactly because $\alpha_G = -1/G(1)$. Hence
$\varphi$ is a left eigenvector of $DT(G)$ with eigenvalue $\lambda_K$, for ANY
$G$ with $G(0) = 1$. From this immediately: $X_{0} = \ker \varphi$ is invariant
(this is what ``$T$ preserves the normalization'' means), in the decomposition
$X = X_{0} \oplus \operatorname{span}(1)$ the operator is block-triangular, and
$\spec_X(DT) = \spec(DT|_{X_{0}}) \cup \{\lambda_K\}$ --- exactly, not a
``structural observation''. The right eigenvector exists and is unique up to a
scalar factor, since $\lambda_K \notin \spec(DT|_{X_{0}})$ --- and this has been
measured (float64, not an interval certificate; \S\ref{sec:9}). For the outline just given it is necessary and sufficient that the
contribution of this direction to $DF$ is closed by $\Delta T_5 \le
1.93\times 10^{-17}$.

\begin{remark}[A claim about an eigenvector, presented as a fact]
An earlier version of the text here said ``the eigenvalue of $DT(\Gbar)$
\textbf{on} $K = \Gbar - 3t\Gbar'$ is $3|\Gbar'(1)|/|\Gbar(1)|$'' --- a claim
about an eigenvector, presented as a fact, in a section carrying the proof. It
could not be verified: the certificate contains only $\lambda_K$, $\Gbar(1)$
and $\Gbar'(1)$. The same class of defect as the one corrected for $\lambda_K$
itself, just below (substitution of $\norm{\Gbar'}_X$ in place of $\Gbar'(1)$): the printed
result was correct, but the justification underneath it was absent.
\end{remark}

In an earlier version the numerator here read $2.1085438\ldots$, i.e. $\norm{\Gbar'}_X$
--- the norm on the whole X-disk, not the value at the point $t = 1$. The
printed result of 7.16329 was correct; only the substitution was wrong: there
was no way to verify it, because $\Gbar'(1)$ was nowhere recorded. Now it is
--- this is what closes the return of the error.

\textbf{Remark on the ``honesty'' of L.} In the first version
$L = 1.04\times 10^{18}$ was chosen CIRCULARLY as $L_{\max}/10$ (from the
condition $\Zhat<1$) --- suspicious. Here L is \textbf{rigorously derived};
the margin on the gates $\Zhat<1$ is preserved enormous ($\approx$ 7 orders of
magnitude).


\subsection{Pointwise hyperbolicity certificate at $\ell = 3$}
\label{sec:5.2}

\subsubsection{Statement}
\label{sec:5.2.1}

\begin{theorem}[Hyp-5-W6-L3]
$DT(\Gs)|_{X_{0}}$ has the spectral decomposition
\begin{multline*}
DT(\Gs)|_{X_{0}} = {}\\
(\delta^{\star} / \langle \varphi^{\star}, v\rangle) \cdot v
\otimes \varphi^{\star} + Q^{\star} ,
\end{multline*}
where $\delta^{\star} = 6.084689074262122$ (Krawczyk box width
$6.226\times 10^{-19}$,
\texttt{03-\allowbreak CERTIFICATES/\allowbreak pointwise/\allowbreak task\_\allowbreak A\_\allowbreak krawzyk\_\allowbreak result.json\allowbreak ::\allowbreak K\_\allowbreak delta\_\allowbreak width})
is a simple eigenvalue, $v \in X$ ($\norm{v}_X \le 1.695$), $\varphi^{\star} \in
X^{*}$ ($\norm{\varphi^{\star}}_\infty \le 1.000000006$ --- the normalization
$\norm{\widehat{\varphi}}_\infty = 1$ is imposed on the approximation, and the
exact $\varphi^{\star}$ is at distance at most $5.782\times 10^{-9}$ from it,
\texttt{q\_\allowbreak perturb\_\allowbreak patch.json::\allowbreak caps.phi\_\allowbreak star}), $\langle \varphi^{\star},
v\rangle \in [1.085888346, 1.085888367]$ (min $\ge 1.085 > 0.5$), and
$Q^{\star} := DT(\Gs)|_{X_{0}} - (\delta^{\star}/\langle \varphi^{\star},v\rangle)
\cdot v\otimes \varphi^{\star}$ kills the unstable direction: $Q^{\star} v = 0$
(exact spectral deflation).
$\rho(Q^{\star}) \le \norm{(Q^{\star})^{16}}_1^{1/16} \le (0.2905947)^{1/16} \le
0.92567 < 1$ (patch \S\ref{sec:5.2.2}). Since $(\delta^{\star}, \varphi^{\star})$ is the left
eigenpair ($\varphi^{\star}\circ DT(\Gs)|_{X_{0}} = \delta^{\star}\cdot
\varphi^{\star}$, certified in \texttt{eigen\_\allowbreak l3\_\allowbreak fixed.json\allowbreak ::\allowbreak gates.\allowbreak gate\_\allowbreak 2\_\allowbreak B\_\allowbreak inclusion}
(80/80) and \texttt{task\_\allowbreak B\_\allowbreak krawzyk\_\allowbreak result.json}), the invariant complement
$Y^{\star} := \ker \varphi^{\star}$ --- a subspace of codimension 1 --- is an
invariant subspace on which $DT(\Gs)$ is rigorously hyperbolic.
\end{theorem}

\subsubsection{Sources and logic}
\label{sec:5.2.2}

\begin{table}[htbp]
\centering
\small
\caption{Certificate sources for the spectral decomposition of the hyperbolicity certificate at $\ell=3$.}
\label{tab:5.2.2-sources}
\begin{tabular}{p{0.30\textwidth} p{0.35\textwidth} p{0.30\textwidth}}
\toprule
Step & File & Key item \\
\midrule
Certification of $(\delta^{\star}, v, \varphi^{\star})$ &
\texttt{03-\allowbreak CERTIFICATES/\allowbreak pointwise/\allowbreak eigen\_\allowbreak l3\_\allowbreak fixed.json} &
5/5 gates PASS (Moore Krawczyk) \\
Uniform certification over the ball &
\texttt{03-\allowbreak CERTIFICATES/\allowbreak pointwise/\allowbreak hyperb\_\allowbreak l3.json} &
Gates 1: uniform Krawczyk A, B (80/80, \texttt{K\_max\_width} $\approx$ 4e-16) \\
Spectral deflation, $\norm{\hat{Q}}_1$ ($\hat{Q}$ is the certified computable
analogue of $Q^{\star}$, $E=\norm{Q^{\star}-\hat{Q}}_1$ --- Appendix~\ref{sec:A}) &
\texttt{03-\allowbreak CERTIFICATES/\allowbreak pointwise/\allowbreak hyperb\_\allowbreak l3.json} &
\texttt{majority\_\allowbreak computation.\allowbreak defl\_\allowbreak coef\_\allowbreak value} = 5.60342044786\ldots \\
$\norm{Q^n}_1$, $n\in \{1,2,4,8,16,32,64\}$ (by doubling) &
\texttt{03-\allowbreak CERTIFICATES/\allowbreak pointwise/\allowbreak norms\_\allowbreak q\_\allowbreak n\_\allowbreak rigorous\_\allowbreak l3.json} &
$\norm{\hat{Q}^{16}}_1 \le$ 2.8995e-01 (rigorous arb ball) \\
Spectral gap (transfer patch E) &
\texttt{03-\allowbreak CERTIFICATES/\allowbreak pointwise/\allowbreak q\_\allowbreak perturb\_\allowbreak patch.json} &
$E \le 2.682\cdot 10^{-7} \Rightarrow \norm{(Q^{\star})^{16}}_1 \le 0.2906$ \\
\bottomrule
\end{tabular}
\end{table}

\textbf{On the two endpoints of $\langle \varphi^{\star},v\rangle$.} Both
appear in the package: the deflation certifier divides $\delta^{\star}$ by the
UPPER endpoint (\texttt{hyperb\_\allowbreak l3.json\allowbreak ::\allowbreak majority\_\allowbreak computation.\allowbreak defl\_\allowbreak coef\_\allowbreak value =
5.603420447862\ldots}, line 31 of Appendix~\ref{sec:A} --- this is the lower endpoint of
the ratio), while the transfer patch divides by the LOWER
(\texttt{q\_\allowbreak perturb\_\allowbreak patch.json\allowbreak ::\allowbreak caps.delta\_\allowbreak hat\_\allowbreak over\_\allowbreak ip\_\allowbreak hat =
5.603420547256\ldots}, the upper endpoint of the ratio). There is no
discrepancy here: the spread between the endpoints, $9.9394\times 10^{-8}$, is
not thrown away but enters $E$ as a separate term
$\mathtt{caps.term\_a} = (\mathrm{spread})\cdot \norm{v^{\star}}\cdot
\norm{\varphi^{\star}} = 1.6844\times 10^{-7}$. Earlier this was not
stated anywhere, and the two different numbers $5.6034204\ldots$ and
$5.6034205\ldots$ looked like a typo of each other.

The operator $DF|_{\Gs}$ depends analytically on $G$ with Lipschitz constant
$L = 2.065\times 10^{3}$ and with $\eps_{\mathrm{DT}} = L\cdot 1.75\cdot \hat{\rho}' = L\cdot\hat{\rho}_X = 3.6136\times 10^{-17}$
(\texttt{03-\allowbreak CERTIFICATES/\allowbreak pointwise/\allowbreak hyperb\_\allowbreak l3.json\allowbreak ::\allowbreak majority\_\allowbreak computation.\allowbreak eps\_\allowbreak DT}). The majorant
matrix $M$ of size $81\times 81$ (head block + 1 super-coord, $\pm
\eps_{\mathrm{DT}}$ inflation) is built in
\texttt{04-\allowbreak CODE/\allowbreak build\_\allowbreak majorant\_\allowbreak rigorous.py} (this script is included in the
package as the ORIGIN of the constants, not as a verification step: its inputs
point to an unsaved working tree, and it is not rerun from the package --- the
same is said in \texttt{05-\allowbreak DATA-\allowbreak INDEX.md}; what is verified here are the
values, not their recomputation); the sequence $\norm{\hat{Q}^n}_1$ for
$n\in \{1,2,4,8,16,32,64\}$ is computed in arb by matrix powering and saved as
\textbf{rigorous upper bounds}.

One step in that inflation deserves to be written out, because the two norms
involved are not the same. $\eps_{\mathrm{DT}}$ bounds
$\norm{\Delta\, \DT}_{X\to X}$, whereas $M$ is inflated \textbf{entry by entry} in
the b-basis, and $\norm{\hat Q^n}_1 \le \norm{M^n}_1$ rests on entrywise
domination, not on column sums. The entrywise bound costs one more factor of
$1.75$:
\[
  |(\Delta\,\DT)_{jk}| \;\le\; \sum_{j}|(\Delta\,\DT)_{jk}|
  \;=\; \norm{\Delta\,\DT\, b_k}_{b} \;\le\; \norm{\Delta\,\DT\, b_k}_{X}
  \;\le\; \norm{\Delta\,\DT}_{X\to X}\norm{b_k}_{X}
  \;\le\; 1.75\,\eps_{\mathrm{DT}} ,
\]
where $\norm{v}_{b}\le\norm{v}_{X}$ holds because the basis is normalized
($\norm{e_k}_X = 1$, \S\ref{sec:3.1}), so the $X$-norm of $v$ is the
$\ell^1$-norm of its $e$-coefficients and carries the extra $e_0$ term.
The deposited certificate inflated by $\eps_{\mathrm{DT}}$ and not by
$1.75\,\eps_{\mathrm{DT}}$; the script cannot be rerun from the package (its
inputs are the unsaved working tree named above), so the effect is bounded here
instead of recomputed. The shortfall per entry is $0.75\,\eps_{\mathrm{DT}} =
2.71\ldots\times 10^{-17}$, over a column of $81$ entries $2.20\times 10^{-15}$, and
$\norm{M^{16}}_1 \le 16\,\norm{M}_1^{15}\cdot 81\cdot 0.75\,\eps_{\mathrm{DT}}$
larger, that is $5.29\times 10^{-12}$: the published
$\norm{M^{16}}_1 = 0.2899488032200199$ becomes at most $0.2899488032253063$
against the gate threshold $0.5$, and $\rho \le \norm{M^{16}}_1^{1/16}$ moves in
the twelfth digit. \textbf{No gate changes its verdict}, and the load-bearing
count ``exactly one eigenvalue outside'' does not run through
$\eps_{\mathrm{DT}}$ at all --- it goes through the Stein--Taussky--Wimmer
inertia of \S\ref{sec:4.7}, of which \S\ref{sec:5.2} is an independent second
confirmation.

\textbf{Sanity (not rigorous).} Computing by the UPPER endpoints of the arb
balls (not midpoint, despite the name printed in the code) the eigenvalues of
$\hat{Q}$: spectral radius $\approx$ 0.2371 (stored in
\texttt{03-\allowbreak CERTIFICATES/\allowbreak pointwise/\allowbreak hyperb\_\allowbreak l3.json\allowbreak ::\allowbreak gates.\allowbreak gate\_\allowbreak 4\_\allowbreak sanity\_\allowbreak midpoint\_\allowbreak spectrum}),
top-5 $\approx \{0.2371, -0.1985, -0.1149, 0.0733, -0.0382\}$ (printed to
stdout by \texttt{04-\allowbreak CODE/\allowbreak build\_\allowbreak majorant\_\allowbreak rigorous.py:262-264}, not persisted
to JSON) --- roughly confirms that the rigorous upper bound $\rho \le 0.92567$
is not vacuous.

\textbf{Uniformity over the ball.} The Krawczyk operators with $\pm
\eps_{\mathrm{DT}}$ inflation described above confirm: the eigenpair
$(\delta^{\star},v,\varphi^{\star})$ exists uniformly over the whole ball
$\norm{G-\Gbar}_X \le 1.75\cdot \hat{\rho}'$ (\texttt{K\_max\_width} $\approx
4\times 10^{-16}$ against $w = 10^{-10}$). The fact of hyperbolicity does not
depend on the exact point $\Gs$.


\subsection{The case $\ell$ = 5/2: rigorous residual and existence certificate}
\label{sec:5.3}

\subsubsection{Rigorous residual (from \texttt{03-\allowbreak CERTIFICATES/\allowbreak pointwise/\allowbreak xresid\_\allowbreak l25.json})}
\label{sec:5.3.1}

\begin{table}[htbp]
\centering
\small
\caption{Rigorous residual quantities for $\ell=5/2$.}
\label{tab:5.3.1-residual}
\begin{tabular}{p{0.45\textwidth} p{0.45\textwidth}}
\toprule
Quantity & Value \\
\midrule
\textbf{$\norm{F(\Gbar)}_X$ (rigorous)} & \textbf{2.538195012737449 $\times$ $10^{-33}$} \\
$\norm{\bar{F}}_X$ (poly) & 9.631 $\times$ $10^{-46}$ \\
$\norm{\bar{F}}_X$ (tail) & 2.539 $\times$ $10^{-33}$ \\
$\norm{h_W}$ (Cauchy tail of $\Wbar$) & 4.387 $\times$ $10^{-34}$ \\
$R$ (large circle) & 1.20 \\
$\NW$ (Miller order) & 192 \\
$\Real U_{\min}$ on $D_R$ & $\ge$ 0.566555 at the gate $\Real U > 0$ (rigorous,
\S\ref{sec:5.3.2} and \texttt{04-\allowbreak CODE/\allowbreak lemma\_\allowbreak domains.py}); source ---
\texttt{03-\allowbreak CERTIFICATES/\allowbreak pointwise/\allowbreak reu\_\allowbreak min\_\allowbreak l25.json\allowbreak ::\allowbreak ReU\_\allowbreak min\_\allowbreak lower} \\
\bottomrule
\end{tabular}
\end{table}

Target gate: $\norm{\bar{F}}_X \le 10^{-25}$. Achieved \textbf{8 orders of
magnitude better}.

\subsubsection{Branch-uniqueness argument for Miller}
\label{sec:5.3.2}

Raising to the power $\Ubar^{(5/2)}$ does not single out a branch uniquely on
$\CC\setminus (-\infty, 0]$. For the Miller recurrence to coincide with the
Taylor coefficients of $\Wbar = \Ubar^{(5/2)}$
(\texttt{03-\allowbreak CERTIFICATES/\allowbreak pointwise/\allowbreak xresid\_\allowbreak l25.json}):

\begin{enumerate}
\item $\Ubar$ is holomorphic on $D_R$ ($R=1.20$) --- it is a polynomial of
degree 64, holomorphic everywhere.
\item $U(D_R) \subset \{\Real>0\}$: rigorously $\Real U \ge 0.566555$ on the
whole circle $|t-c| = R$, and hence on the disk ($U$ is entire, $\Real U$ is
harmonic, the minimum is on the boundary --- the minimum principle for
harmonic functions, the same trick as \S\ref{sec:10.2}; see also
\texttt{04-\allowbreak CODE/\allowbreak lemma\_\allowbreak domains.py}); from this also $|U| \ge 0.566555$. The
count is interval over 2048 arcs in \texttt{arb}/\texttt{acb}, generated by
\texttt{04-\allowbreak CODE/\allowbreak tests/\allowbreak \_\allowbreak reu\_\allowbreak min\_\allowbreak l25.py}, certificate
\texttt{03-\allowbreak CERTIFICATES/\allowbreak pointwise/\allowbreak reu\_\allowbreak min\_\allowbreak l25.json}. The lower bound is
rounded DOWN.
\item $u_0 = 0.853 > 0.5 \Rightarrow$ the $m$-th power branch from $u_0 = u_0^m >
0$ (no phase).
\item $\Wbar = \Ubar^{(5/2)}$ is holomorphic on $D_R$ as a composition of the
holomorphic $\Ubar$ and $z^{(5/2)}$.
\item The Taylor coefficients are defined uniquely; the Miller recurrence with
$w_0 = u_0^{(5/2)}$ (the correct initial condition) yields the same expansion
(both are solutions of the same ODE $W'\cdot U = (5/2)\cdot W\cdot U'$ with one
initial condition).
\end{enumerate}

Tests in \texttt{04-\allowbreak CODE/\allowbreak tests/\allowbreak sanity\_\allowbreak miller.py}: $\ell$=2 --- 129/129
coefficients matched the direct arb convolution of $\Ubar^{2}$; $\ell$=3 ---
193/193 with $\Ubar^{3}$; $\ell$=2.5 --- evaluation of $\Wbar$ at 10 points of
$[0,1]$ versus mpmath dps=500 --- enclosure in all 10 points (with tolerance
1e-60, bounded by the accuracy of the seed a64).

\subsubsection{Existence certificate at $\ell = 5/2$}
\label{sec:5.3.3}

The mirror of the $\ell = 3$ certificate of \S\ref{sec:5.1} for non-integer $\ell$=2.5: $\Wbar = \Ubar^{(5/2)}$
(Miller, poly-part up to $\NW=192$ + Cauchy tail),
$DW\cdot v = (5/2)\cdot \Ubar^{3/2}\cdot (v\circ \bar{s} + (\Gbar'\circ
\bar{s})(Ds\cdot v))$ --- the second non-integer power $\Ubar^{3/2}$ by the
same Miller; the DF formula is validated by finite differences up to
$9.97\times 10^{-57}$ (a high-precision diagnostic run, not a certificate:
no gate tests this agreement, and no statement depends on its value).

\begin{theorem}[existence, $\ell$=5/2; \texttt{03-\allowbreak CERTIFICATES/\allowbreak pointwise/\allowbreak kant\_\allowbreak l25.json}, all 6 gates PASS as recorded there; $Y_0$ and $\rho^{\star}$ below are reassembled from the honest tail --- see ``A superseded field'']
At $\ell$=5/2, $c=3/5$, $r=4/5$, $\hat{\rho}'=1\times 10^{-15}$ there exists a
unique $\Gs \in \Aff$ ($\Gs(0)=1$) in the ball $\norm{G-\Gbar}_X \le
1.75\times 10^{-15}$, analytic on $D(c,r) \supset [0,1]$, solving $F(G)=0$;
$\norm{\Gs-\Gbar}_X \le 1.75\cdot \rho^{\star} \le 4.4419\times 10^{-33}$
(exactly: 4.4418419387\ldots$\times 10^{-33}$).
\end{theorem}

The single-valuedness of the branch $U_G^{5/2}$ (avoidance of the cut
$(-\infty,0]$, $\Real U>0$) over the whole ball $\norm{G-\Gbar}_X\le
1.75\times 10^{-15}$ is not certified by a separate gate in
\texttt{kant\_\allowbreak l25.json} (there are 6 gates there: \texttt{B1}, \texttt{f0\_eq\_0}, \texttt{Z1\_lt\_1},
\texttt{Z\_hat\_lt\_1}, \texttt{gate5\_root}, \texttt{rho\_star} --- Re U is not among them), but is
supported by an order-of-magnitude estimate: the shift
$|U_G-\Ubar|\le \kappa_G\cdot 1.75\times 10^{-15}\le 5.74\times 10^{-15}$
against the margin $\Real \Ubar\ge 0.566555$ (for $\Gbar$, \S\ref{sec:5.3.2}) --- an
enormous margin, but as a separate premise it is silent.

\begin{table}[htbp]
\centering
\small
\caption{Rigorous constants of the existence certificate at $\ell=5/2$.}
\label{tab:5.3.3-constants}
\begin{tabular}{p{0.20\textwidth} p{0.20\textwidth} p{0.50\textwidth}}
\toprule
Constant ($\ell$=2.5) & Value & Source \\
\midrule
$Y_0$ & 2.5382$\times 10^{-33}$ & \texttt{03-\allowbreak CERTIFICATES/\allowbreak pointwise/\allowbreak kant\_\allowbreak l25.json\allowbreak ::\allowbreak Y0.\allowbreak Y0\_\allowbreak part1\_\allowbreak head} $+$ \texttt{.Y0\_\allowbreak part2\_\allowbreak tail} $+$ \texttt{xresid\_\allowbreak l25.json\allowbreak ::\allowbreak tail\_\allowbreak F} \\
$Z_1$ & 1.50$\times 10^{-7}$ & \texttt{03-\allowbreak CERTIFICATES/\allowbreak pointwise/\allowbreak kant\_\allowbreak l25.json\allowbreak ::\allowbreak gates.Z1\_\allowbreak lt\_\allowbreak 1.value.mid} \\
L (rigorous chain) & 1.54$\times 10^{3}$ & \texttt{03-\allowbreak CERTIFICATES/\allowbreak pointwise/\allowbreak kant\_\allowbreak l25.json\allowbreak ::\allowbreak gates.Z\_\allowbreak hat\_\allowbreak lt\_\allowbreak 1.L.mid} \\
$\Zhat$ & 1.50$\times 10^{-7}$ & \texttt{03-\allowbreak CERTIFICATES/\allowbreak pointwise/\allowbreak kant\_\allowbreak l25.json\allowbreak ::\allowbreak gates.Z\_\allowbreak hat\_\allowbreak lt\_\allowbreak 1.value.mid} \\
$\rho^{\star}$ & 2.5382$\times 10^{-33}$ & $Y_0/(1-\Zhat)$ from the two entries above \\
$\norm{\Ahat}_1$ & 2.992 &
\texttt{03-\allowbreak CERTIFICATES/\allowbreak pointwise/\allowbreak kant\_\allowbreak l25.json\allowbreak ::\allowbreak A\_\allowbreak tilde.norm\_\allowbreak 1.mid} \\
\bottomrule
\end{tabular}
\end{table}

\noindent\textbf{A superseded field.} The stored \texttt{Y0.value} and
\texttt{gates.\allowbreak rho\_\allowbreak star} of
\texttt{kant\_\allowbreak l25.json} hold $1.56\times10^{-33}$, a tail constant
carried over from an earlier seed. It cannot be a bound on $Y_0$: on the tail
$\Ahat = I$, so $Y_0$ is at least the tail of $\norm{\bar F}_X$ itself, and
that tail --- recomputed and gated by the acceptance step
\texttt{tests/\allowbreak \_\allowbreak xresid\_\allowbreak tail\_\allowbreak F\_\allowbreak l25.py}
--- is $2.53819501273648598\times10^{-33}$, larger by a factor $1.63$. The
table and the theorem above use the honest value, assembled from the head and
inner-tail parts of the same certificate and that recomputed tail; the two
stored fields are superseded and are used nowhere. Nothing else moves with
them: the enclosure is claimed in the ball of radius $1.75\times10^{-15}$,
eighteen orders wider, and no step of Theorems A--D passes through this
pointwise certificate (\S\ref{sec:5.3.4}).

\FloatBarrier

Outside that table, and not a rigorous constant: $\mathrm{cond}(M') = 14.5$, a
float64 diagnostic of the conditioning of the Newton matrix. Its source is not
persisted in the package --- the value was obtained by a single run and not
written into the certificates --- and no step of the argument uses it.

\subsubsection{Hyperbolicity certificate at $\ell = 5/2$}
\label{sec:5.3.4}

\begin{theorem}[Hyp-5-W9-L25]
$DT(\Gs)|_{X_{0}}$ at $\ell$=5/2 has a unique eigenvalue outside the unit disk:
$\delta^{\star} = 5.412736346083841$ (uniform interval
$[5.412736346069764, 5.412736346094575]$ with inflation
$\eps_{\mathrm{DT}} = 2.70\cdot 10^{-12}$ over the whole uniqueness ball
$\norm{G-\Gbar}_X \le 1.75\cdot 10^{-15}$), simple, with certified $v$ ($\norm{v}_1 \le
1.501$) and $\varphi^{\star}$ ($\langle \varphi^{\star},v\rangle \in
[1.0680404529, 1.0680404734]$); on the complement $Y=\ker \varphi^{\star}$:
$\norm{(Q^{\star})^{8}}_1 \le 0.1457517 < 1$ rigorously (majorant +
E-transfer deflator patch), i.e. the spectral radius $\le
0.1457517^{1/8} \le 0.78607 < 1$. Sanity (not rigorous): midpoint radius $\approx 0.2026 =
B^{+}(2.5)$ (stored in \texttt{03-\allowbreak CERTIFICATES/\allowbreak
pointwise/\allowbreak hyperb\_\allowbreak l25.json\allowbreak ::\allowbreak
gates.\allowbreak gate\_\allowbreak 6\_\allowbreak spectrum\_\allowbreak
sanity.\allowbreak max\_\allowbreak modulus}, with the top five in
\texttt{spectrum\_\allowbreak extra\_\allowbreak top5\_\allowbreak
midpoint\_\allowbreak NOT\_\allowbreak RIGOROUS} of the same file;
$B^{+}(\ell)$ is the second-largest-magnitude eigenvalue of
$DT(\Gs(\ell))|_{X_{0}}$, denoted $\lambda_2$ below; see \S\ref{sec:5.3.5}).
[Remark: in intermediate notes this value read 0.917812. That is \textbf{not}
the value of $\norm{(Q^{\star})^{8}}_1^{1/8}$: the same estimate, taken from a
looser intermediate constant, gives 0.917812, and as an upper bound on the
spectral radius it is valid ($0.9178 < 1$), but weaker. From the
$\norm{(Q^{\star})^{8}}_1 \le 0.1457517$ given here one obtains 0.78607. In the
text we use 0.78607; 0.917812 is mentioned only so that a reader
cross-checking with intermediate notes does not mistake the discrepancy for an
error.]
\end{theorem}

Sources: \texttt{03-\allowbreak CERTIFICATES/\allowbreak pointwise/\allowbreak eigen\_\allowbreak l25.json} (pairs:
$K(X)\subset \operatorname{int}(X)$ 80/80 both, $\delta$-boxes $\mathrm{width}=0.0$
when printing in double, $\eps_v=$2.7169e-15, $\eps_\varphi$=6.7800e-9),
\texttt{03-\allowbreak CERTIFICATES/\allowbreak pointwise/\allowbreak hyperb\_\allowbreak l25.json} (majorant: the
$\norm{Q^{2^m}}$ series, $n=8$ first $< 0.4$; $\norm{(Q^{\star})^{8}}$ patch;
uniform Krawczyk with $\pm \eps_{\mathrm{DT}}$). Remark on the gates: the
target $\delta$ interval of width 1e-11 from the brief was calibrated for
$\ell$=3 (where $\eps_{\mathrm{DT}}=$3.6e-17); for $\ell$=2.5 the inflation box
is wider (2.5e-11) --- it does not affect the theorem (only $|\delta|>1$ and
uniqueness are needed).

The same remark on the branch as in \S\ref{sec:5.3.3} (the
single-valuedness of $U_G^{5/2}$/$\Ubar^{3/2}$ --- avoidance of the cut,
$\Real U>0$ --- is not certified by a separate gate) also applies to the
hyperbolicity certificate just stated: neither \texttt{eigen\_\allowbreak l25.json} nor \texttt{hyperb\_\allowbreak l25.json} contains a
gate on $\Real U$ / the cut; the numerical margin is the same ($\sim$12 orders of
magnitude).

\subsubsection{Certified $\lambda_2$ (universal constant of the class, $\ell$=3)}
\label{sec:5.3.5}

\begin{proposition}[\texttt{03-\allowbreak CERTIFICATES/\allowbreak pointwise/\allowbreak lam2\_\allowbreak l3.json}]
The second (dominant stable) eigenvalue of $DT(\Gbar)|_{X_{0}}$ at $\ell$=3:
$\lambda_2 = B^{+} \in [0.23714611165001134 \pm 1.6\cdot 10^{-18}]$
(mean-value-form Krawczyk, inclusion 80/80, agreement with the non-rigorous
pos branch up to 15 digits). Here $1.6\cdot 10^{-18}$ is a rounded-up bound on the widest of the 80 component K-boxes (field \texttt{K\_max\_width} $= 1.5231\cdot 10^{-18}$ in \texttt{03-\allowbreak CERTIFICATES/\allowbreak pointwise/\allowbreak lam2\_\allowbreak l3.json}); the K-box of $\lambda_{2}$ itself is narrower --- its width (field \texttt{K\_lam2\_width}) prints as 0.0 when output in double. The gate of the certificate is passed against a much WIDER target interval, $[0.2371461116,\, 0.2371461117]$ (field \texttt{gate2\_\allowbreak lam2\_\allowbreak box\_\allowbreak subset\_\allowbreak target}), that is against a half-width of $5\cdot 10^{-11}$; the enclosure printed above is therefore the stronger statement, not the gate itself.
\end{proposition}

It is the subleading spectrum of the period-doubling operator that governs the corrections to scaling studied in the physics literature \cite{mao1987corrections, reick1992universal}; the value certified above is not compared with those non-rigorous estimates.

\textbf{Pointwise nature of the certificate (important for reading the
heading).} Unlike $\delta^{\star}$ (\S\ref{sec:5.2}), $\lambda_2$ here is
certified at a SINGLE point $\Gbar$, $\ell$=3, WITHOUT inflation over the
uniqueness ball $\norm{G-\Gbar}_X \le 1.75\cdot \hat{\rho}'$ --- ``universal
constant of the class'' in the heading means only that this is the same
quantity for the whole Feigenbaum--Cvitanovi\'{c} class at the given point
$\ell$, not that the value is certified uniformly over the solution
$\Gs(3)$ (as $\delta^{\star}$ is). According to \S\ref{sec:4.11}, $\lambda_2(\ell)$ as a
function of $\ell$ is not certified at all (count in float) --- this single
pointwise certificate at $\ell$=3 is not part of any premise of Theorems A--D
and is used only diagnostically (\S\ref{sec:5.3.4}, \S\ref{sec:7.1}, \S\ref{sec:7.4}).



\section{Analytic dependence on $\ell$ (Theorem C)}
\label{sec:6}

Both previous theorems give a statement ``for each $\ell$'' and say nothing about
how the solution and the constant depend on $\ell$. This section closes the
question: the dependence is real-analytic, and all derivation constants are
computed.


\subsection{Statement}
\label{sec:6.1}

We recall the setup. The space

\[
X=\Bigl\{h(t)=\textstyle\sum_{k\ge0} h_k(t-c)^k\;:\;\|h\|_X=\sum_{k\ge0}|h_k|\,r^k<\infty\Bigr\},
\qquad c=\tfrac35,\quad r=\tfrac45,
\]

--- a Banach algebra of functions analytic in the disk $D(c,r)\supset[0,1]$;
$X_0=\{h\in X:h(0)=0\}$ is the linear tangent subspace to the normalization
$G(0)=1$ (the derivatives $DF|_{X_0}$, $DT|_{X_0}$ live there, not the
normalized function itself); the affine slice
$\Aff:=1+X_0=\{G\in X:G(0)=1\}$ --- here $P(\ell)$ and $\Gs(\ell)$
live (see \S\ref{sec:2}). The Feigenbaum--Cvitanović operator in the $t$-coordinate
($t=|x|^{\ell}$) has the form

\begin{gather*}
F(G,\ell)(t)=G(t)+\alpha_G\cdot G(W_G(t)),
\qquad
\alpha_G=-\tfrac1{G(1)},\\
W_G=U_G^{\ell},\quad
U_G(t)=G(s_G(t)),\quad
s_G(t)=t\cdot\alpha_G^{-\ell}.
\end{gather*}

Theorem A gives, for each $\ell\in[2,4]$, a solution
$\Gs(\ell)\in\Aff$ (unique in the ball $1.75\cdot10^{-7}$ around
the center $P(\ell)$); Theorem B gives that $DT(\Gs(\ell))|_{X_0}$ has
exactly one eigenvalue $\delta(\ell)$ outside the unit disk, simple and real,
and the rest of the spectrum is rigorously inside.

\begin{theoremCrestated}[restated]
The maps
\[
\Phi:\ell\mapsto \Gs(\ell)\in\Aff,\qquad
\Delta:\ell\mapsto\delta(\ell)\in\RR
\]
are real-analytic on $[2,4]$. More precisely, for each
$\ell_0\in[2,4]$ there exist a complex neighbourhood
$\mathcal U_{\ell_0}\subset\CC$ of $\ell_0$ and holomorphic maps
$\widetilde\Phi:\mathcal U_{\ell_0}\to\Aff_{\CC}$ and
$\widetilde\Delta:\mathcal U_{\ell_0}\to\CC$, coinciding with
$\Phi,\Delta$ on $\mathcal U_{\ell_0}\cap[2,4]$.
\end{theoremCrestated}

Here $\Aff_{\CC}=1+X_{0,\CC}$, and $X_{0,\CC}$ is the
complexification of the tangent $X_0$ (\S\ref{sec:6.3}); holomorphicity of
$\widetilde\Phi$ as a map into $\Aff_{\CC}$ is equivalent to
holomorphicity of $\widetilde\Phi-1$ as a map into $X_{0,\CC}$ --- the
linear shift by a constant changes neither analyticity nor derivatives. Thus
the union $\bigcup_{\ell_0\in[2,4]}\mathcal U_{\ell_0}$ is an open
neighbourhood of $[2,4]$ in $\CC$, on which $\widetilde\Phi$ and
$\widetilde\Delta$ glue into single-valued holomorphic functions (\S\ref{sec:6.7},
\S\ref{sec:6.8}); that is, real-analyticity is understood in the strong sense ``analytic
on an open set containing the closed interval $[2,4]$''.


\subsection{Premises and their categorization}
\label{sec:6.2}

Each premise is labeled: \textbf{(a)} --- certified in the package (with the
file and field indicated), \textbf{(b)} --- proved in this section,
\textbf{(c)} --- accepted as a hypothesis. \textbf{Theorem C has exactly one
premise of category (c), and it is the package-wide one:} A4 of the registry
in \S\ref{sec:10.11} --- the correctness of the interval arithmetic
(\texttt{python-flint}/\texttt{arb}) through which every premise labeled (a)
below is obtained. It is external to this work and nothing here establishes
it (\S\ref{sec:9}); if it fails, every (a) label falls at once. Earlier
versions of this registry declared category (c) empty, which was true only of
the premises the registry itself lists.


\subsection*{(a) Certified in the package}
\label{sec:6.2.a}

The addresses below are written as \texttt{file::field}; the \texttt{field}
part may name a path inside the JSON object. The only wildcard used is the
per-box one, as in \texttt{box\_\allowbreak g*.json}: the existence certificates are one file
per box, named by the label of that box, so the wildcard runs over the 46
labels \texttt{g00}--\texttt{g43}, \texttt{ganchor} and \texttt{gtop} --- and likewise
\texttt{box\_\allowbreak t*.json} over the labels of the 706 tube boxes, one
file each (there the label carries a suffix \texttt{L}/\texttt{R}/\texttt{b}
wherever a box was split); the branch data \texttt{branch\_\allowbreak g*.npz}
follow the same 46 labels. Where a wildcard
would stand for anything other than such a per-box family, the files are named
individually instead.

{\small
\begin{longtable}{@{}c>{\raggedright\arraybackslash}p{0.44\linewidth}>{\raggedright\arraybackslash}p{0.40\linewidth}@{}}
\caption{Premises of Theorem C certified in the rigorous package.}
\label{tab:certified-premises}\\
\toprule
\# & Statement & Source \\
\midrule
\endfirsthead
\multicolumn{3}{@{}l}{\textit{Table \thetable\ (continued)}}\\
\toprule
\# & Statement & Source \\
\midrule
\endhead
\bottomrule
\endlastfoot
a1 & Existence: for every $\ell$ of the box the zero $\Gs(\ell)$ lies in
$B_X(P(\ell),\,1.75\rho^{\star})$, $\rho^{\star}\le4.85\cdot10^{-13}$ (in the
$X$-norm $1.75\rho^{\star}\le8.49\cdot10^{-13}$); uniqueness in the ball
$1.75\cdot10^{-7}$; the union of the 46 boxes covers $[2.0000,4.0075]$ &
\texttt{03-\allowbreak CERTIFICATES/\allowbreak existence/\allowbreak box\_\allowbreak g*.json::\allowbreak best.rho\_\allowbreak star},
\texttt{glue.json} \\
a2 & $\|I-A(\ell)M'(\ell)\|_1\le Z_1(\ell)\le 0.1501$ uniformly in $\ell$ over
the boxes & \texttt{box\_\allowbreak g*.json::\allowbreak Z1}, worst case in \texttt{box\_g43.json} \\
a3 & $\|A(\ell)\|_1\le B_A(\ell)$; worst case $B_A\le3.62$ &
\texttt{box\_\allowbreak g*.json::\allowbreak B\_\allowbreak A} (range $[2.865,\,3.617]$),
worst case in \texttt{box\_gtop.json} \\
a4 & Lipschitz constant $\norm{DT(G_1)-DT(G_2)}_{X}\le L\,\norm{G_1-G_2}_{X}$
in the $X$-norm at both ends;
$L\le4.26\cdot10^4$ & \texttt{box\_\allowbreak g*.json::\allowbreak best.L} \\
a5 & $\lvert G(1)\rvert$ is separated from zero: $\min\lvert G(1)\rvert\ge
m_1\ge0.393$ over all 46 boxes & \texttt{04-\allowbreak CODE/\allowbreak tests/\allowbreak \_\allowbreak nozeros\_\allowbreak check.py}
(interval lower bound $\min\lvert G(1)\rvert\ge0.3962$ over all 46 boxes) and
\texttt{04-\allowbreak CODE/\allowbreak tests/\allowbreak \_\allowbreak kappaG\_\allowbreak bound.py} (the same quantity under a different
subdivision, $\ge0.3987$); consistent with the pointwise
\texttt{03-\allowbreak CERTIFICATES/\allowbreak pointwise/\allowbreak kant\_\allowbreak l3.json::\allowbreak m1} at $\ell=3$ \\
a6 & Domain lemma: on the boundary $\partial D(c,r)$ one has
$\lvert U_{P(\ell)}\rvert\ge U_{\min}\ge0.626>0.5$, \texttt{margin\_a} $> 0$,
\texttt{margin\_c} $> 0$, the cut $U=(-\infty,0]$ is not touched --- and in fact
rigorously more than that: the gate requires $\operatorname{Re}U>0$ on both
circles (flag \texttt{nonzero\_ok}) &
\texttt{box\_\allowbreak g*.json::\allowbreak lemma1\_\allowbreak \{U\_\allowbreak min,margin\_\allowbreak a,margin\_\allowbreak c,nonzero\_\allowbreak ok,ReU\_\allowbreak min\_\allowbreak r\_\allowbreak val,ReU\_\allowbreak min\_\allowbreak R\}};
code \texttt{04-\allowbreak CODE/\allowbreak lemma\_\allowbreak domains.py}; the gate is \texttt{tube\_\allowbreak lib.py:1418}
on \texttt{nonzero\_ok}, and the same flag is written into the json as
\texttt{lemma1\_\allowbreak nonzero\_\allowbreak ok} by \texttt{rig2fix3.py:829} \\
a7 & Hyperbolicity: $\sup_{\lvert z\rvert=1}\|(DT_0-z)^{-1}\|_1\le R\le3.7961$,
$\|\Delta\|\le0.280$, $R\,\|\Delta\|\le0.9896$, the full test of \S\ref{sec:4.6}
$\le0.9919<1$; 706 boxes, their union covers $[2.0000,4.0012]$ (all 706 are nested inside their own existence box; \texttt{t0695} is a degenerate grid remainder beyond the right edge, and although its intersection with $[2,4]$ is a single point, it lies wholly inside \texttt{gtop}, so the \texttt{extra} term is justified on it too, see \S\ref{sec:4.8}) &
\texttt{03-\allowbreak CERTIFICATES/\allowbreak tube/\allowbreak box\_\allowbreak t*.json::\allowbreak \{R,delta\_\allowbreak norm,prod\}},
\texttt{tube/\allowbreak \_\allowbreak accept.out} \\
a8 & Spectrum of $DT(\Gs(\ell))\rvert_{X_0}$: exactly one value
$\delta(\ell)$ outside $\{\lvert z\rvert\le1\}$, simple and real, the
remaining spectrum lying in the open unit disk; in particular the spectrum
does not touch $\{\lvert z\rvert=1\}$ & Theorem B, \S\ref{sec:4} \\
a9 & Radius of holomorphy: all the premises of \S\ref{sec:6.5} (separation of
$\lvert G(1)\rvert\ge0.396279$ from zero, $\lvert\arg\alpha_G\rvert<\pi/2$,
$s_G(D) \subset D$ against \texttt{margin\_a},
$\kappa_U=\|U_G-1\|_X\le0.363184<1$ --- $X$-valued holomorphy of $W_G$ (not
merely the pointwise (C1), see item 5 of \S\ref{sec:6.5}), $\hat q\le0.955131<1$) hold
uniformly on the affine ball
$P(\ell)+B_{X_{0,\CC}}(0,\,1.75\cdot10^{-7})\subset\Aff_{\CC}$,
whence $R_G\ge1.75\cdot10^{-7}$ &
\texttt{box\_\allowbreak g*.json::\allowbreak \{lemma1\_\allowbreak margin\_\allowbreak a\}} +
\texttt{04-\allowbreak CODE/\allowbreak tests/\allowbreak \_\allowbreak holo\_\allowbreak radius.py} (46/46, field \texttt{kap\_U}) \\
\end{longtable}
}

All quantities in (a) are \textbf{uniform upper bounds in $\ell$ over the box}
(the supremum over the closed $\ell$-interval, not over a grid of nodes), which
is exactly what the premises of the implicit function theorem require.

On 7 boxes (\texttt{g37}\ldots\texttt{g43}, where the subdivision is
\texttt{lemma\_\allowbreak bases\_\allowbreak K=2}) all six fields are published from a SINGLE
sub-piece --- the one on which \texttt{margin\_c} is minimal
(\texttt{rig2fix3.py:481-500}) --- and not as a separate minimum for each
quantity; for \texttt{margin\_a}/\texttt{U\_min}/\texttt{ReU\_min\_R} this is
anti-conservative, and by an amount that differs per field (both $\ell$-halves of all seven boxes recomputed): the true minimum over the box is below the published one by $\le 6.67\cdot10^{-5}$ out of $\approx0.166$ for \texttt{margin\_a} ($0.041\,\%$, worst \texttt{g37}), by $\le 4.45\cdot10^{-4}$ out of $\approx0.663$ for \texttt{U\_min} ($0.068\,\%$, worst \texttt{g43}) and by $\le 5.59\cdot10^{-4}$ out of $\approx0.571$ for \texttt{ReU\_min\_R} ($0.098\,\%$, worst \texttt{g43}).
The three bounds above ARE reproducible from the package: \texttt{fp\_\allowbreak lib.\allowbreak solve\_\allowbreak fixed\_\allowbreak point} on the seed of \texttt{hiprec\_\allowbreak seeds.\allowbreak json}, then \texttt{rig\_\allowbreak lib.\allowbreak build\_\allowbreak branch}, \texttt{rig\_\allowbreak lib.\allowbreak P\_\allowbreak at\_\allowbreak interval} and \texttt{tube\_\allowbreak lib.\allowbreak verify\_\allowbreak lemma\_\allowbreak U} on each $\ell$-half, the fields taken as the lower end of the \texttt{arb} ball exactly as the engine prints them; the arithmetic is rigorous --- interval, \texttt{flint.arb} at \texttt{ctx.prec = 256}, and on the boxes with \texttt{lemma\_\allowbreak bases\_\allowbreak K=1} the recomputation reproduces the published fields to the last printed digit. The knock-on effect on $\kappa_G$ was measured separately, by an isolated run of \texttt{\_\allowbreak kappaG\_\allowbreak bound.py::\allowbreak kappa\_\allowbreak on}, and THAT run is not reproducible from the package: its \texttt{main()} walks all boxes with its own subdivision --- the same status as the batteries of \S\ref{sec:3.3} and \S\ref{sec:4.10}: the
effect on $\kappa_G$ is of order $4\cdot10^{-9}$, nine orders of magnitude
below the quantity $4.130319$ itself (and more than two orders below its last
printed digit) --- the published $\kappa_G\le4.14$ does not change; the gate
$\Delta s < \mathtt{margin\_a}$ is likewise out of risk (a margin of 5-6 orders
of magnitude). Each sub-piece separately passes all four conditions of the
lemma (\texttt{margin\_a} $> 0$, \texttt{margin\_c} $> 0$, \texttt{nonzero\_ok},
$U_{\min}\ge0.5$) --- the verdict of the box is unaffected; the matter
concerns only the printed numbers.


\subsection*{(b) Proved in this section}
\label{sec:6.2.b}

\begin{enumerate}
\item[1.] \textbf{Absence of zeros of $U$ in the disk} --- an honest bridge from
  the boundary certificate a6 to the statement on all of $D(c,r)$ (\S\ref{sec:6.4}).
  This replaces the incorrect ``minimum modulus principle without checking
  for zeros''.
\item[2.] \textbf{Holomorphy of $F$ in $(G,\ell)$} in a complex neighbourhood of
  every point $(\Gs(\ell_0),\ell_0)$ (\S\ref{sec:6.5}). That the neighbourhood
  covers $\Gs$ itself, and not merely the centre $P$, is not an appeal
  to continuity: the radius in $G$ is bounded below by a number (a9), and
  $\|\Gs-P\|_X$ is $2.06\cdot10^{5}\times$ smaller.
\item[3.] \textbf{Invertibility of $D_GF(\Gs(\ell_0),\ell_0)$} (transfer
  from the centre to the solution) and application of the analytic implicit
  function theorem (\S\ref{sec:6.6}).
\item[4.] \textbf{Gluing} of the local holomorphic branches into a single
  analytic function on a neighbourhood of $[2,4]$ (\S\ref{sec:6.7}).
\item[5.] \textbf{Analyticity of $\delta(\ell)$} via the Riesz projector on the
  \textbf{unit circle} (\S\ref{sec:6.8}); the incorrect contour $|z-\delta|=\rho$ of the
  draft has been eliminated.
\end{enumerate}


\subsection*{(c) Hypotheses}
\label{sec:6.2.c}

A4 (\S\ref{sec:10.11}): the correctness of \texttt{python-flint}/\texttt{arb} interval
arithmetic, common to every (a) premise of this section and to the whole
package. There are no further hypotheses specific to Theorem~C.


\subsection{Complexification of $X$ and of the operator $F$}
\label{sec:6.3}

\subsubsection{The space $X_{\CC}$}
\label{sec:6.3.1}

Let

\[
X_{\CC}=\Bigl\{h(z)=\textstyle\sum_{k\ge0}h_k(z-c)^k:\;h_k\in\CC,\
\|h\|_X=\sum_{k\ge0}|h_k|r^k<\infty\Bigr\}
\]

--- a complex Banach space with the same norm $\|\cdot\|_X$; $X$ embeds into it
as a closed $\RR$-linear subspace (real coefficients). The slice

\[
X_{0,\CC}=\{h\in X_{\CC}:h(0)=0\}
\]

is closed; $X_0=X\cap X_{0,\CC}$. Since $\|\cdot\|_X$ is a weighted
$\ell^1$ norm with weights $r^k$, the space $X_{\CC}$ is a Banach
algebra under convolution, and in particular the multiplication
$(g,h)\mapsto g\cdot h$ is continuous and holomorphic in each variable. For
the \textbf{composition} of series the algebra structure alone is not enough:
one also needs the inner argument to map the disk strictly into its interior.
Quantitatively this is the condition $q<1$; it is introduced and certified in
\S\ref{sec:6.5} (item 6) and is not yet used here.
$F$ is nonlinear in $G$ (through $\alpha_G=-1/G(1)$ and the composition
$G\circ W_G$); what is linear is not the operators of the article themselves
but their \textbf{derivatives at a fixed base point}: for a given $G$ the map
$h\mapsto D_GF(G)\,h$ (and, correspondingly, $DT(G)$) is $\RR$-linear in
the direction $h$ and extends to $h\in X_{0,\CC}$ by
$\CC$-linearity (the complex directional derivative) --- this and only
this is used below. Since the matrices $A(\ell)$, $M'(\ell)$ in the
$b$-basis (precisely these derivatives at a real centre $\ell$) have real
coefficients, their column $\ell^1$-norm over $\CC$ coincides with the
column norm over $\RR$; hence the quantities $Z_1$, $B_A$, $R$,
$\|\Delta\|$ of (a), taken at the real centre, carry over to the
complexification unchanged. For the same reason the remaining two quantities
of the Kantorovich gate --- $Y_0$ and $L$ --- carry over as well (they are
needed in \S\ref{sec:6.7}, step 1, where the uniqueness part of Theorem A is applied on
the complexified ball $P(\ell)+B_{X_{0,\CC}}(0,1.75\cdot10^{-7})$ and not
only on its real slice): $Y_0=\|\tilde A\cdot\tilde F(P)\|_1$ is a NUMBER
computed at the fixed real centre $P(\ell_0)$, and complexifying $h$ does not
change it; $L$ is derived in \texttt{tube\_\allowbreak lib.\_\allowbreak compute\_\allowbreak L} by a chain of
majorants of the form $\sum|\text{coef}|\cdot r^k$ and by the triangle
inequality, no step of which appeals to the realness of the increment $h$ ---
the derivation is valid verbatim for $h\in X_{0,\CC}$ as well. The
holomorphy of $F$ itself in $(G,\ell)$ (that is, the correctness of
complexifying the base point and not only the direction) is proved
separately, block by block, below in \S\ref{sec:6.3.2}--\S\ref{sec:6.5} --- through the explicit
branches $\exp(\ell\,\mathrm{Log}(\cdot))$, not through linearity. Details on
the $b$-basis and on the equivalence of norms are in \S\ref{sec:4.3}.

\subsubsection{What changes for complex $\ell$}
\label{sec:6.3.2}

``Real-analytic'' means extendability to a complex neighbourhood of $\ell$.
Therefore every construction containing a power with exponent $\ell$ (namely
$\alpha_G^{-\ell}$ and $U^{\ell}$) must make sense for complex $\ell$ near the
real axis. The only reasonable way is to define them as holomorphic branch
functions:

\[
\alpha_G^{-\ell}:=\exp(-\ell\,\mathrm{Log}\,\alpha_G),\qquad
U^{\ell}:=\exp(\ell\,\mathrm{Log}\,U).
\]

The correctness of these definitions is exactly what is verified in
\S\ref{sec:6.4}--\ref{sec:6.6}: $\alpha_G$ lies in the right half-plane (it does not cross the
cut $(-\infty,0]$), and $U(t)$ is separated from the cut on all of $D(c,r)$.
Then $\mathrm{Log}$ is taken to be the principal branch and is holomorphic;
for real $\ell$ and real values these definitions coincide with the original
$U^{\ell}$, $W_G$ (branch normalization along the real axis, as in \S\ref{sec:5.3.2}). The non-integrality of $\ell$ thereby \textbf{introduces no}
singularities: the only singularity of the original $|x|^{\ell}$ sits at
$x=0$ and is absorbed by the substitution $t=|x|^{\ell}$ (\S\ref{sec:1.1}); here $\exp$ and $\mathrm{Log}$ are holomorphic wherever they are
defined.


\subsection{The domain lemma: from the boundary to the disk (correction of error 1)}
\label{sec:6.4}

\subsubsection{What is actually certified}
\label{sec:6.4.1}

The code \texttt{04-\allowbreak CODE/\allowbreak lemma\_\allowbreak domains.py::\allowbreak verify\_\allowbreak lemma\_\allowbreak 1} computes
\texttt{margin\_a} (via the triangle inequality, correctly \textbf{on the
whole disk}) and \texttt{U\_min}, \texttt{margin\_c}, \texttt{cut\_ok} ---
\textbf{only on the circle} $|t-c|=r$ (moreover \texttt{cut\_ok} itself is not
written into the json and does not enter any gate: both premise a6 and the
code rely on the rigorously stronger \texttt{nonzero\_ok},
$\operatorname{Re}U>0$ on both circles) (the comment
\verb!# Compute on |t-c| = r!, lines 118--133), although the docstring (line
86) claims a minimum over the closed disk. The quantity \texttt{margin\_c} is
nevertheless correct on the disk as well: $W-c$ is analytic where it is
defined, and the maximum modulus principle gives
$\max_{D(c,r)}|W-c|=\max_{\partial D(c,r)}|W-c|$. But $U_{\min}\ge0.626$ on
the \textbf{boundary} does \textbf{not} by itself imply $|U|\ge0.626$ in the
disk: the minimum modulus principle applies only to a function without zeros
inside, and the absence of zeros is precisely what has to be proved.
(Counterexample from the brief: $f(z)=z$ on the unit disk --- the minimum on
the boundary is $1$, with a zero at the centre.) The transfer ``boundary
$\Rightarrow$ disk'' in the draft closed on itself in a circle.

\subsubsection{An honest derivation: $|U-1|<1$ in the disk}
\label{sec:6.4.2}

The centre $P(\ell)$ is a \textbf{polynomial} in $t$ of degree $\le N_p=80$
with coefficients depending holomorphically (in fact polynomially) on $\ell$;
its $0$-centred coefficients $a_k(\ell)$ ($a_0=1$, i.e. $G(0)=1$) are
extracted from \texttt{branch\_\allowbreak g*.npz::\allowbreak c\_\allowbreak a}
(same per-box convention, \S\ref{sec:6.2.a}). Then

\begin{align*}
U_{P(\ell)}(t)&=P(\ell)\bigl(t\,\alpha(\ell)^{-\ell}\bigr)\\
&=\sum_{k=0}^{N_p} a_k(\ell)\,\alpha(\ell)^{-k\ell}\,t^k,
\end{align*}

and this equality is \textbf{exact} (a polynomial, not the limit of a
series), hence

\begin{align*}
|U_{P(\ell)}(t)-1|
&=\Bigl|\sum_{k=1}^{N_p}a_k(\ell)\,\alpha(\ell)^{-k\ell}\,t^k\Bigr|\\
&\le\sum_{k=1}^{N_p}|a_k(\ell)|\,\alpha(\ell)^{-k\ell}\,(c+r)^k
\qquad(|t|\le c+r=\tfrac75).
\end{align*}

Here $\alpha(\ell)=-1/P(\ell)(1)>0$. The sign does not follow from a5 (which
certifies only the modulus $|G(1)|$), but neither is it taken from reference
points: \texttt{04-\allowbreak CODE/\allowbreak tests/\allowbreak \_\allowbreak nozeros\_\allowbreak check.py::\allowbreak box\_\allowbreak sup\_\allowbreak piece}
constructs $\alpha(\ell)$ in interval arithmetic on EVERY piece of EVERY one
of the 46 boxes and halts with an error if $\alpha>0$ is not guaranteed on
that piece --- that is, positivity is verified by machine and uniformly over
the box, not at isolated points. The substantive explanation of why the sign
cannot change comes from the continuity of $\ell\mapsto P(\ell)(1)$ on each
box, from its negativity at the reference points ($G(1)=-0.39954$ at
$\ell=2$; $\bar G(1)=-0.51876$ at $\ell=3$, Appendix~\ref{sec:A}), and from the fact
that $|P(1)|\ge0.39>0$ (a5) leaves the sign no room to change inside a box or
across the seams (each centre lies within $\rho^{\star}_{X}\le 8.49\cdot 10^{-13}$
of $\Gs(\ell)$ by \S\ref{sec:3.6}, so two neighbouring centres differ by at
most $1.7\cdot 10^{-12}$ in the norm $X$, and $1.7\cdot 10^{-12}\ll0.39$); thus
$\alpha(\ell)^{-k\ell}=\exp(-k\ell\log\alpha(\ell))$ is a real number,
strictly decreasing in $k$. The sum on the right is finite and is
\textbf{rigorously evaluated in arb interval arithmetic} over all 46 boxes
(the supremum over $\ell$ in the box; $\ell$ is substituted as an interval,
and the coefficients $a_k(\ell)$ as interval Chebyshev bounds). The result:

\[
\sup_{\ell\,\text{of the box}}\ \sum_{k=1}^{N_p}|a_k(\ell)|\,\alpha(\ell)^{-k\ell}\,(c+r)^k
\ \le\ 0.379\ <\ 1
\]

(worst box \texttt{g06}: $0.3648$; script
\texttt{04-\allowbreak CODE/\allowbreak tests/\allowbreak \_\allowbreak nozeros\_\allowbreak check.py}, subdivision of the box into $K=2$
parts). The article quotes the round $0.379$ rather than $0.3648$ deliberately: it is a
slack bound fixed once and used in a dozen places, so that a re-run with a different
subdivision cannot invalidate a printed number. Three figures for this one quantity
$\sup|U_{P(\ell)}-1|$ therefore appear in the text, in decreasing order of slack:
$0.379$ (quoted throughout), $0.3648$ (the worst box at $K=2$, above), and
$\kappa_{U}\le 0.363184$ (\S\ref{sec:6.5} item~5, sharper still). They do not disagree.
Consequently

\begin{equation}
|U_{P(\ell)}(t)-1|\le 0.379<1\quad\text{for all } t\in D(c,r),\ \ell\in[2,4].
\tag{C1}
\end{equation}

From this it follows at once that $U_{P(\ell)}$ \textbf{has no zeros} in
$D(c,r)$ and, moreover, that $\mathrm{Re}\,U_{P(\ell)}(t)\ge 1-0.379=0.621>0$
--- the image of the disk lies in the right half-plane, far from the cut
$(-\infty,0]$.

\textbf{Remark on consistency.} Formally (C1) is a premise of category (b):
it is not stored as a ready-made field in the certificate but is derived from
the already certified coefficients $c_a$ by a finite interval sum, without
new analytic assumptions. The margin is enormous ($0.379$ against the
threshold $1$), which makes the conclusion robust to further perturbations
(see \S\ref{sec:6.4.3}).

\subsubsection{Transfer to the solution and to a complex neighbourhood}
\label{sec:6.4.3}

In Theorem C the separation of $U_G$ from zero is needed not only for the
centre $P(\ell)$ but also for $G=\Gs(\ell)$ and for complex $\ell$. This
follows from (C1) by continuity, and here the margin $0.621$ absorbs all the
perturbations.

\textbf{(a) In $G$.} The map $G\mapsto U_G=G\circ s_G$
($s_G(t)=t\alpha_G^{-\ell}$, $\alpha_G=-1/G(1)$) is smooth (indeed
holomorphic) in a neighbourhood of $P(\ell)$, since $G(1)\neq0$ (a5).
Consequently it is Lipschitz with a finite constant $\kappa_G$, and for
$\|G-P(\ell)\|_X\le1.75\cdot10^{-7}$ (the uniqueness ball, a1)

\begin{align*}
|U_G(t)-1|&\le|U_{P(\ell)}(t)-1|+\kappa_G\|G-P(\ell)\|_X\\
&\le0.379+\kappa_G\cdot1.75\cdot10^{-7}<1.
\end{align*}

\textbf{The value of $\kappa_G$ has been computed} (script
\texttt{04-\allowbreak CODE/\allowbreak tests/\allowbreak \_\allowbreak kappaG\_\allowbreak bound.py}, a rigorous arb bound over all 46
boxes). The norm of the differential of the map $G\mapsto U_G$ at the point
$G$ in the direction $h$, $\|h\|_X=1$, is estimated as follows: from
$|h(u)|\le\|h\|_X$ for $|u-c|\le r$, $|h(1)|\le\|h\|_X$ (since
$|1-c|=0.4<r$), $|G'(s)|\le\sum_k k|g_k|r^{k-1}$ and $|t|\le c+r$ it follows
that

\begin{align*}
\|D_GU_G\,h\|_\infty&\le 1+\Bigl(\sum_k k|g_k|r^{k-1}\Bigr)\,(c+r)\,\ell\,
\frac{|\alpha_G|^{-\ell-1}}{|G(1)|^{2}}\\
&= 1+\Bigl(\sum_k k|g_k|r^{k-1}\Bigr)\,(c+r)\,\ell\,|G(1)|^{\ell-1}.
\end{align*}

The second equality is the identity $\alpha_G=-1/G(1)$: the factor
$|\alpha_G|^{-\ell-1}/|G(1)|^{2}=|G(1)|^{\ell+1}/|G(1)|^{2}$ depends on
\textbf{one} quantity $G(1)$, and to bound its two factors separately
($|\alpha_G|^{-\ell-1}\le1$ on its own and $1/|G(1)|^{2}$ on its own) means
losing a factor $|G(1)|^{-(\ell+1)}$ for no reason at all. The script takes
the right-hand expression: coarsening to the left-hand one is legitimate but
costs roughly $10.4\times$ in $\kappa_G$ itself ($43.07$ instead of $4.14$; the
$43.07$ is a one-off diagnostic of the rejected variant, not persisted in the
package).

This is a bound at a \textbf{single point} $G$; the Lipschitz constant on the
segment $[P(\ell),P(\ell)+\eta]$, however, requires the supremum over the
whole segment, not the value at its endpoint. Hence all three $G$-dependent
quantities are inflated over the ball $\|\eta\|_X\le\rho=1.75\cdot10^{-7}$:

\begin{itemize}
\item $|G(1)|\ge|P(1)|-\rho$ --- since $|\eta(1)|\le\|\eta\|_X$;
\item $|G(1)|\le|P(1)|+\rho<1$ --- whence, first, $|\alpha_G|>1$ (an order on
  $\CC$ is not defined; this is needed for the bound on the shift
  $\Delta s$ below) and, second, the base of the power in $|G(1)|^{\ell-1}$
  is less than one, so that the power \textbf{increases} in the base and
  \textbf{decreases} in the exponent; the upper bound is therefore taken at
  the largest base $|P(1)|+\rho$ and the smallest exponent $\ell_{\min}-1>0$
  (the positivity of the exponent is verified separately --- for $\ell<1$
  the monotonicity in the base would reverse sign);
\item $|G'(s)|\le\|P'\|+\rho/\bigl(r(1-\theta)^2\bigr)$, where
  $\theta=(r-\mathtt{margin\_a}+\Delta s)/r$ --- a Cauchy estimate inside the disk, taken on the circle that the image of $s_{G}$ actually reaches for every $G$ of the ball, not merely the image of $s_{P}$ (the shift $\Delta s$ is the one bounded in the next paragraph; the gate $\Delta s<\mathtt{margin\_a}$ that appears there is exactly what keeps $\theta<1$):
  for $|u-c|\le\theta r$ one has
  $\sum_k k|\eta_k|(\theta r)^{k-1}\le\|\eta\|_X\,r^{-1}\sum_k k\theta^{k-1}
  =\|\eta\|_X/\bigl(r(1-\theta)^2\bigr)$.
\end{itemize}

In addition it is verified that the image $s_G(D(c,r))$ stays inside
$D(c,r)$ for \textbf{all} $G$ of the ball: the shift does not exceed
$(c+r)\,\ell\,\Delta\alpha$ with $\Delta\alpha\le\rho/(|P(1)|(|P(1)|-\rho))$,
and it is required to be smaller than the certified \texttt{margin\_a} of
the domain lemma (field \texttt{lemma1\_\allowbreak margin\_\allowbreak a}, \S\ref{sec:6.4.1}). On all 46 boxes this
holds with a margin of at least four orders of magnitude. Without this check
the first of the three inequalities above ($|h(u)|\le\|h\|_X$ on the image)
would have no force.

In the draft there was no inflation: the script substituted only the centre
$P(\ell)$, and the declared ``constant on the ball'' was the norm of the
differential at the centre. The difference is numerically small ($4.130310$
against $4.130319$; the first of the two comes from the superseded draft script
and is not persisted in the package --- it is quoted only to show how small the
gap was), but without it the quantitative transfer to the ball is
not derived.

Here $g_k$ are the coefficients of the expansion \textbf{about the centre
$c$}: it is precisely these that enter the norm $\|h\|_X=\sum_k|h_k|r^k$, and
for them $G(1)=\sum_k g_k(1-c)^k$. This is essential: the parametric centre is
stored in the monomial basis $t^k$, and substituting the monomial
coefficients into these formulas gives a different quantity --- the value of
$P$ at the point $0.4$ instead of $G(1)$. Check: at $\ell=2$ one obtains
$G(1)=-0.39954$, that is $\alpha=2.5029078750959$, which agrees with the
published Feigenbaum constant to 14 digits (\S\ref{sec:3.9}).

The supremum over the box is taken \textbf{in interval arithmetic}: the box
is split into eight subintervals, and $\ell$ and the coefficients are
substituted as intervals --- sampling at nodes is inadmissible here, since
between the nodes $|P(1)|$ and $\|P'\|$ are not bounded by anything. The
maximum over all 46 boxes equals $\kappa_G\le4.14$ (in fact $4.130319$,
worst box \texttt{g43}); along the way $\min|P(1)|\ge0.398720225464$ (the
number is rounded DOWN), whence over the ball (subtracting
$\rho_X=1.75\cdot10^{-7}$) $|G(1)|\ge0.398720050464$, which is consistent
with (a5). Hence $\kappa_G\cdot1.75\cdot10^{-7}=7.23\cdot10^{-7}$ against the
margin $0.621$ --- with a margin factor of $8.5\cdot10^{5}$. The transfer
from the centre to the solution is thereby justified by a number, not by an
order-of-magnitude estimate.

The value has been cross-checked against a second, independently written
computation of the same constant (\texttt{04-\allowbreak CODE/\allowbreak layer/\allowbreak layer\_\allowbreak kappa.py},
\S\ref{sec:10.7}): under the same coarse subdivision $K=1$ they give $4.415912$ and
$4.4197$ on \texttt{g43} (a discrepancy of $0.09\,\%$), while the published
$4.130319$ is the same computation at $K=8$. The sample boxes of the table in
\S\ref{sec:10.7} agree as well: \texttt{gtop} $4.269531$ against $4.2709$, \texttt{g42}
$4.371065$ against $4.3748$; the worst box \texttt{g43} is one and the same
for both computations.

In the draft this place carried the estimate ``$\kappa_G$ is a quantity of
order unity'', supported only by the fact that both corrections --- the one
from $\alpha_G$ via $d\alpha_G=dG(1)/G(1)^2$ and the one from the shift of
the composition argument --- are of order $10^{-7}$. That is not enough:
``order unity'' is not a number, and there is nothing by which to check such
a claim. The finite arb bound given above replaces it.

\textbf{(b) In $\ell$ (complex).} For $\ell'$ in a small complex
neighbourhood of $\ell_0$ the coefficients $a_k(\ell')$ and the quantities
$\alpha(\ell')^{-\ell'}=\exp(-\ell'\log\alpha(\ell'))$ are holomorphic in
$\ell'$ (\S\ref{sec:6.3.2}), hence the sum in (C1) is continuous in $\ell'$; for
$|\ell'-\ell_0|$ small enough it remains $<1$. That is, there exists
$\varepsilon_b>0$ for which $|U_{P(\ell')}(t)-1|<1$ for all
$|\ell'-\ell_0|<\varepsilon_b$ and $t\in D(c,r)$. Together with (a) this
yields the holomorphy of $W_G$ in a complex neighbourhood (\S\ref{sec:6.5}).

\textbf{Conclusion.} Error 1 has been eliminated: instead of the unjustified
transfer of the boundary $U_{\min}$ to the disk, an honest majorant (C1)
$|U-1|<1$ is supplied, from which both the absence of zeros and the
separation from the cut follow, and moreover uniformly in $\ell$ and stably
in $G$. The appeal to the minimum modulus principle in the statement of
Theorem C is no longer needed.



\subsection{Holomorphy of $F$ in $(G,\ell)$}
\label{sec:6.5}

Fix $\ell_0\in[2,4]$ and set
$\Gamma=\bigl(P(\ell_0)+B_{X_{0,\mathbb C}}(0,R_G)\bigr)\times B_{\mathbb C}(\ell_0,R_\ell)$
for sufficiently small $R_G,R_\ell>0$ --- in $G$ this is an AFFINE ball in
$\mathcal A_{\mathbb C}=P(\ell_0)+X_{0,\mathbb C}$ (the linear ball
$B_{X_{0,\mathbb C}}(0,R_G)$ translated to the centre $P(\ell_0)$), and not a
linear ball in $X_{0,\mathbb C}$: $P(\ell_0)\notin X_{0,\mathbb C}$, since
$P(\ell_0)(0)=1\neq0$. We analyse $F$ block by block; each block is holomorphic
on $\Gamma$.

\begin{enumerate}
\item[1.] \textbf{$G(1)$} is a $\mathbb C$-linear functional. By a5,
   $|P(\ell_0)(1)|\ge0.393$, so for $R_G$ small enough
   $|G(1)|\ge m_1/2>0$ on $P(\ell_0)+B_{X_{0,\mathbb C}}(0,R_G)$, and
   $\alpha_G=-1/G(1)$ is holomorphic in $G$.

\item[2.] \textbf{$\alpha_G^{-\ell}=\exp(-\ell\,\mathrm{Log}\,\alpha_G)$}.
   Since $\alpha_G>0$ is real at the centre and, by item 1, $\alpha_G$ varies
   continuously and holomorphically in $G$, for small $R_G$ we have
   $\mathrm{Re}\,\alpha_G>0$, i.e. $\alpha_G$ does not cross the cut. Then
   $\mathrm{Log}\,\alpha_G$ (the principal branch) and $\alpha_G^{-\ell}$ are
   holomorphic in $(G,\ell)$. The sign condition $\mathrm{Re}\,\alpha_G>0$
   ensures that $|\arg\alpha_G|<\pi/2$ and that the branch of the logarithm does
   not jump.

\item[3.] \textbf{$s_G(t)=t\,\alpha_G^{-\ell}$} is holomorphic in $(t,G,\ell)$.
   The certified $\mathrm{margin\_a}>0$ (a6) means
   $s_{P(\ell)}(D(c,r))\Subset D(c,r)$ for real $\ell$ of the box. This
   carries over to the ball not ``by continuity'' but with a number: the shift of
   the image under the passage $P\to G$ is estimated as
   $|\Delta s|\le(c+r)\,\ell\,|\alpha^{-\ell-1}|\Delta\alpha$ with
   $\Delta\alpha\le\rho_X/(|P(1)|(|P(1)|-\rho_X))$, and for
   $\rho_X=1.75\cdot10^{-7}$ over all 46 boxes
   $|\Delta s|\le3.20\cdot10^{-6}$ against $\mathrm{margin\_a}\ge0.151$
   (a margin of $4.7\cdot10^{4}$, the quotient of the two numbers just printed,
   rounded down). Consequently $s_G(D(c,r))\subset D(c,r)$ for
   all $G$ of the ball, and the composition below is well defined.

   \textbf{In complex \boldmath$\ell$.} As for (C1) and $\kappa_U$: $\Delta s$ is built from a finite number of factors -- $(c+r)$, $\ell$, $|\alpha^{-\ell-1}|$, $\Delta\alpha$ -- each holomorphic in $\ell$ (the same building blocks of \S\ref{sec:6.3.2}). Since $\Delta s(\ell_0)\approx3.2\cdot10^{-6}\ll \mathrm{margin\_a}(\ell_0)\ge0.151$ at the real center, by continuity there exists $\varepsilon>0$ such that the same inequality holds for all $|\ell'-\ell_0|<\varepsilon$ -- by the same device as in \S\ref{sec:6.4.3}(b) for (C1) and in item 5 for $\kappa_U$.

\item[4.] \textbf{$U_G(t)=G(s_G(t))$} is a composition of holomorphic functions, hence holomorphic in $(t,G,\ell)$ (the image of $s_G$ stays in $D(c,r)$ by item 3).

\item[5.] \textbf{$W_G(t)=U_G(t)^{\ell}=\exp(\ell\,\mathrm{Log}\,U_G(t))$}. The branch $\mathrm{Log}\,U_G$ is chosen by restricting the principal branch to the right half-plane -- and it is single-valued on $D(c,r)$ precisely because the entire image of $U_G$ lies there: \S\ref{sec:6.4} (C1) gives $|U_{P(\ell)}(t)-1|<1$ uniformly over $\ell$ in the box and \textbf{stably in \boldmath$G$} (not only at the center $P$, but on the whole ball $P(\ell_0)+B_{X_{0,\mathbb C}}(0,R_G)$, by the same argument as in item 2 for $\alpha_G$) -- that is, $\mathrm{Re}\,U_G>0$, $U_G$ does not cross the branch cut, and $D(c,r)$ is simply connected. (The normalization ``$U_G>0$ on $[0,1]$, since $G>0$ on $[0,1]$'' does not apply here: on this ball $G$ is complex rather than real positive -- real-valuedness is needed only at the center $\ell_0$ as the starting point, from which $\mathrm{Re}\,U_G>0$ extends to the whole ball by continuity via (C1).) Hence $W_G$ is holomorphic in $(t,G,\ell)$ \textbf{pointwise} -- as a function of $t$ taking values in $\mathbb C$ at each fixed $t$.

   \textbf{X-valued holomorphy --- a separate premise.} Further on (item 6, \S\ref{sec:6.5.1}) $W_G$ is used not pointwise, but as an ELEMENT of the Banach algebra $X_{\mathbb C}$ -- through the norm $\|W_G-c\|_X$ and through $D_GW_G$ as a bounded operator on $X_{0,\mathbb C}$. This requires that $\mathrm{Log}\,U_G$ and $W_G=\exp(\ell\,\mathrm{Log}\,U_G)$ THEMSELVES be representable as a power series in $(t-c)$ convergent in the norm $\|\cdot\|_X$ -- not merely pointwise bounded. Scalar (C1) does not give this: $|U_G(t)-1|<1$ at each point does NOT imply $\|U_G-1\|_X<1$ (the $X$-norm is the sum of MODULI of the coefficients, the supremum of values does not majorize it; for a power series the inequality goes the other way: $\sup_D|h|\le\|h\|_X$, not conversely).

   The needed premise is available -- it has already been computed. $X_{\mathbb C}$ is a commutative Banach algebra (\S\ref{sec:6.3.1}); the holomorphic functional calculus in a Banach algebra gives: if $\|h\|_X<1$, then $\mathrm{Log}(1+h)=\sum_{n\ge1}\frac{(-1)^{n+1}}{n}h^n$ converges ABSOLUTELY in $X$ (majorant $\sum\|h\|_X^n/n<\infty$), and the composition with $\exp$ -- by the same device. We apply it at $h=U_G-1$: the quantity $\kappa_U:=\|U_G-1\|_X$ is computed and gated by \texttt{04-\allowbreak CODE/\allowbreak tube\_\allowbreak lib.py::\allowbreak \_\allowbreak compute\_\allowbreak L} (\texttt{if not (kap\_\allowbreak U < arb(1)): return None}) on the SAME ball $P(\ell)+B_{X_{0,\mathbb C}}(0,\rho_X)$ as $\hat q$ in item 6 below -- this is the same chain, whose intermediate result was not previously published as a certificate. Over all 46 boxes: $\kappa_U\le0.363184$ (worst box \texttt{g05}), with a large margin to $1$. It is published in \texttt{03-\allowbreak CERTIFICATES/\allowbreak pointwise/\allowbreak holo\_\allowbreak radius.json::\allowbreak kap\_\allowbreak U\_\allowbreak max} (per box -- the field \texttt{kap\_U}; generator \texttt{04-\allowbreak CODE/\allowbreak tests/\allowbreak \_\allowbreak holo\_\allowbreak radius.py}).

   \textbf{In complex \boldmath$\ell$.} The certificate above was computed on the real $\ell$-box. As for (C1) (\S\ref{sec:6.4.3}(b)), we need continuity in complex $\ell'$: $\kappa_U$ is built from a FINITE number of summands (the degree of $U_{P(\ell)}$ in $t$ is at most $N_p$, the perturbation $U_{\mathrm{diff}}$ is likewise a finite sum of quantities $m_1$, $\|P'\|$, $\overline{\alpha}_{\mathrm{lo}}$), each of which is holomorphic in $\ell$ (the same building blocks of \S\ref{sec:6.3.2} as for $a_k(\ell)$, $\alpha(\ell)^{-\ell}$ in (C1)) -- hence the sum itself is continuous in $\ell'$ WITHOUT any theorem on uniform convergence (which would be needed for an infinite series, but not for a sum of $\le N_p{+}O(1)$ summands). Since $\kappa_U(\ell_0)\le0.363184\ll1$ at the real center, by continuity there exists $\varepsilon_c>0$ such that $\kappa_U(\ell')<1$ for all $|\ell'-\ell_0|<\varepsilon_c$ -- by the same device as in \S\ref{sec:6.4.3}(b) for (C1) and in item 2 of this subsection for $\alpha_G$.

   Separately we need the image: $\mathrm{margin\_c}>0$ (a6) gives $W_{P(\ell)}(D(c,r))\Subset D(c,r)$ for real $\ell$. On the ball this follows already without appeal to continuity -- from $\hat q<1$ of item 6: for $|t-c|\le r$ we have $|W_G(t)-c|\le\sum_k|w_k|r^k=\|W_G-c\|_X\le\hat qr<r$, that is, $W_G(D(c,r))\subset D(c,\hat qr)\Subset D(c,r)$ for all $G$ in the ball at once.

\item[6.] \textbf{$F(G,\ell)(t)=G(t)+\alpha_G\,G(W_G(t))$}. Here a careful argument is needed, not a reference to finite degree: the center $P(\ell)$ is a polynomial of degree $\le N_p$, but an arbitrary element of $X$ is not a polynomial, and ``the coefficients are polynomial expressions in the coefficients of $G$'' does not hold for it.

   Holomorphy gives what is needed: $X$ with norm $\|h\|_X=\sum_k|h_k|r^k$ is a \textbf{Banach algebra} ($\|fg\|_X\le\|f\|_X\|g\|_X$), together with the certified inequality
   \[
   q\;:=\;\frac{\|W_G-c\|_X}{r}\;<\;1 .
   \]
   The quantity $q$ is stored in each certificate (\texttt{box\_\allowbreak g*.json::\allowbreak bases\_\allowbreak q}, flag \texttt{bases\_ok}) and over all 46 boxes does not exceed $0.9911$. Then the composition series converges in the $X$ norm:
   \[
   \|G\circ W_G\|_X\le\sum_k|g_k|\,\|(W_G-c)^k\|_X
   \le\sum_k|g_k|\,(qr)^k\le\sum_k|g_k|\,r^k=\|G\|_X ,
   \]
   and with a margin: a perturbation of $W$ by $\varepsilon<r(1-q)$ (in the worst box $7.1\cdot10^{-3}$) leaves the series convergent. Hence $(G,W)\mapsto G\circ W$ is a convergent power series in $W$ with coefficients linear in $G$, i.e., a holomorphic map on a neighborhood. The remaining links are simpler: $G\mapsto G(1)$ is a bounded linear functional (since $|1-c|=0.4<r$), $\alpha_G=-1/G(1)$ is holomorphic where $G(1)\ne0$ (a5 gives $|G(1)|\ge0.393$), and $s_G$ and $W_G=U_G^{\ell}$ are holomorphic in $(\alpha_G,\ell)$ on the right half-plane, into which the image falls by (C1).

   It is exactly $q<1$, and nothing else, that bounds \emph{this} estimate --- the composition-tail estimate realized through $q$ --- from above in $\ell$: when $q\ge1$ the composition series diverges, and with it the $\hat\beta$-majorant and the Cauchy tails (\S\ref{sec:3.6}) collapse.

   The bound is on that realization only, and not on the method in general: \S\ref{sec:10.4} shows a weaker sufficient condition --- \texttt{margin\_c} $> 0$ --- for a related estimate in lemma 2, where $q\ge1$ is no longer an obstacle.

   \textbf{In complex \boldmath$\ell$.} By the same device as in \S\ref{sec:6.4.3}(b) and in item 5: $\hat q=(\|W_P-c\|_X+\Delta W)/r$ is a finite sum of summands holomorphic in $\ell$, hence continuous in $\ell'$. Since $\hat q(\ell_0)\le 0.955131\ll1$, the margin $1-0.955131$ gives $\varepsilon>0$ with $\hat q(\ell')<1$ for all $|\ell'-\ell_0|<\varepsilon$.
\end{enumerate}

\textbf{Conclusion: \boldmath$R_G$ is quantitative.} All premises of items 1--6 have been verified not at the center, but uniformly on the affine ball $P(\ell)+B_{X_{0,\mathbb C}}(0,\rho_X)$ with $\rho_X=1.75\cdot10^{-7}$ -- the very uniqueness ball of Theorem A on which $L$ is already computed in \S\ref{sec:6.6}. Over all 46 boxes:

\begin{table}[ht]
\centering
\small
\caption{Premises verified uniformly on the affine ball, with the worst values over the boxes.}
\label{tab:ball-premises}
\begin{tabular}{lll}
\toprule
premise & what is verified on the ball & worst value \\
\midrule
item 1 & $|G(1)|\ge |P(1)|-\rho_X>0$ & $\ge0.396279$ \\
item 2 & $|\arg\alpha_G|\le\arcsin(\rho_X/|P(1)|)<\pi/2$ & $\le4.42\cdot10^{-7}$ \\
item 3 & $|\Delta s|<\mathrm{margin\_a}$ & $3.20\cdot10^{-6}$ against $0.151$ \\
item 5 ($\kappa_U$, a separate premise) & $\kappa_U=\norm{U_G-1}_X<1$ --- $X$-valued holomorphy of $W_G$ & $\le0.363184$ \\
items 5--6 & $\hat q=(\|W_P-c\|_X+\Delta W)/r<1$ & $\le0.955131$ \\
\bottomrule
\end{tabular}
\end{table}

The quantity $\hat q$ is not the same as $q$ in item 6, although both are obtained by the same function \texttt{tube\_\allowbreak lib.\_\allowbreak compute\_\allowbreak L}, only with different arguments: $q$ -- on the box as a whole, without correction for the shift over the ball; $\hat q$ in this table -- from \texttt{04-\allowbreak CODE/\allowbreak tests/\allowbreak \_\allowbreak holo\_\allowbreak radius.py}, on $K=2$ (fixed, not varied) narrower sub-intervals in $\ell$, with the addition $\Delta W$ on the ball $B_X(\hat a,\rho_X)$ of holomorphy radius and fixed $\rho'=10^{-7}$. Hence $\hat q\ne q$ numerically: on 39 of 46 boxes $\hat q<q$ (worst discrepancy $-0.1127$ on \texttt{g36}), at the upper edge the discrepancy is at most $5\cdot10^{-6}$ (the largest, $+4.61\cdot10^{-6}$, on \texttt{g43}). The comparison is box by box between \texttt{03-\allowbreak CERTIFICATES/\allowbreak pointwise/\allowbreak holo\_\allowbreak radius.json} (the field \texttt{q\_\allowbreak hat} of each element of \texttt{boxes}) and \texttt{03-\allowbreak CERTIFICATES/\allowbreak existence/\allowbreak box\_\allowbreak g*.json\allowbreak ::\allowbreak bases\_\allowbreak q} (the per-box convention of \S\ref{sec:6.2.a}). The box in the certificate gates exactly $q$ (\texttt{bases\_q}, \texttt{bases\_ok}) -- a rigorous check \texttt{beta\_\allowbreak bound\_\allowbreak lib.bases\_\allowbreak check} before the Lipschitz chain starts; \texttt{tube\_\allowbreak lib.\_\allowbreak compute\_\allowbreak L} is called once more inside the certificate gate itself, on sub-boxes of different width and with $\rho'$ chosen by enumeration -- this is also a $\hat q$-like quantity, nowhere published separately, and its failure there only discards the current enumeration point, not the box.

Therefore we may take $R_G=1.75\cdot10^{-7}$, and then
\[
\|\Gs(\ell_0)-P(\ell_0)\|_X\le\rho^{\star}_X\le8.49\cdot10^{-13}
\;<\;1.75\cdot10^{-7}=R_G
\]
with a margin of $2.06\cdot10^{5}$, i.e., $\Gs(\ell_0)\in\Gamma$, and
the analytic implicit function theorem (IFT, \S\ref{sec:6.6}) applies legitimately at the point $(\Gs(\ell_0),\ell_0)$.

The quantity $R_\ell$ remains existential, and this is not a gap: $\ell_0$ is the \textbf{center} of the complex disk in $\ell$, nothing stands off it, and there is nothing to compare $R_\ell$ with. The quantification was required precisely for the radius in $G$, because $\Gs(\ell_0)$ does stand off the center $P(\ell_0)$.

The draft at this place carried a comparison $\|\Gs(\ell_0)-P(\ell_0)\|_X\le8.49\cdot10^{-13}\ll R_G$ without any lower estimate on $R_G$, and the next revision -- a claim that no lower estimate is needed at all, since the IFT itself builds a branch inside $\Gamma$. Both variants are wrong: the first compares a number with a quantity introduced by an existential quantifier; the second presupposes what is to be proved -- to apply the IFT at $\Gs$, one must first know that $F$ is holomorphic in a neighborhood of $\Gs$. The replacement is Table~\ref{tab:ball-premises}.

\subsubsection{Holomorphy of $D_GF$ and $DT$}
\label{sec:6.5.1}

The differential
\[
D_GF(G,\ell)\,h
=h+\alpha_G\,h(W_G)+\tfrac{\partial\alpha_G}{\partial G}h\cdot G(W_G)
+\alpha_G\,G'(W_G)\cdot D_GW_G(G,\ell)\,h
\]
is a sum of compositions of holomorphic (linear/bilinear) operators, hence $D_GF$ is holomorphic in $(G,\ell)$ with values in $\mathcal B(X_{0,\mathbb C})$. Since $F=I-T$, we have $DT=I-D_GF$; in particular $\ell\mapsto DT(\Gs(\ell),\ell)=I-D_GF(\Gs(\ell),\ell)$ is holomorphic in $\ell$ (composition of the holomorphic maps $\ell\mapsto \Gs(\ell)$ and $(G,\ell)\mapsto D_GF$), which is what is required in \S\ref{sec:6.8}.


\subsection{Invertibility of $D_GF$ and the analytic implicit function theorem}
\label{sec:6.6}

\subsubsection{Invertibility at the center}
\label{sec:6.6.1}

From (a2): $\|I-A(\ell_0)M'(\ell_0)\|_1\le Z_1(\ell_0)<1$, where $M'(\ell_0)=D_GF(P(\ell_0),\ell_0)$. The Neumann series gives
\[
\|D_GF(P(\ell_0),\ell_0)^{-1}\|_1
\le\frac{\|A(\ell_0)\|_1}{1-Z_1(\ell_0)}
\le\max_{b}\frac{B_A^{(b)}}{1-Z_1^{(b)}}
\le4.24.
\]
The maximum is over the boxes, each with \textbf{its own} pair $(B_A^{(b)},Z_1^{(b)})$:
it equals $4.2372$ and is attained at \texttt{g43}. It must not be assembled from
the two global maxima, which sit at different boxes --- $B_A\le 3.6162$ at
\texttt{gtop} and $Z_1\le 0.1501$ at \texttt{g43}: their quotient, ${\approx}\,4.255$,
is a strictly weaker bound, attained on no box (\S\ref{sec:6.6.2} makes the
same point for the transferred estimate).

\subsubsection{Transfer to the solution}
\label{sec:6.6.2}

By (a4) $\|DT(\Gs(\ell_0))-DT(P(\ell_0))\|_X\le L\,\|\Gs(\ell_0)-P(\ell_0)\|_X$. The solution $\Gs(\ell_0)$ lies inside the uniqueness ball of radius $1.75\cdot10^{-7}$ (a1), i.e., $\|\Gs(\ell_0)-P(\ell_0)\|_X\le1.75\cdot10^{-7}$ (in fact it lies in the ball of radius $1.75\rho^{\star}\le8.49\cdot10^{-13}$, but conservatively we take the uniqueness ball). Then $\|DT(\Gs)-DT(P)\|_X\le L\cdot1.75\cdot10^{-7}$, which over all 46 boxes does not exceed $4.2518\cdot10^4\cdot1.75\cdot10^{-7}<7.5\cdot10^{-3}$.

The bound of \S\ref{sec:6.6.1} is in the column norm $\|\cdot\|_1$ and this one is in $\|\cdot\|_X$, so the two may not be multiplied as they stand. On the slice $\|v\|_1\le\|v\|_X\le1.75\|v\|_1$ (\S\ref{sec:3.1}), whence $\|T\|_{1\to1}\le1.75\,\|T\|_{X\to X}$ for every operator: the input vector costs a factor $1.75$ and the output costs nothing. Carrying it,
$\|D_GF(\Gs)-D_GF(P)\|_1\le1.75\cdot7.5\cdot10^{-3}<1.32\cdot10^{-2}$, and hence
\[
D_GF(\Gs(\ell_0),\ell_0)
=D_GF(P(\ell_0),\ell_0)\bigl(I+\underbrace{D_GF(P)^{-1}\cdot(D_GF(\Gs)-D_GF(P))}_{\|\cdot\|_1\ \le4.24\cdot1.32\cdot10^{-2}<5.6\cdot10^{-2}<1}\bigr)
\]
is invertible, and
\[
\|D_GF(\Gs(\ell_0),\ell_0)^{-1}\|_1
\le\frac{B_A/(1-Z_1)}{1-\bigl(B_A/(1-Z_1)\bigr)\,1.75\,L\cdot1.75\cdot10^{-7}}.
\]

The estimate is taken \textbf{for each box with its own \boldmath$B_A,Z_1,L$}; the maximum over all 46 boxes is $4.3414<4.35$ (attained at \texttt{g43}). One must not combine global maxima of different quantities here -- they are attained at different boxes, and such a combination would yield a strictly weaker number.

The margin is huge: even the rough estimate $L\cdot1.75\cdot10^{-7}$ gives a perturbation $\sim10^{-3}$, whereas the invertibility threshold is $\sim1$.

\subsubsection{Analytic implicit function theorem (IFT)}
\label{sec:6.6.3}

\begin{proposition}[Complex analytic IFT]
Let $U,V$ be complex Banach spaces, $\Omega\subset U\times V$ open, $F:\Omega\to U$ holomorphic, $F(x_0,y_0)=0$ and $D_xF(x_0,y_0)$ a linear homeomorphism. Then there exist a neighborhood $\mathcal U'\subset V$ of $y_0$ and a neighborhood $\mathcal V\subset U$ of $x_0$ and a holomorphic $\tilde x:\mathcal U'\to\mathcal V$, $\tilde x(y_0)=x_0$, such that for $y\in\mathcal U'$ the equation $F(\tilde x(y),y)=0$ has in $\mathcal V$ a unique solution $\tilde x(y)$.
\end{proposition}

This is the standard complex analog of the implicit function theorem (its real version is in any text on functional analysis; the holomorphic variant is obtained by the same contraction $x\mapsto x-D_xF(x_0,y_0)^{-1}F(x,y)$, which is holomorphic in $y$ and applies the contraction mapping principle layer by layer; holomorphy of the limit follows from the uniform convergence of holomorphic iterations). The proposition is a theorem about \textbf{linear} Banach spaces $U,V$; we apply it not to $F$ directly (its domain $\mathcal A_{\mathbb C}\times\mathbb C$ in $G$ is affine, not linear), but to the shift by the constant $\widehat F(h,\ell):=F(1+h,\ell)$, $h\in X_{0,\mathbb C}$ -- shifting by the fixed element $1\in X$ is holomorphic (an affine change of variable) and does not change either the holomorphy of $F$ or $D_GF=D_h\widehat F$ (the directional derivative is invariant under a constant shift of the base point). The hypotheses are verified for $\widehat F$: $U=X_{0,\mathbb C}$, $V=\mathbb C$, $\widehat F$ is holomorphic wherever $F$ is (\S\ref{sec:6.5}), invertibility of $D_h\widehat F=D_GF$ from \S\ref{sec:6.6.2}.

\begin{corollary}[Local branches at real points]
Let $\ell^{\ast}\in[2,4]$ and $h^{\star}(\ell):=\Gs(\ell)-1\in X_{0,\mathbb C}$ be the coordinates of the solution relative to the normalization center (not to be confused with the coordinates relative to $P(\ell)$ from \appref{sec:5} and \S\ref{sec:6.2}). The hypotheses of the proposition hold for $\widehat F$ at the point $(h^{\star}(\ell^{\ast}),\ell^{\ast})$: $h^{\star}(\ell^{\ast})\in X_{0,\mathbb C}$ by construction (unlike $\Gs(\ell^{\ast})\in\mathcal A_{\mathbb C}$ itself, which does not lie in $X_{0,\mathbb C}$ -- $\Gs(0)=1\ne0$), $\widehat F$ is holomorphic in its neighborhood (\S\ref{sec:6.5}), and $D_h\widehat F=D_GF$ is invertible with a bound $\|D_GF^{-1}\|_1\le4.35$ uniform over the entire segment (\S\ref{sec:6.6.2}). Hence there exist an open disk $D_{\ell^{\ast}}\subset\mathbb C$ centered at a \textbf{real} point $\ell^{\ast}$, a ball $W_{\ell^{\ast}}\subset X_{0,\mathbb C}$ around $h^{\star}(\ell^{\ast})$, and a holomorphic $\widehat\Phi_{\ell^{\ast}}:D_{\ell^{\ast}}\to W_{\ell^{\ast}}$ with $\widehat F(\widehat\Phi_{\ell^{\ast}}(\ell),\ell)=0$, and in $W_{\ell^{\ast}}$ this is the unique solution. Returning to $G$-coordinates: $\Phi_{\ell^{\ast}}:=1+\widehat\Phi_{\ell^{\ast}}:D_{\ell^{\ast}}\to V_{\ell^{\ast}}:=1+W_{\ell^{\ast}}\subset\mathcal A_{\mathbb C}$ is holomorphic (shift by a constant), $F(\Phi_{\ell^{\ast}}(\ell),\ell)=0$, and in $V_{\ell^{\ast}}$ this is the unique solution -- the same fact as for $\widehat F$, transported back by the shift.
\end{corollary}

The proposition does not estimate the disk radius, and we use it nowhere -- what matters is only that it is positive. Additionally we \textbf{shrink} the disk so that three conditions hold (each achievable by shrinking the radius):

\begin{itemize}
\item (i) $D_{\ell^{\ast}}\cap\mathbb R\cap[2,4]$ lies inside a single (closed) existence box $b$. This is possible for any point of the segment, including the endpoints $\ell=2,4$ (at $\ell=2$ itself $D_{\ell^{\ast}}\cap\mathbb R$ does not fit strictly inside the box \texttt{g00}$=[2.0,2.05]$ -- that is an endpoint, not an interior point; intersection with $[2,4]$ suffices, since step 2 does not require $D_{\ell^{\ast}}\cap\mathbb R$ to lie entirely inside $[2,4]$): 46 boxes cover $[2,4]$ and neighboring boxes overlap on a non-empty open interval (\texttt{glue.json} carries the 46 box records and the strict verdict \texttt{seams\_ok}; the 45 seams themselves are computed from those records by \texttt{glue.py}, which prints them rather than depositing a table --- 42 of the 45 seams join neighbouring boxes of width $0.05$ spaced at a step of $0.045$, which therefore overlap by $0.005$ in $\ell$, that is by $0.1$ of a width; the other three overlap by more --- \texttt{g21}/\texttt{ganchor} by $0.020$, \texttt{ganchor}/\texttt{g22} by $0.035$, and \texttt{g43}/\texttt{gtop} by $0.0075$, which is $0.25$ of the width $0.03$ of \texttt{gtop} --- so $0.005$ is the minimum over all 45, not the value at every seam), so every point of the segment is interior to at least one box (taking into account intersection with $[2,4]$ at the endpoints);
\item (ii) $\Phi_{\ell^{\ast}}(D_{\ell^{\ast}})\subset B_X(\Gs(\ell^{\ast}),\eta)$ with $\eta:=8\cdot10^{-8}$ --- a condition on the \textbf{image} of the disk, not on the ball $V_{\ell^{\ast}}$ of the corollary. The distinction matters: $V_{\ell^{\ast}}$ lives in the function space and does not shrink when the disk does, whereas the image contracts to the single point $\Phi_{\ell^{\ast}}(\ell^{\ast})=1+h^{\star}(\ell^{\ast})=\Gs(\ell^{\ast})$, because $\Phi_{\ell^{\ast}}$ is holomorphic and hence continuous at $\ell^{\ast}$. Nothing below needs more: \S\ref{sec:6.7}, step~1 uses only that the \emph{value} $\Phi_{\ell^{\ast}}(\ell)$ lies in that ball;
\item (iii) $\|P(\ell)-P(\ell^{\ast})\|_X<\eta$ for real $\ell\in D_{\ell^{\ast}}$. This is possible, since inside the box $P(\ell)$ is a polynomial in $\ell$ with coefficients from $X$, hence continuous in the $X$-norm (quantitatively -- the quantity \texttt{deltaP} of the certificate).
\end{itemize}

Conditions (i)--(iii) are compatible with the SAME shrinkage radius: (i) depends only on the position of $D_{\ell^{\ast}}$ relative to the box grid (a discrete condition, independent of $\eta$), while (ii) and (iii) both tend monotonically to $0$ as the disk shrinks --- (ii) because the image $\Phi_{\ell^{\ast}}(D_{\ell^{\ast}})$ contracts to the point $\Gs(\ell^{\ast})$, (iii) by continuity of $P$ --- hence the minimum of the three radii found individually, all three being radii of one and the same disk $D_{\ell^{\ast}}\subset\CC$, works for all three at once.


\subsection{Gluing branches over $[2,4]$}
\label{sec:6.7}

\textbf{Step 1 (the branches coincide with \boldmath$\Gs$ on the real trace).} Let $\ell\in D_{\ell^{\ast}}\cap\mathbb R\cap[2,4]$. Then
\begin{align*}
\|\Phi_{\ell^{\ast}}(\ell)-P(\ell)\|_X
&\le\underbrace{\|\Phi_{\ell^{\ast}}(\ell)-\Gs(\ell^{\ast})\|_X}_{<\,\eta\ \text{by (ii)}}
+\underbrace{\|\Gs(\ell^{\ast})-P(\ell^{\ast})\|_X}_{\le\,8.49\cdot10^{-13}\ \text{(a1)}}\\
&\quad+\underbrace{\|P(\ell^{\ast})-P(\ell)\|_X}_{<\,\eta\ \text{by (iii)}}
<1.7\cdot10^{-7},
\end{align*}
that is, $\Phi_{\ell^{\ast}}(\ell)$ lies in the uniqueness ball of Theorem A of radius $1.75\cdot10^{-7}$ around $P(\ell)$. Uniqueness here is needed not on the real slice, but on the entire COMPLEXIFIED ball $P(\ell)+B_{X_{0,\mathbb C}}(0,1.75\cdot10^{-7})$ -- $\Phi_{\ell^{\ast}}(\ell)$ is not required a priori to be real-valued; real-valuedness is what this step must establish, not what can be assumed in it; the transfer of uniqueness of Theorem A to this ball is justified in \S\ref{sec:6.3.1}: $Y_0$ is a number computed at a fixed real center, and the derivation of $L$ in \texttt{tube\_\allowbreak lib.\_\allowbreak compute\_\allowbreak L} nowhere appeals to the real-valuedness of the increment, so both Kantorovich gate quantities transfer to $X_{0,\mathbb C}$ verbatim. The domain on which $L$ is certified is an X-ball of radius $3.0625\rho'+dP$ around the sub-box center (\S\ref{sec:3.1}), and it contains the uniqueness ball $B(P(\ell),1.75\rho')$, $1.75\rho'=1.75\cdot10^{-7}$. In this ball the solution is unique, and $\Phi_{\ell^{\ast}}(\ell)$ is a solution. Hence $\Phi_{\ell^{\ast}}(\ell)=\Gs(\ell)$.

\textbf{Step 2 (finite subcover).} The sets $I_{\ell^{\ast}}:=D_{\ell^{\ast}}\cap\mathbb R$ are open in $\mathbb R$ and contain $\ell^{\ast}$, so $\{I_{\ell^{\ast}}\}_{\ell^{\ast}\in[2,4]}$ is an open cover of the compact set $[2,4]$. By Heine--Borel we choose a finite subcover with centers $\ell_1,\dots,\ell_n\in[2,4]$; we denote the corresponding disks by $D_1,\dots,D_n$ and the branches by $\Phi_1,\dots,\Phi_n$.

\textbf{Step 3 (any two intersecting disks share a real point).} Here we use the fact that all centers are real. Let $D_i=D(\ell_i,r_i)$, $D_j=D(\ell_j,r_j)$ intersect, i.e., $|\ell_i-\ell_j|<r_i+r_j$. The point
\[
\mu:=\frac{\ell_i r_j+\ell_j r_i}{r_i+r_j}\in\mathbb R
\]
lies between $\ell_i$ and $\ell_j$, hence $\mu\in[2,4]$, and at the same time $|\mu-\ell_i|=r_i|\ell_i-\ell_j|/(r_i+r_j)<r_i$, similarly $|\mu-\ell_j|<r_j$. Hence $\mu\in D_i\cap D_j\cap\mathbb R$, and this intersection is an open set containing $\mu$, i.e., it contains a whole real interval around $\mu$.

\textbf{Step 4 (coincidence on the intersection).} On this interval both branches equal $\Gs$ by step 1, hence they equal each other. The intersection of two disks is convex (intersection of convex sets), hence connected and open, and the coincidence set has an accumulation point in it. By the uniqueness theorem for holomorphic functions with values in a Banach space (it reduces to the coordinate-wise version: $\ell\mapsto(\Phi_i-\Phi_j)(t_k)$ is holomorphic in $\mathbb C$ and vanishes on an interval) we obtain $\Phi_i\equiv\Phi_j$ on all of $D_i\cap D_j$.

\textbf{Step 5 (conclusion).} Hence the formula $\widetilde\Phi(\ell):=\Phi_j(\ell)$ for any $j$ with $\ell\in D_j$ correctly defines a holomorphic function on the open set $\mathcal U:=\bigcup_{j=1}^{n}D_j\supset[2,4]$, extending $\Phi$: by step 1 $\widetilde\Phi(\ell)=\Gs(\ell)$ for real $\ell\in[2,4]$. Since $\mathcal U$ contains an open complex neighborhood of the segment, $\Phi=\widetilde\Phi|_{[2,4]}$ is real-analytic on $[2,4]$.

\emph{A remark on what is NOT being done here.} Branches are constructed at real
points, not ``by boxes,'' and the overlap of existence boxes is used only in
condition (i) --- so that for every point of the segment there is a box
containing it together with some neighborhood \textbf{within \boldmath$[2,4]$} (at the
endpoints $\ell = 2, 4$ this is a one-sided neighborhood: the point 2 itself is
the left edge of box \texttt{g00}, and condition (i) is formulated in \S\ref{sec:6.6.3}
precisely via the intersection with $[2,4]$). Overlap of the boxes does NOT
imply overlap of the complex IFT-neighborhoods: the proposition does not
estimate their radii, and such an inference would be incorrect. Non-emptiness
of the real intersection is taken from step 3, i.e., from the realness of the
centers, not from anything about the boxes.


\subsection{Analyticity of $\delta(\ell)$ (correction of error 2)}
\label{sec:6.8}

\subsubsection{The draft's incorrect contour}
\label{sec:6.8.1}

The draft took the Riesz projector along the contour $|z-\delta|=\rho$ with
$\rho<0.004$, justifying this by the claim that ``$1-0.9919=0.0081$ is the gap
between $\delta(\ell)$ and the unit circle.'' This is a misreading: $0.9919$
is the value of the full test of \S\ref{sec:4.6} (the product in the Rouché inequality
after passing from the projection to the operator; $R\,\|\Delta\|$ itself
does not exceed $0.9896$), i.e. the margin for transferring the eigenvalue
\emph{count}, not a spectral gap. In fact $\delta(3)=6.0847$, roughly five units
away from the unit circle, but no certificate \textbf{states} that there is no
other eigenvalue near $\delta$ within radius $0.004$ (the second-largest in
modulus, $\lambda_2\approx0.24$, is inside the disk, but on the opposite side
from $\delta$). Relying on such a contour is unjustified.

\subsubsection{The correct contour --- the unit circle}
\label{sec:6.8.2}

What correctly separates $\delta$ from the rest of the spectrum is precisely
$\Gamma=\{|z|=1\}$: by Theorem B (a8), $\delta(\ell_0)$ lies outside
$\Gamma$, the entire remaining spectrum lies strictly inside, and $\Gamma$
contains no spectrum. Set $T(\ell)=DT(\Gs(\ell))|_{X_{0,\mathbb C}}$
(the operator at the TRUE solution, not at the center $P(\ell)$). Then on
$\Gamma$ the resolvent $R(z,\ell_0)=(zI-T(\ell_0))^{-1}$ exists and is
bounded: this follows from a8 (the spectrum of $T(\ell_0)$ does not touch
$\Gamma$, the resolvent is continuous away from the spectrum, and $\Gamma$ is
compact) --- and NOT from the numerical $R\le3.7961$ in a7, which was proved
for a different operator, $DT_0=DT(P(\ell_0),\ell_0)$ at the box center
(\S\ref{sec:4.3}). This distinction does not affect the conclusion: below we use only
the finiteness of $R_{\mathrm{full}}=\sup_{|z|=1}\|(zI-T(\ell_0))^{-1}\|$ and
the holomorphy of $T(\ell)$ in $\ell$ --- the concrete numerical value of
$R_{\mathrm{full}}$ is never substituted below. Define the Riesz projector
of the ``inner'' part:
\[
P_{\mathrm{in}}(\ell)=\frac1{2\pi i}\oint_{|z|=1}(zI-T(\ell))^{-1}\,dz
\]
($\Gamma$ oriented counterclockwise). At $\ell=\ell_0$ the integral exists as
a Bochner integral in $\mathcal B(X_{0,\mathbb C})$ (the integrand is
continuous and bounded on the compact $\Gamma$ by a8). For complex $\ell$ in
a small neighborhood of $\ell_0$:
\[
(zI-T(\ell))^{-1}=(zI-T(\ell_0))^{-1}\sum_{n\ge0}\bigl[(T(\ell)-T(\ell_0))(zI-T(\ell_0))^{-1}\bigr]^n
\]
converges uniformly in $z\in\Gamma$ provided
$\|T(\ell)-T(\ell_0)\|\cdot\sup_{|z|=1}\|(zI-T(\ell_0))^{-1}\|<1$. Since
$T(\ell)$ is holomorphic in $\ell$ (\S\ref{sec:6.5.1}), $\|T(\ell)-T(\ell_0)\|\le
C_T|\ell-\ell_0|$, and for $|\ell-\ell_0|<1/(C_T\,R_{\mathrm{full}})$ (where
$R_{\mathrm{full}}=\sup_{|z|=1}\|(zI-T(\ell_0))^{-1}\|<\infty$) the resolvent
exists. Hence $P_{\mathrm{in}}(\ell)$ is holomorphic in $\ell$ in this
neighborhood (differentiation under the integral sign).

The complementary projector
\[
P_{\mathrm{out}}(\ell):=I-P_{\mathrm{in}}(\ell)
\]
--- the projector onto the spectral subspace of the ``outer'' part of the
spectrum. At $\ell=\ell_0$ this is exactly the root subspace of the isolated
eigenvalue $\delta(\ell_0)$; since $\delta(\ell_0)$ is simple (a8),
$\mathrm{rank}\,P_{\mathrm{out}}(\ell_0)=1$. The rank of a finite-dimensional
projector depending continuously on a parameter is locally constant, so
$\mathrm{rank}\,P_{\mathrm{out}}(\ell)=1$ for $\ell$ in a neighborhood.

Finally,
\begin{equation*}
\delta(\ell)=\mathrm{tr}\bigl(T(\ell)P_{\mathrm{out}}(\ell)\bigr).
\tag{C2}
\end{equation*}

Indeed, $T(\ell)P_{\mathrm{out}}(\ell)$ is a rank-1 operator, and its unique
nonzero eigenvalue is $\delta(\ell)$ (on $\mathrm{ran}\,P_{\mathrm{out}}(\ell)$
the operator $T(\ell)$ acts as multiplication by $\delta(\ell)$). Since
$T(\ell)$ and $P_{\mathrm{out}}(\ell)$ are holomorphic in $\ell$, the rank-1
operator-function $T(\ell)P_{\mathrm{out}}(\ell)$ is holomorphic, and the
trace of a holomorphic finite-rank family is holomorphic (explicitly:
$P_{\mathrm{out}}(\ell)=v(\ell)\otimes\varphi(\ell)$ with holomorphic
$v(\ell)\in X_{0,\mathbb C}$, $\varphi(\ell)\in X_{0,\mathbb C}^{\star}$, and
$\delta(\ell)=\varphi(\ell)\bigl(T(\ell)v(\ell)\bigr)$ --- a holomorphic
function). This is standard perturbation theory of an isolated eigenvalue
(see Kato \cite{kato1976perturbation},
Ch. III \S6.4 (``Separation of the spectrum''), \S6.5
(``Isolated eigenvalues'') and \S4.3 (``Degenerate operators. The trace
and determinant''), together with Ch. VII \S1.3 (``Separation of the
spectrum and finite systems of eigenvalues''); we do not reproduce the exact theorem
number --- see \S\ref{sec:6.11}).

\textbf{Realness and agreement with Theorem B.} For real $\ell$ the operator
$T(\ell)$ is real, so $\delta(\ell)$ from (C2) is real (a non-real value
would come paired with its conjugate, giving an outer count $\ge2$ --- a
contradiction with a8). For real $\ell$, the value (C2) coincides with
$\delta(\ell)$ of Theorem B (both are the unique eigenvalue outside the
disk). Hence $\widetilde\Delta:=\delta$ from (C2) is the analytic
continuation of the function $\Delta$ of Theorem C, and $\Delta$ is real-analytic on $[2,4]$. The gluing of $\widetilde\Delta$ over the whole of $\bigcup_{\ell_0}\mathcal U_{\ell_0}$ promised in \S\ref{sec:6.1} repeats steps 2--5 of \S\ref{sec:6.7} verbatim and needs nothing new: the local neighbourhoods are disks with \textbf{real} centers $\ell_0\in[2,4]$, two intersecting such disks share a real interval, on it both local branches equal the $\delta(\ell)$ of Theorem~B (the paragraph ``Realness and agreement'' above), and the uniqueness theorem for holomorphic functions on the connected intersection does the rest.


\subsection{The endpoints of the segment $[2,4]$}
\label{sec:6.9}

The point $\ell=2$ is the left edge of box \texttt{g00} ($[2.0,2.05]$), $\ell=4$ is
an interior point of box \texttt{gtop} ($[3.9775,4.0075]$); for hyperbolicity
$[2.0000,4.0012]\supset[2,4]$ (a1, a7). For the analytic IFT it does not
matter whether $\ell_0$ is an interior point of a \emph{box}: what is needed is
only the holomorphy of $F$ near $(\Gs(\ell_0),\ell_0)$ and the
invertibility of $D_GF$ there (\S\ref{sec:6.5}, \S\ref{sec:6.6}) --- both hold at the points $2$ and
$4$, since these are points of the segment $[2,4]$ where Theorems A and B
apply. Consequently the IFT yields a complex neighborhood $\mathcal U_2\ni2$
and $\mathcal U_4\ni4$, and the branches extend past the endpoints (into
$\mathcal U_2,\mathcal U_4$). Hence the statement is for the \textbf{closed
segment \boldmath$[2,4]$}, not $(2,4)$. To avoid confusion, let us stress
separately: the extension past the endpoints is a bare IFT consequence of
the holomorphy of $F$ exactly at the points $2,4$, nothing more; it does NOT
extend the certified region (the box covering ends at $4.0075$/$4.0012$) and
asserts nothing about the uniqueness of this branch among possible
continuations, nor about hyperbolicity, nor about any connection to the
global picture outside $[2,4]$.


\subsection{Consequences}
\label{sec:6.10}

\begin{enumerate}
\item[1.]
\textbf{Smoothness implies everything needed for this section.} $\Phi,\Delta$ are
   infinitely differentiable, with Taylor series converging at every point
   of $[2,4]$. In particular $\delta'(\ell)$ is given by the standard
   formula from the perturbation theory of an isolated eigenvalue (Rellich;
   in physics, Hellmann--Feynman)
   $\delta'(\ell)=\varphi(\ell)\,(D_\ell T(\ell))\,v(\ell)$ (the derivative
   of an isolated eigenvalue with respect to a parameter), where $D_\ell T$
   is the derivative of $DT$ with respect to $\ell$, computed via the
   differential of $F$ with respect to $\ell$ (\S\ref{sec:6.5.1}). The observations of
   \S\ref{sec:7.3} (strict monotonicity and constant sign of convexity of
   $\delta,\alpha,|\lambda_2|$ over 101 points) can now be raised from the
   status of an ``observation'' to a ``provable statement'' --- provided that
   $|\delta'(\ell)|\ge c>0$ and the sign of $\delta''$ are certified
   intervally; analyticity alone does not give these.

\item[2.] \textbf{\boldmath$\delta$ spreads apart the class constants.} $\delta(2)=4.669\ldots$
   and $\delta(4)=7.284\ldots$ (\S\ref{sec:7.1}) are values of one and the same
   real-analytic function; statements of the form ``$\delta=4.6692\ldots$''
   are accurate only with the qualification $\ell=2$.

\item[3.] \textbf{Methodologically}, Theorem C requires no new boxes: all of its
   premises are quantities already certified by Theorems A and B, plus a
   finite interval sum (C1) over the center's coefficients.
\end{enumerate}


\subsection{What is not proved in this section}
\label{sec:6.11}

\begin{itemize}
\item \textbf{Continuation beyond \boldmath$[2,4]$.} As a global statement, this is not proved
  (the certificate covering ends at $4.0075$/$4.0012$; nothing is asserted
  about existence or hyperbolicity outside $[2,4]$). The local branch on
  $\mathcal U_2,\mathcal U_4$ is an analytic continuation, but its status as
  a ``solution of the Feigenbaum--Cvitanović equation'' outside $[2,4]$ is not
  certified.
\item \textbf{The layer is fixed} ($c=3/5$, $r=4/5$): like Theorems A and B, the
  entire statement pertains to the operator on this space of analytic
  functions. Independence from the choice of $(c,r)$ within the
  \textbf{exhibited} family $\Ffam$ is proved separately (Theorem D, \S\ref{sec:10};
  what limits the result is the size of $\Ffam$, not gaps inside it --- \S\ref{sec:10.9}); in
  particular, the analyticity of $\delta(\ell)$ in $\ell$ is
  automatically independent of the layer, since this is a statement about the
  same function $\delta(\ell)$ (\S\ref{sec:10.5}, Consequence 2).
\item \textbf{Kato --- without a theorem number.} We deliberately do not reproduce the
  number (Ch. IV/VII, etc.): in the draft the references to Kato were
  blurred and partly invented. Here we use only the general fact that a
  simple isolated eigenvalue of a holomorphic family of bounded operators
  depends holomorphically on the parameter, with a complete treatment of
  its conditions in our (non-Hilbert, weighted $\ell^1$) space: the Riesz
  projector $P_{\mathrm{in}}$ is defined as a Bochner integral over $\Gamma$
  and is holomorphic, and the finite-rank condition is provided by Theorem
  B, not by compactness. The reference is deliberately given only down to
  chapter and section (Ch. III \S6.4 (``Separation of the spectrum''), \S6.5
  (``Isolated eigenvalues'') and \S4.3 (``Degenerate operators. The trace
  and determinant''), together with Ch. VII \S1.3 (``Separation of the
  spectrum and finite systems of eigenvalues'')), without a theorem number: a number we had not
  checked against the book ourselves would look more precise than it
  actually is --- exactly the defect present in the draft. (In the book's own
  terminology: Ch. I \S5 is titled ``The eigenvalue problem'' --- ``Riesz
  projectors'' there is not the heading but our paraphrase; the projector
  $P_{\mathrm{in}}$ is there called the eigenprojection --- ``eigenprojection,''
  not ``Riesz projection,'' which does not occur in the book at all; Ch. II \S1
  is titled ``Analytic perturbation of eigenvalues'' --- this is the chapter's
  first section, not the whole chapter. Earlier versions of this work cited
  exactly those two chapters. That address is wrong, and wrong in a way
  worth naming: Ch. I is ``Operator theory in finite-dimensional vector
  spaces'' and Ch. II is ``Perturbation theory in a finite-dimensional
  space'', whereas $X_{0,\mathbb C}$ here is an infinite-dimensional weighted
  $\ell^1$. The Banach-space statements are the ones cited above.) This does
  not affect the argument itself: the conditions are worked out above from first principles, not
  taken from the reference. Beyond this case, the remaining 19 entries of
  \texttt{06-\allowbreak REFERENCES.md} were additionally reread while preparing the final
  review for similar drift (title/year/journal versus how the source is
  used in the text). Two further instances have since been found and
  corrected: the pair $\alpha(3)$/$\delta(3)$ was taken from a tertiary
  source while the primary one is in the bibliography, and the number of
  digits attributed to BOT was that of an intermediate bound rather than of
  their Table~II. For item 18 it has
  already been separately recorded (in the same place) that part of the
  data was checked only against a secondary source, not the primary text.
\end{itemize}

\emph{Numerical verification script: \texttt{04-\allowbreak CODE/\allowbreak tests/\allowbreak \_\allowbreak nozeros\_\allowbreak check.py} (rigorous
arb-sum (C1) over 46 boxes, worst value $0.3648$ on \texttt{g06}; the value is
checked automatically in \texttt{tests/\allowbreak \_\allowbreak numbers\_\allowbreak check.py}).}



\section{Constants as functions of $\ell$: what is visible over the whole segment}
\label{sec:7}

The parametric certificate yields a byproduct that the march could not give
in principle: the center $P(\ell)$ is known \textbf{analytically} on every box, so the
spectrum of the derivative along the branch can be computed on any grid in $\ell$
practically for free. Below is an exploratory scan over $101$ points with step
$0.02$ (\texttt{05-\allowbreak DATA/\allowbreak delta\_\allowbreak scan.py}).

\textbf{This is not a certificate.} Everything is computed in float64; rigorous
values exist only at two points ($\ell=3$ and $\ell=5/2$). The scan's value lies
elsewhere: it shows where to look for trouble, and --- through the control
points --- checks the branch itself far from the interpolation nodes.


\subsection{Relation between the spectra and six control points}
\label{sec:7.1}

$M'$ is the derivative of the \textbf{residual} $F(G) = G + \alpha_G\cdot G\circ W$, not of the
renormalization operator. Their spectra are related by $\mu = 1 - \lambda$. Checking
this relation is itself the first control point, and it turned out to be the
most informative one:

\begin{table}[ht]
\centering
\caption{Control points of the scan against independently certified reference
values. The discrepancy is \emph{relative} ($|\mathrm{scan}-\mathrm{reference}|
/ |\mathrm{reference}|$) and is computed from the unrounded scan value, while
the scan column is rounded for printing --- so it is reproduced from
\texttt{05-\allowbreak DATA/\allowbreak delta\_\allowbreak scan.json}, not by
subtracting the two printed columns. Each discrepancy is rounded UP, as an error measure must be.}
\label{tab:control-points}
{\small\setlength{\tabcolsep}{4pt}
\begin{tabular}{ll>{\raggedright\arraybackslash}p{0.463\textwidth}l}
\toprule
control & scan & reference & discrepancy \\
\midrule
$\delta(3)$ & $6.084689074262$ & $6.084689074262122$ --- our Krawczyk certificate & $2.59\cdot 10^{-16}$ \\
$\lambda_{2}(3)$ & $0.23714611165001$ & $0.23714611165001134$ --- our certificate & $9.26\cdot 10^{-15}$ \\
$\delta(5/2)$ & $5.412736346084$ & $5.412736346083841$ --- our certificate & $3.94\cdot 10^{-16}$ \\
$\delta(2)$ & $4.669201609103$ & $4.669201609102990$ --- \textbf{the classical Feigenbaum $\delta$} & $3.71\cdot 10^{-15}$ \\
$\alpha(2)$ & $2.502907875096$ & $2.502907875095892$ --- \textbf{the classical Feigenbaum $\alpha$} & $2.71\cdot 10^{-14}$ \\
$\delta(4)$ & $7.284686217073$ & $7.28468621707334336431$ --- the known quartic case, proven to $325$ digits by BOT & $6.68\cdot 10^{-15}$ \\
\bottomrule
\end{tabular}}
\end{table}
\FloatBarrier

The first three check internal consistency (our own certificates), the last
three check agreement with external knowledge independent of us, moreover at
\textbf{both ends} of the segment. The branch, constructed for non-integer $\ell$,
lands on the classical integer cases to double precision. The quartic reference is proven to $325$ digits
(BOT, Table~II), so the row measures the scan and not the truncation of the
reference: the float64 scan agrees with it to $14$ significant digits. An
earlier version of this table quoted the reference truncated to eleven digits and
printed a discrepancy of $3.66\cdot 10^{-12}$ --- $548$ times the true one, and a
measure of the truncation rather than of the scan.


\subsection{Box seams are invisible in the data}
\label{sec:7.2}

The scan crosses the boundaries between neighboring boxes dozens of times in
a row (the exact number depends on how seams are counted and is not detailed
here). If neighboring boxes converged to different solutions, jumps would
appear at the seams. In fact the difference in $\delta$ at every seam equals the
ordinary step in $\delta$ between neighbouring scan points: the differences at the $17$ seam crossings lie in $2.24\cdot 10^{-2}\dots 3.01\cdot 10^{-2}$, while the ordinary step over all $100$ intervals of the grid ranges over $2.24\cdot 10^{-2}\dots 3.16\cdot 10^{-2}$ with median $2.58\cdot 10^{-2}$ --- that
is, the seam is \textbf{not visible at all} in the data.

This is an independent, but NOT rigorous (a float64 scan, not a certificate),
confirmation of the gluing --- for a different aspect than \S\ref{sec:3.6}: there it was
rigorously verified that the zero of each box lies in the neighbor's
uniqueness ball ($\mathrm{dist}(P_i,P_j)(\ell^{\star}) + 1.75\cdot \rho^{\star}_i < 1.75\cdot \rho'_j$; mere
intersection of the balls is NOT sufficient for this, see \S\ref{sec:3.6}); here it is
seen illustratively that the derivative of the solution is continuous across
the seam. ``Stronger'' here does not mean ``more rigorous'': \S\ref{sec:3.6} remains the
sole rigorous argument about the gluing, and this scan supplements it rather
than surpassing it in reliability.


\subsection{What is visible}
\label{sec:7.3}

\begin{table}[ht]
\centering
\caption{$\delta$, $\alpha$, $|\lambda_{2}|$ and $|\lambda_{3}|$ at $\ell = 2, 3, 4$, with the trend of each across the segment.}
\label{tab:visible}
{\footnotesize\setlength{\tabcolsep}{4pt}
\begin{tabular}{lrrrr}
\toprule
quantity & $\ell=2$ & $\ell=3$ & $\ell=4$ & behavior \\
\midrule
$\delta$ (doubling rate) & $4.669202$ & $6.084689$ & $7.284686$ & strictly increasing, concave \\
$\alpha$ (scaling) & $2.502908$ & $1.927691$ & $1.690303$ & strictly decreasing, convex \\
$|\lambda_{2}|$ (leading stable) & $0.159628$ & $0.237146$ & $0.291838$ & strictly increasing, concave \\
$|\lambda_{3}|$ & $0.123653$ & $0.198494$ & $0.255665$ & strictly increasing \\
\bottomrule
\end{tabular}}
\end{table}

All four quantities are monotonic, and the sign of the second difference is
constant across the whole segment for $\delta$, $\alpha$ and
$|\lambda_{2}|$. \textbf{The spectrum is everywhere real}: not
a single complex pair among the leading eigenvalues at the $101$ points --- that
is, on $[2,4]$ there is neither eigenvalue collision, nor resonance, nor
spectral bifurcation. There are no surprises in the $\ell$-dependence on this
segment, and this too is a result: it means that the hyperbolicity tube will
encounter no qualitative obstructions on $[2,4]$, only a quantitative
deterioration of the margin.

A local approximation of $\delta$ by the quadratic polynomial
$\delta \approx 1.2249 + 1.9407\cdot\ell - 0.10682\cdot\ell^{2}$ matches the scan with error $\le 0.0099$
over all of $[2,4]$; the logarithmic model $a + b\cdot\ln\ell$ gives error $0.084$,
i.e.\ almost $8.5$ times worse --- the growth of $\delta$ is closer to linear than to
logarithmic.


\subsection{The one worrying trend: $\lambda_{2}$ and $\lambda_{3}$ are converging}
\label{sec:7.4}

The hyperbolicity margin $1/|\lambda_{2}|$ falls from \textbf{6.26} at $\ell=2$ to \textbf{3.43} at
$\ell=4$ --- monotonically, with no collapses. The upper end of the segment is
nearly twice as hard as the lower one, and it is precisely there that the hyperbolicity tube pays the most for it: to hold the same margin the box width falls from $0.007$ at $\ell=2$ to $0.0012$ at $\ell=4$, so the last tenth of the segment, $[3.8,4]$, carries $168$ boxes by center against $39$ in the first (both counts with
closed ends: exactly one tube box has its center at $\ell_{0} = 2.2$, and
the half-open convention $[2, 2.2)$ would give $38$). The delivered margin itself is worst not at the end but at $\ell_{0}=3.6976$ (box \texttt{t0464}, full test $0.9919$, i.e.\ a margin $1-\text{test}=0.0082$ to the threshold); at $\ell=4$ the full test is $0.5420$ (exactly $0.541939231356$, rounded OUTWARD as a majorant of a quantity that must stay below $1$), a margin of $0.4580$ --- wider by a factor of $56$, not by a factor of two.

Separately noteworthy is the ratio $|\lambda_{2}|/|\lambda_{3}|$: it too falls monotonically,
$1.2909$ ($\ell=2$) $\to$ $1.1415$ ($\ell=4$). Linear extrapolation of this data gave a merger
(ratio $= 1$) around \textbf{$\ell \approx 5.8$} --- beyond our segment, but not infinitely far
(the slope was taken by least squares over the full scan
\texttt{05-\allowbreak DATA/\allowbreak delta\_\allowbreak scan.json}, not from the two endpoints alone: from the
endpoints alone it would come out $5.89$). If the merger were real, at larger
$\ell$ $\lambda_{2}$ would cease to be a simple real eigenvalue and would move into a
complex pair, and statements of the form ``$\lambda_{2}$ is a universal constant of the
class'' would require a qualification on the range of $\ell$.

The caveat that extrapolating with a linear model $1.8$ units beyond the edge
of the data is a \textbf{hypothesis, not a result} was mandatory, and the section
itself proposed checking it with the same scan further out in $\ell$. \textbf{The
check has been done, and the prediction was not confirmed.} It is \textbf{NOT
RIGOROUS}: float64/mpmath, without interval certificates; no theorem of the
article and no acceptance run relies on it --- it refutes the prediction of
this section rather than proving anything. The ratio was recomputed on
$[4,12]$ with two different engines --- layer-free collocation in a monomial
basis (mpmath, $N=24$, dps=$30$, a fine, not everywhere uniform step --- the union
of two exploratory runs) and the package's $u$-coordinate engine ported to
mpmath (dps=$50$); at the common points they agree to four decimal places.
Both were seeded by a third, previously computed dataset (independent real
collocation, mpmath dps=$50$, $N=40/48$), with which the result agrees at seven
points $\ell = 4.0 \dots 7.0$ to no worse than $3\cdot 10^{-12}$ --- this is a check of
truncation sufficiency, not an independent confirmation. The exploratory
data themselves are not part of the package: the section is an observation,
no statement of the article relies on it, and there is no need to look for
it among the certificates.

\begin{table}[ht]
\centering
\caption{The ratio $|\lambda_{2}|/|\lambda_{3}|$ recomputed on $[4,12]$, outside the scan range (not rigorous).}
\label{tab:ratio-extended}
\begin{tabular}{lrrrrrrr}
\toprule
$\ell$ & $4$ & $5$ & $5.8$ & $6.5$ & $8$ & $10$ & $12$ \\
\midrule
ratio & $1.1415$ & $1.1078$ & $1.0890$ & $1.0764$ & $1.0573$ & $1.0415$ & $1.0314$ \\
\bottomrule
\end{tabular}
\end{table}

At the predicted point $\ell \approx 5.8$ the ratio equals $1.0890$, not $1$. The decrease
is monotonic across all $281$ computed points (union and deduplication of two
exploratory runs --- data outside the package, sweep\_mp\_l2l3 +
sweep\_mp\_bonus, hence the non-integer relation to the range/step of the
single run above), but it \textbf{slows down}: the average slope falls from
$-0.034$ on $[4,5]$ to $-0.0075$ on $[7,12]$ (both from the $281$-point set, not
from the seven tabulated values) --- a more than fourfold change on this one
interval, whereas the linear model kept it constant. The imaginary parts of
$\lambda_{2}$ and $\lambda_{3}$ over all of $[4,12]$ do not exceed $3\cdot 10^{-25}$ --- this is rounding
noise; the spectrum remains real, no complex pair arises. Where the ratio
goes as $\ell \to \infty$ is not settled by this data: a slowing decrease is equally
consistent with a limit strictly greater than $1$ (no merger at all) and with
a merger far beyond $\ell = 12$. The convergence trend thus remains --- but
without a demonstrated merger point and without grounds to expect one near
the segment.

\textbf{Independent cross-check of the scan against the tube.} The scan computes
the spectrum from the fixed-point solution at each point
(\texttt{fp\_\allowbreak lib.solve\_\allowbreak fixed\_\allowbreak point} $\to$ \texttt{build\_Mprime}); the tube computes it from the
parameterized center $P(\ell)$ via the \texttt{tube\_lib} package. These are different
code paths and different routes to the matrix. Comparison of the $101$ scan
points with the $706$ tube centers (\texttt{04-\allowbreak CODE/\allowbreak tests/\allowbreak \_\allowbreak scan\_\allowbreak vs\_\allowbreak tube.py}): where a
scan point coincides with a box center, $\delta$ and $\lambda_{2}$ agree \textbf{to $10^{-15}$}
(e.g.\ $\ell=3$: $8.9\cdot 10^{-16}$; $\ell=3.9$: $4.4\cdot 10^{-15}$); at the other points the difference
does not exceed $1.7\cdot 10^{-6}$ and is entirely accounted for by linear
interpolation between centers. This is an independent confirmation of both
the parameterized centers and the scan itself.

\emph{Data: \texttt{05-\allowbreak DATA/\allowbreak delta\_\allowbreak scan.py}, \texttt{05-\allowbreak DATA/\allowbreak delta\_\allowbreak scan.json},
\texttt{05-\allowbreak DATA/\allowbreak delta\_\allowbreak scan.md}, \texttt{04-\allowbreak CODE/\allowbreak tests/\allowbreak \_\allowbreak scan\_\allowbreak vs\_\allowbreak tube.py}.}


\section{A side observation: the family $\alpha^{(-k\ell)}$ in the spectrum at even $\ell$}
\label{sec:8}

This is \textbf{an observation, not a theorem}: the numbers are reproducible and
stable, but there is no rigorous certificate for them, and the mechanism has
been established only partially. The section is included because the
observation looks structural rather than accidental, and because it
delineates the boundary of applicability of the t-trick.


\subsection{What is observed}
\label{sec:8.1}

At \textbf{even integer} $\ell$, the spectrum of $\DT$ contains a whole family of
numbers $\alpha(\ell)^{(-k\ell)}$, $k \in \{1,2,3,4,6\}$ (k$=$5 is not tabulated separately),
with a miss at the level of float64 error:

\begin{table}[ht]
\centering
\caption{The family $\alpha^{-k\ell}$ at even $\ell$: miss against the observed spectrum.}
\label{tab:family-even}
{\footnotesize\setlength{\tabcolsep}{4pt}
\begin{tabular}{crrrrrrr}
\toprule
$\ell$ & $\alpha$ & $k=1$ & $k=2$ & $k=3$ & $k=4$ & $k=6$ \\
\midrule
$2$ & $2.502907875096$ & $2.14\cdot 10^{-14}$ & $4.84\cdot 10^{-15}$ & $2.89\cdot 10^{-15}$ & $2.56\cdot 10^{-15}$ & $5.38\cdot 10^{-15}$ \\
$4$ & $1.690302971405$ & $7.20\cdot 10^{-15}$ & $2.83\cdot 10^{-15}$ & $9.79\cdot 10^{-16}$ & $1.02\cdot 10^{-14}$ & $2.16\cdot 10^{-14}$ \\
\bottomrule
\end{tabular}}
\end{table}
\FloatBarrier

\textbf{The positions in the spectrum are exactly $k\ell$}: at $\ell=2$ the value
$\alpha^{(-2k)}$ occupies position $2k$ when ranked by decreasing modulus, at $\ell=4$ ---
position $4k$.

At $\ell$ that is not an even integer ($2.25$, $2.5$, $3.0$, $3.5$) the nearest eigenvalue is separated
from $\alpha^{(-\ell)}$ by $3\cdot 10^{-2} \dots 7\cdot 10^{-2}$. The number simply is not there.


\subsection{What has been checked for stability}
\label{sec:8.2}

\begin{itemize}
\item \textbf{Not a truncation artifact:} the miss does not change as $N_{p}$
  ranges from $20$ to $200$ (except for $N_{p} = 20$ at $\ell=4$, where the matrix is
  shorter than $4k$).
\item \textbf{Not an artifact of the choice of layer:} the miss does not
  depend on $c \in \{0.4 \dots 0.7\}$ and $r \in \{0.6 \dots 1.0\}$.
\item \textbf{Not an artifact of center imprecision:} as $P$ is refined, the
  miss drops ($2.14\cdot 10^{-14} \to 4.16\cdot 10^{-15}$ at $\ell=2$), i.e.\ the equality is
  genuine, not an approximate coincidence.
\end{itemize}


\subsection{Mechanism: what was confirmed and what was refuted}
\label{sec:8.3}

The natural guess is that these are ``gauge modes'' of the coordinate
$t = |x|^{\ell}$: renormalization contains the scaling $x \mapsto x/\alpha$, to which
the linear map $t \mapsto \alpha^{(-\ell)}t$ corresponds in the t-coordinate, and for a
composition with a linear map the eigenvalues on the monomials $t^{k}$
equal $\alpha^{(-k\ell)}$. The guess was tested in its strong form and
\textbf{refuted}: the eigenvector for $\alpha^{(-k\ell)}$ is not the monomial $t^{k}$.
The cosine with $t^{k}$ lies in $0.05 \dots 0.47$; the scaling equation
$v(\alpha^{(-\ell)}t) = \alpha^{(-k\ell)}v(t)$ is violated by $2$--$3$ orders of magnitude;
$M'\cdot t^{k}$ is not proportional to $t^{k}$.

However, \textbf{a causal link to scaling is confirmed by a direct test}.
In the derivative of the operator, scaling enters through the linear
shift $s(u) = \alpha^{(-\ell)}u + c(\alpha^{(-\ell)} - 1)$. If $\alpha^{(-\ell)}$ is replaced by $1$
in it, the family $\alpha^{(-k\ell)}$ \textbf{disappears} from the spectrum entirely.
If it is replaced by $\alpha^{(-\ell/2)}$ or $\alpha^{(-2\ell)}$, the family also
disappears, rather than shifting to $\alpha^{(-k\ell/2)}$. That is, scaling is
necessary but enters nonlinearly, and there is no closed form for the
eigenvector of these modes.


\subsection{Why only even $\ell$}
\label{sec:8.4}

Three explanations were tested; they are distinguished by the
behavior at $\ell = 2 + \varepsilon$:

\begin{table}[ht]
\centering
\caption{Three explanations of the behaviour at $\ell = 2+\varepsilon$, and how each fared.}
\label{tab:variants}
\begin{tabular}{p{0.4\textwidth} p{0.25\textwidth} p{0.25\textwidth}}
\toprule
variant & prediction & result \\
\midrule
(a) the mode $x^{k\ell} = t^{k}$ is non-analytic at non-integer $\ell$ and leaves the space & the miss grows as $O(\varepsilon)$, the spectrum stays real & \textbf{consistent}: miss $\approx 0.114\cdot\varepsilon$ at $\ell=2+\varepsilon$ and $\approx 0.062\cdot\varepsilon$ at $\ell=4-\varepsilon$ \\
(b) it is due to the $X_{0}$ normalization & gives no quantitative prediction & nothing to test against \\
(c) the eigenvalue splits into a complex pair & complex pairs appear nearby & \textbf{refuted}: the spectrum stays real at non-integer $\ell$ \\
\bottomrule
\end{tabular}
\end{table}
\FloatBarrier

Conclusion: the crossing is \textbf{transversal} --- the value $\alpha^{(-\ell)}$
crosses the spectral curve at a nonzero angle exactly at the even
points, rather than being tangent to it.


\subsection{Prediction at $\ell = 6$: confirmed}
\label{sec:8.5}

The observation was obtained at $\ell = 2$ and $\ell = 4$. The real test is to
predict and hit a target where no one has looked yet. The prediction
was fixed in advance: at $\ell = 6$ the spectrum must contain $\alpha(6)^{(-6k)}$
\textbf{at positions $6k$}.

Reaching $\ell = 6$ by continuation in $\ell$ with this engine does not
succeed: the continuation breaks off around \textbf{$\ell \approx 5.3$} (step $0.1$ ---
at $5.3$, step $0.01$ --- at $5.35$).

The workaround is a jump $\ell = 4 \to 6$ by Newton iteration in \texttt{mpmath}
with dps$=$50: convergence is quadratic, the residual is
$\norm{\tilde{F}}_{\infty} = 4.8\cdot 10^{-45}$.

\textbf{The break-off is a property of the continuation method, not of the
branch.} This must be stated explicitly, because it is tempting ---
and wrong --- to read the break-off as a singularity of the solution:
what the jump establishes is that at $\ell = 6$ a solution with a residual of
$10^{-45}$ exists at all; whether it is the continuation of the same branch is
a separate question, and it is answered below --- not rigorously, but by three
numerical arguments of different kinds: the direct check, the dependence of the
break-off on the truncation, and the layer-free sweep of \S\ref{sec:7.4}. Direct check (NOT RIGOROUS, float64): the solution computed
by an independent layer-free engine is substituted into the same
layer $(c, r) = (0.6, 0.8)$ with the same $N_{p} = 80$ --- the very one
that breaks off at $5.3$ --- and the residual grows smoothly, with no
kink:

\begin{table}[ht]
\centering
\caption{Residual of the layer representation, $\norm{\tilde{F}}_{\infty}$ and $\norm{W-c}_{X}/r$, at $\ell = 4.0 \dots 7.0$.}
\label{tab:residual-layer}
{\footnotesize\setlength{\tabcolsep}{4pt}
\begin{tabular}{lrrrrrr}
\toprule
$\ell$ & $4.0$ & $5.0$ & $5.5$ & $6.0$ & $6.5$ & $7.0$ \\
\midrule
$\norm{\tilde{F}}_{\infty}$ & $3.9\cdot 10^{-16}$ & $6.0\cdot 10^{-16}$ & $1.1\cdot 10^{-15}$ & $6.8\cdot 10^{-15}$ & $1.9\cdot 10^{-13}$ & $3.2\cdot 10^{-12}$ \\
$\norm{W-c}_{X}/r$ & $0.839$ & $0.996$ & $1.061$ & $1.117$ & $1.168$ & $1.213$ \\
\bottomrule
\end{tabular}}
\end{table}

The layer represents the solution both at $\ell = 6$ and at $\ell = 7$; what
breaks is not the solution but the small-step Newton continuation
over it. Independently: layer-free engines pass through $[4,12]$ in
one continuous sweep without a single failure (\S\ref{sec:7.4}), and the layer
continuation itself, if a Newton system of size $N = 64$, $60$, or $40$ is
taken instead of $N_{p} = 80$, passes $\ell \approx 5.3$ smoothly and reaches
$\ell = 6$ (residual $2.9\cdot 10^{-13}$, $7.5\cdot 10^{-14}$, and $3.0\cdot 10^{-15}$
respectively --- the SMALLER the truncation, the better). That is,
exactly one value of the truncation parameter breaks off, and the
break-off moves together with it, which cannot happen for a
singularity of the solution. We do not have an exact mechanism for
the break-off. The second row of the table gives a candidate:
$\norm{W-c}_{X}/r$ --- an X-norm estimate, hence certainly not underestimated,
of the same self-embedding $W(D) \subseteq D$. This estimate is the quantity
$q$, which WAS condition (D4) in
\S\ref{sec:10.2} (later dropped as redundant --- see \S\ref{sec:10.2}, \S\ref{sec:10.4}); the self-embedding
itself was never dropped. The ratio crosses
unity around $\ell \approx 5.0$, so the cheap guarantee of disk invariance
disappears before the break-off. But this cause alone is not
sufficient: on four different layers the moment of this crossing
differs substantially (for one of them the ratio is already above
unity at $\ell = 4$), while all of them break off at roughly the same
place.
For the observation of this section this does not matter: what matters
is that at $\ell = 6$ a solution exists and that the numerical evidence above
identifies it with the same one.

The resulting fixed point gives $\alpha(6) = 1.467742450319899$
(monotonicity holds: $\alpha(2) = 2.5029 > \alpha(4) = 1.6903 > \alpha(6) = 1.4677$).
The spectrum of $\DT$ at this point:

\begin{table}[ht]
\centering
\caption{Confirmation of the predicted spectrum of $\DT$ at $\ell = 6$.}
\label{tab:spectrum-ell6}
\begin{tabular}{clll}
\toprule
$k$ & $\alpha(6)^{(-6k)}$ & miss & position from maximum \\
\midrule
$1$ & $1.000232\cdot 10^{-1}$ & $2.0\cdot 10^{-15}$ & \textbf{6} \\
$2$ & $1.000465\cdot 10^{-2}$ & $5.3\cdot 10^{-15}$ & \textbf{12} \\
$3$ & $1.000697\cdot 10^{-3}$ & $7.1\cdot 10^{-14}$ & \textbf{18} \\
$4$ & $1.000929\cdot 10^{-4}$ & $5.7\cdot 10^{-13}$ & \textbf{24} \\
$5$ & $1.001162\cdot 10^{-5}$ & $2.4\cdot 10^{-12}$ & \textbf{30} \\
$6$ & $1.001395\cdot 10^{-6}$ & $3.4\cdot 10^{-11}$ & \textbf{36} \\
\bottomrule
\end{tabular}
\end{table}
\FloatBarrier

The positions $6k$ are exactly as predicted. The stability is the
same as at $\ell = 2$, $4$: the miss does not change for $N_{p}$ from $40$ to
$120$ and for $(c, r)$ from four pairs.

\textbf{The prediction was confirmed.} For an observation this carries
more weight than an explanation: the numbers $\alpha(6)$ and the positions
were known before the computation --- there is nothing to tune.


\subsection{What has not been done}
\label{sec:8.6}

\begin{itemize}
\item There is no rigorous (interval) certificate that $\alpha^{(-k\ell)}$ belongs
  to the spectrum --- all numbers were obtained in float64.
\item There is no closed form for the eigenvector; it is known only that
  it is analytic and decays rapidly (convergence in $N_{p}$ from
  15--30).
\end{itemize}

\emph{Numerical data for the section and verification logs ---
\texttt{05-\allowbreak DATA-\allowbreak INDEX.md}.}

---

\emph{Numbers in the article's main tables are checked by an independent
run of the corresponding script (\texttt{02-\allowbreak REPRODUCE.md}). Exceptions are
cross-checked against the
certificate values, but are not recomputed from the package:
\notrecomputedroll{}
Precision is not uniform across the package, and no single figure covers it:
the working precisions run from prec=$128$ to prec=$1536$. The engines that
produce and check the certificates occupy the lower part of that range --- the
rigorous majorant libraries and their driver at $512$, the existence engine and
the Rouch\'e test of the tube gate at $256$, the
acceptance pass that recomputes finished boxes at $128$ --- while the two
highest settings belong to test scripts at $\ell = 5/2$. The full map,
script by script, is in \S\ref{sec:10.7}.
The difference $0.28995$ (without the patch) versus $0.29060$ (with the
E-transfer patch) is explained in \S\ref{sec:5.2.2}: it is a rigorous upper
estimate accounting for the box certification of eigenpairs.}



\section{Caveats and Open Issues}
\label{sec:9}

\textbf{Invariant class.} The result holds on the invariant class of
functions of $t=|x|^{\ell}$ --- the working layer $D(3/5,4/5)$, a space of the type
Lanford used in 1982 \cite{lanford1982computer} but not his disk (\S\ref{sec:1.1}) --- not
on the full space of unimodal pairs $G-Y$.

\textbf{The fixed layer $(c, r)$} --- \textbf{done (\S\ref{sec:10})}, but not for an
arbitrary layer. Theorems A--C are formulated for the operator
restricted to the space $X = X(3/5, 4/5)$. \S\ref{sec:10}{} proves (Theorem D)
that the nonzero spectrum of $\DT(\Gs(\ell))$ --- and hence the count
``exactly one outside'', and $\delta(\ell)$ itself --- does not depend on the
choice of $(c, r)$ within the explicitly exhibited and certified family of
layers $\Ffam$, on all of which the transfer $P\to\Gs$ of \S\ref{sec:10.7}
is certified. This removes the
concern that $\delta(\ell)$ is an artifact of the normalization $c = 3/5$,
$r = 4/5$.

What remains is a limitation of \textbf{coverage}, not of correctness:
the exhibited members of $\Ffam$ are the rectangle $\Qcell = (0.56, 0.64) \times (0.76, 0.84)$
(a continuum of layers) plus 14 individual layers scattered over
$[0.55, 0.70] \times [0.70, 0.90]$, plus a connected corridor of 7 cells
between the base layer and one of the farthest of the exhibited ones (the
corridor joins the two by an explicitly certified connected chain, and the transfer
$P\to\Gs$ of \S\ref{sec:10.7}{} is certified on its cells as well,
see \S\ref{sec:10.13}). That the exhibited coverage reaches an arbitrary
admitted layer is not proved, and near the upper edge
$\ell \to 4$ the family is indeed thin (\S\ref{sec:10.9}).

The numerical picture that was the sole evidence before \S\ref{sec:10}{} remains
valid as an independent illustration: 16 runs --- all 14 pointwise layers of the
family $\Ffam$ plus two variations of $N_p$ (the base layer itself at
$N_p = 80$ is not a separate entry: it is the reference against which
the others are compared), $c$ from 0.55 to 0.70, $r$ from 0.70 to
0.90, $N_p$ from 60 to 100 --- at four values of $\ell$ (2, 5/2, 3, 4) give
the same twenty leading eigenvalues to a tolerance of $10^{-11}$, which
is what the test checks; over the fifteen layer records of the probe file ---
the base layer and the fourteen others, without the $N_p$ variations --- the
total spread of $\delta$ does not
exceed $1.51\cdot 10^{-14}$ (Appendix~\ref{sec:A}, row 57). The count ``exactly one
outside'' does not change anywhere (\texttt{04-\allowbreak CODE/\allowbreak tests/\allowbreak \_\allowbreak layer\_\allowbreak invariance.py}).

The structural reason this was to be expected: the truncation
operates on polynomials of degree $\le N_p$, and this subspace does not
depend on $(c, r)$ --- polynomials of degree $\le N_p$ in $(t-c)$ and in
$(t-c')$ form the same set. A change of layer acts on it as a change
of basis: diagonal in $r$ (weights $r^k$) and triangular in $c$ (the
binomial theorem). The proof in \S\ref{sec:10}{} takes a different route --- not
comparing the truncated matrices at all, but working at the level of
the operator: the fixed point and the eigenfunctions are continued
from one layer to another by the equation itself.

\textbf{The horseshoe and the combinatorics of renormalization} (the
symbolic/combinatorial structure of the orbits of the renormalization
operator along the certified family $\Ffam$) --- outside the scope of this
work, not addressed.

\textbf{The tube over $\ell \in [2,4]$} --- \textbf{done} (\S\ref{sec:3}, \S\ref{sec:4}). The pilot (summary ---
\texttt{05-\allowbreak DATA/\allowbreak COVERAGE.md}, machine-readable --- \texttt{05-\allowbreak DATA/\allowbreak coverage.json})
closed 55 disconnected components on $[2.07,\, 3.5808]$ (0.303 in total in
$\ell$, see \S\ref{sec:3.8}) by the march; the full interval was closed by a
different method --- the parametric center --- and cost 42.2 core-minutes
instead of the extrapolated $\approx 422$ core-hours of the march.

\textbf{The existence certificate at $\ell=2.5$} --- \textbf{done}
(\texttt{03-\allowbreak CERTIFICATES/\allowbreak pointwise/\allowbreak kant\_\allowbreak l25.json}), and it further COEXISTS
with the general certificate on [2,4] as an independent pointwise
check (the identification of the zeros of $\Gs(2.5)$ and $\Gbar$ follows the
same way as at $\ell = 3$, \S\ref{sec:4.7}: the same computation from
\texttt{branch\_\allowbreak g11.npz} and the deposited seed gives
$\|\Gbar - P(2.5)\|_{X} \le 7.88\cdot 10^{-15}$, seven orders below the
uniqueness radius $1.75\cdot 10^{-7}$, so the pointwise zero lies in the
uniqueness ball of Theorem~A and is therefore $\Gs(2.5)$).

\textbf{All rigor rests on} Arb via python-flint 0.9.0, the package chain,
and high-precision seeds ($\sim 10^{-80}$), whose precision bounds the
admissible Miller sanity check (1e-60). The correctness of that interval
arithmetic is premise A4 of the registry in \S\ref{sec:10.11}: it is external to
this work, common to the field, and nothing here establishes it.

\textbf{Limits recorded elsewhere in the article and repeated here, because this
is the list a reader consults for them.}

\begin{itemize}
\item \textbf{Regeneration is not deterministic.} $Y0$, $\tau_{F}$,
  $\rho^{\star}_{b}$, the $\mathrm{extra}$ that inherits from it, and also
  $d_{\mathrm{chk}}$ and through it $Z_{\mathrm{conv}}$, are built
  through a float seed that varies between library builds --- $Y0$ by tens of
  percent (\S\ref{sec:3.6}, \S\ref{sec:3.7}). A rerun of box \texttt{g01}
  with the shipped code reproduces 38 of the 51 fields bit for bit and moves
  seven: $Y0$ by $\times 2.30$, $\tau_{F}$ by $\times 7.96$,
  $d_{\mathrm{chk}}$ by $\times 1.88$, $\rho^{\star}$ by $\times 2.30$.
  Note the direction: the guard criterion rests on $d_{\mathrm{chk}}$, and its
  share $B_{A}d_{\mathrm{chk}}/Y0$ in that rerun becomes \emph{worse} than in
  the deposited run ($0.0133$ against $0.0108$, threshold $0.5$) --- the
  non-determinism does not move it only in the safe direction. The high-precision seeds named
  just above are those of the deposited run; they do not make a fresh
  regeneration reproduce it digit for digit. What must survive a regeneration is
  the verdict, not the digits (\emph{Reproducibility}).

\item \textbf{The printed fields \texttt{margin\_a}, \texttt{U\_min} and
  \texttt{ReU\_min\_R} on \texttt{g37}--\texttt{g43} are anti-conservative} by
  $\le 6.67\cdot10^{-5}$, $\le 4.45\cdot10^{-4}$ and $\le 5.59\cdot10^{-4}$ respectively (\S\ref{sec:6.2.a}); the conclusions hold with the
  honest values, but the printed ones are not majorants.
\item \textbf{(E1) and (E2) are not gated by the certifier}
  \texttt{one\_box} (\S\ref{sec:10.2}). They are, however, machine-checked:
  they are elementary inequalities in the layer's own $(c', r')$, and the
  numbers gate now reads those constants from every deposited layer and cell
  (\texttt{tests/\allowbreak \_\allowbreak numbers\_\allowbreak check.py},
  block K4: $40$ layers and cells, worst margin $0.05$ at the layer
  $(0.70, 0.75)$; for a cell the worst corner is taken). What remains true is
  that a layer never deposited is not covered by anything.
  \texttt{D\_margin\_a} and
  \texttt{D\_margin\_c} have no second, independently written implementation
  (\S\ref{sec:10.9}).
\item \textbf{No automatically run check assembles $\DT$ a second time at a
  non-integer interior $\ell$.} The cross-assembly of \S\ref{sec:3.9} is run by
  the driver at $\ell = 3$ only. The two further points are not reproduced by
  hand, though: they are the same gated script with one command-line argument
  (\texttt{tests/\allowbreak \_\allowbreak dt\_\allowbreak fd\_\allowbreak cross.py 2.5}),
  which returns a verdict and an exit code --- at $\ell = 2.5$ it reports an
  elementwise discrepancy of $7.36\cdot 10^{-14}$ against its own threshold
  $10^{-6}$ in $98$ seconds. What is missing is two more steps in the driver,
  not a rigorous computation.
\item \textbf{A published table disagrees at $\ell = 3$ in $\delta$ and agrees
  in $\alpha$.} The value $\delta(3) = 5.9679687037\ldots$ of Briggs
  \cite{briggs1991precise} differs in the second significant figure from the
  enclosure certified here; the three grounds given in \S\ref{sec:2} place the
  discrepancy on the side of the table, but auditing another author's
  computation is not something this work undertakes. The $\alpha(3)$ of that
  same table agrees with the float64 scan to $14$ significant digits, and with
  the certified $\alpha(\Gs(3))$ to $18$ (\S\ref{sec:3.9}) --- which is the
  one external comparison this work has at an interior point (\S\ref{sec:2}).
\end{itemize}

\textbf{The resolved question of the K-mode.} The original gate
``$M\cdot \tilde K \approx 0$'' rested on the erroneous hypothesis that the rescaling
family yields the kernel of DF. In fact the closure
$\alpha_G = -1/G(1)$ breaks the invariance of the family. The eigenvalue
$\lambda_K = \ell\cdot \Gbar'(1)/\Gbar(1) = 7.16329$ of $\DT$ on the full $X$ does genuinely
exist --- this is derived in \S\ref{sec:5.1.2}{} via the LEFT eigenvector
$\varphi(h) = h(0)$, without recourse to artifacts. However, the
identification of the right eigenvector with $K = \Gbar - \ell\cdot t\Gbar'$ is only
a hypothesis; it is not used here. The accompanying observation that
``$1-\lambda_K = -6.1633$ is exactly the `extra' number in the spectrum of
$M$'' is a leftover from that earlier period and is \textbf{not confirmed as
an artifact} in the package: the full $(81\times 81)$ matrix of the operator
on $X$, in the available basis, is not a representation of it (the
numbers spread out to $10^{8}$, the spectrum is meaningless --- this is
stated directly in \texttt{04-\allowbreak CODE/\allowbreak tests/\allowbreak \_\allowbreak kmode.py}), so there is nothing to
check $-6.1633$ against. Below, only what is checkable is checked.
There is no kernel of DF; the $X_{0}$-slice is used as the normalization
$G(0)=1$, not to bypass a degeneracy.

The number 7.16329 is greater than unity, and it is worth
distinguishing right away from Theorem B: \textbf{it is not in the spectrum
of $\DT|_{X_{0}}$}. It is also worth stating the converse without
softening: on the FULL $X$, $\DT$ has \textbf{two} eigenvalues outside the
unit circle, and the leading one among them is not $\delta$ but precisely
$\lambda_K$ ($7.16329$ versus $\delta(3) = 6.08469$). The statement ``the
unstable direction is one-dimensional'' is a statement about the
normalization slice $X_{0}$, not about $\DT$ on $X$; the reduction to the
slice itself is legitimate (by the block triangularity of \S\ref{sec:5.1.2},
$\Aff$ is invariant), but a reader who concludes that the transversal
direction is ``small and harmless'' will be wrong about the leading
mode. Direct measurement at $\ell=3$ (\texttt{04-\allowbreak CODE/\allowbreak tests/\allowbreak \_\allowbreak kmode.py},
acceptance step \texttt{kmode}): the distance from 7.16329 to the nearest
eigenvalue of $\DT|_{X_{0}}$ is \textbf{no less than 1.078} (measured 1.0786),
whereas $\deltastar = 6.084689074262122$ lies in this spectrum to a precision
of $2.7\cdot 10^{-15}$. The distance is published ROUNDED DOWN, the precision
ROUNDED UP: the opposite rounding would make both statements stronger
than what was measured.
\textbf{Both numbers here are float64 measurements}, not interval enclosures:
the matrix is assembled in arb at \texttt{prec=256} and then cast to float64
before the eigenvalues are taken, and the conservative rounding just stated is
what makes them safe to quote. Nothing in Theorems A--D rests on them --- the
count ``exactly one outside'' is settled by the Rouch\'e test, not here. The same place also measures why a K-mode
cannot be in this spectrum: $K(0) = \Gbar(0) = 1 \ne 0$, that is,
$K \notin X_{0}$. The K-mode lives in the full space and is removed by
exactly that normalization; the count ``exactly one outside'' pertains
to the slice $X_{0}$ and does not conflict with it.

\textbf{Remarks on rigor.} $\Ahat$ is taken as a matrix of exact arb numbers
(obtained from the float64 inversion of the midpoint) --- this is
rigorous: $Z_{1a}$ honestly measures its residual by interval
arithmetic, and a separate ``inversion-error ball'' is not needed.
$L = 2.065\cdot 10^3$ is a rigorous UPPER bound (the chain in \S\ref{sec:5.1.2}, item
5). The margin in $\Zhat$ shows that $\hat\rho'$ can be raised to $\sim 10^{-13}$ if one
wishes to enlarge the uniqueness ball.



\section{Independence of the Spectrum from the Layer (Theorem D)}
\label{sec:10}

\S\ref{sec:9} names the fixed layer $(c, r)$ as the work's principal
open point. Here that point is closed for the explicitly exhibited
family of layers.

Each premise below is tagged with one of three labels:

\begin{itemize}
\item \textbf{[A]} --- certified elsewhere in the package (Theorems A/B/C, the
  acceptance runs \texttt{glue.py}, \texttt{tube\_\allowbreak accept.py});
\item \textbf{[C]} --- certified here, by interval computation (code
  \texttt{04-\allowbreak CODE/\allowbreak layer/\allowbreak }, logs \texttt{03-\allowbreak CERTIFICATES/\allowbreak layer/\allowbreak });
\item \textbf{[P]} --- proved here analytically.
\end{itemize}


\subsection{Setup}
\label{sec:10.1}

The package proves Theorems A/B/C on \textbf{one} layer $D_0 = D(c_0, r_0)$,
$c_0 = 3/5$, $r_0 = 4/5$, in the space $X(c_0,r_0)$ of functions analytic
in $D_0$ with norm
\[
\norm{h}_{X(c,r)} = \sum_k |h_k| \cdot r^k , \qquad h(t) = \sum_k h_k \cdot (t - c)^k .
\]

Formally this means that $\delta(\ell)$ is an eigenvalue of \textbf{our} operator
on \textbf{our} space. \S\ref{sec:9} records this as the principal open point: it is
not proved that the count of eigenvalues outside the unit circle, and
$\delta(\ell)$ itself, do not depend on the choice of layer.

Here this is closed for an explicit family of layers. The strategy
sidesteps the dead end named in \S\ref{sec:9} (the truncation is defined not
only by the image but also by the complement, and the complement
depends on $c$): all statements are made \textbf{at the level of the
operator}, and the truncated matrices are not compared at all.


\subsection{The family of layers $\Ffam$}
\label{sec:10.2}

A layer $(c', r')$ belongs to $\Ffamz$ \textbf{precisely when} four groups of
conditions hold (E, D, I1, I2) --- there used to be five: see below regarding
(D4), dropped as redundant. The four groups are thus not merely sufficient for
membership: they define $\Ffamz$, and every layer of $\Ffamz$ satisfies them.

\textbf{Elementary (E).}

\begin{itemize}
\item[(E1)] $|c'| < r'$ --- then $0 \in D'$, and the slice $X_0 = \{h(0) = 0\}$
  is meaningful;
\item[(E2)] $|1 - c'| < r'$ --- then $t = 1 \in D'$, evaluation at the point $1$
  is bounded: $|h(1)| \le \sum|h_k|\, |1-c'|^k \le \norm{h}_{X'}$,
  which is needed for $\alpha_G = -1/G(1)$ and
  $d\alpha(h) = h(1)/G(1)^2$.
\end{itemize}

(A caveat on the scope of applicability of the code: (E1)/(E2) are
not gated by the certifier \texttt{one\_box} itself. They are checked
separately, by the numbers gate, for every deposited layer and cell
(\texttt{tests/\allowbreak \_\allowbreak numbers\_\allowbreak check.py}, block
K4; for a cell the worst corner is taken): $40$ objects, worst margin $0.05$
at the layer $(0.70, 0.75)$, where $0.70 < 0.75$. Until that gate was added
these two conditions were verified by hand, and code applied to a layer
outside the deposited set still does not check them.)

\textbf{Certificate D (the layer on its own).} For each $\ell$-box of the
package:

\begin{itemize}
\item[(D1)] $\mathrm{margin}_a = r' - \max_{|t-c'|\le r'} |t\cdot\alpha^{-\ell} - c'| > 0$ --- $s(D')\subseteq D'$;
\item[(D2)] $\mathrm{margin}_c = r' - \max_{|t-c'|\le r'} |W - c'| > 0$ --- $W(D')\subseteq D'$;
\item[(D3)] $\min_{|t-c'|\le r'} |U| \ge 1/2$ and $\operatorname{Re} U > 0$ on $|t-c'| \in \{r', 6/5\}$ --- $U$ has no zeros, the principal branch of $\operatorname{Log} U$ is single-valued, $W = U^\ell$ is analytic in $D'$.
\end{itemize}

(There used to also be (D4): $q = \norm{W-c'}_{X'}/r' < 1$,
$q_s = \norm{s-c'}_{X'}/r' < 1$ --- a numerical check of compactness. It has
been removed from $\Ffamz$: in \S\ref{sec:10.4} the same compactness is derived from
(D1)+(D2) via the spectral radius of the Banach algebra, without a
separate number. This is not a waiver of the check but the
replacement of one sufficient check by another, manifestly weaker
one (\S\ref{sec:10.4}) --- and, as it turned out in practice (\S\ref{sec:10.6}), (D4) was
precisely the one real (not merely theoretical) bottleneck on wide
cells.)

\textbf{Certificate I1 (transfer from $D'$ to $D_0$).}

\begin{itemize}
\item[(I1a)] $r_0 - \big( |c'\cdot\alpha^{-\ell} - c_0| + r'\cdot|\alpha^{-\ell}| \big) > 0$ --- $s(D')\subseteq D_0$;
\item[(I1b)] $r_0 - \max_{|t-c'|=r'} |W - c_0| > 0$ --- $W(D')\subseteq D_0$;
\item[(I1c)] $\operatorname{Re} U > 0$ on $|t - c'| = r'$.
\end{itemize}

\textbf{Certificate I2 (in reverse, from $D_0$ to $D'$).} The same with
$(c_0,r_0) \leftrightarrow (c',r')$ swapped.

\paragraph{The certified family $\Ffam$.} Passing the four groups above admits a
layer as a candidate; it is not yet everything the conclusions about the
solution $\Gs(\ell)$ require. Write
\[
  \Ffam \;=\; \bigl\{\,(c',r')\in\Ffamz \;:\;
     \text{the transfer of \S\ref{sec:10.7} is certified there as well}\,\bigr\}
  \;\subseteq\; \Ffamz ,
\]
that is, those admitted layers for which in addition $c'+r' \le 1.53$ and the
margins listed above --- together with the minimum of $\operatorname{Re} U$ on the
inner circle $|t-c'| = r'$ of (D3) and on the circles of (I1c) --- dominate the
transfer perturbations
$\Delta U/\Delta W/\Delta s$. The outer circle $|t-c'| = 6/5$ of (D3) is not
among them: the transfer of \S\ref{sec:10.7} is certified for
$|t| \le \max(c'+r',\,1.4)$ and does not reach it, and no lemma behind
Theorem D uses it (\S\ref{sec:10.9}). All of these are properties of the pair
$(c',r')$ alone, so the definition is not circular.
Every layer of the presented family --- the fourteen pointwise layers of
\S\ref{sec:10.6}, the corridor cells of \S\ref{sec:10.13}, and every point of $\Qcell$ --- lies in
$\Ffam$; the statements below about how far the admitted set reaches --- its
geometry, its boundary, and the version of Theorem D for the interpolant
$P(\ell)$ --- are statements about $\Ffamz$, and say so.

\paragraph{A remark on boundaries.} In (D2), (I1b), (I2b) the maximum is taken
over the \textbf{circle}, whereas the statement is needed on the disk.
This is legitimate: $U = G\circ s$ is a composition of an entire function
(the polynomial $P$, and after Lemma 1 the continued $\Gs$) with the
affine $s$, so $\operatorname{Re} U$ is harmonic, and its minimum over the closed
disk is attained on the boundary; by (D3), $\operatorname{Re} U > 0$ inside, hence
$W = U^\ell$ is analytic in the disk, and $|W - c|$ is the modulus of an
analytic function, whose maximum is likewise attained on the
boundary. \textbf{[P]} Numerical control: inside the disk the worst
$|W - c_0|$ = $0.636$ / $0.521$ / $0.624$ against the certified boundary
values $0.683$ / $0.541$ / $0.667$ (boxes \texttt{g42} / \texttt{g00} / \texttt{gtop}, layer
$(0.55, 0.80)$) --- no violations. \textbf{[C]}

Separately, on $\min |U| \ge 1/2$ from (D3). Here the same ``boundary
versus disk'' question is sharper: the minimum of the MODULUS of an
analytic function, unlike the maximum, is in general not attained on
the boundary --- the minimum-modulus principle requires that the
function have no zero inside. The code
(\texttt{04-\allowbreak CODE/\allowbreak layer/\allowbreak layer\_\allowbreak lemma\_\allowbreak u.py:31-40}) computes $U_{\min}$ by arc
subdivision precisely of the CIRCLE $|t-c'| = r'$. The bridge is
built in the following order, and it is not circular:

\begin{itemize}
\item[1.] $\operatorname{Re} U$ is harmonic (see above), so $\operatorname{Re} U > 0$ on the circle
   implies $\operatorname{Re} U > 0$ on the whole closed disk --- this is a statement
   about $\operatorname{Re} U$, and it uses neither $|U|$ nor the absence of zeros;
\item[2.] $\operatorname{Re} U > 0$ on the disk implies $U \ne 0$ on the disk;
\item[3.] only now can the minimum-modulus principle be applied: $U$ is
   analytic and has no zeros in the disk, so $\min |U|$ over the
   closed disk is attained on the boundary, and the certified
   boundary minimum suffices.
\end{itemize}

The order matters: if step 3 were applied before step 2, the argument
would be exactly the circular one that \S\ref{sec:6.4.1} recognizes as
insufficient.


\subsection{Lemma 1 (continuation of the fixed point)}
\label{sec:10.3}

\begin{lemmaone}
Let $(c',r') \in \Ffam$ and let $\ell$ belong to one of the
$\ell$-boxes. Then $\Gs(\ell)$, constructed in \textbf{[A]} on $D_0$, continues
analytically to $D'$, and the continuation lies in $X(c',r')$.
\end{lemmaone}

\begin{proof}
\textbf{[P]} From $F(\Gs) = 0$ follows the identity
\[
\Gs(t) = -\alpha \cdot \Gs(W(t)) , \qquad \alpha = -1/\Gs(1), \quad W = U^\ell, \quad U = \Gs\circ s .
\]

(The certificates (I1a)--(I1c), as well as (D1)--(D3) and (I1) in Lemma 2, and
(I1)/(I2) in Lemma 3, are computed numerically in \S\ref{sec:10.6} for the
interpolant $P(\ell)$, not for the true $\Gs(\ell)$, which is used here and
in Lemmas 2--4; the transfer from $P(\ell)$ to $\Gs(\ell)$ is justified in
\S\ref{sec:10.7} uniformly and WITHOUT splitting into individual premises: the
certificates of \S\ref{sec:10.6} are transferred there in one step, without
division into embedded layers and the rest,
with a margin $\ge 225.9\times$ on the 14 pointwise layers
and $\ge 71.9\times$ over the whole cell $\Qcell$, and $\ge 128\times$ on every corridor cell
of \S\ref{sec:10.13} --- this entire step of the argument, in
each of Lemmas 1--4, must be read together with \S\ref{sec:10.7}.)

By (I1b), $W(D') \subseteq D_0$, where $\Gs$ is already defined, and by (I1c)
and the Remark in \S\ref{sec:10.2}, $W$ is analytic in $D'$. Hence the right-hand side
is analytic in $D'$ and continues $\Gs$ there; on $D' \cap D_0$ the
continuation coincides with the original function by uniqueness of
analytic continuation. Estimate:
\[
\sup_{D'} |\Gs| \le |\alpha| \cdot \sup_{D_0} |\Gs| \le |\alpha| \cdot \norm{\Gs}_{X(c_0,r_0)} .
\]

Finiteness of the norm of the continuation. We need $\varepsilon > 0$ such that
$(c', r'+\varepsilon)$ also satisfies (I1a)--(I1c) (so that the inclusions above
already hold for $D(c',r'+\varepsilon)$); then, by Cauchy's formula on this
larger disk, the coefficients of $\Gs$ about $c'$ decay as
$(r'/(r'+\varepsilon))^k$, and $\norm{\Gs}_{X(c',r')} < \infty$. For a \textbf{cell} $(c',r')$
this $\varepsilon$ is built into the interval certification itself: an interior
point of the cell by construction has a margin to its edge. For a layer given
as a single \textbf{point} --- the 14 \textbf{pointwise} layers of
\S\ref{sec:10.6} among them --- there is no such margin in the
definition. It is supplied without a new computation: for a fixed
$\ell$-box, each of (I1a), (I1b), (I1c) is continuous
in $(c',r')$: (I1a) is a closed algebraic expression in $(c',r')$; (I1b) is
$r_0$ minus the maximum over the circle $|t-c'|=r'$ of a function depending
continuously on $(c',r',\theta)$ --- a composition of the
polynomial $P$, the affine $s$, and the power $U^\ell$; and (I1c) is the
minimum of $\operatorname{Re} U$ over that same circle. The maximum of a
function continuous in $(c',r',\theta)$, over compact $\theta\in[0,2\pi)$, is
continuous in $(c',r')$ (a standard fact: a consequence of uniform
continuity on a compact set). Since (I1a)--(I1c) hold with strict inequality
at \textbf{every} point of $\Ffamz$ --- that is what the conditions themselves
say: (I1a) and (I1b) are written as $>0$, and the minimum of the continuous
$\operatorname{Re} U$ over the compact circle of (I1c) is therefore positive as
well --- by continuity they
remain strict also in some open neighborhood of that point in the
space $(c',r')$; in particular along the ray $r' \to r'+\varepsilon$ for
sufficiently small $\varepsilon > 0$. Nothing here is special to the fourteen:
the step covers every layer to which Lemma 1 is applied.
This gives the existence of $\varepsilon$ without
explicitly computing its value --- exactly what is needed for the
estimate above. (The elementary conditions (E1)--(E2), meanwhile, only improve as $r'$ grows for fixed $c'$.)

One step is still owed here, and it is worth spelling out. The continuity argument just given is about the \textbf{interpolant}: (I1b) and (I1c) are written through $P$, the affine $s$ and the power $U^{\ell}$ with $U=U_P$, and the strictness is taken from (I1a)--(I1c) holding at every point of $\Ffamz$ --- a $P$-condition (\S\ref{sec:10.6}). What the chain above consumes is the $\Gs$-version of (I1b)/(I1c) on the enlarged layer $(c',r'+\varepsilon)$. The bridge is the transfer of \S\ref{sec:10.7}, and it covers that layer as well: the perturbations $\Delta U/\Delta W/\Delta s$ are explicit algebraic functions of $R_t=\max(c'+r',1.4)$ and of the $\ell$-box, monotone in $c+r$, and the majorant they rest on, $\sup|U_P-1|\le0.400$, is certified not at the exhibited points but on the \textbf{whole} domain $c+r\le1.53$ (\texttt{04-\allowbreak CODE/\allowbreak tests/\allowbreak \_\allowbreak nozeros\_\allowbreak check.py}, constants \texttt{CR\_\allowbreak WIDE}$\,=1.53$ and \texttt{THRESH\_\allowbreak WIDE}$\,=0.400$), whereas the widest exhibited layer reaches only $c+r=1.522$. Since the $P$-margins dominate $\Delta U/\Delta W/\Delta s$ with the factors already recorded above ($\ge225.9\times$ on the 14 pointwise layers, $\ge71.9\times$ over $\Qcell$, $\ge128\times$ on every corridor cell) and are themselves continuous in $(c',r')$, for small enough $\varepsilon$ --- with $c'+r'+\varepsilon\le1.53$ --- they still dominate at $(c',r'+\varepsilon)$, and the $\Gs$-version of (I1a)--(I1c) holds there. That is the $\varepsilon$ the estimate needs. (A route avoiding $\varepsilon$ altogether also exists: $\Gs\in X(c',r')$ follows from (I1b) by the same Gelfand argument used three paragraphs below in Lemma 2, since $\rho(W-c_0)=\max_{\bar D(c',r')}|W-c_0|=r_0-\texttt{I1\_margin\_W}<r_0$ makes $\sum_k b_k (W-c_0)^k$ converge in $X(c',r')$.) \qedhere
\end{proof}

\paragraph{Consequence.} Each of the 14 pointwise layers of \S\ref{sec:10.6} is in fact the center
of some (implicitly) small open region on which (I1a)--(I1c) continue to hold, rather than an
isolated point with no margin. What this does \emph{not} yet give is fourteen
small regions of \emph{admitted} layers: admission asks for (E1)/(E2), (D1)--(D3)
and (I2) as well, and \S\ref{sec:10.10} asks about the admitted set. Its argument
therefore continues to rest on the fourteen certified points themselves, far
removed from one another, and not on regions around them.
Computing the size of this region explicitly is not required ---
Lemma 1 only needs the existence of $\varepsilon$, not its value.


\subsection{Lemmas 2--4}
\label{sec:10.4}

\begin{lemmatwo}[Compactness on every layer]
For $(c',r') \in \Ffam$, the operator
$\DT(\Gs(\ell))$ is bounded and compact on $X(c',r')$, and its essential spectrum
is $\{0\}$.
\end{lemmatwo}

\begin{proof}
\textbf{[P]} The operator is a sum of three types of terms:

\begin{itemize}
\item[1.] \textbf{finite rank:} the term entering through $d\alpha(h) = h(1)/\Gs(1)^2$ is a
   bounded functional (by (E2)), multiplied by a fixed function; rank $1$;
\item[2.] \textbf{multiplication} by the fixed functions $\Gs\circ W$, ${\Gs}'\circ W$, $\ell U^{\ell-1}$,
   ${\Gs}'\circ s$ --- bounded, since the norm $\norm{\cdot}_{X(c,r)}$ is submultiplicative, and
   the functions themselves lie in $X(c',r')$ by Lemma 1 and (I1). For ${\Gs}'$
   this is not just $\Gs$ from Lemma 1: the same construction yields more ---
   $\Gs$ is extended to a STRICTLY LARGER disc $D(c',r'+\varepsilon)$ with coefficients
   $a_k$ decaying as $(r'/(r'+\varepsilon))^k$. From here a standard Cauchy estimate for
   the derivative of a power series: $k\cdot(r'/(r'+\varepsilon))^k \to 0$ exponentially, so it
   is majorized by a constant $C(\varepsilon)$, and $\sum_k k|a_k| r'^{k-1} \le C(\varepsilon)\cdot \norm{\Gs}_{X(c',r'+\varepsilon)} < \infty$
   --- by the same margin $\varepsilon$, without additional computation;
\item[3.] \textbf{compositions} $h \mapsto h\circ W$ and $h \mapsto h\circ s$.
\end{itemize}

The list is exhaustive in a way that matters: \textbf{no term multiplies $h$ itself.} A bare multiplication operator $h\mapsto m\cdot h$ on $X(c',r')$ is bounded but NOT compact, so had item 2 acted on $h$ directly, the conclusion below would not follow from items 1--3. It does not: the identity part of $D_{G}F$ cancels in $\DT = I - D_{G}F$, and every fixed function listed in item 2 stands as a factor at $h\circ W$, at $h\circ s$, or at the rank-one term of item 1.


For the composition $h \mapsto h\circ W$: $X(c',r')$ is a commutative Banach algebra with
unity under \textbf{pointwise} multiplication (submultiplicativity of the norm has
already been used above, for the multiplication terms). Let $g = W - c'$ be an element of this algebra, $z = t - c'$ be the generator. (That $W = U^{\ell}$ and $U^{\ell-1}$ belong to $X(c',r')$ at NON-INTEGER $\ell$ is neither a composition nor a product, and does not follow from analyticity on $D'$ alone; it follows from the holomorphic functional calculus in this same algebra, by the identification of the maximal ideal space made just below: the spectrum of $U$ in the algebra equals $U(\bar D(c',r'))$, and $\operatorname{Re} U > 0$ there --- (D3) certifies this on the circle $|t-c'| = r'$, and it follows on the whole closed disc because $\operatorname{Re} U$ is harmonic and takes its minimum on the boundary (the remark on boundaries in \S\ref{sec:10.2}) --- so the spectrum misses the cut $(-\infty,0]$ and $U^{a} = \exp(a\,\Log U)$ lies in $X(c',r')$ for every real $a$.) Polynomials in $z$
are dense in $X(c',r')$ (the tail $\sum_{k>N}|h_k| r'^k \to 0$ as $N\to\infty$), and
$\norm{z} = r'$ exactly (coefficient $1$ at $k=1$), hence the spectral radius of
$z$ in the algebra equals $r'$, and every continuous character $\varphi$ (homomorphism
into $\CC$) is uniquely determined by the value $\lambda = \varphi(z)$, $|\lambda| \le r'$; conversely,
each such $\lambda$ yields a character --- evaluation at a point, $\varphi_\lambda(h) = h(c'+\lambda)$.
Hence the maximal ideal of the algebra is exactly the closed disc $\bar D(c',r')$
(a classical fact for the Wiener weighted algebra $\sum|h_k| r'^k$), and the Gelfand
formula for the spectral radius gives
\begin{align*}
\rho(g) = \lim_k \norm{g^k}^{1/k} = \max_\varphi |\varphi(g)| &= \max_{\bar D(c',r')} |W - c'| \\
&= r' - \mathrm{margin}_c < r' ,
\end{align*}
the last equality being exactly the definition of $\mathrm{margin}_c$ (D2) (as in
Lemma 1, numerically certified in \S\ref{sec:10.6} for $P(\ell)$ and transferred to $\Gs(\ell)$
in \S\ref{sec:10.7}), and it is $>0$.
Since $\rho(g) < r'$ strictly, for any $\delta$ with $\rho(g) < r'-\delta < r'$ there exists
$K = K(\delta)$ such that $\norm{g^k} \le (r'-\delta)^k$ for all $k \ge K$. Hence, if $\Pi_N$
is the projection onto $k \le N$ (on the \textbf{input}, i.e. in the expansion of
$h$ itself), and $C_W : h \mapsto h\circ W$ is the composition operator, then for
$N \ge K$
\begin{align*}
\norm{C_W (\Id - \Pi_N) h} = \norm{\textstyle\sum_{k>N} h_k g^k} &\le \sum_{k>N} |h_k| (r'-\delta)^k \\
&\le \left(\frac{r'-\delta}{r'}\right)^N \sum_{k>N} |h_k| r'^k \le \left(\frac{r'-\delta}{r'}\right)^N \norm{h} ,
\end{align*}
that is, $\norm{C_W(\Id-\Pi_N)} \to 0$ as $N\to\infty$ (not the other way round: truncating
$h\circ W$ itself in its \textbf{own} coefficients would mix contributions of all $k$,
since $W(c')-c' \ne 0$ in general, and such a formula would be incorrect).
Hence $C_W$ is the norm limit of the finite-rank operators $C_W \Pi_N$, i.e.
it is compact; the same holds for $C_s : h\mapsto h\circ s$ through $\mathrm{margin}_a > 0$ (D1)
instead of $\mathrm{margin}_c$. The sum of compact and finite-rank operators is compact,
so the essential spectrum is $\{0\}$, and the nonzero spectrum consists of
isolated eigenvalues of finite algebraic multiplicity. \qedhere
\end{proof}

\paragraph{Remark (why this is better than the previous version).} The numerical quantity
$q = \norm{W-c'}_{X'}/r'$ (the previous certificate D4 --- see \S\ref{sec:10.2}) estimates
$\norm{g^k}$ more crudely, via submultiplicativity itself ($\norm{g^k} \le \norm{g}^k = (qr')^k$),
which yields the same conclusion but requires the stronger condition $q < 1$.
Since for a power series the maximum modulus on the disc is always no greater
than the sum of moduli of the coefficients ($\rho(g) \le \norm{g}$, triangle inequality),
the spectral radius is a strictly weaker and therefore broader condition:
$\mathrm{margin}_c > 0$ can hold where $q \ge 1$. On narrow cells both paths yield the
same set of $46$ boxes (there $q<1$ holds a fortiori), but on wide cells (\S\ref{sec:10.6})
$q$ temporarily exceeds $1$ for edge boxes ($\ell\to 4$), while $\mathrm{margin}_c$ remains
strictly positive --- i.e. (D4) turned out to be not just mathematically
redundant but a real (not merely theoretical) bottleneck of the computation.

\begin{lemmathree}[Bootstrap of eigenfunctions]
Let $(c',r') \in \Ffam$, $\lambda \ne 0$.
\begin{enumerate}[label=(\roman*)]
\item If $v \in X(c_0,r_0)$, $v(0) = 0$ and $\DT v = \lambda v$ on $X_0(c_0,r_0)$, then
$v$ extends analytically to $D'$ and the extension lies in $X_0(c',r')$,
remaining an eigenfunction with the same $\lambda$.
\item Conversely, from $X_0(c',r')$ to $X_0(c_0,r_0)$.
\item The same holds for Jordan chains: if $(\DT - \lambda)w = v$, then $w$
extends together with $v$.
\end{enumerate}
\end{lemmathree}

\begin{proof}
\textbf{[P]} Rewrite the equation as an extension formula:
\[
v = \lambda^{-1} \cdot \DT v ,
\]
and note that the right-hand side at point $t$ uses $v$ at the points
$W(t)$, $s(t)$ and --- through the rank-one term $d\alpha(h) = h(1)/\Gs(1)^2$ --- at the
point $1$, while everything else is fixed functions that are analytic
on $D'$ (Lemma 1). By (I1a) $s(D') \subseteq D_0$ and by (I1b) $W(D') \subseteq D_0$, while
$1 \in D_0$ and $1 \in D'$ by (E2) for the base and for the target layer respectively, so the
right-hand side is well-defined and analytic for all $t \in D'$: the formula
itself extends $v$ from $D_0$ to $D'$. Finiteness of the norm --- as in Lemma 1
(via the $\varepsilon$-margin of the cell, or, for a pointwise layer, by the
continuity argument given there). The value at zero is preserved, since it is
the same function: $v(0) = 0$.

Why $\lambda$ remains an eigenvalue. The operator on both layers is given by \textbf{one
and the same} expression (see \S\ref{sec:10.1}), so the function $\DT v - \lambda v$ is analytic
on $D_0 \cup D'$ (the set is connected: both discs contain $0$ by (E1)) and
vanishes identically on $D_0$. By uniqueness of analytic continuation it
vanishes on $D'$ too, i.e. the extended $v$ is an eigenfunction of the operator
on $X_0(c',r')$ with the same $\lambda$.

Item (ii) is the same argument with (I2) instead of (I1). Item (iii):
$w = \lambda^{-1}(\DT w - v)$, and by induction along the chain the right-hand side is
already extended. \qedhere
\end{proof}

\begin{lemmafour}[Coincidence of algebraic multiplicities]
For $(c',r') \in \Ffam$
and $\lambda \ne 0$, the root subspaces of $\DT$ in $X_0(c_0,r_0)$ and in $X_0(c',r')$
are isomorphic; in particular the algebraic multiplicities coincide.
\end{lemmafour}

\begin{proof}
\textbf{[P]} By Lemma 2, $\lambda \ne 0$ is isolated in both spectra and the root
subspaces are finite-dimensional. The extension $D_0\to D'$ (Lemma 3(i),(iii)) and
the converse extension $D'\to D_0$ (Lemma 3(ii); coincides with the literal
restriction only when $D'\subseteq D_0$, i.e. on $\Qcell_{\mathrm{in}}$, \S\ref{sec:10.5}) are linear maps between
the root subspaces; their
compositions in both directions are identical by uniqueness of analytic
continuation (both functions are analytic and coincide on the nonempty open
set $D_0 \cap D' \ne \emptyset$, since both discs contain $0$ by (E1)). Hence they are
mutually inverse bijections. \qedhere
\end{proof}

Lemmas 1--4 are stated with the premise $(c',r')\in\Ffam$, and their proofs rely
on the transfer of \S\ref{sec:10.7} (the margins of $D$/$I$, transferred from $P(\ell)$
to $\Gs(\ell)$), which is certified on $c+r\le 1.53$. Every layer exhibited in
\S\ref{sec:10.6} and \S\ref{sec:10.13} has $c+r\le 1.522$ --- the maximum, attained on the
corridor --- so that domain contains all of them. More: $c+r\le 1.53$ is one of the two
conditions defining $\Ffam$, so the domain contains every member of $\Ffam$, exhibited
or not, and the chain of lemmas is closed on the whole of $\Ffam$. For a layer of $\Ffamz$
that has not been shown to satisfy the two further conditions, they have to be
checked directly before Theorem D applies to it; the remark at the end of \S\ref{sec:10.5} says so.


\subsection{Theorem D}
\label{sec:10.5}

\begin{namedtheorem}[Independence of layer]
\label{thm:D}
Let $(c',r') \in \Ffam$. Then for every
$\ell \in [2,4]$:
\begin{itemize}
\item[1.] $\Gs(\ell) \in X(c',r')$ and $\DT(\Gs(\ell))$ is compact on $X_0(c',r')$;
\item[2.] the nonzero spectra of $\DT(\Gs(\ell))$ on $X_0(c',r')$ and on $X_0(c_0,r_0)$
   coincide together with the algebraic multiplicities;
\item[3.] in particular, outside the unit disc there is exactly one eigenvalue, it
   is simple, real, and equals $\delta(\ell)$ from Theorem B.
\end{itemize}
In particular, $\delta(\ell)$ is the same on all layers of the family and is not
an artifact of the normalization $c = 3/5$, $r = 4/5$.
\end{namedtheorem}

\emph{(The formulation ``$\delta(\ell)$ is a property of the problem, not of the layer''
is deliberately not used: independence of $(c,r)$ inside $\Ffam$ is proved, not
independence of the choice of function space altogether --- see \S\ref{sec:10.9}, item 2.)}

\paragraph{Remark (the version for $P(\ell)$).} Items 1--2 remain true, for every
$(c',r')\in\Ffamz$ and not only on $\Ffam$, if one replaces
$\Gs(\ell)$ everywhere by the interpolant $P(\ell)$: for Lemma 1 the extension to
$D'$ is trivial, since $P(\ell)$ is a polynomial (an entire function; the
argument of \S\ref{sec:10.3} via the identity $F(\Gs)=0$, valid for the true solution, is
not needed for the polynomial itself); Lemmas 2--4 in this replacement use, of the
numerically certified conditions, only (D1)--(D3)/(I1)/(I2), computed in \S\ref{sec:10.6}
specifically for $P(\ell)$, together with the elementary (E1)/(E2); they do not
require the transfer of \S\ref{sec:10.7}.
It is precisely this $P$-version
that is referenced by \S\ref{sec:10.11}, \S\ref{sec:10.12} and \S\ref{sec:10.10} when speaking
of ``items 1--2 about $P(\ell)$''. About $\delta(\ell)$ (item 3), unlike the $\Gs$-version,
it says nothing.

\begin{proof}
(1) --- Lemmas 1 and 2 (\S\ref{sec:10.3}, \S\ref{sec:10.4}). (2) --- Lemma 4 (\S\ref{sec:10.4}), applied to each $\lambda \ne 0$. (3) ---
from (2) and Theorem B \textbf{[A]}, which gives on the base layer exactly one
eigenvalue outside the disc. \qedhere
\end{proof}

\paragraph{Consequence 1 (the family is nonempty and explicit).} $\Ffam$ contains the
open rectangle $\Qcell = (0.56, 0.64) \times (0.76, 0.84)$ --- a continuum of layers,
twice as wide in half-width as the previous one (which was blocked by the
spurious, now removed certificate D4 --- see \S\ref{sec:10.2}, \S\ref{sec:10.4}, \S\ref{sec:10.6}) --- and, besides it,
14 separate layers listed in \S\ref{sec:10.6}, scattered over $[0.55, 0.70] \times [0.70, 0.90]$,
and the seven-cell corridor of \S\ref{sec:10.13} joining the base layer to one
of the farthest of them.
In particular the statement of Theorem D is substantive: it links spectra on
layers whose truncated matrices differ by 9.8--118\,\% (\S\ref{sec:10.8}; that
range is a float64 illustration of how far apart the layers are, not an
interval certificate --- nothing in Theorem D depends on its value).

\paragraph{Explicit decomposition of $\Qcell$ by nesting} (used in \S\ref{sec:10.12}). The condition
$D(c,r)\subseteq D_0$ is $|c-0.6|+r\le 0.8$, piecewise-linear; inside $\Qcell$ it cuts out the
triangle
\begin{align*}
\Qcell_{\mathrm{in}} &= \{ (c,r) \in \Qcell : 0.76 \le r \le 0.8 - |c-0.6| \} \\
&= \Qcell \cap \conv\{ (0.56,0.76), (0.64,0.76), (0.60,0.80) \} ,
\end{align*}
exactly $25\,\%$ of the area of $\Qcell$ ($0.0016$ out of $0.0064$) --- the simplest,
``one-line'' part of the statement: there the transfer of \S\ref{sec:10.7} is trivial
(not even the argument via $s(D')\subseteq D_0$ is needed, a direct Cauchy estimate
without reference to (I1a) suffices).
$\Qcell_{\mathrm{out}} = \Qcell\setminus \Qcell_{\mathrm{in}}$ --- the remaining $75\,\%$, where the transfer is also rigorous,
but requires the argument via $s(D')\subseteq D_0$, not the trivial embedding of discs.
Both parts are \textbf{[C]}, the decomposition is purely descriptive and does not
change the resulting coverage of $\Qcell$.

\paragraph{Consequence 2 (matching with Theorem C)} The analyticity of $\delta(\ell)$ in $\ell$
--- that conclusion of Theorem C which concerns $\delta$ --- is automatically independent
of the layer throughout the exhibited family, on which the transfer of \S\ref{sec:10.7} is
certified: this is a statement about one and the same function $\delta(\ell)$.
The other conclusion of Theorem C --- analyticity of $\Phi:\ell\mapsto\Gs(\ell)$ as a map
into $\mathcal A(c_0,r_0)$ --- is NOT claimed here for a new layer, and is not needed
anywhere: Theorem D gives $\Gs(\ell)\in X(c',r')$ pointwise in $\ell$, whereas holomorphy
of $\ell\mapsto\Gs(\ell)$ in the norm of the new layer would require a uniform-in-$\ell$
bound on $\norm{(W_\ell-c_0)^k}_{X(c',r')}$, which is not certified in the package.

\paragraph{Remark on the nature of the statement.} The condition $(c',r') \in \Ffamz$ is not
abstract: it is a finite set of interval inequalities (\S\ref{sec:10.2}), verifiable for
any given layer in a matter of minutes of computation. But verification of
$\Ffamz$ alone is not enough for the proof: the transfer of center certificates to
the solution (\S\ref{sec:10.7}) additionally relies on the dominance of the margin over
the transfer perturbations and on $c+r\le 1.53$ --- these two facts are verified
as a property of the PRESENTED family, and are not
part of the definition of $\Ffamz$: they are precisely what distinguishes $\Ffam$
from $\Ffamz$. For a layer OUTSIDE the presented family that
passes \S\ref{sec:10.2}, they must be verified separately. Within the presented family nothing is left over:
its widest point, the corridor cell reaching $c+r = 1.522$, lies inside the certified
domain $1.53$, and the margin dominance is exhibited there as well (\S\ref{sec:10.7}).
In this sense Theorem D is
conditional in exactly the same sense as Theorems A/B/C of the package are
conditional, and no further.



\subsection{What exactly is numerically certified}\label{sec:10.6}

\textbf{Grid in $\ell$.} The 46 boxes of the package, whose union covers $[2, 4]$
entirely (the extreme box \texttt{gtop} extends beyond 4: $[3.9775, 4.0075]$ --- this
is covering the endpoint, not extending the statement). \textbf{[C]}

\textbf{Pointwise layers.} 14 layers: the cross $(0.55,0.80)$, $(0.70,0.80)$,
$(0.60,0.70)$, $(0.60,0.90)$; the corners $(0.55,0.70)$, $(0.55,0.90)$,
$(0.70,0.75)$, $(0.65,0.85)$; and $(0.575,0.80)$, $(0.625,0.80)$, $(0.65,0.80)$,
$(0.60,0.75)$, $(0.60,0.85)$, $(0.65,0.75)$. A total of $14 \times 46 = 644$
certificates, all passed: $D$, $I1$, $I2$ with positive margins. \textbf{[C]}
Part of the boxes require subdivision of the $\ell$-box into $K$ parts;
subdivision is not curve-fitting: it only narrows the interval of $P(\ell)$
coefficients inside the box, and the statement is still made for the whole
box as the union of parts. A complete run of all 644 pairs by the certifier
\texttt{04-\allowbreak CODE/\allowbreak layer/\allowbreak layer\_\allowbreak cert.py} (driver \texttt{layer\_\allowbreak recheck\_\allowbreak all.py}) shows that
subdivision was needed in only one place --- 9 boxes of the layer $(0.55,0.90)$
at $K=2$; all the remaining 635 pairs pass at $K=1$.

Bottlenecks --- the worst value for each gate separately over all 644 pairs
(\texttt{04-\allowbreak CODE/\allowbreak layer/\allowbreak layer\_\allowbreak worst\_\allowbreak by\_\allowbreak gate.py}; the overall ``worst margin'' is the
minimum over all six gates at once, so for the table one needs the minimum
for each separately):

\begin{table}[ht]
\centering
\caption{Worst value of each gate over the 644 pointwise-layer pairs.}
\label{tab:worst-by-gate}
{\small\setlength{\tabcolsep}{4pt}
\begin{tabular}{>{\raggedright\arraybackslash}p{0.571\textwidth}ll}
\hline
gate & worst value & where \\
\hline
\texttt{D\_margin\_a} (self-embedding $s(D')\subseteq D'$) & $2.44 \cdot 10^{-2}$ & $(0.70,0.75)$, \texttt{g00} \\
\texttt{D\_margin\_c} (self-embedding $W(D')\subseteq D'$) & $5.28 \cdot 10^{-2}$ & $(0.55,0.80)$, \texttt{g43} \\
\texttt{I1\_margin\_s} ($s(D')\subseteq D_{0}$) & $0.126$ & $(0.55,0.90)$, \texttt{g00} \\
\texttt{I1\_margin\_W} ($W(D')\subseteq D_{0}$) & $1.95 \cdot 10^{-3}$ & $(0.55,0.90)$, \texttt{g43} \\
\texttt{I2\_margin\_s} ($s(D_{0})\subseteq D'$) & $\mathbf{1.09 \cdot 10^{-3}}$ & $(0.70,0.75)$, \texttt{g00} \\
\texttt{I2\_margin\_W} ($W(D_{0})\subseteq D'$) & $6.80 \cdot 10^{-3}$ & $(0.55,0.70)$, \texttt{g43} \\
$q$ (not required, see \S\ref{sec:10.2}/\S\ref{sec:10.4} --- kept for memory) & $1.2352$ & $(0.70,0.80)$, \texttt{g43} \\
\hline
\end{tabular}}
\end{table}
\FloatBarrier

\textbf{The global minimum over all 6 gates and all 644 pairs is \texttt{I2\_margin\_s} $= 1.09\cdot 10^{-3}$ at \texttt{(0.70,0.75)/\allowbreak g00}} --- the same number quoted as the ``worst
margin over 644'' in \S\ref{sec:10.9} and \S\ref{sec:10.11}
(\S\ref{sec:10.7} names a different number, \texttt{I1\_margin\_W} $= 1.954\cdot 10^{-3}$ at the
pair \texttt{(0.55,0.90)/\allowbreak g43}, and it is not a competing answer to this one: what is
minimised there is the ratio margin\,/\,perturbation, not the margin itself, and
$1.954\cdot 10^{-3}$ is the numerator of that worst ratio.) All thin margins improve under subdivision of the
$\ell$-box: actually $K=2$ was needed only by 9 boxes of the layer $(0.55,0.90)$
(above), after which they exceed $10^{-3}$. \textbf{[C]}

\textbf{Interval cells (continuum statement).} Certificates are computed at
\textbf{interval} $c' \in [c_{lo}, c_{hi}]$, $r' \in [r_{lo}, r_{hi}]$: ball arithmetic then
gives the statement at once for all layers of the cell. \textbf{[C]}

The first run (half-widths below) showed that the degradation with growing
cell is explained not by the ``irreducible'' conservatism of interval
arithmetic in general, but by a specific place in the code: the engine
\texttt{tube\_\allowbreak lib.verify\_\allowbreak lemma\_\allowbreak U} (via \texttt{lemma\_\allowbreak domains.verify\_\allowbreak lemma\_\allowbreak 1}), before checking
self-embedding, \textbf{re-expands} the monomial polynomial around the \textbf{interval}
center $c'$ ($G_c = P(c' + (t-c'))$ via the \texttt{arb\_poly} composition), and then
evaluates $W - c'$ on this re-expansion. The re-expansion coefficients and the
evaluation point depend on \textbf{the same} $c'$, but \texttt{arb} arithmetic treats them
as independent intervals --- the classical dependency problem, doubling the
width. Diagnostics on box g22: as the half-width grows from 0 to $5\cdot 10^{-3}$
this margin (\texttt{D\_margin\_c}) drops from 0.125 to 0.010, whereas the \texttt{I1}/\texttt{I2}
margins ($s$- and $W$-image against the \textbf{fixed} base layer $(3/5,4/5)$),
computed by Horner's scheme directly on the monomial coefficients without
re-expansion,
lose over the same range only 0.01--0.02 --- the same data, a
different computation path, and a loss smaller by a factor of five to eleven.
All four figures in this paragraph are diagnostics of that superseded path,
measured before the fix described next: they are not certificates, they are not
persisted in the package, and no step of any proof rests on them.

Fix (\texttt{04-\allowbreak CODE/\allowbreak layer/\allowbreak layer\_\allowbreak lemma\_\allowbreak u.py}): compute the self-embedding
$W(D') \subseteq D'$ by the same Horner scheme on the monomial coefficients already
used for \texttt{I1}/\texttt{I2} --- without re-expansion around $c'$. The \texttt{margin\_a} formula
(part (a) of the lemma, $s(D')\subseteq D'$) was not changed: it does not re-expand
the polynomial and did not degrade. Result on the \textbf{same} cell $\pm 3\cdot 10^{-3}$,
where previously only the margin $8.72\cdot 10^{-4}$ held: \textbf{[C]}

\begin{table}[ht]
\centering
\caption{Cell margins under the historical certifier modes (with gate D4).}
\label{tab:cell-margins-d4}
{\small\setlength{\tabcolsep}{4pt}
\begin{tabular}{>{\raggedright\arraybackslash}p{0.208\textwidth}llll}
\hline
cell (half-width around $(3/5, 4/5)$) & boxes passed & worst margin & \texttt{q\_max} & first failed gate \\
\hline
$\pm 3\cdot 10^{-3}$, old \texttt{lemma\_U} & \textbf{46 / 46} & $8.72\cdot 10^{-4}$ (\texttt{g28}) & $0.9112$ & --- \\
$\pm 3\cdot 10^{-3}$, monomial \texttt{lemma\_U}, STILL with gate D4 & \textbf{46 / 46} & $0.145$ (\texttt{g00}) & $0.9954$ & --- \\
$\pm 5\cdot 10^{-3}$ & \textbf{46 / 46} & $0.140$ (\texttt{g00}) & $0.9933$ & --- \\
$\pm 1\cdot 10^{-2}$ & \textbf{46 / 46} & $0.130$ (\texttt{g00}) & $0.9942$ & --- \\
$\pm 1.5\cdot 10^{-2}$ & \textbf{46 / 46} & $0.116$ (\texttt{gtop}) & $0.9994$ & --- \\
$\pm 2\cdot 10^{-2}$ & \textbf{46 / 46} & $\mathbf{0.103}$ (\texttt{gtop}) & $0.9993$ & --- \\
$\pm 3\cdot 10^{-2}$, STILL with gate D4 & 41 / 46 & $7.43\cdot 10^{-2}$ (\texttt{g39}) & $1.0305$ & \texttt{bases\_q}, \texttt{g40} \\
\hline
\end{tabular}}
\end{table}

(margins rounded DOWN, \texttt{q\_max} --- UP; protocol ---
\texttt{03-\allowbreak CERTIFICATES/\allowbreak layer/\allowbreak cells3\_\allowbreak hist.json}, driver ---
\texttt{04-\allowbreak CODE/\allowbreak layer/\allowbreak layer\_\allowbreak cells\_\allowbreak hist.py}, modes set by \texttt{LAYER\_\allowbreak PKG\_\allowbreak LEMMA} /
\texttt{LAYER\_\allowbreak GATE\_\allowbreak D4})

(The drivers for the table of this subsection are
\texttt{04-\allowbreak CODE/\allowbreak layer/\allowbreak layer\_\allowbreak cells\_\allowbreak hw.py} (search of the boundary by half-width),
\texttt{layer\_\allowbreak cells\_\allowbreak bisect.py} (three control half-widths at the very boundary of
\Qcell) and \texttt{layer\_\allowbreak cells\_\allowbreak offcenter.py} (off-center cell); protocols ---
\texttt{cells3\_\allowbreak hw.json}, \texttt{cells3\_\allowbreak bisect.json},
\texttt{cells3\_\allowbreak offcenter.json} and the common \texttt{cells3.log}.)

The upper table was obtained in two historical modes of the certifier, and
both are reproduced by the same file \texttt{04-\allowbreak CODE/\allowbreak layer/\allowbreak layer\_\allowbreak cert.py} through
environment variables: \texttt{LAYER\_\allowbreak PKG\_\allowbreak LEMMA=1} turns on the engine
\texttt{tube\_\allowbreak lib.verify\_\allowbreak lemma\_\allowbreak U} (re-expansion around the interval $c'$, $M=64$) ---
this is the line ``old \texttt{lemma\_U}''; \texttt{LAYER\_\allowbreak GATE\_\allowbreak D4=1} returns gate D4 as a
failure condition --- these are the lines ``STILL with gate D4''. By default
both modes are off.
Control: under the mode \texttt{LAYER\_\allowbreak PKG\_\allowbreak LEMMA=1 LAYER\_\allowbreak GATE\_\allowbreak D4=1}
the box \texttt{g28} on the cell $\pm 3\cdot 10^{-3}$ gives $\mathrm{margin}_c = 8.7204\cdot 10^{-4}$, i.e.
exactly the published $8.72\cdot 10^{-4}$; with \texttt{LAYER\_\allowbreak PKG\_\allowbreak LEMMA}
switched off the same box gives \texttt{0.1825}. Both runs are the two
\texttt{half\_\allowbreak width}$\,=\,$0.003 records of
\texttt{03-\allowbreak CERTIFICATES/\allowbreak layer/\allowbreak cells3\_\allowbreak hist.json}
(\texttt{mode}$\,=\,$\texttt{pkg\_\allowbreak lemma+D4} and \texttt{mode}$\,=\,$\texttt{D4});
in each the number is
\texttt{boxes.\allowbreak g28.\allowbreak worst.\allowbreak D\_\allowbreak margin\_\allowbreak c}.

The lines above are the certification still \textbf{with} the numerical gate D4
($q<1$, $q_s<1$); it stopped the growth precisely at $\pm 3\cdot 10^{-2}$. Below this
gate is removed as mathematically redundant (\S\ref{sec:10.2}, \S\ref{sec:10.4}) --- the table continues
without it, by the same code (\texttt{layer\_\allowbreak cert.py}, \texttt{M\_ARCS=128}, the same $K$
search), $q$ is still computed and written to the protocol, but no longer a
failure condition:

\begin{table}[ht]
\centering
\caption{Cell margins after removing gate D4.}
\label{tab:cell-margins-nod4}
{\small\setlength{\tabcolsep}{4pt}
\begin{tabular}{>{\raggedright\arraybackslash}p{0.090\textwidth}lll>{\raggedright\arraybackslash}p{0.288\textwidth}}
\hline
half-width, WITHOUT gate D4 & boxes passed & worst margin & \texttt{q\_max} & first failed gate \\
\hline
$\pm 2\cdot 10^{-2}$ & \textbf{46 / 46} & $5.49\cdot 10^{-2}$ (\texttt{g43}) & $1.224$ & --- \\
$\pm 3\cdot 10^{-2}$ & \textbf{46 / 46} & $9.33\cdot 10^{-3}$ (\texttt{g43}) & $1.297$ & --- \\
$\pm 3.5\cdot 10^{-2}$ & \textbf{46 / 46} & $\mathbf{3.30\cdot 10^{-3}}$ (\texttt{g39}) & $1.270$ & --- \\
$\pm 4\cdot 10^{-2}$ & \textbf{46 / 46} & $6.32\cdot 10^{-4}$ (\texttt{g42}), thin & $1.223$ & --- \\
$\pm 4.5\cdot 10^{-2}$ & 44 / 46 & --- & $1.178$ & \textbf{\texttt{lemma\_U}} (negative \texttt{margin\_c}: $-3.4\cdot 10^{-4}$ and $-1.8\cdot 10^{-4}$; \texttt{margin\_a} still $\approx 8.3\cdot 10^{-2}$), \texttt{g43} and \texttt{gtop} \\
\hline
\end{tabular}}
\end{table}

\noindent The \emph{worst margin} column of this table is not comparable with
the same column of Table~\ref{tab:cell-margins-d4} row by row. At the two
half-widths the two tables share, they report different worst boxes and
different margins: at $\pm 2\cdot 10^{-2}$, $0.103$ at \texttt{gtop} there
against $5.49\cdot 10^{-2}$ at \texttt{g43} here; at $\pm 3\cdot 10^{-2}$,
$7.43\cdot 10^{-2}$ at \texttt{g39} there against $9.33\cdot 10^{-3}$ at
\texttt{g43} here. The gate changes which boxes fail, and hence which ones the
$K$ search subdivides, so the column is not a like-for-like measurement across
the two tables; each shows the boundary within its own mode.

The boundary has grown exactly twice in half-width ($\pm 2\cdot 10^{-2}\to\pm 4\cdot 10^{-2}$, four
times by area) and, more importantly, the \textbf{character} of the limit has
changed: for the first time not a numerical gate (\texttt{lemma\_U} in the old sense
--- re-expansion, or \texttt{bases\_q} --- redundancy) fails, but the very geometry of
the family --- \texttt{margin\_c} actually vanishes at \texttt{g43}/\texttt{gtop} (boxes
$\ell \to 4$) on the cell $\pm 4.5\cdot 10^{-2}$, \texttt{margin\_a} staying there
at $\approx 8.3\cdot 10^{-2}$, two orders of magnitude away from zero. This is the same gate code as before
(\texttt{lemma\_\allowbreak U\_\allowbreak direct}, \texttt{ok\_D}), only now it is the sole and last line of defense,
not one of two. The \texttt{q\_max} values (1.17--1.30, i.e.\ definitely $>1$) in the
new table are \textbf{intentionally kept}: they show that the old gate would
forbid everything that is allowed here --- a vivid confirmation that it was
spurious, not just ``not interfering yet''.

\FloatBarrier

\textbf{Return to the mode STILL-WITH-gate-D4} (not Table~\ref{tab:cell-margins-nod4}, the ``without D4'' one ---
there at $\pm 3\cdot 10^{-2}$ the gate is not failed, because there $q<1$ is no longer
a failure condition). The margin on the former cell $\pm 3\cdot 10^{-3}$ has grown by
$\approx 166$ times. The cell holds up to $\pm 2\cdot 10^{-2}$ --- 6.6 times wider than the former
limiting $\pm 3\cdot 10^{-3}$ --- and \textbf{fails not on \texttt{lemma\_U}}, but on \texttt{bases\_q}
(compactness): starting from $\pm 3\cdot 10^{-2}$ the boxes at the upper edge
$\ell \to 4$ (\texttt{g40}--\texttt{g43}, \texttt{gtop}) fail $q < 1$. This is exactly the same place
that was earlier masked by an earlier \texttt{lemma\_U} failure, and is independently
confirmed by the off-center run (see below).

Caveat on the cause of the $166\times$ growth: it decomposes into two contributions,
and re-expansion is not the only difference of \texttt{layer\_\allowbreak lemma\_\allowbreak u} from the engine
\texttt{tube\_\allowbreak lib.verify\_\allowbreak lemma\_\allowbreak U}. The engine version rigidly fixes the number of
arcs $M=64$ (\texttt{M\_\allowbreak ARCS\_\allowbreak LEMMA} in \texttt{tube\_lib.py}), whereas \texttt{layer\_\allowbreak lemma\_\allowbreak u} is
called with $M=128$ (the same value already used for \texttt{I1}/\texttt{I2}). The mere
transition $M: 64\to 128$ on the original, re-expanded formula raises \texttt{margin\_c}
for g28 on the cell $\pm 3\cdot 10^{-3}$ from $8.72\cdot 10^{-4}$ to $0.0297$ ($\approx 34\times$; the $0.0297$ is a one-off run of the superseded formula at an arc count the shipped code does not use for it, and is not persisted in the package) --- i.e.\ a
noticeable part of the gain is this doubling of the number of arcs, not only
the elimination of re-expansion. This does not undermine correctness (Horner
with arcs is a valid estimate for any $M$), but means that ``166 times'' is
the effect of \textbf{both} edits together, not the pure effect of eliminating
interval dependence. The second clarification is more accurate than the first
draft (also caught by verification): \texttt{margin\_a} is not the only limit over
the entire range $\pm 3\cdot 10^{-3}$\ldots$\pm 2\cdot 10^{-2}$, but only on its narrow half:

\begin{table}[ht]
\centering
\caption{Comparison of \texttt{margin\_a} and \texttt{margin\_c} across the certified range.}
\label{tab:margin-a-vs-c}
\begin{tabular}{lllll}
\hline
half-width & worst box & \texttt{margin\_a} & \texttt{margin\_c} & limiting \\
\hline
$\pm 3\cdot 10^{-3}$\ldots$\pm 1\cdot 10^{-2}$ & \texttt{g00} & $0.130$--$0.145$ & $0.252$--$0.274$ & \texttt{margin\_a} \\
$\pm 1.5\cdot 10^{-2}$ & \texttt{gtop} & $0.140$ & $0.116$ & $\mathbf{\texttt{margin\_c}}$ \\
$\pm 2\cdot 10^{-2}$ & \texttt{gtop} & $0.132$ & $\mathbf{0.103}$ & $\mathbf{\texttt{margin\_c}}$ \\
\hline
\end{tabular}
\end{table}

That is, \texttt{margin\_c} (the very quantity that was being fixed) is not ``fully
removed as a bottleneck'' --- it continues to degrade with width ($0.27\to 0.10$) and
on the widest certified cell \Qcell{} again limits the result.
It is only true that on the larger part of the range the limit temporarily passes
to \texttt{margin\_a} (a quantity that never involved re-expansion), before \texttt{margin\_c}
again overtakes it. The numbers (\Qcell, \texttt{0.103}, the boundary by halves) do
not change from this --- only the mechanism has been clarified. Note on
\texttt{q\_max}: it grows non-monotonically (0.911 for the old \texttt{lemma\_U} $\to$ 0.995 for
the new on the same cell) --- not a regression, but a consequence of the fact
that previously some boxes passed self-embedding only after subdividing the
$\ell$-box ($K > 1$, independently confirmed from the log: e.g.\ \texttt{g39}/\texttt{gtop}
previously required $K=4$, now $K=2$), which concurrently also tightened
$q$; now self-embedding passes at smaller $K$, and $q$ is shown at this
coarser but already sufficient level of subdivision.

That is, \Ffam{} contains the open rectangle

\[
\Qcell = (0.56, 0.64) \times (0.76, 0.84) ,
\]

taken on the rigorous (not merely ``comfortable'') boundary --- all 46 boxes pass
at $\pm 4\cdot 10^{-2}$, including the thin case $\mathrm{margin} = 6.32\cdot 10^{-4}$. A reader who needs
a margin rather than the limiting width should take the sub-cell $\pm 3.5\cdot 10^{-2}$
($(0.565,0.635)\times(0.765,0.835)$), where the worst margin is $3.30\cdot 10^{-3}$ --- five times
thicker. By \S\ref{sec:10.7} (the transfer is rigorously closed for the entire \Qcell),
the worst margin of the full \Qcell{} ($6.32\cdot 10^{-4}$) exceeds the worst
transfer perturbation at the most unfavorable corner of \Qcell{} (\texttt{(0.64,0.84)},
$\Delta W = 8.787\cdot 10^{-6}$) by more than $\mathbf{71\times}$; on the comfortable sub-cell $\pm 3.5\cdot 10^{-2}$
the margin is more than $\mathbf{375\times}$. Both margins are \textbf{rigorously [C]}.

Further extension of the cell beyond \Qcell{} now runs into the geometry
of the family itself --- \texttt{margin\_c} (D2) actually vanishes for the
top-edge boxes ($\ell \to 4$, \texttt{g43}/\texttt{gtop}) at $\pm 4.5\cdot 10^{-2}$, while
\texttt{margin\_a} (D1) stays at $\approx 8.3\cdot 10^{-2}$ there
(Table~\ref{tab:cell-margins-nod4}) --- and not
into an artifact of the counting method, as was the case with \texttt{lemma\_U} (re-expansion)
and $q$ (sub-multiplicativity in place of the spectral radius). Extending \Qcell{} further
is possible only by a genuine geometric argument (for example, by proving \texttt{margin\_c}
for $\ell \to 4$ boxes separately, outside the $\ell$-uniform cell) --- this is no longer the
work of this step.

The off-center cell $c \in [0.55, 0.625]$, $r \in [0.70, 0.80]$ (center
\texttt{(0.5875, 0.75)}, outside \Qcell) has been recomputed by the same code without the D4 gate:
it was \texttt{31 / 46} (\texttt{q\_max} $= 1.037$, first failures at \texttt{g30}--\texttt{g32} on \texttt{bases\_q}), it became
\texttt{39 / 46} --- the removal of D4 honestly returned 8 boxes (\texttt{g30}--\texttt{g37}) that had previously
failed precisely on it. The remaining 7 (\texttt{g38}--\texttt{g43}, \texttt{gtop}) now fail on
\texttt{margin\_a}/\texttt{margin\_c} (gate \texttt{lemma\_U}) --- that is, shifting the center closer to the
top edge of $\ell \to 4$ by itself exhausts the geometric margin (D1/D2)
faster than mere expansion of the half-width around \texttt{(3/5,4/5)}. The cell as a whole
still does not belong to \Ffamz, but stops precisely where it should have:
on genuine geometry, not on the $q$ gate. \textbf{[C]}

\textbf{Check of the cut $(-\infty,0]$ --- as an explicit protocol field.} The condition ``$U$ does not
encroach on the cut of the principal branch of \texttt{Log}'' (part of D3) has been exposed as
a separate Boolean field \texttt{D\_cut\_ok} in the protocol of each certificate (\texttt{layer\_\allowbreak cert.py}),
rather than being checked only implicitly through the general gate \texttt{lemma\_U}. \textbf{Mutation test}
(\texttt{04-\allowbreak CODE/\allowbreak layer/\allowbreak layer\_\allowbreak mutation1.py}):
on a pair that a priori passes (base, \texttt{g00}) the Horner scheme has been deliberately
corrupted (shift of $\operatorname{Re}$ by $-10$, simulating ``encroached on the cut'') --- the gate correctly
caught the corruption (\texttt{pass=False}, \texttt{D\_\allowbreak cut\_\allowbreak ok=False}, \texttt{gate='lemma\_\allowbreak U'}), and after
restoration it again passes cleanly. So the check is not a ``rubber stamp'':
it is genuinely capable of failing a priori bad input. \textbf{[C]}

\textbf{Mutation test --- gates D1/D2} (\texttt{04-\allowbreak CODE/\allowbreak layer/\allowbreak layer\_\allowbreak mutation2.py}).
The first mutation test corrupts \texttt{horner} entirely, and therefore touches
all quantities at once (\texttt{margin\_a}, \texttt{margin\_c}, the cut) --- it does not show that
the gates are independent of one another. Targeted corruption: a wrapper around
\texttt{lemma\_\allowbreak U\_\allowbreak direct} that replaces ONLY \texttt{margin\_a} (or ONLY \texttt{margin\_c}) with an a priori
negative value, with \texttt{OK} recomputed by the same rule as in the original. The result
on the pair (base, \texttt{g00}): corrupting \texttt{margin\_a} yields \texttt{pass=False}, in the protocol
\texttt{D\_\allowbreak margin\_\allowbreak a=-0.001}, while \texttt{D\_\allowbreak margin\_\allowbreak c=0.2965} and \texttt{D\_\allowbreak cut\_\allowbreak ok=True} --- untouched; symmetrically
for corruption of \texttt{margin\_c} (\texttt{D\_\allowbreak margin\_\allowbreak c=-0.001}, \texttt{D\_\allowbreak margin\_\allowbreak a=0.166}, \texttt{D\_\allowbreak cut\_\allowbreak ok=True}).
The neighboring values here are higher than the base ones (\texttt{0.2842} and \texttt{0.1511} at $K=1$)
because after a failure \texttt{one\_box} escalates the subdivision to $K=8$, and the splitting
of the $\ell$-box narrows the intervals and raises the margin; this is the protocol of the last,
finest attempt. So the gates D1 (\texttt{margin\_a}), D2 (\texttt{margin\_c}), and the cut check (\texttt{D\_cut\_ok})
are three independent fields, each catching the corruption of precisely its own quantity,
without confusing itself with its neighbors and without failing/passing by association. \textbf{[C]}


\subsection{Transfer of certificates from $P(\ell)$ to $\Gs(\ell)$}\label{sec:10.7}

The certificates of \S\ref{sec:10.6} were computed for the interpolation center $P(\ell)$, while
Theorem D speaks of $\Gs(\ell)$. The gap is closed by the uniqueness ball of Theorem A,
$\norm{G^{\star} - P}_X \le 1.75\cdot 10^{-7}$, and Lipschitz estimates:

\begin{gather*}
\widehat{G}_{1} := |P(\ell)(1)| + 1.75\cdot 10^{-7}, \qquad M_{G'} := \norm{P'}_{X(c_{0},r_{0})}, \\
R_t := \max(c+r, c_{0}+r_{0}), \\
\kappa_G \le 1 + M_{G'} \cdot R_t \cdot \ell \cdot \widehat{G}_{1}^{\ell-1}, \\
\Delta U \le \kappa_G \cdot 1.75\cdot 10^{-7}, \\
\Delta W \le \ell \cdot (\sup|U|)^{\ell-1} \cdot \Delta U, \qquad \sup|U| \le 1 + 0.400 + \Delta U = 1.400 + \Delta U, \\
\Delta s \le R_t \cdot \ell \cdot \widehat{G}_{1}^{\ell-1} \cdot 1.75\cdot 10^{-7}.
\end{gather*}

(The formula for $\Delta W$ is a majorant of $|U_{1}^{\ell}-U_{2}^{\ell}|$ along the segment between $U_{1}$ and $U_{2}$.
For the interpolant, $\operatorname{Re} U > 0$ on the relevant circles is certified by (D3)/(I1c), \S\ref{sec:10.2}.
For the solution it follows from that same certificate together with $\Delta U$, which is derived
above without any appeal to the sign of $\operatorname{Re} U$: over the whole family the certified
minimum of $\operatorname{Re} U$ on the circles the transfer covers is $\ge 0.6025$ (fields \texttt{D\_\allowbreak ReU\_\allowbreak r\_\allowbreak val}, \texttt{I1\_\allowbreak ReU}, \texttt{I2\_\allowbreak ReU} of the layer certificates: the minimum over those three fields is $0.6025186598$ on the corridor cell reaching $c+r = 1.522$, box \texttt{g03}, and $0.6084424753$ over the 644 pointwise pairs; the field \texttt{D\_\allowbreak ReU\_\allowbreak R} belongs to the
outer circle $|t-c'| = 6/5$, which the transfer does not reach and no lemma
uses --- see the Remainder below), while the worst $\Delta U$ over all pairs ``box $\times$ layer'' is $8.26\cdot 10^{-7}$
(the figure recorded below in this same subsection), so
$\operatorname{Re} U_{\Gs} \ge 0.6025 - 8.26\cdot 10^{-7} > 0$ --- more than five orders of margin. (Earlier editions printed $0.5152$ here. That number is the minimum of \texttt{D\_\allowbreak ReU\_\allowbreak R}, the outer circle which this very sentence excludes, and not of the three fields named; over the 644 pointwise pairs that field descends to $0.5064163295$ on the pair $(0.70,0.75)/\texttt{g03}$, with the same value on $(0.70,0.80)/\texttt{g03}$, since that circle has the fixed radius $6/5$ and so depends on $c$ alone, i.e. below $0.5152$, so the old figure was not a bound for it either. Neither fact touches the conclusion: the outer circle enters no lemma.)
Both endpoints therefore lie in the convex set
$\{\operatorname{Re}>0\}$, and so does the segment --- this is enough to estimate
the derivative of the holomorphic $z^{\ell}$ by integration along the segment, without the assumption
that $U$ is real and positive. Separately about the constant $\sup|U|$: the value $1.379$ follows from
(C1), and (C1) was proved on the base circle $|t|\le c_{0}+r_{0}=1.4$, while the target layers reach
$c'+r'$ up to $1.522$ (the corridor of \S\ref{sec:10.13}; the pointwise layers and \Qcell{} stop at or below $1.50$).
Transferring the base number there is NOT ALLOWED: the sum
$\sum|a_{k}|\cdot \alpha^{-k\ell}\cdot (c+r)^{k}$ consists of positive terms and grows monotonically with $(c+r)$.
Therefore the same sum is computed by a second pass directly on the widest domain
$c+r = 1.53$ (\texttt{04-\allowbreak CODE/\allowbreak tests/\allowbreak \_\allowbreak nozeros\_\allowbreak check.py}, constants \texttt{CR\_WIDE}/\texttt{THRESH\_WIDE}):
$\sup|U-1| \le 0.399346$ on the worst box \texttt{g06} and $< 0.400$ on all 46.
The value $1.53$ is a round
number above the maximum of $c+r$ over the exhibited family, so by the monotonicity just stated this
single bound is conservative for every layer with $c+r \le 1.53$, hence for every member of
$\Ffam$ --- the pointwise layers, every point of
\Qcell, and every cell of the corridor among them.
The $+ \Delta U$ term in the formula is not for show: $0.400$ majorizes $|U-1|$ for the
INTERPOLANT $P(\ell)$ (this is exactly what \texttt{\_\allowbreak nozeros\_\allowbreak check.py} computes; its docstring says so
directly), whereas the segment along which the derivative is taken goes from $U_{P(\ell)}$ to $U_{\Gs}$,
and its second endpoint is at distance at most $\Delta U$ from the first. This is exactly how the code
computes it: \texttt{04-\allowbreak CODE/\allowbreak layer/\allowbreak layer\_\allowbreak kappa.py}, $\sup U = \text{arb}('1.400') + dU$. For the same reason and in
the same block, the base $\widehat{G}_{1}$ is $|P(\ell)(1)|$ PLUS the radius of the ball, not $|P(\ell)(1)|$ itself:
$\alpha = -1/G(1)$, hence $\alpha^{-\ell} = (-G(1))^{\ell}$, and $\kappa_G$ together with $\Delta s$ are also majorants of the
mean value along the segment from $P(\ell)$ to $\Gs(\ell)$. The base here is less than one ($|G(1)| \le 0.60$),
and the exponent $\ell-1$ is positive, so $x^{\ell-1}$ grows with the base --- omitting the radius would
underestimate the majorant rather than overestimate it. The same trick with the same justification
sits in \texttt{04-\allowbreak CODE/\allowbreak tests/\allowbreak \_\allowbreak kappaG\_\allowbreak bound.py} (line 209), which computes the same $\kappa_G$ on the base layer.
Numerically this changes nothing: the worst $\Delta U$ over all pairs ``box $\times$ layer'' equals $8.26\cdot 10^{-7}$
(same place where the worst $\Delta W$ is --- the corridor cell reaching $c+r = 1.522$ on box \texttt{g43}), that is,
the correction is of relative order $10^{-6}$. Justifying it instead through \texttt{margin\_c} (D2) for $\Gs$
would be circular: D2 for $\Gs$ is itself transferred from $P$ via $\Delta W$, which is precisely what
is being derived here.)

(Why the formula uses precisely $\norm{P'}$ --- the norm of the derivative of the INTERPOLANT,
and not of the true $\Gs$ (which is not computed explicitly). This is not carelessness,
but necessity: Theorem A's ball gives $\norm{G-P}\le 1.75\cdot 10^{-7}$ in the Wiener norm, but this norm
does not control the norm of the DIFFERENCE OF DERIVATIVES $\norm{G'-P'}$ (the sum
$\sum k|d_{k}|r^{k-1}$ is not bounded by $\sum|d_{k}|r^{k}$). The formula is saved not through $G'$,
but through the direct expansion
$|G(s_{G}(t))-P(s_{P}(t))| \le |(G-P)(s_{G}(t))| + |P(s_{G}(t))-P(s_{P}(t))| \le \norm{G-P} + \sup|P'|\cdot \Delta s$
--- in the second term it is precisely $\norm{P'}$ that appears.)

\textbf{Why \texttt{R\_t}, and not simply $(c+r)$.} The argument $t$ in these majorants
does not range over just one disk: for (D), (I1) the input lies in the target $D'$
($|t|\le c'+r'$), but for the reverse transfer (I2) the input lies in the BASE $D_{0}$
($|t|\le c_{0}+r_{0}=1.4$) --- the legality comes precisely from the self-embedding of the base layer
$s_{G}(D_{0})\subseteq D_{0}$ (certificate (D1), $\mathrm{margin}_a \ge 0.15$), and not from (I1a), which gives
an embedding for a different input circle; both certificates of \S\ref{sec:10.2} are transferred by THESE
SAME formulas (the ones written out just above). The scalar $(c+r)$ of the target layer is suitable only for one
of the two directions; for the six labels of the family carried in this
comparison with $c'+r'<1.4$ --- five of the 14 pointwise layers and the
corner \texttt{(0.56,0.76)} of $\Qcell$ (the worst ---
\texttt{(0.55,0.70)}, $c+r=1.25$) --- it would underestimate $\Delta U$/$\Delta W$/$\Delta s$ precisely where they are
compared with the I2 margins. $R_t := \max(c+r, c_{0}+r_{0})$ covers both circles at once.
For the global worst case (wide layers, $c+r\ge 1.4$) this makes no difference: there
$R_t=c+r$ coincides with the old expression literally.

where $(c,r)$ are the parameters of the \textbf{target} layer (enter purely algebraically, without
any $\Gs$-norm), and $M_{G'}$ is computed \textbf{once} on the base layer
(see below, why this is legitimate for any layer of the family, not only for one embedded in $D_{0}$).
Fixed here means fixed with respect to the \textbf{target layer}, and \emph{not} across the $\ell$-boxes:
$M_{G'}$ is built from $P(\ell)$ and therefore carries its own value in each box, exactly as
$\kappa_G$ does in Table~\ref{tab:kappa-samples}. The check this affords: $(\kappa_G-1)/\Delta s$ must be constant
along the layers of a single box --- and over all $1058$ records of \texttt{kappa\_\allowbreak full.json}
it is, to within a relative $6\cdot 10^{-10}$ (worst box \texttt{gtop}) --- while between the extreme boxes it ranges over a factor
of ${\approx}\,1.61$.

Honest interval values (supremum over the entire $x \in [-1,1]$, not over three points),
computed at $(c,r) = (3/5,4/5)$ --- i.e., so far only $1\times M_{G'}$, without the target-layer factor: \textbf{[C]}

\begin{table}[ht]
\centering
\caption{Sample boxes of $\kappa_G$ and transfer perturbations on the base layer.
The entries are the certified majorants of the formulas above, each rounded UP;
they are bounds, not rounded equalities.}
\label{tab:kappa-samples}
\begin{tabular}{lllll}
\hline
box & $\ell$ & $\kappa_G$ & $\Delta U$ & $\Delta W$ \\
\hline
\texttt{g00} & $2.025$ & $3.0094$ & $5.27\cdot 10^{-7}$ & $1.54\cdot 10^{-6}$ \\
\texttt{ganchor} & $3.000$ & $3.6066$ & $6.32\cdot 10^{-7}$ & $3.78\cdot 10^{-6}$ \\
\texttt{g42} & $3.915$ & $4.3748$ & $7.66\cdot 10^{-7}$ & $8.12\cdot 10^{-6}$ \\
\texttt{gtop} & $3.993$ & $4.2709$ & $7.48\cdot 10^{-7}$ & $8.24\cdot 10^{-6}$ \\
\hline
\end{tabular}
\end{table}

$\kappa_G$ grows with $\ell$ in general, but not strictly monotonically: by the full 46-box
recomputation on the base layer the maximum is attained on \textbf{\texttt{g43}} (\texttt{4.4197}), not on
\texttt{gtop} (\texttt{4.271}) and not on \texttt{g42} (\texttt{4.3748}); Table~\ref{tab:kappa-samples} gives only four sample boxes.
The honest bound on the base layer is $\mathbf{\kappa_G \le 4.42}$ ($4.4197 < 4.42$). $\Delta W > \Delta U$
(the factor $\ell|U|^{\ell-1} > 1$) --- verified and confirmed.

\textbf{Transfer between layers.} There is a subtlety here that is easy to miss:
if the formula required $\norm{G'}$ in the LAYER norm $X(c,r)$, one would need
to know $\norm{\Gs-P}_{X(c,r)}$, and the uniqueness ball of Theorem A gives this only in the
\textbf{base} norm $X(c_{0},r_{0})$ --- and the transfer would be justified only for
$D(c,r) \subseteq D_{0}$, that is, on 4 layers out of 14. In fact, the formula needs the norm $\norm{P'}$
(not $\norm{G'}$ --- see the clarification above) not on the whole $D(c,r)$, but only at points
$s(t)$, $t \in D(c,r)$ --- that is, on $s(D')$, and by certificate \textbf{(I1a)}
$s(D') \subseteq D_{0}$ for \textbf{any} $(c',r') \in \Ffamz$ without exceptions (not only for embedded ones).
It therefore suffices

\[
M_{G'} = \norm{P'}_{X(c_{0},r_{0})} \ge \sup_{D_{0}} |P'| \ge \sup_{s(D')} |P'|
\]

--- \textbf{one} quantity, computed \textbf{on the base layer by monomial re-expansion around the fixed point $c_{0} = 3/5$}
(not an interval --- re-expansion around a point does not create the interval-dependence problem,
unlike \texttt{lemma\_U} in \S\ref{sec:10.6}). The $(c+r)$ in the formula are parameters of the target layer, enter only
algebraically, and require no $\Gs$ norm. The condition $D(c,r) \subseteq D_{0}$ is \textbf{no longer needed anywhere} ---
the formula is valid for any exhibited layer $(c',r')$ (structurally: $M_{G'}$ is taken once on the base layer,
\texttt{(c,r)} enters algebraically; and the numerical substitution $\sup|U| \le 1.400 + \Delta U$ is certified on
$c+r \le 1.53$, which covers the exhibited family entirely --- its maximum is $1.522$, attained on the corridor).
The same for $\Delta s$
(it did not depend on the layer norm before either: $|G(1)|$ is a base quantity,
$|h(1)| \le \norm{h}_{X(c_{0},r_{0})}$ by (E2)).

Honest recomputation (\texttt{04-\allowbreak CODE/\allowbreak layer/\allowbreak layer\_\allowbreak kappa.py}) over all 14 layers, the base
layer \texttt{(0.60, 0.80)}, the 4 corners of \Qcell, the far corners of the last four of the
seven corridor cells of \S\ref{sec:10.13}, and 46 boxes
($46 \times 23 = 1058$ pairs --- that is exactly how many records sit in \texttt{kappa\_\allowbreak full.json}):
$\kappa_G$ grows linearly with $(c+r)$, as it should; the worst value over all 23 labels is at
its widest point, the corridor cell $[0.588,0.612]\times[0.880,0.910]$ with $c+r = 1.522$, box \texttt{g43}:
$\kappa_G = 4.7177$, $\Delta W = 8.983\cdot 10^{-6}$ --- growth from the base $\kappa_G(g43) = 4.4197$ by only
$+6.7\,\%$. Among the pointwise layers the worst is $c+r=1.50$ (\texttt{(0.70,0.80)}, \texttt{(0.60,0.90)},
\texttt{(0.65,0.85)} --- give ONE and the same number, as expected, since only the sum $c+r$ appears in
the formula): $\kappa_G = 4.6640$, $\Delta W = 8.880\cdot 10^{-6}$. For specific pairs needed in
\S\ref{sec:10.5}/\S\ref{sec:10.6}:
$\Delta W((0.60,0.90), g34) = 6.361\cdot 10^{-6}$; $\Delta W$ in the angle \Qcell{} \texttt{(0.64,0.84)} on \texttt{g43}
equals $8.787\cdot 10^{-6}$ --- this is the worst corner of \Qcell, since $c+r = 1.48$ is maximal there.

Summary: the certificates of \S\ref{sec:10.6} are transferred from $P(\ell)$ to $\Gs(\ell)$ \textbf{rigorously [C]} on
\textbf{all} of the presented family --- the 14 exhibited layers, the \textbf{entire} cell \Qcell, and every cell of the
corridor --- without division into ``embedded'' and ``others''. The worst ``margin / perturbation''
ratio: on pointwise layers ---
$1.954\cdot 10^{-3} / 8.648\cdot 10^{-6} \approx 225.9\times$ (pair \texttt{(0.55,0.90)/\allowbreak g43}; minimum over all 644 pairs, by
collapsing \texttt{full\_\allowbreak recheck\_\allowbreak 644.log} with \texttt{kappa\_\allowbreak full.json});
on \Qcell{} --- $6.32\cdot 10^{-4} / 8.787\cdot 10^{-6} \approx 71.9\times$ on the rigorous boundary $\pm 4\cdot 10^{-2}$ and
$3.30\cdot 10^{-3} / 8.787\cdot 10^{-6} \approx 375\times$ on the comfortable sub-cell $\pm 3.5\cdot 10^{-2}$ (\S\ref{sec:10.5}, \S\ref{sec:10.6});
on the corridor --- the last four cells of the table of \S\ref{sec:10.13}, the four with the narrowest margins, give
$1.06\cdot 10^{-2} / 8.890\cdot 10^{-6} \approx 1192\times$, $4.57\cdot 10^{-3} / 8.983\cdot 10^{-6} \approx 508\times$,
$1.58\cdot 10^{-3} / 8.946\cdot 10^{-6} \approx 176\times$ and $1.14\cdot 10^{-3} / 8.839\cdot 10^{-6} \approx 128\times$
(cell margins from the table of \S\ref{sec:10.13}, worst gate against worst perturbation).

\textbf{Narrow layers separately.} Six labels of the family carried in this
comparison have $c'+r'<1.4$
(\texttt{(0.55,0.70)}, \texttt{(0.60,0.70)}, \texttt{Q corner (0.56,0.76)}, \texttt{(0.55,0.80)},
\texttt{(0.60,0.75)}, \texttt{(0.575,0.80)}; worst $c+r=1.25$). After $R_t=\max(c+r,1.4)$
their $\kappa_G$, $\Delta U$, $\Delta W$, $\Delta s$ on each box coincide numerically with the values of
the BASE layer \texttt{(0.60,0.80)} for the same box (both points give $R_t=1.4$). The worst
case among the narrow layers is box \texttt{g43}: $\Delta W=8.415\cdot 10^{-6}$ versus
\texttt{I2\_margin\_W}$=6.80\cdot 10^{-3}$ --- margin $808.0\times$ (was $896.3\times$ from an erroneous record,
that is, the number itself was anti-conservative, and not only its
justification; that record is not in the package and the $896.3\times$ is not
reproducible from it). For $\Delta s$ the maximum, as before, sits on the widest point --- the
narrow layers do not approach it, and the \texttt{I2\_margin\_s} margin does not change.

Separately --- the per-layer number for $\Delta s$, which was previously mentioned only in words.
$\Delta s$ on each layer is computed by the EXACT version of the estimate of \S\ref{sec:6.5} item 3
(with the explicit factor $|G(1)|^{\ell-1}$, which \S\ref{sec:6.5} does not expose conservatively --- the same
``naive/exact'' fork as for $\kappa_G$, \S\ref{sec:6.4.3}), so the number does not coincide with the
number of \S\ref{sec:6.5} ($3.20\cdot 10^{-6}$), although both estimates are valid; over all 46 boxes and 23
layer labels (field \texttt{dS} in \texttt{kappa\_\allowbreak full.json}) the maximum is $\Delta s \le 2.399\cdot 10^{-7}$ on the
widest corridor cell $\times$ \texttt{g43}; among the pointwise layers it is $2.364\cdot 10^{-7}$ on
\texttt{(0.70,0.80)/\allowbreak g43}.
Collapsing with the certificate's s-margins over all 644 pairs gives the worst ratio
$\mathbf{\approx 5083\times}$ on the pair \texttt{(0.70,0.75)/\allowbreak g00} --- the same place where the global
minimum of the margin \texttt{I2\_margin\_s} is attained, and where $\Delta s \le 2.157\cdot 10^{-7}$. The ratio is
formed from the certificate's own $1.0965\cdot 10^{-3}$ (\S\ref{sec:10.9} item 5; row 46 of
Appendix~\ref{sec:A}), not from the three-digit $1.09\cdot 10^{-3}$ that \S\ref{sec:10.6} prints for the
same pair --- the latter would give $5053\times$. That is,
in the $s$-direction the worst ratio, $5083\times$, is some seventy times the worst ratio in
$W$ ($71.9\times$ on \Qcell), and $\Delta W$ remains the limiting one.

\textbf{Remainder (not closed, smaller in scale).} Certificate D3 numerically checks
$\operatorname{Re} U > 0$ not only on $|t-c'|=r'$, but also on the fixed external circle
$|t-c'|=6/5$ (\S\ref{sec:10.2}). Verified: the analytic argument that is ACTUALLY used in the text
(``Remark on the boundaries'', \S\ref{sec:10.2}, and everywhere below) --- the harmonicity of $\operatorname{Re} U$ and the
minimum principle --- rests \textbf{only} on the circle $|t-c'|=r'$; the external \texttt{6/5} is not
cited anywhere in the proofs of Lemma 1, Lemma 2, or Theorem D. In the package itself this
circle is NOT redundant: \texttt{04-\allowbreak CODE/\allowbreak lemma\_\allowbreak domains.py} takes on it $\mathtt{M\_W\_R} = \max |U^{\ell}|$ --- the input
of the Cauchy tail in $W$ --- and $\operatorname{Re} U > 0$ on it is needed for the uniqueness of the principal
branch of $U^{\ell}$ already on the larger circle $D(c,R)$, without which this estimate would not
make sense. For Theorem D, neither one nor the other is required. Numerically, the check is
nowhere limiting (over the $966$ cell records $\texttt{D\_ReU\_R} \ge 0.5152$, and over the $644$ pointwise pairs the same field descends to $0.5064163295$, as recorded earlier in this subsection --- either way with a large margin, and neither figure enters a proof. The field is \texttt{D\_\allowbreak ReU\_\allowbreak R} inside \texttt{worst} of each element of \texttt{boxes}, in the five files \texttt{03-\allowbreak CERTIFICATES/\allowbreak layer/\allowbreak cells3\_\allowbreak bisect.json}, \texttt{cells3\_\allowbreak hist.json}, \texttt{cells3\_\allowbreak hw.json}, \texttt{cells3\_\allowbreak offcenter.json}, \texttt{corridor.json} --- $966$ records in all --- and in \texttt{full\_\allowbreak recheck\_\allowbreak 644.json}), it does not affect the
resulting \Qcell{} and has not been removed from \texttt{ok\_D} --- but the dependence of the proof of
Theorem D on it has been removed.

\textbf{Connection with the premise of Theorem C.} The same Lipschitz constant $\kappa_G$ is also
computed by \texttt{04-\allowbreak CODE/\allowbreak tests/\allowbreak \_\allowbreak kappaG\_\allowbreak bound.py} --- by the formula of \S\ref{sec:6.4.3}(a), on the base layer.
This is a second, independently written computation of the same quantity, and serves as a cross-check
for the formula above.

Both expressions coincide algebraically: $\alpha_{G}^{-\ell-1}/|G(1)|^{2} = |G(1)|^{\ell-1}$,
since $\alpha_{G} = -1/G(1)$. Estimating the factors separately ($\alpha_{G}^{-\ell-1} \le 1$ separately
and $1/|G(1)|^2$ separately) means losing the factor $|G(1)|^{-(\ell+1)}$,
that is, by about a factor of ten; therefore both computations take $|G(1)|^{\ell-1}$ in full,
with an explicit check that the base is less than one --- otherwise the direction of
monotonicity with respect to the exponent would change.

The implementations are nevertheless different: here $M_{G'}$ is taken once as
$\norm{P'}_{X(c_{0},r_{0})}$, while in \texttt{\_\allowbreak kappaG\_\allowbreak bound.py} the quantity $\sum k|g_{k}| r^{k-1}$
is recomputed on each sub-interval of the $\ell$-box, inflated by the uniqueness ball
($\rho = 1.75\cdot 10^{-7}$) and accompanied by a check $\Delta s < \mathrm{margin}_a$. With the same coarse subdivision
$K = 1$ they yield \texttt{4.415912} (output of \texttt{\_\allowbreak kappaG\_\allowbreak bound.py}) and \texttt{4.4197} (output of
\texttt{layer\_\allowbreak kappa.py} via \texttt{kappa\_\allowbreak full.json}) on \texttt{g43} --- discrepancy \texttt{0.09 \%}; the sample boxes of
Table~\ref{tab:kappa-samples} are reproduced in the same way (\texttt{gtop}: \texttt{4.269531} versus \texttt{4.2709};
\texttt{g42}: \texttt{4.371065} versus \texttt{4.3748}). The worst box \texttt{g43} is the same in both computations.
The published $\kappa_G \le 4.14$ (in fact \texttt{4.130319}) in \S\ref{sec:6.4.3} --- is the same computation at
$K = 8$, where the interval inflation in $\ell$ is smaller. \textbf{[C]}

The bound $\kappa_G \le 4.42$ is confirmed in four ways: by direct interval computation
(\texttt{04-\allowbreak CODE/\allowbreak layer/\allowbreak layer\_\allowbreak kappa.py}), by manual substitution into the formula
\S\ref{sec:6.4.3}(a), by the algebraic identity $|\alpha^{-\ell-1}|/|G(1)|^{2} = |G(1)|^{\ell-1}$ that
connects both expressions, and by the second, independently written computation
\texttt{04-\allowbreak CODE/\allowbreak tests/\allowbreak \_\allowbreak kappaG\_\allowbreak bound.py}, agreeing with it within \texttt{0.09 \%} under equal
subdivision.

\textbf{Precision and computation parameters.} The code \texttt{04-\allowbreak CODE/\allowbreak layer/\allowbreak } everywhere uses
$\texttt{ctx.prec} = 256$. There is no uniform precision across the package, and it would be wrong
to say ``in the rest of the package 512'' --- the actual map is as follows:
\texttt{1536} --- \texttt{tests/\allowbreak \_\allowbreak xresid\_\allowbreak tail\_\allowbreak F\_\allowbreak l25.py:67} (\texttt{PREC}, tail of $F$ at $\ell = 5/2$);
\texttt{1024} --- \texttt{tests/\allowbreak \_\allowbreak reu\_\allowbreak min\_\allowbreak l25.py:53} (\texttt{PREC}, $\min \operatorname{Re} U$ at $\ell = 5/2$);
\texttt{512} --- libraries \texttt{rig\_\allowbreak lib.py:40}, \texttt{tube\_\allowbreak lib.py:45} (\texttt{PREC}),
\texttt{beta\_\allowbreak bound\_\allowbreak lib.py:36}, driver \texttt{build\_\allowbreak majorant\_\allowbreak rigorous.py:36}, watchdogs
\texttt{tests/\allowbreak \_\allowbreak numbers\_\allowbreak check.py:2569} and \texttt{tests/\allowbreak sanity\_\allowbreak miller.py:17};
\texttt{300} --- \texttt{tests/\allowbreak \_\allowbreak tau\_\allowbreak test.py:25};
\texttt{256} --- \texttt{hyp\_\allowbreak rouche.py:70}, \texttt{kmode\_\allowbreak cert.py:49}, \texttt{tube\_\allowbreak offset.py:49}, all of
\texttt{04-\allowbreak CODE/\allowbreak layer/\allowbreak }, and fourteen of the nineteen tests that set the precision
explicitly --- thirteen of them fix it for the whole module, the fourteenth ---
\texttt{tests/\allowbreak \_\allowbreak arb\_\allowbreak sanity\_\allowbreak battery.py}, whose module-level \texttt{PREC} is also this, although
the F sub-tests deliberately change it in both directions (that is the point of the sub-test);
the twentieth file, \texttt{tests/\allowbreak pure\_\allowbreak interval.py}, takes the precision as a parameter;
\texttt{128} --- \texttt{tube\_\allowbreak accept.py:26} (cross-check of finished tube boxes).
This map says where a precision is \emph{set} in the source. What a given run
actually used is recorded in the certificate it produced, and the two need not
agree: \texttt{rig\_\allowbreak lib.py:40} sets $512$ at import time, but the
existence engine \texttt{rig2\allowbreak fix3.py} takes its own \texttt{prec}
(default $256$) and reassigns \texttt{ctx.prec} on entry, so every one of the
$46$ certificates \texttt{03-\allowbreak CERTIFICATES/\allowbreak existence/\allowbreak box\_\allowbreak g*.json}
records \texttt{"prec": 256}. Where the two could differ, the certificates are
what the statements rest on.
The diversity is real, but threatens no conclusion:
the \texttt{arb} arithmetic is interval-based, and if \texttt{256} bits were not enough, this would
manifest itself as an inflated (not just inaccurate) radius of the ball and a failure of
the certificate itself, not as a silent error; on the scale of the considered margins
($10^{-3}$\ldots$10^{-4}$) \texttt{256} bits ($\approx 77$ decimal digits) --- a margin of many orders of magnitude.
The same argument also covers the \texttt{128} bits of \texttt{tube\_\allowbreak accept.py}
($\approx 38$ decimal digits): this step does not build boxes, but cross-checks already built ones,
and the checked quantities are separated from the thresholds by tens of orders of magnitude
beyond the resolution of such arithmetic.
\texttt{M\_\allowbreak ARCS\_\allowbreak CAUCHY} of the package (external $R=6/5$) equals \texttt{128} and coincides with the
choice here; the \texttt{64} versus \texttt{128} disagreement concerned only the internal circle
$r'$ (\S\ref{sec:10.6}). $N_{p} = 80$ coincides.



\subsection{Numerical illustration: invariance is not tautological}
\label{sec:10.8}

It has been verified separately that changing the layer actually changes the object, and is not
ignored by the engine. On the anchor box (\texttt{ganchor}, $\ell = 3$) the matrix $\DT$ has been
assembled in the layer basis on four layers; the matrices themselves were compared, not only
the spectrum (float64 illustration, not an interval certificate):

\begin{table}[h]
\centering
\caption{Four sample layers and the corresponding $\DT$ matrices.}
\label{tab:10.8-layers}
{\footnotesize\setlength{\tabcolsep}{3pt}
\begin{tabular}{cc>{\centering\arraybackslash}p{0.107\textwidth}>{\centering\arraybackslash}p{0.138\textwidth}cc}
\toprule
layer & $\|M\|_F$ & $\|M - M_0\|_F / \|M_0\|_F$ & \texttt{max} element difference & $\delta$ & outside circle \\
\midrule
$(0.60,0.80)$ & 7.991 & 0 & 0 & 6.084689074262122 & 1 \\
$(0.65,0.85)$ & 8.121 & 0.121 & 0.459 & 6.084689074262116 & 1 \\
$(0.55,0.90)$ & 7.529 & 0.290 & 0.779 & 6.084689074262118 & 1 \\
$(0.70,0.75)$ & 13.908 & 1.178 & 1.982 & 6.084689074262119 & 1 \\
\bottomrule
\end{tabular}}
\end{table}
\FloatBarrier

The radii of the interval balls of the elements are $\leq 1.2\cdot 10^{-73}$, so the discrepancy of the
matrices is not numerical noise. Table~\ref{tab:10.8-layers} shows 4 sample layers; a full run over all 14
exhibited layers (\texttt{04-\allowbreak CODE/\allowbreak layer/\allowbreak layer\_\allowbreak dt\_\allowbreak probe.py}) gives the maximum difference of the matrices
\textbf{118}\,\% --- the same pair \texttt{(0.70,0.75)}, extending the sample found nothing worse --- and the
smallest non-zero difference \textbf{9.8}\,\% on \texttt{(0.575,0.80)} (these two ends are what the range
``9.8--118\,\%'' quoted elsewhere in the text refers to); the maximum
deviation of $\delta$ from the base layer over all 14 layers is $8.9\cdot 10^{-15}$ (exactly: $8.882\cdot 10^{-15}$;
this is NOT the same as the ``full spread'' $1.51\cdot 10^{-14}$ from \S\ref{sec:9}/Appendix~\ref{sec:A} --- that quantity is
$\max-\min$ between layers, this one is the maximum deviation from ONE, the base layer); on each of the
14 layers (really, not by extrapolation) $n_{\mathrm{out}} = 1$ has been verified. The top spectrum over
all 14 layers coincides to 12+ digits (the same float64 illustration, not a
certificate). That is the agreement observed; the threshold the acceptance test
actually requires is looser --- a tolerance of $10^{-11}$ on the twenty leading
eigenvalues (\S\ref{sec:9}).


\subsection{What is NOT proved}
\label{sec:10.9}

\begin{enumerate}
\item The statement concerns the \textbf{certified family} $\Ffam$, not an arbitrary
   layer: outside $\Ffam$ nothing is asserted \emph{about the solution} $\Gs$
   (the $P(\ell)$-version of items 1--2 reaches further, over all of $\Ffamz$ ---
   the remark in \S\ref{sec:10.5}). Certified continuously is the
   rectangle $\Qcell = (0.56, 0.64) \times (0.76, 0.84)$ (half-width $4\cdot 10^{-2}$)
   around $(3/5, 4/5)$; besides it, 14 separate layers scattered over
   $[0.55, 0.70] \times [0.70, 0.90]$ and the 7-cell corridor of
   \S\ref{sec:10.13} (which reaches $r = 0.915$, outside that rectangle,
   and contains the widest exhibited point) are exhibited in $\Ffam$.
   The exhibited set is not connected: nine of those fourteen layers lie
   outside $\Qcell$ and outside every corridor cell. What is joined by a
   connected certified region is the base layer and one of the farthest
   exhibited ones --- $\Qcell$ together with the corridor. \emph{(Below is the
   history of two refinements of this bound in the course of the work on
   the section, given for reproduction context; it does not affect the
   final status stated above.)} The size of $\Qcell$ grew twice in the
   course of the work on the section --- both times by removing excess
   conservativeness of the computation, not by changing the problem.
   First (\S\ref{sec:10.6}) --- the engine's \texttt{lemma\_U}, which re-expanded the monomial
   polynomial around the interval centre $c'$ (the interval dependency
   problem): direct monomial evaluation raised $\Qcell$ from $\pm 3\cdot 10^{-3}$ to
   $\pm 2\cdot 10^{-2}$ (6.6$\times$). Then --- the numerical gate D4
   ($q<1$, $q_s<1$, \S\ref{sec:10.2}): a rigorous argument via the spectral radius of
   a Banach algebra (\S\ref{sec:10.4}) showed that the already available
   \texttt{margin\_a}>0/\texttt{margin\_c}>0 (D1/D2) suffice for compactness and that the
   gate can be removed --- and this widened $\Qcell$ by a further factor of
   two, to $\pm 4\cdot 10^{-2}$. Further widening ($\pm 4.5\cdot 10^{-2}$ no longer passes,
   \S\ref{sec:10.6}) is blocked not by a computational artefact but by the geometry
   itself: \texttt{margin\_c} genuinely approaches zero at the boxes of
   the upper edge $\ell \to 4$ (\texttt{g43}, \texttt{gtop}) under a cell expansion uniform
   over all of $\ell \in [2,4]$, while \texttt{margin\_a} there is still of order
   $8\cdot 10^{-2}$. This remains a limitation of the \textbf{exhibited}
   enclosure, not an assertion that such layers are inadmissible (every
   separate layer that has been checked passes) --- but closing the whole
   rectangle beyond $\pm 4.5\cdot 10^{-2}$ is now possible only through a separate
   geometric argument for the $\ell \to 4$ boxes (\S\ref{sec:10.10}, path 3), not by
   correcting the computational method.

\item There is no claim that ``$\delta(\ell)$ is independent of any reasonable space''.
   What is proved is independence of $(c,r)$ within a single family of
   weighted $\ell^1$-spaces. Other norms (for example $H^2$, or spaces with a
   different weight) are not considered.

\item $\ell$ --- only $[2,4]$: outside this segment there are neither branches
   $P(\ell)$ nor Theorem A.

\item At the upper edge $\ell \to 4$ the family is thin even without the $q$
   gate (\S\ref{sec:10.2}, \S\ref{sec:10.4}), and without subdivision of the
   $\ell$-box ($K = 1$) \textbf{nine of the 644 pointwise pairs do not pass at
   all}: boxes \texttt{g36}--\texttt{g43} and \texttt{gtop} of the layer
   $(0.55, 0.90)$. The gate that fails there is \texttt{I1\_margin\_W}, which
   goes negative, from $-4.73\cdot 10^{-4}$ on \texttt{g36} to
   $-3.80\cdot 10^{-2}$ on \texttt{g43} ($-7.85\cdot 10^{-3}$ on \texttt{gtop}).
   \texttt{margin\_c} is not the pre-subdivision bottleneck for the pointwise
   layers at all: over all 644 pairs at $K = 1$ its minimum is
   $2.22\cdot 10^{-2}$, on that same pair $(0.55,0.90)/\texttt{g43}$. At
   $K \le 2$ all 644 pass, worst margin $1.0965\cdot 10^{-3}$
   (\S\ref{sec:10.6}); none of these $K = 1$ figures is among the stored
   certificates, which hold the accepted values. The same thinness surfaces at the very
   boundary of the continuous $\Qcell$ ($\pm 4\cdot 10^{-2}$: $\mathrm{margin} = 6.32\cdot 10^{-4}$, the
   same order of magnitude, item 1). The statement remains rigorous, but
   the region there is narrow, and this is now the \textbf{only} limit --- a
   geometric one (\texttt{margin\_c} genuinely approaches zero; \texttt{margin\_a} stays
   at $\approx 8.3\cdot 10^{-2}$), not one of numerical method.

\item All 644 pointwise certificates have been recomputed by the same
   certifier (\texttt{layer\_\allowbreak cert.one\_\allowbreak box}, monomial \texttt{lemma\_U}, K-escalation
   $K \in \{1,2,4,8\}$): \textbf{644/644 passed, not a single failure}, the worst
   margin being $1.0965\cdot 10^{-3}$ at the pair \texttt{(0.7,0.75)/\allowbreak g00}. This is a
   repeat run of \textbf{one and the same code}, not a second independent
   implementation: it confirms independence of the PACKAGE engine and of
   the old \texttt{lemma\_U}, but not of the section's own code. A separately
   written second implementation of the transfer margins \texttt{I1}/\texttt{I2} (its
   own traversal of arcs and subintervals, without \texttt{layer\_\allowbreak lemma\_\allowbreak u} and
   without the \texttt{one\_box} structure --- \texttt{04-\allowbreak CODE/\allowbreak layer/\allowbreak layer\_\allowbreak crosscheck.py})
   agrees with the certifier on 6 pairs with a relative discrepancy of
   $\leq 4.16\cdot 10^{-3}$. In it the polynomial is summed by powers rather
   than by the Horner scheme, so that the agreement is not a consequence
   of shared code. A second separation was intended --- $A = (-1/G(1))^{-\ell}$
   taken via $\exp(-\ell \cdot \log(\cdot))$ rather than by the power operator
   --- but it is not one: \texttt{arb} implements the power operator as exactly
   that composition whenever the exponent is not an exact number, and here
   $\ell$ is an interval. On every branch entering this reconciliation the
   two values of $A$ agree bit for bit, in midpoint and in radius; the
   numbers gate checks this (\code{tests/\allowbreak \_numbers\_check.py}, item I12c). The direction of the discrepancy is
   NOT one-sided, and it would be an error to claim otherwise: the Horner
   scheme and direct summation have different interval dependency trees,
   neither of which majorises the other identically. Measurement from
   \texttt{crosscheck.json}: on five of the six pairs \texttt{alt} is more conservative
   than the certifier, while on the pair \texttt{(0.55,0.80)/\allowbreak g00} it is, on the
   contrary, thinner than it by $3.05\cdot 10^{-6}$ and $2.22\cdot 10^{-6}$ in the two
   \texttt{\_W}-margins. The reconciliation checks the MAGNITUDE of the
   discrepancy ($\leq 10^{-2}$), not its sign; the sign is only measured and
   written into the protocol. The margins \texttt{I1\_margin\_s} / \texttt{I2\_margin\_s},
   into which the polynomial does not enter at all, are for that reason
   \textbf{not covered} by the second implementation: with $A$ shared, the two
   sides are the same expression over the same inputs, so the $0$ recorded
   for them in \texttt{crosscheck.json} is an identity and not a measurement.
   What this reconciliation measures is the pair \texttt{I1\_margin\_W} /
   \texttt{I2\_margin\_W}, where the polynomial is summed by two different schemes
   and the discrepancy is real.

   \textbf{What this reconciliation does not cover.} In \texttt{crosscheck.json}
   four quantities are listed (the field \texttt{key} in \texttt{rows}) ---
   \texttt{I1\_margin\_s}, \texttt{I1\_margin\_W}, \texttt{I2\_margin\_s}, \texttt{I2\_margin\_W} ---
   of which only the two \texttt{\_W} margins are genuinely compared, for the
   reason just given. The margins
   \texttt{D\_margin\_a} and \texttt{D\_margin\_c}, that is, the very gates of self-embedding
   of a layer (\S\ref{sec:10.2}), are \textbf{not cross-checked at all} by the second
   implementation: they have gates and mutation tests
   (\texttt{layer\_\allowbreak mutation1/\allowbreak 2}), but there is no independent second computation.
   This is a known gap in the enclosure, not something claimed and left
   undone: the statement above speaks only of the transfer margins
   \texttt{I1}/\texttt{I2}.

   The limits of this reconciliation. It is \textbf{not} a check of the
   arithmetic itself: both implementations rely on one library
   (\texttt{python-flint}/\texttt{arb}), one \texttt{rig\_\allowbreak lib.P\_\allowbreak at\_\allowbreak interval} and --- essentially ---
   on LITERALLY the same function for enclosing an arc in an \texttt{acb} box
   (\texttt{my\_arc} was copied without modification), so the arc enclosure is not
   checked here at all. For independence of the library, and not only of
   the computational scheme, the certificates would have to be recomputed
   with a second interval engine (for example \texttt{mpmath.iv}) --- this has not
   been done.

\end{enumerate}


\subsection{Anticipated reviewer questions}
\label{sec:10.10}

\textbf{``Why is only a rectangle of half-width $4\cdot 10^{-2}$ closed continuously? Is the
region of admissible layers really that small?''}

No, and this is evident from our own data: 14 separate layers scattered over the whole of $[0.55, 0.70] \times [0.70, 0.90]$ pass all 46 boxes each (644 certifications, \texttt{03-\allowbreak CERTIFICATES/\allowbreak layer/\allowbreak full\_\allowbreak recheck\_\allowbreak 644.json}), and across all of them the self-embedding margins never drop below $2.44 \cdot 10^{-2}$ for \texttt{D\_margin\_a} and $5.28 \cdot 10^{-2}$ for \texttt{D\_margin\_c} (\texttt{layer/\allowbreak worst\_\allowbreak by\_\allowbreak gate.json}; these two gates specifically, not all six --- for \texttt{I1}/\texttt{I2} the worst margins are noticeably narrower, see \S\ref{sec:10.6}), that is, orders of magnitude above zero. The size of $\Qcell$ is a property of the \textbf{enclosure}, not of the problem, and
during the work on the layer family this was corrected twice. Initially (the draft
version) $\Qcell$ was $\pm 3\cdot 10^{-3}$, and the degradation was attributed to ``interval
enclosure inherently loses accuracy''. That turned out to be inaccurate: there
were \textbf{two} culprits, and both have been removed. The first --- the engine's
\texttt{lemma\_U} (\texttt{tube\_\allowbreak lib.verify\_\allowbreak lemma\_\allowbreak U}), which re-expanded the monomial
polynomial around the interval $c'$ (the interval dependency problem); a
direct Horner scheme on the monomial coefficients (\S\ref{sec:10.6}) raised $\Qcell$ to
$\pm 2\cdot 10^{-2}$. The second --- the numerical gate D4 ($q<1$, $q_s<1$, \S\ref{sec:10.2}): for
compactness of the composition operator it required a condition stronger than
necessary. A rigorous argument via the spectral radius of a Banach algebra
(\S\ref{sec:10.4}) shows that the already available \texttt{margin\_a}>0/\texttt{margin\_c}>0 (D1/D2)
suffice --- the gate can be removed, and in practice this eliminated a real
(not merely theoretical) bottleneck: $\Qcell$ grew by a further factor of two, to
$\pm 4\cdot 10^{-2}$ (twice in half-width and fourfold in area relative to
$\pm 2\cdot 10^{-2}$; $13\times$ in half-width and $177\times$ in area relative to the
original $\pm 3\cdot 10^{-3}$).

\textbf{``Then why not improve the enclosure all the way?''}

Both gates found initially have been removed --- this is precisely ``all the
way'' in the sense of the known defects of the computational method. The
remaining limit (\S\ref{sec:10.6}, \S\ref{sec:10.9} items 1/4) is no longer an artefact: at
$\pm 4.5\cdot 10^{-2}$ the margin \texttt{margin\_c} \textbf{genuinely} vanishes at the
boxes of the upper edge $\ell \to 4$ (\texttt{g43}, \texttt{gtop}) --- \texttt{margin\_a}
does not, it stays at $\approx 8.3\cdot 10^{-2}$ --- that is, the geometry of the
family $\Ffamz$ itself (the condition $s(D')\subseteq D'$/$W(D')\subseteq D'$) ceases to hold there
under such a cell expansion uniform over $\ell \in [2,4]$. There are three standard
ways of widening the enclosure beyond this point:

\begin{enumerate}
\item \textbf{Direct evaluation in the original basis, without re-expansion.} Already
   carried out --- for \texttt{lemma\_U} (\S\ref{sec:10.6}) and, in effect, by the discovery of
   redundancy for $q$ (\S\ref{sec:10.2}, \S\ref{sec:10.4}). Both sources of excess conservativeness
   have been removed; there is nothing left to widen along this path ---
   the limit now lies in the problem itself.

\item \textbf{Affine arithmetic / Taylor models in $(c', r')$} instead of naive
   interval arithmetic --- not done; it could shift $\pm 4.5\cdot 10^{-2}$ slightly, but
   it will not change the fact that the $\ell \to 4$ boxes are inherently pressed
   closer to the boundary of $\Ffamz$.

\item \textbf{A separate argument for the $\ell \to 4$ boxes (not uniform over all of
   $\ell \in [2,4]$)} --- for example, a cell narrower in $r'$ precisely there, since
   the thinness is concentrated in 4--5 specific boxes of the upper edge (\S\ref{sec:10.6},
   \S\ref{sec:10.9} item 4) --- not done; probably the cheapest way to widen $\Qcell$ further,
   should that be needed.
\end{enumerate}

\textbf{``Are 14 layers and one rectangle enough to speak of
independence?''}

The statement of Theorem D does not rest on the number of layers exhibited:
items 1--2 (compactness, coincidence of the non-zero spectra of $P(\ell)$) hold
for \textbf{any} two layers of $\Ffamz$, i.e.\ satisfying the four condition groups of \S\ref{sec:10.2} --- a finite set of
interval inequalities that can be checked for a given layer in minutes (the certifier
gates D/I1/I2; (E1)/(E2) are the hand check named there). But
for the full conclusion about the solution $\Gs(\ell)$ (not only about the
interpolant $P(\ell)$) this is not sufficient: the transfer of \S\ref{sec:10.7} is also
needed, and it is certified on $c+r\leq 1.53$ --- a domain that contains the whole of the
exhibited family, but not every layer one might write down (see \S\ref{sec:10.5}, the remark after
Consequence 2).
The exhibited $\Qcell$, the 14 layers and the seven-cell corridor are needed
only in order to show that the condition is non-vacuous, that it holds on layers
far removed from one another (the matrices differ by 9.8--118 \%; that range is a
float64 illustration, not an interval certificate, \S\ref{sec:10.8}), and that the
base layer and one of the farthest of the exhibited layers are joined by a connected
certified region --- $\Qcell$ together with the corridor. The exhibited set as a
whole is not connected: nine of the fourteen pointwise layers lie outside
$\Qcell$ and outside every corridor cell.

\textbf{``What happens at the upper edge $\ell \to 4$?''}

There the method is indeed at its limit, and this is now a limit of the
problem itself rather than of the way it is computed: without subdivision
($K = 1$) nine of the 644 pointwise pairs fail outright, on
\texttt{I1\_margin\_W} rather than on \texttt{margin\_c} (\S\ref{sec:10.9},
item 4; after $K \le 2$ all of them hold, and that pre-subdivision picture is a diagnostic of an intermediate state
and is not among the stored certificates), while on the continuous $\Qcell$ the same thinness shows up at $\pm 4\cdot 10^{-2}$
($\mathrm{margin}=6.32\cdot 10^{-4}$) and the margin \texttt{margin\_c} vanishes at
$\pm 4.5\cdot 10^{-2}$, while \texttt{margin\_a} stays there at $\approx 8.3\cdot 10^{-2}$. All the inequalities are
strict, but the margin in $r'$ in this zone is small, and widening the region
precisely there now requires a separate (non-uniform in $\ell$) geometric
argument --- see the previous question, path 3.


\subsection{Summary table: what is proved, what is certified, what is assumed}
\label{sec:10.11}

A single map of the section --- for the reader who needs one table instead of
the whole of this section.

\textbf{Registry of assumptions (what the section takes as given).}

{\small
\begin{longtable}{p{0.04\textwidth}p{0.32\textwidth}p{0.22\textwidth}p{0.30\textwidth}}
\caption{Registry of assumptions underpinning the section.}
\label{tab:10.11-assumptions}\\
\toprule
\# & Assumption & Source & What happens if it is false \\
\midrule
\endfirsthead
\multicolumn{4}{@{}l}{\textit{Table \thetable\ (continued)}}\\
\toprule
\# & Assumption & Source & What happens if it is false \\
\midrule
\endhead
\bottomrule
\endlastfoot
A1 & Existence and local uniqueness of $\Gs(\ell)$ on $D_0$, for all 46 boxes; the section uses $\norm{\Gs-P}\leq 1.75\cdot 10^{-7}$, the \emph{uniqueness} radius (the certified distance is far smaller, $\rhostar_X \le 8.49\cdot 10^{-13}$, \S\ref{sec:3.6}) & Theorem A of the package \textbf{[A]} & The whole edifice of the section collapses --- it is its foundation \\
A2 & Exactly one eigenvalue outside the disc on the base layer, simple, real, $=\delta(\ell)$ & Theorem B of the package \textbf{[A]} & Theorem D item 3 (transfer of $\delta(\ell)$) loses its meaning, items 1--2 (compactness/multiplicities) do not \\
A3 & Analyticity of $\delta(\ell)$ in $\ell$ & Theorem C of the package \textbf{[A]} & Consequence 2 falls away, the rest does not \\
A4 & The \texttt{python-flint}/\texttt{arb} arithmetic is correct (interval, with upward rounding) & Common to the whole field, not checked here & All \texttt{[C]} labels lose force simultaneously --- a systemic, not a local risk \\
\end{longtable}
}

\textbf{Map of the proof.}

{\small
\begin{longtable}{p{0.30\textwidth}p{0.19\textwidth}p{0.23\textwidth}p{0.19\textwidth}}
\caption{Map of the proof: status and dependencies of each block.}
\label{tab:10.11-proof}\\
\toprule
Block & Status & Relies on & Where \\
\midrule
\endfirsthead
\multicolumn{4}{@{}l}{\textit{Table \thetable\ (continued)}}\\
\toprule
Block & Status & Relies on & Where \\
\midrule
\endhead
\bottomrule
\endlastfoot
(D1)--(D3), (I1), (I2) for 644 pointwise pairs, $\Qcell$ and the seven corridor cells of \S\ref{sec:10.13} & \textbf{[C]} certified, interval arithmetic & A1, A4 & \S\ref{sec:10.2}, \S\ref{sec:10.6} \\
(E1)--(E2) for the same layers, the corridor cells among them & \textbf{[P]} elementary inequalities in $(c',r')$; not gated by \texttt{one\_box}, but machine-checked for every deposited layer and cell by the numbers gate (block K4) & --- & \S\ref{sec:10.2} \\
D4 ($q<1$, $q_s<1$) & \textbf{removed} --- redundant (not a certificate but a historical stage) & --- & \S\ref{sec:10.2}, \S\ref{sec:10.4} \\
Lemma 1 (continuation of $\Gs$ to $D'$, including pointwise layers) & \textbf{[P]} proved analytically & (I1a)--(I1c) \textbf{[C]}, (E2) \textbf{[P]}, continuity of the margin in $(c,r)$ & \S\ref{sec:10.3} \\
Lemma 2 (compactness, spectral radius) & \textbf{[P]} proved analytically & (D1), (D2), (D3), (I1) \textbf{[C]}, (E2) \textbf{[P]}, Lemma 1 & \S\ref{sec:10.4} \\
Lemma 3 (bootstrap of eigenfunctions) & \textbf{[P]} proved analytically & (I1), (I2) \textbf{[C]}, (E1), (E2) \textbf{[P]}, Lemma 1 & \S\ref{sec:10.4} \\
Lemma 4 (multiplicities) & \textbf{[P]} proved analytically & (E1) \textbf{[P]}, Lemmas 2, 3 & \S\ref{sec:10.4} \\
Theorem D (items 1--2: compactness and coincidence of the non-zero spectra) & \textbf{[P]+[C]} & Lemmas 1--4 & \S\ref{sec:10.5} \\
Theorem D (item 3: $\delta(\ell)$ specifically) & \textbf{[P]+[A]} & items 1--2, A2 & \S\ref{sec:10.5} \\
Transfer $P(\ell)\to\Gs(\ell)$ ($\kappa_G/\Delta U/\Delta W/\Delta s$) --- the exhibited family: the 14 layers, all of $\Qcell$, every corridor cell & \textbf{[P]+[C]} & A1, (I1a)--(I1c), (D1), (D2), (D3), (I2), (E2) & \S\ref{sec:10.7} \\
$R=6/5$ in D3 & not used by the proof of Theorem D (in the package it is needed for the Cauchy tail in $W$) & --- & \S\ref{sec:10.7} \\
Non-vacuousness of $\Ffam$: $\Qcell$ + 14 layers + the seven-cell corridor & \textbf{[C]} & all of the above & \S\ref{sec:10.5}, \S\ref{sec:10.6} \\
\end{longtable}
}

\textbf{What remains open (for the full list see \S\ref{sec:10.9}; only the headings here).}

\begin{table}[h]
\centering
\caption{Open items (headings only; for full descriptions see \S\ref{sec:10.9}).}
\label{tab:10.11-open}
\begin{tabular}{p{0.04\textwidth}p{0.66\textwidth}p{0.22\textwidth}}
\toprule
\# & Open item & Severity \\
\midrule
1 & The whole rectangle $[0.55,0.70]\times[0.70,0.90]$ is NOT certified in its entirety --- only $\Qcell$, 14 points scattered over that rectangle, and the corridor of \S\ref{sec:10.13}, whose cells reach $r = 0.915$ and so leave the rectangle & A limitation of the enclosure, not of correctness \\
2 & Independence only within the family of weighted $\ell^1$-norms, not of an arbitrary space & A limitation of the formulation \\
3 & $\ell$ only $[2,4]$ & A limitation of the formulation \\
4 & At $\ell \to 4$ the method is at its limit (margin $\approx 6\cdot 10^{-4}$) & Numerical thinness, not a defect \\
5 & All 644 certificates have been recomputed by a repeat run of the same certifier (\textbf{644/644}, worst margin $1.09\cdot 10^{-3}$); a cross-check by a different computational route --- only on 6 pairs; recomputation by a second engine independent of \texttt{python-flint}/\texttt{arb} --- not done & Low: repeatability of the algorithm is closed, independence of the implementation and of the arithmetic library is open \\
\bottomrule
\end{tabular}
\end{table}

The bottom line in one sentence: \textbf{the main result (Theorem D, items 1--2 on
$P(\ell)$) is [P]+[C], with no gaps inside the exhibited family $\Ffam$; the full
conclusion about the solution $\Gs(\ell)$ (the transfer of \S\ref{sec:10.7}) has no gaps
inside it either --- the certified domain $c+r\leq 1.53$ contains that family whole,
its widest point being the corridor cell at $1.522$; the remaining open
items are all about the enclosure (how large $\Ffam$ is) or about the formulation
(in which class of spaces), and none about the correctness of what has already
been exhibited.}

\textbf{What has been checked, and by what.} The verification of this article is
not uniform. Some of the statements are re-checked by machine at every build
of the package: for these a canary has been set up --- a deliberately corrupted
copy on which the corresponding check is obliged to go red, and this has been
verified, not merely asserted. Some are a one-off numerical sweep: an
independent recomputation performed once outside the package code and recorded
with a date and a report. Some have been read by a human, usually in two
independent passes, including the final blind focused readings, which knew
nothing of the previous night's edits. The statements of the theorems, the
decisive thresholds and the load-bearing prose are locked by a hash ratchet ---
an editorial safeguard over the manuscript, not one of the verification scripts
and not part of the deposited package:
it does not understand the meaning of what is written, but it does not allow a
statement to be silently weakened or substituted without an explicit entry in
the edit log. The numbers of the article have been reconciled with the
certificates they refer to, and addresses of the form ``\S N'' with the body of
the section named. A few places (the direction of an inequality in a geometric
or numerical argument, the sufficiency of an appeal to ``continuity'', the
rigour of the construction of a series) remain a consciously accepted boundary
of the method: checking that an argument is correct, and not merely that the
text has not been changed, requires a human understanding of the proof --- this
is open, not concealed, and is discussed separately.


\subsection{Where the statement is strongest}
\label{sec:10.12}

Theorem D (items 1--2, on the interpolant $P(\ell)$) is proved at once for the
whole of the exhibited $\Ffam$. The full conclusion about the solution $\Gs(\ell)$
additionally requires the transfer of \S\ref{sec:10.7}, certified on $c+r\leq 1.53$ (see
\S\ref{sec:10.5}, the remark after Consequence 2) --- a domain that contains the exhibited
family whole, corridor included. It is nevertheless useful to distinguish three parts by the
strength of their grounding, because they can be accepted separately. The
\textbf{14 pointwise layers}, scattered over $[0.55,0.70]\times[0.70,0.90]$, are the
most unconditional part: there (D1)--(D3), (I1), (I2) are certified box by
box, while (E1)--(E2) are checked by the numbers gate for each deposited layer
rather than by the certifier (they are not gated
automatically by \texttt{one\_box} --- see the caveat in \S\ref{sec:10.2});
Lemma 1 provides an $\varepsilon$-margin from the
continuity of the margin (\S\ref{sec:10.3}), and the transfer of $\kappa_G$ is rigorous for any
layer of the presented family $\Ffam$ without exception (\S\ref{sec:10.7}). A refinement about the widths: the widest of the
14 layers --- \texttt{(0.70,0.80)}, \texttt{(0.60,0.90)}, \texttt{(0.65,0.85)} --- have $c+r=3/2$, and the
widest exhibited point overall is the corridor cell reaching $1.522$. Both sit inside the
domain on which the majorant $\sup|U-1|\leq 0.400$ is certified
(\texttt{CR\_\allowbreak WIDE} $=1.53$, the actual value being \texttt{0.399346}), so the layers, $\Qcell$
and the corridor are covered alike;
at the same time the $\DT$ matrices on these layers
differ by 9.8--118\,\% (a float64 illustration, not an interval certificate, \S\ref{sec:10.8}), so the coincidence of the non-zero spectra is
meaningful.
The word ``unconditional'' here refers to the completeness of the certification
of the conditions, not to the number of implementations: the section has a
single certifier (\texttt{04-\allowbreak CODE/\allowbreak layer/\allowbreak }), and there is no independent second
computation of the margins \texttt{D\_margin\_a}/\texttt{D\_margin\_c} at all (\S\ref{sec:10.9}, item 5).
The \textbf{triangle $\Qcell_{\mathrm{in}}$} (25\,\% of the area of $\Qcell$, where $D(c,r) \subseteq D_0$) is the
most conservative continuous part: there the transfer of \S\ref{sec:10.7} follows from the
direct inclusion of the discs, without the argument via $s(D') \subseteq D_0$. The
\textbf{remaining 75\,\% of $\Qcell$} ($\Qcell_{\mathrm{out}}$) rest on the same apparatus of
certificates, but the transfer there proceeds via (I1a). All three parts are
\textbf{[P]+[C]}, none of them hangs on an unproved step; the only incompleteness
of the section is the enclosure (what part of the rectangle
$[0.55,0.70]\times[0.70,0.90]$ is exhibited continuously), not the proof (\S\ref{sec:10.9},
item 1).


\subsection{The corridor: a connected chain instead of isolated points}
\label{sec:10.13}

The first of the three parts distinguished in \S\ref{sec:10.12} rests on 14
separate pointwise layers. Formally this is sufficient (the certificates of \S\ref{sec:10.2} are checked for
each layer independently), but a reviewer may still be left with the question
``why exactly these 14 points --- were they not chosen so as to skirt the
problematic places of the family?''. The answer is to construct an explicit
certified CONNECTED chain joining the base layer to one of the separately
exhibited ones, rather than merely displaying scattered points.

An L-shaped corridor of 7 overlapping cells has been constructed (with the
same \texttt{layer\_cert} engine as $\Qcell$; driver \texttt{04-\allowbreak CODE/\allowbreak layer/\allowbreak layer\_\allowbreak corridor.py}):
from the base
layer $(0.60,0.80)$ vertically up to $(0.60,0.90)$, then horizontally to
$(0.55,0.90)$ --- one of the three furthest layers equidistant in the
Euclidean metric (on a par with $(0.55,0.70)$, $(0.70,0.75)$), and the only
one where K-escalation is needed already at a single point ($K=2$, \S\ref{sec:10.6}). The
worst margin among all 644 pairs is not here, however, but at
\texttt{(0.7,0.75)/\allowbreak g00} --- yet it is precisely \texttt{(0.55,0.90)} that is thinnest with
respect to the subdivision requirement, which is why the corridor leads to it.
The path deliberately passes next to the already exhibited \texttt{(0.60,0.85)} and
\texttt{(0.60,0.90)} --- the corridor does not skirt them, but connects them. \textbf{[C]}

\textbf{Enclosure (important for reading Theorem D along the corridor).}
\texttt{layer\_\allowbreak corridor.py} certifies D1/D2/D3/I1/I2 for each cell; together with
(E1)/(E2), which that program does not gate and which the numbers gate checks
(\S\ref{sec:10.2}), this is their membership in $\Ffamz$. On these cells the hand check is
immediate: $c \le 0.612 < 0.800 \le r$ gives (E1), and $|1-c| \le 0.452 < 0.800 \le r$ gives (E2). The transfer $P(\ell)\to\Gs(\ell)$ (\S\ref{sec:10.7}, $\kappa_G/\Delta U/\Delta W/\Delta s$) is tabulated separately, in \texttt{kappa\_\allowbreak full.json}: 18 explicitly exhibited pointwise layers and corners of $\Qcell$, the base layer, and the far corners of the last four of the seven corridor cells of Table~\ref{tab:10.13-corridor} --- 23 \texttt{layer} labels in the file, $46\times 23 = 1058$ records.
Since only the sum $c+r$ enters the formula and $\kappa_G$ is monotone in it, the far corner of a cell majorizes the whole cell, and a cell is covered as soon as the list holds a label at a $c+r$ at least as large. Run that over the seven. The exhibited layers reach $c+r = 1.50$, which already covers cells one to three (far corners $1.442$, $1.462$, $1.482$) and the seventh ($1.491$). Cells four and six ($1.502$ and $1.514$) exceed $1.50$, so covering them takes a label from the corridor itself; the fifth cell's, at $c+r = 1.522$, is the widest label in the file, and it covers both of them --- and the fifth cell, which nothing else reaches. Keeping that one label therefore covers all seven, so three of the four corridor labels are redundant, not one. The file carries all four because they are simply the last four rows of Table~\ref{tab:10.13-corridor}; keeping the three redundant ones costs nothing, while dropping the label at $c+r = 1.522$ would leave the fifth cell uncovered. The widest point of the corridor, $c+r=1.522$, lies inside the domain $c+r\leq 1.53$ on which the wide majorant $\sup|U-1|\leq 0.400$ is certified, so the transfer covers the corridor as it covers the rest of the exhibited family; the corridor therefore does two things at once --- it joins the base layer to one of the farthest of the exhibited layers by a certified connected chain, and it carries the full conclusion of Theorem D (about $\Gs$, not only about $P$) along its own cells. The margins are not tight anywhere along it: the worst cell margins $1.06\cdot 10^{-2}$, $4.57\cdot 10^{-3}$, $1.58\cdot 10^{-3}$, $1.14\cdot 10^{-3}$ stand each against the perturbation $\Delta W$ of its own cell --- $8.890$, $8.983$, $8.946$ and $8.839$, all $\cdot 10^{-6}$ (the second of them is at the same time the maximum of $\Delta W$ over the 23 labels, row 54 of Appendix~\ref{sec:A}) --- that is, factors of $1192\times$, $508\times$, $176\times$ and $128\times$.

\begin{table}[h]
\centering
\caption{Seven corridor cells, all passing 46/46 boxes; worst margin narrows monotonically.}
\label{tab:10.13-corridor}
\begin{tabular}{cccc}
\toprule
Cell ($c\times r$) & pass & worst margin & worst box \\
\midrule
$[0.588,0.612]\times[0.800,0.830]$ & 46/46 & $7.94\cdot 10^{-2}$ & g43 \\
$[0.588,0.612]\times[0.820,0.850]$ & 46/46 & $6.24\cdot 10^{-2}$ & g43 \\
$[0.588,0.612]\times[0.840,0.870]$ & 46/46 & $3.67\cdot 10^{-2}$ & g43 \\
$[0.588,0.612]\times[0.860,0.890]$ & 46/46 & $1.06\cdot 10^{-2}$ & g43 \\
$[0.588,0.612]\times[0.880,0.910]$ & 46/46 & $4.57\cdot 10^{-3}$ & g39 \\
$[0.571,0.599]\times[0.885,0.915]$ & 46/46 & $1.58\cdot 10^{-3}$ & g35 \\
$[0.548,0.576]\times[0.885,0.915]$ & 46/46 & $1.14\cdot 10^{-3}$ & g31 \\
\bottomrule
\end{tabular}
\end{table}

\textbf{In total: 7/7 cells passed, worst margin $1.14\cdot 10^{-3}$} (the last cell, the
one covering the end point $(0.55,0.90)$), total computation time $2091$~s
($\approx 35$ min). The margin narrows honestly and monotonically as one moves towards
the far corner --- exactly the behaviour that was expected from the separately
exhibited points (\S\ref{sec:10.6}, \S\ref{sec:10.9} item 4): not an artefact, but the geometry of the
approach to the boundary of $\Ffamz$ at $\ell \to 4$, visible along the whole path. At no
point did the margin cross zero --- the corridor is nowhere broken.

Significance for \S\ref{sec:10.12}: step 1 (``14 separate layers'') can now be
strengthened --- at least the base and one of the most distant of the 14 layers are
joined by an explicitly exhibited continuous certified path, and not merely by
the coincidence of certificates at isolated points. The same device would apply
to any other pair among the 14 layers where the geometry permits (the
off-center cell of \S\ref{sec:10.6} shows that it need not: it fails on
\texttt{margin\_a}/\texttt{margin\_c}, not on the withdrawn $q$ gate), though no corridor for another pair has
been computed --- the corridor was built as a demonstration
of the method, not as an exhaustive covering of the whole of
$[0.55,0.70]\times[0.70,0.90]$ (that remains open item \S\ref{sec:10.9}, item 1). \textbf{[C]}




\clearpage
\appendix

\section{All constants with their provenance}
\label{sec:A}

\textbf{Rounding convention (in force throughout the text).} Published numbers are
rounded OUTWARD from the certified fact, not to nearest: a majorant (``$\le X$'')
upward, a minorant (``$\ge X$'') downward, both ends of a range or of an enclosing
interval --- spreading apart. The printed digits therefore often do NOT agree
with round-to-nearest (for example, the range $B_A$ is printed as
$[2.865, 3.617]$ while the fact is $[2.86551, 3.61620]$): this is not a
rounding error but deliberate conservativeness. The opposite direction (a
bound rounded inward) is a defect. In the prose of \S3 and \S4 it is caught by \texttt{\_numbers\_check.py} (block F), in the table below by \texttt{\_appendix\_check.py} (\texttt{check\_direction}). A truncated value is distinguished from a bound by an ellipsis ($1.1385\ldots\times 10^{-42}$): it makes no claim about a side, and the check does not demand one.

The table below lists the constants of the \textbf{pointwise certificates} of \appref{sec:5}
($\ell = 3$ and $\ell = 5/2$): each of them has a single value, and it is convenient to
exhibit it as a list. The constants of the parametric certificates of \S\ref{sec:3} and
\S\ref{sec:4} are not numbers but ranges over 46 and 706 boxes; they are collected in the
tables of \S\ref{sec:3.6} and \S\ref{sec:4.8}, and box by box they reside in
\texttt{03-\allowbreak CERTIFICATES/\allowbreak existence/\allowbreak box\_\allowbreak g*.json} and \texttt{03-\allowbreak CERTIFICATES/\allowbreak tube/\allowbreak box\_\allowbreak t*.json}
(the meaning of each field is given in \texttt{05-\allowbreak DATA-\allowbreak INDEX.md}). The test
\texttt{04-\allowbreak CODE/\allowbreak tests/\allowbreak \_\allowbreak numbers\_\allowbreak check.py} reconciles both ends of every range in the
package's copy of the article against the actual minimum and maximum in the
certificates, so there
is no point in duplicating every range here as a list. The few parametric
constants that do carry a statement of their own --- the worst $\hat q$ and
$\rho^{\star}_{X}$ over the 46 boxes, and the factor at which the gate breaks
with respect to $\tau$ --- are listed at the end of the table, together with
the constants of the layer family.

{\footnotesize\setlength{\tabcolsep}{2pt}
\begin{longtable}{>{\raggedright\arraybackslash}p{0.065\textwidth}>{\raggedright\arraybackslash}p{0.275\textwidth}>{\raggedright\arraybackslash}p{0.345\textwidth}>{\raggedright\arraybackslash}p{0.205\textwidth}}
\toprule
\textbf{\#} & \textbf{Name} & \textbf{Value} & \textbf{Source file} \\
\midrule
\endhead
1 & $c$ & $3/5$ & \texttt{03-\allowbreak CERTIFICATES/\allowbreak pointwise/\allowbreak kant\_\allowbreak l3.json::\allowbreak c} \\
2 & $r$ & $4/5$ & \texttt{03-\allowbreak CERTIFICATES/\allowbreak pointwise/\allowbreak kant\_\allowbreak l3.json::\allowbreak r} \\
3 & $N_{p}$ & $80$ & \texttt{03-\allowbreak CERTIFICATES/\allowbreak pointwise/\allowbreak kant\_\allowbreak l3.json::\allowbreak N\_\allowbreak p} (method parameter; introduced in \S\ref{sec:2}, operative in \S\ref{sec:4.6}) \\
4 & $\bar{\alpha}$ (seed) & $1.92769096384764\ldots$ & \texttt{03-\allowbreak CERTIFICATES/\allowbreak pointwise/\allowbreak kant\_\allowbreak l3.json::\allowbreak abar} \\
5 & $m_{1} = |\bar{G}(1)|$ & $0.5187553496666\ldots$ & \texttt{03-\allowbreak CERTIFICATES/\allowbreak pointwise/\allowbreak kant\_\allowbreak l3.json::\allowbreak m1} \\
6 & $\norm{\bar{G}'}_{X}$ & $2.108543823369\ldots$ & \texttt{03-\allowbreak CERTIFICATES/\allowbreak pointwise/\allowbreak kant\_\allowbreak l3.json::\allowbreak Gprime\_\allowbreak norm} \\
7 & $\norm{t}_{X}$ & $1.4$ & \texttt{03-\allowbreak CERTIFICATES/\allowbreak pointwise/\allowbreak kant\_\allowbreak l3.json::\allowbreak t\_\allowbreak norm} \\
8 & $\hat{\rho}'$ & $1.0\times 10^{-20}$ & \texttt{03-\allowbreak CERTIFICATES/\allowbreak pointwise/\allowbreak kant\_\allowbreak l3.json::\allowbreak rho\_\allowbreak prime} \\
9 & $\hat{\rho}_{X}$ & $1.75\times 10^{-20}$ & \texttt{03-\allowbreak CERTIFICATES/\allowbreak pointwise/\allowbreak kant\_\allowbreak l3.json::\allowbreak rho\_\allowbreak X} \\
10 & $\norm{\bar{U}}_{X}$ (strict) & $1.0487118131237275\ldots$ & \texttt{03-\allowbreak CERTIFICATES/\allowbreak pointwise/\allowbreak fixL\_\allowbreak result.json::\allowbreak norm\_\allowbreak U\_\allowbreak strict\_\allowbreak upper} \\
11 & $\hat{q}$ & $0.6903152621075212\ldots$ & \texttt{03-\allowbreak CERTIFICATES/\allowbreak pointwise/\allowbreak fixL\_\allowbreak result.json::\allowbreak q\_\allowbreak hat} \\
12 & $\kappa_{s}$ & $0.6279201508804733\ldots$ & \texttt{03-\allowbreak CERTIFICATES/\allowbreak pointwise/\allowbreak fixL\_\allowbreak result.json::\allowbreak kappa\_\allowbreak s} \\
13 & CD2 & $10.42702451618060\ldots$ & \texttt{03-\allowbreak CERTIFICATES/\allowbreak pointwise/\allowbreak fixL\_\allowbreak result.json::\allowbreak CD2} \\
14 & $\Delta\alpha_2$ & $2.508\times 10^{-19}$ & \texttt{03-\allowbreak CERTIFICATES/\allowbreak pointwise/\allowbreak fixL\_\allowbreak result.json::\allowbreak Delta\_\allowbreak alpha2} \\
15--18 & $T_2, T_3, \Delta T_4, \Delta T_5$ (upper bounds, rounded upward) & $1.61\times 10^{-18}, 4.98\times 10^{-18}, 1.04\times 10^{-17}, 1.93\times 10^{-17}$ & \texttt{03-\allowbreak CERTIFICATES/\allowbreak pointwise/\allowbreak fixL\_\allowbreak result.json::\allowbreak T2, T3, Delta\_\allowbreak T4, Delta\_\allowbreak T5} \\
19 & $\mathbf{L}$ & $\mathbf{2.0648725990526113 \times 10^{3}}$ & \texttt{03-\allowbreak CERTIFICATES/\allowbreak pointwise/\allowbreak fixL\_\allowbreak result.json::\allowbreak L\_\allowbreak strict} \\
20 & $Y_{0}$ & $1.138567618961918\ldots\times 10^{-42}$ & \texttt{03-\allowbreak CERTIFICATES/\allowbreak pointwise/\allowbreak kant\_\allowbreak l3.json::\allowbreak Y0} \\
21 & $Z_{1}$ & $\mathbf{5.20010896072677\ldots\times 10^{-8}}$ (rigorous closed-form $\hat{\beta}$); $1.47421154471446\ldots\times 10^{-7}$ --- the earlier estimate with a float fit & \texttt{03-\allowbreak CERTIFICATES/\allowbreak pointwise/\allowbreak Z1c\_\allowbreak l3\_\allowbreak strict.json::\allowbreak Z1\_\allowbreak strict} + \texttt{03-\allowbreak CERTIFICATES/\allowbreak pointwise/\allowbreak fixL\_\allowbreak result.json::\allowbreak Z1} \\
22 & $\Zhat_{\mathrm{new}}$ & $5.2001090\times 10^{-8}$ (with the rigorous closed-form $\hat{\beta}$); $1.4742115467\ldots\times 10^{-7}$ --- with the earlier $Z_{1}$ (a historical field, not updated) & \texttt{03-\allowbreak CERTIFICATES/\allowbreak pointwise/\allowbreak fixL\_\allowbreak result.json::\allowbreak \{Z\_\allowbreak hat\_\allowbreak new\_\allowbreak strict,Z\_\allowbreak hat\_\allowbreak new\}} \\
23 & $\rhostar_{\mathrm{new}}$ & $1.1385677868\ldots\times 10^{-42}$ (from the historical \texttt{Z\_hat\_new}, see row 22; with the rigorous \texttt{Z\_\allowbreak hat\_\allowbreak new\_\allowbreak strict} it would be $1.1385676782\times 10^{-42}$, i.e.\ slightly SMALLER; since $\rhostar$ is an enclosure radius, the larger of the two is the conservative one, and it is the one printed here) & \texttt{03-\allowbreak CERTIFICATES/\allowbreak pointwise/\allowbreak fixL\_\allowbreak result.json::\allowbreak rho\_\allowbreak star\_\allowbreak new} \\
24 & $\beta_{81}$ & $3.45728670541315\ldots\times 10^{-8}$ & \texttt{03-\allowbreak CERTIFICATES/\allowbreak pointwise/\allowbreak df\_\allowbreak package.json::\allowbreak beta\_\allowbreak 81} \\
24a & $\hat{\beta}_{N_{p}+1}$ (the tail estimate in the Kantorovich chain --- a different quantity, not row 24) & $3.45728350665437\ldots\times 10^{-8}$ & \texttt{03-\allowbreak CERTIFICATES/\allowbreak pointwise/\allowbreak kant\_\allowbreak l3.json::\allowbreak beta\_\allowbreak Np1} \\
25 & $\norm{\Ahat}_{1}$ & $3.1399814582478575$ & \texttt{03-\allowbreak CERTIFICATES/\allowbreak pointwise/\allowbreak kant\_\allowbreak l3.json::\allowbreak A\_\allowbreak norm1} \\
26 & $\deltastar$ & $6.084689074262122$ (Krawczyk box width $6.226\times 10^{-19}$) & \texttt{03-\allowbreak CERTIFICATES/\allowbreak pointwise/\allowbreak eigen\_\allowbreak l3\_\allowbreak fixed.json::\allowbreak delta\_\allowbreak box\_\allowbreak intersect} + \texttt{03-\allowbreak CERTIFICATES/\allowbreak pointwise/\allowbreak task\_\allowbreak A\_\allowbreak krawzyk\_\allowbreak result.json::\allowbreak K\_\allowbreak delta\_\allowbreak width} \\
27 & $\varepsilon_{v,\mathrm{loose}}$ & $2.0507\times 10^{-12}$ & \texttt{03-\allowbreak CERTIFICATES/\allowbreak pointwise/\allowbreak eigen\_\allowbreak l3\_\allowbreak fixed.json::\allowbreak eps\_\allowbreak v\_\allowbreak loose} \\
28 & $\varepsilon_{\varphi,\mathrm{loose}}$ & $5.6820\times 10^{-9}$ & \texttt{03-\allowbreak CERTIFICATES/\allowbreak pointwise/\allowbreak eigen\_\allowbreak l3\_\allowbreak fixed.json::\allowbreak eps\_\allowbreak phi\_\allowbreak loose} \\
29 & $\langle\varphi^{\star},v\rangle$, ball & lower end $1.0858883467601657$, upper end $1.0858883660218137$ (the division in the next row starts from the UPPER end --- that is how the certifier computed it) & \texttt{03-\allowbreak CERTIFICATES/\allowbreak pointwise/\allowbreak eigen\_\allowbreak l3\_\allowbreak fixed.json::\allowbreak \{inner\_\allowbreak product.abs\_\allowbreak min, inner\_\allowbreak product.hi\}} \\
30 & $\varepsilon_{\mathrm{DT}}$ & $3.613527048343\times 10^{-17}$ & \texttt{03-\allowbreak CERTIFICATES/\allowbreak pointwise/\allowbreak hyperb\_\allowbreak l3.json::\allowbreak majority\_\allowbreak computation.eps\_\allowbreak DT} \\
31 & $\delta/\langle\varphi,v\rangle$, division by the UPPER end of row 29 $\then$ this is the LOWER end of the ratio; the upper end (division by the lower one, $5.6034205472563325\ldots$) resides in \texttt{q\_\allowbreak perturb\_\allowbreak patch.json::\allowbreak caps.delta\_\allowbreak hat\_\allowbreak over\_\allowbreak ip\_\allowbreak hat}, while the spread between the ends, $9.9394\times 10^{-8}$ (stored under its own name, \texttt{caps.delta\_\allowbreak over\_\allowbreak ip}), is the quantity from which $E$'s term \texttt{term\_a} is formed --- \texttt{term\_a} itself holds $1.6844\times 10^{-7}$, not the spread & $5.603420447862033716\ldots$ & \texttt{03-\allowbreak CERTIFICATES/\allowbreak pointwise/\allowbreak hyperb\_\allowbreak l3.json::\allowbreak majority\_\allowbreak computation.defl\_\allowbreak coef\_\allowbreak value} \\
32 & $\norm{\hat{Q}}_{1}$ & $1.3969171932439115$ & \texttt{03-\allowbreak CERTIFICATES/\allowbreak pointwise/\allowbreak q\_\allowbreak perturb\_\allowbreak patch.json::\allowbreak floats.qhat\_\allowbreak 1norm} \\
33 & $\norm{\hat{Q}^{16}}_{1}$ & $0.2899488032200199$ & \texttt{03-\allowbreak CERTIFICATES/\allowbreak pointwise/\allowbreak q\_\allowbreak perturb\_\allowbreak patch.json::\allowbreak floats.Q16\_\allowbreak HAT} \\
34 & \textbf{$\mathbf{\norm{{Q^{\star}}^{16}}_{1}}$ (patched)} & $\mathbf{0.2905947}$ & \texttt{03-\allowbreak CERTIFICATES/\allowbreak pointwise/\allowbreak q\_\allowbreak perturb\_\allowbreak patch.json::\allowbreak summary.sum\_\allowbreak q16} \\
35 & $\rho(Q^{\star})$ upper --- the sixteenth root of the row above; it is not stored in the json as a separate number & $0.9256685\ldots$ & \texttt{03-\allowbreak CERTIFICATES/\allowbreak pointwise/\allowbreak q\_\allowbreak perturb\_\allowbreak patch.json::\allowbreak summary.sum\_\allowbreak q16} \\
36 & $E = \norm{Q^{\star}-\hat{Q}}_{1}$ & $2.681787\times 10^{-7}$ & \texttt{03-\allowbreak CERTIFICATES/\allowbreak pointwise/\allowbreak q\_\allowbreak perturb\_\allowbreak patch.json::\allowbreak summary.E} \\
37 & $u_{0}$ ($\ell=2.5$) & $0.8531715032047797\ldots$ & \texttt{03-\allowbreak CERTIFICATES/\allowbreak pointwise/\allowbreak xresid\_\allowbreak l25.json::\allowbreak u\_\allowbreak 0} \\
38 & $\norm{\bar{F}}_{X}$ ($\ell=2.5$) & $2.538195012738\times 10^{-33}$ (rounded upward; exactly $2.538195012737448990\ldots\times 10^{-33}$) & \texttt{03-\allowbreak CERTIFICATES/\allowbreak pointwise/\allowbreak xresid\_\allowbreak l25.json::\allowbreak F\_\allowbreak norm\_\allowbreak strict} \\
39 & $\norm{h_{W}}$ ($\ell=2.5$) & $4.386547958464\times 10^{-34}$ (rounded upward; exactly $4.386547958463239236\ldots\times 10^{-34}$) & \texttt{03-\allowbreak CERTIFICATES/\allowbreak pointwise/\allowbreak xresid\_\allowbreak l25.json::\allowbreak tail\_\allowbreak hW} \\
40 & $\norm{W_{\mathrm{ball}}-c}_{X}$ ($\ell=2.5$) & $0.5282423818964933\ldots$ & \texttt{03-\allowbreak CERTIFICATES/\allowbreak pointwise/\allowbreak xresid\_\allowbreak l25.json::\allowbreak Wc\_\allowbreak norm} \\
41 & $\norm{\Ahat}_{1}$ ($\ell=2.5$) & $2.9917507063308273$ (midpoint of the ball; its radius is $3.1\cdot 10^{-153}$, rounded up --- exactly $3.09\cdot 10^{-153}$ --- so the midpoint and either endpoint agree far beyond every digit printed) & \texttt{03-\allowbreak CERTIFICATES/\allowbreak pointwise/\allowbreak kant\_\allowbreak l25.json::\allowbreak A\_\allowbreak tilde.norm\_\allowbreak 1.mid} + \texttt{03-\allowbreak CERTIFICATES/\allowbreak pointwise/\allowbreak kant\_\allowbreak l25.json::\allowbreak A\_\allowbreak tilde.norm\_\allowbreak 1.rad} \\
\bottomrule
\end{longtable}}

\noindent\emph{On the derived values in the tables of this appendix.} Three
rows print
a value that is computed rather than stored --- in one row beside the stored
number, in another beside a second computed one, and in the third on its own ---
and all three are reproducible from the fields cited in the rows named below. In the row for
$\rhostar_{\mathrm{new}}$ the two numbers are $Y_{0}/(1-\Zhat)$ taken with the
historical and with the strict $\Zhat$ of the preceding row, $Y_{0}$ being the
field recorded three rows earlier under its own name. In the row for the
deviation of $\delta$ across
layers, the value is the largest $|\delta - \delta_{\mathrm{base}}|$ over the
non-base layers and the full spread is the difference between the largest and
the smallest $\delta$ over all of them, both from the list of layer records in
the probe file. In the row for the upper bound on $\rho(Q^{\star})$ the number
is the sixteenth root of the row above it, and the row says so itself; both
rows cite the same stored field.

\textbf{Theorem D (\S\ref{sec:10}), independence of the layer.} Rows \textup{42--58} take their values from
\texttt{03-\allowbreak CERTIFICATES/\allowbreak layer/\allowbreak }; margins are published as the lower ends of the
intervals, perturbations as the upper ends. The table then closes with rows
\textup{59--64}: rows \textup{59--61} are the pointwise constants of the K-mode at
$\ell=3$, and rows \textup{62--64} the parametric constants announced above. None of
the six is part of the layer family, and all cite \texttt{03-\allowbreak CERTIFICATES/\allowbreak pointwise/\allowbreak } in their own source column.

{\footnotesize\setlength{\tabcolsep}{3pt}
\begin{longtable}{>{\raggedright\arraybackslash}p{0.070\textwidth}>{\raggedright\arraybackslash}p{0.250\textwidth}>{\raggedright\arraybackslash}p{0.230\textwidth}>{\raggedright\arraybackslash}p{0.290\textwidth}}
\toprule
No. & \textbf{quantity} & \textbf{value} & \textbf{origin} \\
\midrule
\endhead
42 & cell $\Qcell$ (continuum of layers) & $(0.56, 0.64) \times (0.76, 0.84)$ & \texttt{layer/\allowbreak cells3\_\allowbreak bisect.json} (half-width $4\cdot 10^{-2}$) \\
43 & worst margin on $\Qcell$ & $6.32\times 10^{-4}$ (box \texttt{g42}) & \texttt{layer/\allowbreak cells3\_\allowbreak bisect.json} \\
44 & worst margin on the sub-cell $\pm 3.5\cdot 10^{-2}$ & $3.30\times 10^{-3}$ (box \texttt{g39}) & \texttt{layer/\allowbreak cells3\_\allowbreak bisect.json} \\
45 & first failing half-width & $\pm 4.5\cdot 10^{-2}$, 44/46, gate \texttt{lemma\_U} & \texttt{layer/\allowbreak cells3\_\allowbreak bisect.json} \\
46 & worst margin over all 644 pairs & $1.0965\times 10^{-3}$ (\texttt{(0.7,0.75)}/\texttt{g00}) & \texttt{layer/\allowbreak full\_\allowbreak recheck\_\allowbreak 644.json} \\
47 & \texttt{D\_margin\_a}, worst & $2.44\times 10^{-2}$ (\texttt{(0.70,0.75)}/\texttt{g00}) & \texttt{layer/\allowbreak worst\_\allowbreak by\_\allowbreak gate.json} \\
48 & \texttt{D\_margin\_c}, worst & $5.28\times 10^{-2}$ (\texttt{(0.55,0.80)}/\texttt{g43}) & \texttt{layer/\allowbreak worst\_\allowbreak by\_\allowbreak gate.json} \\
49 & \texttt{I1\_margin\_s} / \texttt{I1\_margin\_W}, worst & 0.126 / $1.95\times 10^{-3}$ (\texttt{(0.55,0.90)}) & \texttt{layer/\allowbreak worst\_\allowbreak by\_\allowbreak gate.json} \\
50 & \texttt{I2\_margin\_s} / \texttt{I2\_margin\_W}, worst & $1.09\times 10^{-3}$ / $6.80\times 10^{-3}$ & \texttt{layer/\allowbreak worst\_\allowbreak by\_\allowbreak gate.json} \\
51 & $q$ (gate D4 lifted), maximum & 1.2352 (\texttt{(0.7,0.8)}/\texttt{g43}) --- diagnostic, not a condition & \texttt{layer/\allowbreak worst\_\allowbreak by\_\allowbreak gate.json} \\
52 & $\kappa_{G}$ on the base layer & 4.4197 (box \texttt{g43}), published as $\le 4.42$; this is the coarse subdivision $K = 1$ --- the $\kappa_{G} \le 4.14$ of \S\ref{sec:6.4.3} is the same quantity at $K = 8$, and \S\ref{sec:10.7} reconciles the two & \texttt{layer/\allowbreak kappa\_\allowbreak full.json} \\
53 & $\kappa_{G}$ maximum over the 23 \texttt{layer} labels of the file (18 exhibited layers/corners, the base layer, and the far corners of the last four of the seven corridor cells of \S\ref{sec:10.13}) & 4.7177 (label \texttt{corridor cell (0.612,0.910)}, box \texttt{g43}; $c+r = 0.612+0.910 = 1.522$) & \texttt{layer/\allowbreak kappa\_\allowbreak full.json} \\
54 & $\Delta W$ maximum over the same 23 labels & $8.983\times 10^{-6}$ & \texttt{layer/\allowbreak kappa\_\allowbreak full.json} \\
55 & $\Delta W$ at the worst corner of $\Qcell$ \texttt{(0.64,0.84)}/\texttt{g43} & $8.787\times 10^{-6}$ & \texttt{layer/\allowbreak kappa\_\allowbreak full.json} \\
56 & worst corridor margin (7 cells) & $1.14\times 10^{-3}$ & \texttt{layer/\allowbreak corridor.json} \\
57 & maximum deviation of $\delta$ from the base layer over 14 layers (full spread --- $1.51\times 10^{-14}$) & $8.882\times 10^{-15}$ --- float64 diagnostic, not an interval certificate & \texttt{layer/\allowbreak dt\_\allowbreak layer\_\allowbreak probe.json}: field \texttt{layers[].\allowbreak delta} --- the value is $\max|\delta - \delta_{\mathrm{base}}|$ over the 14 non-base layers, the full spread is $\max - \min$ over all 15 \\
58 & maximum discrepancy of the $\mathrm{DT}$ matrices & $118\,\%$ (\texttt{(0.70,0.75)}) --- float64 diagnostic, not an interval certificate & \texttt{layer/\allowbreak dt\_\allowbreak layer\_\allowbreak probe.json} \\
59 & $\bar{G}(1)$ ($\ell=3$) & $-0.5187553496666370$ & \texttt{03-\allowbreak CERTIFICATES/\allowbreak pointwise/\allowbreak kmode\_\allowbreak l3.json::\allowbreak G\_\allowbreak at\_\allowbreak 1} \\
60 & $\bar{G}'(1)$ ($\ell=3$) & $-1.2386641506999486$ & \texttt{03-\allowbreak CERTIFICATES/\allowbreak pointwise/\allowbreak kmode\_\allowbreak l3.json::\allowbreak Gprime\_\allowbreak at\_\allowbreak 1} \\
61 & $\lambda_{K} = 3\bar{G}'(1)/\bar{G}(1)$ & $7.1632850716389135$ & \texttt{03-\allowbreak CERTIFICATES/\allowbreak pointwise/\allowbreak kmode\_\allowbreak l3.json::\allowbreak lambda\_\allowbreak K} \\
62 & $\hat{q}$ maximum over 46 boxes (premise a9; rounded up to 0.955131 in the text) & $0.9551308853551745$ (\texttt{g43}) & \texttt{03-\allowbreak CERTIFICATES/\allowbreak pointwise/\allowbreak holo\_\allowbreak radius.json::\allowbreak \{q\_\allowbreak hat\_\allowbreak max, q\_\allowbreak hat\_\allowbreak max\_\allowbreak box\}} \\
63 & worst $\rho^{\star}_{X}$ over 46 boxes & $8.482891\times 10^{-13}$ (\texttt{g41}) & \texttt{03-\allowbreak CERTIFICATES/\allowbreak pointwise/\allowbreak holo\_\allowbreak radius.json::\allowbreak \{rho\_\allowbreak X\_\allowbreak max\_\allowbreak field, rho\_\allowbreak X\_\allowbreak max\_\allowbreak box\}} \\
64 & factor at which the gate breaks with respect to $\tau$, worst of the three boxes in \S\ref{sec:3.7} & $4.6\times 10^{5}$ (\texttt{g00}) --- directional, not exact: a single seed-dependent measurement (\S\ref{sec:3.7}) & \texttt{03-\allowbreak CERTIFICATES/\allowbreak pointwise/\allowbreak tau\_\allowbreak share.json::\allowbreak boxes.g00.tau\_\allowbreak break\_\allowbreak factor} \\
\bottomrule
\end{longtable}}